% !TeX root = ../Book2.tex
\BookPart{实变调和分析}{b2:part:04}
% 本编导读（章级标题，不计入编号）
\BookPartGuide

Fourier 侧的形式计算并不能直接给出所有 \(L^p\) 估计。本编先用极大函数描述局部平均与层集，再以好坏分解把 \(L^2\) 有界性推广到非 Hilbert 指数；频率环带与平方函数随后把光滑度分成尺度。最后的选读章研究空间和频率不能同时过度集中的定量与定性原理。

\ProseParagraphBreak
编内主要依赖关系为
\[
\begin{tikzcd}[column sep=large,row sep=large]
 \text{第\ref{b2:ch:17}章：极大与插值}
   \arrow[r] \arrow[dr]
 & \text{第\ref{b2:ch:18}章：奇异积分}\arrow[d]\\
 \text{Fourier 与函数空间}
   \arrow[u] \arrow[r] \arrow[dr]
 & \text{第\ref{b2:ch:19}章：位势与平方函数}\\
 & \text{\chapref{b2:ch:20}：测不准原理}
\end{tikzcd}
\]
从\chapref{b2:ch:17}到\chapref{b2:ch:19}的直接箭头表示 Hedberg 位势估计使用极大算子；从\chapref{b2:ch:18}到\chapref{b2:ch:19}则表示随机环带算子使用 Calderón–Zygmund 理论。极大平均控制位势的局部质量，核抵消控制不同频带叠加后的范数。

\ProseParagraphBreak
\chapref{b2:ch:20}转向集中问题：空间中的脉冲越窄，表示它所需的频率怎样变化？Fourier 反演与 Plancherel 将把这一问题化为位置和频率矩的比较。\chapref{b2:ch:17}的插值、\chapref{b2:ch:18}的零平均坏块与\chapref{b2:ch:19}的环带分解则构成第六编乘子理论的共同基础。

\ProseParagraphBreak
阅读时应把空间、指数与极限方式一起记录。弱型界控制层集而非强范数；极大截断控制全参数主值而非单个截断；齐次平方函数忽略零频，非齐次平方函数则保留低频块。这些区别直接影响后面的应用。
Elias M. Stein 的\cite{Stein1970Singular,Stein1993Harmonic}适合比较实变量抵消与振荡方法，Loukas Grafakos 的\cite{Grafakos2014Classical,Grafakos2014Modern}适合按极大、奇异积分和频率分解主题继续阅读。

% !TeX root = ../../Book2.tex
\chapter{极大函数、弱型估计与插值}
\label{b2:ch:17}
\BookChapterNumber{b2:ch:17}
\begin{chapterabstract}
估计一个算子时，我们往往只能先掌握两个指数：一端来自绝对积分，另一端来自平方范数或本性上确界。中间的指数怎样得到？端点的积分发散以后，还能保留什么？本章从局部平均的极大函数出发，用分布函数记录大值出现的范围，再分别通过大小截断与复参数变化建立两种插值方法。弱空间、Lorentz 空间和 \(K\) 泛函说明了这些估计背后的函数空间结构。
\end{chapterabstract}

\begin{learninggoals}\begin{itemize}
\item\label{b2:item:17-goal-1}比较中心、非中心与二进平均，并完成弱型估计和微分论证。
\item\label{b2:item:17-goal-2}由分布函数计算强、弱与 Lorentz 空间的可积条件。
\item\label{b2:item:17-goal-3}选择大小截断并完成 Marcinkiewicz 的高度积分。
\item\label{b2:item:17-goal-4}用三线估计证明 Riesz--Thorin，并解释 Fourier 变换的延拓。
\end{itemize}\end{learninggoals}

单位球的示性函数有有限积分，但其极大平均在远处只按 \(|x|^{-d}\) 衰减。这个简单例子说明：一个自然算子可能在端点失去可积性，同时仍有精确的超水平集合控制。后面的奇异积分也呈现同样现象，只是核的正平均将换成正负抵消。

\ProseParagraphBreak

本章需要 Lebesgue 积分、Hölder 不等式以及 \(L^p\) 收敛的基本知识。遇到无限求和或换序时，我们会指出非负性、可积控制或有限测度截断在哪里使用。极大函数的微分论证还要用连续紧支集函数的 \(L^1\) 稠密性；Fourier 应用则使用\secref{b2:sec:0401}已经建立的 Plancherel 等距性。复插值所需的复分析知识限于全纯函数的局部性质与最大模原理，\secref{b2:sec:1704}先说明其使用方式，再证明条带上的三线引理。

\ProseParagraphBreak

阅读时可先沿“极大函数的弱端点---分布函数的高度积分---中间强型”这一条路线走完\secref{b2:sec:1701}至\secref{b2:sec:1703}，看清为什么保留输入的大值部分能够消除发散。随后学习\secref{b2:sec:1704}，比较线性结构带来的另一种插值。Lorentz 空间的第二指数和最后的 \(K\) 泛函，可以在掌握这两条证明之后再回读；极大函数及同指数插值的证明并不依赖这两个补充视角。\secref{b2:sec:1701}中的二进立方体分解将在\secref{b2:sec:1802}直接用于 Calderón--Zygmund 分解。

\ProseParagraphBreak

继续阅读可选 Elias M. Stein 的《Singular Integrals and Differentiability Properties of Functions》\cite{Stein1970Singular}，把极大函数与后续奇异积分联系起来；Loukas Grafakos 的《Classical Fourier Analysis》\cite{Grafakos2014Classical}则可用于进一步学习弱型、重排和插值。本章所需的具体估计与证明在正文展开。

% !TeX root = ../../Book2.tex
\section{Hardy–Littlewood 极大算子回顾}
\label{b2:sec:1701}

固定尺度的平均会平滑函数，却可能漏掉比这个尺度窄得多的尖峰。若要同时讨论许多尺度的逼近，就希望用同一个函数控制所有这些平均。极大算子的做法是：先在每一点计算各个尺度上的绝对平均，再取其中的上确界。这个上确界允许观察尺度随位置改变，因而不能简单视为一个固定核的卷积。本节固定欧氏空间的 Lebesgue 测度，\(d\geq1\) 为整数。

\subsection{第一卷结果回顾}
\label{b2:subsec:170101}

对实或复函数 \(f\in L^1_{\mathrm{loc}}(\mathbb R^d)\)，约定
\[
 Mf(x)\coloneqq\sup_{r>0}\dfrac1{|B(x,r)|}\int_{B(x,r)}|f(y)|\,\mathrm dy.
\]
\symidx{\(Mf\)：Hardy--Littlewood 中心极大函数}
这是本卷的\term[hardy-littlewood-jida-hanshu]{Hardy--Littlewood 极大函数}[Hardy--Littlewood maximal function]，相应映射 \(f\mapsto Mf\) 称为极大算子。先取绝对值，是为了不让输入自身的正负抵消遮住误差大小；再除以球体积，是为了使常函数的平均保持原高度。它一般不是线性算子，但满足
\[M(f+g)\leq Mf+Mg,\quad M(cf)=|c|\cdot Mf\quad(c\in\mathbb C).\]
这两条性质分别来自积分三角不等式和绝对齐次性，后面的实插值会利用这种结构。

\ProseParagraphBreak

还须确认无穷多个平均的上确界是可测的。记 \(A_rf(x)=|B(x,r)|^{-1}\int_{B(x,r)}|f|\)。固定 \(r>0\)，当中心从 \(x\) 移到 \(x+h\) 时，两个积分的差不超过
\[\int_{B(x+h,r)\,\triangle\,B(x,r)}|f|.\]
这些球位于一个固定有界区域内，对称差的体积随 \(h\to0\) 趋零；局部可积函数的积分对小测度集合绝对连续，故 \(A_rf\) 对中心连续。固定中心时，令半径趋于某个正数，两个球之间的薄环体积也趋零，且分母体积连续，因此 \(r\mapsto A_rf(x)\) 也连续。于是
\[Mf(x)=\sup_{r\in\mathbb Q,\ r>0}A_rf(x)\]
是可数个可测函数的上确界。另一方面，\(\{Mf>t\}\) 是各固定半径的开超水平集合之并，因而还是开集。局部可积性通过积分的绝对连续性，转化为了平均值对球心与半径的连续性。

\ProseParagraphBreak

一个能直接算完的例子是实直线上的
\[f=\mathbf1_{(0,1)}+\mathbf1_{(2,3)}.\]
在空隙中点 \(x=3/2\)，半径 \(r\leq1/2\) 的平均为零；若 \(1/2<r<3/2\)，球内与两段支集相交的总长度为 \(2r-1\)，故平均为 \(1-1/(2r)\)，随 \(r\) 增大；若 \(r\geq3/2\)，总质量固定为 \(2\)，平均变为 \(1/r\)，随 \(r\) 减小。因此 \(Mf(3/2)=2/3\)，尽管该点附近 \(f=0\)。极大平均会从较大邻域看见质量，不能把它与函数本身在该点的大小混同。

\ProseParagraphBreak

若某个平均超过高度 \(t\)，它所在的球必须含有至少 \(t\) 乘体积的质量。困难在于这些球会重叠，不能直接把它们的积分相加。下面先选出不交的一部分，再用固定倍数的扩大球补回覆盖。

\Needspace{6\baselineskip}
\begin{theorem}[Hardy--Littlewood]\label{b2:thm:17-maximal-weak}
设整数 \(d\geq1\)，\(f\in L^1(\mathbb R^d;\mathbb C)\)，\(t>0\)，并定义 \(Mf(x)=\sup_{r>0}|B(x,r)|^{-1}\int_{B(x,r)}|f|\)，其中球为开球、集合大小取 Lebesgue 测度。则
\[
 |\{Mf>t\}|\leq \dfrac{5^d}{t}\|f\|_1.
\]
\end{theorem}
\begin{proof}
记 \(E=\{Mf>t\}\)，并先固定紧集 \(K\subset E\)。对每个 \(x\in K\)，由上确界严格大于 \(t\)，能选一个以 \(x\) 为中心、平均大于 \(t\) 的球。开球覆盖紧集，所以可取有限子覆盖。

在这有限族中选半径最大的球，将所有与它相交的球删除；在余下的球中重复。得到的不交球记为 \(B_j=B(c_j,r_j)\)。若 \(B(a,s)\) 在选择 \(B_j\) 时被删，则 \(s\leq r_j\)、\(|a-c_j|<s+r_j\)。对其中任意点 \(y\)，有 \(|y-c_j|<s+(s+r_j)\leq3r_j\)。故被删球包含于 \(3B_j\)，所选球的三倍扩大已覆盖原来的有限族，也就覆盖 \(K\)。特别地，可以用 \(5B_j\) 作下列估计：
\[
 |K|\leq5^d\sum_j|B_j|\leq\dfrac{5^d}{t}\sum_j\int_{B_j}|f|\leq\dfrac{5^d}{t}\|f\|_1.
\]
第二步使用每个所选球的平均大于 \(t\)，最后一步才使用球互不相交。对全部紧 \(K\subset E\) 取上确界，Lebesgue 内正则性给出 \(|E|\) 的同一上界，得到结论。即使最初不知道 \(|E|\) 有限，这种内逼近仍然适用。证明只进行了有限族的选择，不需要预设能从无穷族中选出最大半径球。
\end{proof}
注意估计控制的是超水平集合的测度，没有直接断言 \(Mf\) 可积。例如 \(f=\mathbf1_{B(0,1)}\)。在 \(|x|\geq2\) 时，球 \(B(x,|x|+1)\) 包含单位球，故 \(Mf(x)\geq c_d|x|^{-d}\)。反过来，任何碰到单位球的中心球半径至少为 \(|x|-1\)，其所含质量又不超过 \(|B(0,1)|\)，所以 \(Mf(x)\leq C_d|x|^{-d}\)。尾部径向积分于是含有发散的 \(\int_2^\infty r^{-1}\,\mathrm dr\)。这个例子一面满足上述集合估计，一面否定一般的强 \(L^1\) 界，正是下一节引入弱空间的原因。

\subsection{中心/非中心极大函数}
\label{b2:subsec:170102}

对 \(f\in L^1_{\mathrm{loc}}(\mathbb R^d;\mathbb C)\)，若上确界改取所有包含 \(x\) 的开球，所得函数记为 \(M_u f(x)\)，称为\term[feizhongxin-jida-hanshu]{非中心极大函数}[uncentered maximal function]。这里“包含”不要求 \(x\) 是球心。给定 \(x\in B(a,r)\)，有 \(B(a,r)\subset B(x,2r)\)，故
\[\frac1{|B(a,r)|}\int_{B(a,r)}|f|
 \leq2^d\frac1{|B(x,2r)|}\int_{B(x,2r)}|f|.\]
对这类球取上确界，得到
\[
 Mf\leq M_u f\leq2^d Mf.
\]
两种平均的数值并不相同，但任何一者的范数或分布估计都可借此传给另一者，只改变维数常数。非中心函数也可测：\(M_uf(x)>t\) 当且仅当 \(x\) 位于某个平均大于 \(t\) 的开球内，所以该超水平集合是开球的并。对实直线上的区间输入，非中心区间可以把观察点放在端点附近，中心区间却要同时向两侧扩张；两者的差别来自允许怎样选择覆盖球。

\ProseParagraphBreak

对 \(R>0\)、\(f\in L^1_{\mathrm{loc}}(\mathbb R^d)\)，局部版本 \(M_Rf(x)\) 定义为 \(\sup_{0<r<R}|B(x,r)|^{-1}\int_{B(x,r)}|f|\)。若只考察紧集 \(K\)，选一个包含 \(K+B(0,R)\) 的有界球 \(U\)，令 \(F=f\cdot\mathbf1_U\)。则 \(F\in L^1\)，且 \(M_Rf(x)=M_RF(x)\leq MF(x)\) 对 \(x\in K\) 成立。这样才把全空间弱估计合法地用于局部函数。

\ProseParagraphBreak

局部可积本身不保证全尺度极大函数有限。例如在实直线上取 \(f(y)=|y|\)，则以原点为中心、半径 \(r\) 的平均为 \(r/2\)，故 \(Mf(0)=\infty\)。问题来自任意大的尺度；截断半径后，每个固定紧集只使用一个有界邻域内的积分。后面的局部微分论证也会按这个方式先截断位置。

\subsection{二进极大函数}
\label{b2:subsec:170103}

对 \(f\in L^1_{\mathrm{loc}}(\mathbb R^d;\mathbb C)\)，令 \(\mathcal D\) 为半开二进立方体 \(2^{-j}(k+[0,1)^d)\) 的族，其中 \(j\in\mathbb Z\)、\(k\in\mathbb Z^d\)。定义\term[erjin-jida-hanshu]{二进极大函数}[dyadic maximal function]
\[
 M_{\mathcal D}f(x)\coloneqq\sup_{Q\in\mathcal D,\,x\in Q}\dfrac1{|Q|}\int_Q|f|.
\]
\symidx{\(M_{\mathcal D}f\)：二进极大函数}
半开约定使每个尺度恰有一个立方体包含 \(x\)，同一格中的两个立方体不交或一者包含另一者。网格可数，故 \(M_{\mathcal D}f\) 可测。与球族相比，这里的尺度选择少了，却多了唯一的父立方体：每个 \(Q\) 都包含于一个边长加倍、体积为 \(2^d|Q|\) 的二进父块。下一章需要的恰是这个父子结构。

\Needspace{6\baselineskip}
\begin{proposition}[二进超水平集合分解]\label{b2:thm:17-dyadic-level}
设整数 \(d\geq1\)，\(f\in L^1(\mathbb R^d;\mathbb C)\)、\(t>0\)。固定标准半开二进网格 \(\mathcal D=\{\,2^{-j}(k+[0,1)^d):j\in\mathbb Z,\ k\in\mathbb Z^d\,\}\)，定义
\[
 M_{\mathcal D}f(x)=\sup_{Q\in\mathcal D,\,x\in Q}\frac1{|Q|}\int_Q|f|.
\]
则 \(\{M_{\mathcal D}f>t\}\) 是一族互不相交的极大二进立方体 \(Q_j\) 的并，且
\[
 t<\dfrac1{|Q_j|}\int_{Q_j}|f|\leq2^dt,
 \quad \sum_j|Q_j|\leq t^{-1}\|f\|_1.
\]
\end{proposition}
\begin{proof}
对平均大于 \(t\) 的任意二进块 \(Q\)，其第 \(n\) 个祖先 \(Q^{(n)}\) 的平均满足
\[\frac1{|Q^{(n)}|}\int_{Q^{(n)}}|f|
 \leq\frac{\|f\|_1}{2^{nd}|Q|}\longrightarrow0.\]
所以只有有限多个祖先仍可能平均大于 \(t\)；取这些超水平祖先中最大的一个。注意不能仅在第一次遇到平均不超过 \(t\) 的父块时停下，因为更大的祖先可能又包含别处的质量。所有足够大的祖先都不超水平，才保证极大块真正存在。

取全部极大超水平块。二进嵌套性说明两个不同的极大块必不交；网格可数，故它们组成有限或可数族 \(Q_j\)。\(M_{\mathcal D}f(x)>t\) 的每一点都在某个超水平块中，再包含于一个极大块，反向包含也由定义成立。每个 \(Q_j\) 的父块不超水平，因此
\[\int_{Q_j}|f|\leq\int_{\operatorname{parent}(Q_j)}|f|
 \leq t\,2^d|Q_j|.\]
这给出平均上界；下界来自超水平条件。最后由不交性，
\[t\sum_j|Q_j|\leq\sum_j\int_{Q_j}|f|\leq\|f\|_1,\]
得到测度估计。
\end{proof}
这项分解还记录坏区域的平均上界，是下一章分解函数的依据。仅有球覆盖弱型估计并不能直接提供这样的父子关系。

\ProseParagraphBreak

单个固定网格并不等同于全部球平均。对 \(x\in Q\)，以 \(x\) 为中心、半径 \(\sqrt d\,\ell(Q)\) 的球包含 \(Q\)，所以 \(M_{\mathcal D}f\leq C_dMf\)。反向的逐点比较对一个网格一般失败：在实直线上取 \(f=\mathbf1_{(0,1)}\)、\(x=-1\)，包含 \(-1\) 的标准二进区间都留在负半轴，因而 \(M_{\mathcal D}f(-1)=0\)；中心球却可跨过原点取得正平均。二进算子的弱型估计应由它自身的分解证明，而不能把两个算子直接当作数值等价的函数。

\subsection{极大函数微分法}
\label{b2:subsec:170104}

平均能恢复函数，并非因为任意可测代表在每一点都连续，而是坏点的测度可由逼近误差控制。
\Needspace{6\baselineskip}
\begin{proposition}[Lebesgue 微分的极大函数证明]\label{b2:thm:17-differentiation}
设整数 \(d\geq1\)，\(f\in L^1_{\mathrm{loc}}(\mathbb R^d;\mathbb C)\)，并选取一个有限值可测代表。则对几乎处处 \(x\in\mathbb R^d\)，有
\[
 \lim_{r\downarrow0}\dfrac1{|B(x,r)|}\int_{B(x,r)}|f(y)-f(x)|\,\mathrm dy=0.
\]
另令 \(Q_j(x)\) 为标准半开二进网格中唯一包含 \(x\)、边长 \(2^{-j}\) 的立方体，则还对几乎处处 \(x\) 有
\[\lim_{j\to\infty}\frac1{|Q_j(x)|}\int_{Q_j(x)}|f(y)-f(x)|\,\mathrm dy=0.\]
\end{proposition}
\begin{proof}
先设 \(f\in L^1\)。连续紧支集函数在 \(L^1(\mathbb R^d)\) 中稠密（第一卷命题~10.2.4），故对任意 \(\varepsilon>0\)，可取 \(g\in C_c(\mathbb R^d)\)，使 \(\|f-g\|_1<\varepsilon\)。令 \(Df(x)\) 为球平均绝对偏差的上极限。它是可测函数：固定半径的参数积分可测，半径的连续性又允许在有理半径上取上极限。由
\[
 |f(y)-f(x)|\leq |f(y)-g(y)|
       +|g(y)-g(x)|+|g(x)-f(x)|
\]
及 \(g\) 的连续性，得到
\[Df(x)\leq M(f-g)(x)+|f(x)-g(x)|.\]
因此对 \(t>0\)，由极大弱型和 Chebyshev 不等式，
\[
 |\{Df>2t\}|\leq(5^d+1)\varepsilon/t.
\]
左侧集合由 \(f,t\) 固定，不随所选 \(g\) 改变。令 \(\varepsilon\downarrow0\)，可知对每个固定 \(t>0\) 该集合为零测集，再只对 \(t=1/n\) 取可数并，得到 \(Df=0\) 几乎处处。

二进版本令 \(D_{\mathcal D}f\) 为所列平均偏差的上极限。因为 \(\operatorname{diam}Q_j(x)=\sqrt d\,2^{-j}\to0\)，连续函数的平均偏差仍趋零，所以
\[D_{\mathcal D}f\leq M_{\mathcal D}(f-g)+|f-g|.\]
二进超水平分解给出常数为一的弱估计，其余误差集合论证相同。

最后设 \(f\) 仅局部可积。对每个整数 \(N\geq1\)，置 \(f_N=f\cdot\mathbf1_{B(0,N+1)}\in L^1\)。在 \(B(0,N)\) 内，各点的足够小球及足够小二进块都落在 \(B(0,N+1)\) 内，其平均偏差与 \(f_N\) 的相同。于是可在 \(B(0,N)\) 上应用已证结论，再合并整数 \(N\) 对应的可数个零测例外集。
\end{proof}
因此在二进分解的好区域上可以把“所有小立方体平均不超过 \(t\)”转为“\(|f|\leq t\) 几乎处处”。这一步将在\secref{b2:sec:1802}使用。

% !TeX root = ../../Book2.tex
\section{弱 \texorpdfstring{\(L^p\)}{Lᵖ} 空间}
\label{b2:sec:1702}

上一节的极大函数可能不可积，却仍满足 \(|\{Mf>t\}|\leq C\|f\|_1/t\)。这不是没有得到任何控制，而是控制的对象变成了“超过某个高度的区域有多大”。本节先把这种信息写成分布函数，再说明它怎样决定强、弱 \(L^p\) 以及 Lorentz 空间。分布和重排的基本恒等式适用于一般测度空间；局部可积性、卷积和简单函数逼近的具体应用会明确回到 Lebesgue 空间。

\subsection{分布函数}
\label{b2:subsec:170201}

设 \((X,\mathcal A,\mu)\) 为测度空间，\(f:X\to\mathbb C\) 可测，其\term[fenbu-hanshu]{分布函数}[distribution function]定义为
\[d_f(t)\coloneqq\mu\{\,x\in X:|f(x)|>t\,\},\quad t\geq0.\]
\symidx{\(d_f(t)\)：函数的分布函数}
它允许取无穷，且只保留超水平集合的测度，不保留集合的位置。它单调不增；对 \(t_n\downarrow t\)，集合 \(\{|f|>t_n\}\) 递增到 \(\{|f|>t\}\)，所以 \(d_f(t_n)\to d_f(t)\)，即右连续。这里使用从下连续性，无须先假设某个超水平集合有限。

\ProseParagraphBreak

把 \(|f(x)|\) 想成竖直线段的长度，就有
\[|f(x)|=\int_0^\infty\mathbf1_{\{|f(x)|>t\}}\,\mathrm dt.\]
先累加各点的竖直长度，与先累加每一层的水平测度，得到同一个总量。对幂次作相同分层，产生下面的公式。

\Needspace{6\baselineskip}
\begin{proposition}[层蛋糕公式]\label{b2:thm:17-layer-cake}
设 \((X,\mathcal A,\mu)\) 为测度空间，\(h:X\to[0,\infty]\) 可测，\(d_h(t)=\mu\{h>t\}\)，\(0<p<\infty\)。则有允许两端同时为无穷的恒等式
\[\int_Xh^p\,\mathrm d\mu=p\int_0^\infty t^{p-1}d_h(t)\,\mathrm dt.\]
特别地，对实或复可测函数 \(f\) 取 \(h=|f|\)，便得到 \(\int_X|f|^p\,\mathrm d\mu=p\int_0^\infty t^{p-1}d_f(t)\,\mathrm dt\)。
\end{proposition}
\begin{proof}
先设 \(h\) 为非负简单函数，将其正值写在有限个不交可测集合 \(E_j\) 上，记这些值为 \(a_j>0\)。两侧都等于 \(\sum_j a_j^p\mu(E_j)\)：右侧只需对有限个阶梯函数分别积分。若某个正值集合测度无穷，则两侧都为无穷，等式仍成立。

一般情形取非负简单函数 \(h_n\uparrow h\)。对每个 \(t\geq0\)，集合 \(\{h_n>t\}\) 递增到 \(\{h>t\}\)，故 \(d_{h_n}(t)\uparrow d_h(t)\)。分别在 \(X\) 和高度轴上用单调收敛，把简单函数等式传到极限即可。这一证明不要求 \(X\) 为 \(\sigma\)-有限；若改用积空间 Tonelli，必须另外检查所采用版本的假设。
\end{proof}
例如 \(f=a\cdot\mathbf1_E\)、\(a>0\)、\(\mu(E)=m<\infty\) 时，\(d_f(t)=m\) 于 \(0\leq t<a\)，其余为零。公式给出 \(p\int_0^a t^{p-1}m\,\mathrm dt=a^pm\)，与直接积分一致。一般函数可以在高水平有尖峰，也可以在低水平有长尾；在无穷测度空间上，这两端都必须检查。

\subsection{弱型 \texorpdfstring{\(L^p\)}{Lp}}
\label{b2:subsec:170202}

Chebyshev 不等式说明，强 \(L^p\) 函数满足 \(d_f(t)\leq t^{-p}\|f\|_p^p\)。反向若只知道这个逐高度上界，层蛋糕中的被积函数却被估成 \(Ct^{-1}\)，一般无法积分。因而应把“每个高度的集合都不太大”与“所有高度的总积分有限”区分开来。

\ProseParagraphBreak

设 \((X,\mathcal A,\mu)\) 为测度空间，\(f:X\to\mathbb C\) 可测，\(0<p<\infty\)。若
\[\|f\|_{p,\infty}\coloneqq\sup_{t>0}t\cdot d_f(t)^{1/p}<\infty,\]
则称 \(f\) 属于\term[ruo-lp-kongjian]{弱 \(L^p\) 空间}[weak Lp space]，记为 \(L^{p,\infty}(X,\mu)\)。该量等于使 \(d_f(t)\leq(A/t)^p\) 对全部 \(t>0\) 成立的最小 \(A\)。它是拟范数：若 \(A=\|f\|_{p,\infty}\)、\(B=\|g\|_{p,\infty}\)，则超水平集合的包含关系给出
\[\|f+g\|_{p,\infty}\leq2(A^p+B^p)^{1/p}\leq C_p(A+B).\]
常数一的三角不等式确实可能失败。在两个互不相交、测度均为一的集合 \(E,F\) 上，取 \(f=2\cdot\mathbf1_E+\mathbf1_F\)、\(g=\mathbf1_E+2\cdot\mathbf1_F\)。二者的弱 \(L^1\) 拟范数都是 \(2\)，而和等于 \(3\cdot\mathbf1_{E\cup F}\)，弱拟范数为 \(6\)。

\Needspace{6\baselineskip}
\begin{proposition}[弱范数与有限集合积分]\label{b2:thm:17-weak-inclusions}
设 \((X,\mathcal A,\mu)\) 为测度空间，\(f:X\to\mathbb C\) 可测，\(0<p<\infty\)，并令 \(\|f\|_{p,\infty}=\sup_{t>0}t\cdot\mu\{|f|>t\}^{1/p}\)。则 \(\|f\|_{p,\infty}\leq\|f\|_{L^p(X,\mu)}\)。若 \(A=\|f\|_{p,\infty}<\infty\)、\(0<r<p\)，且 \(E\) 可测、\(0<\mu(E)<\infty\)，则
\[\|f\|_{L^r(E)}\leq
 \left(\frac p{p-r}\right)^{1/r}
 A\,\mu(E)^{1/r-1/p}.\]
特别地，\(p>1\) 时有
\[\int_E|f|\,\mathrm d\mu\leq p'\|f\|_{p,\infty}\cdot\mu(E)^{1/p'},
 \quad p'=\frac p{p-1}.\]
反之，设 \(1<p<\infty\)、\((X,\mathcal A,\mu)\) 半有限；若存在 \(C<\infty\)，使每个有限测度可测集 \(E\) 都满足 \(\int_E|f|\,\mathrm d\mu\leq C\mu(E)^{1/p'}\)，则 \(f\in L^{p,\infty}\) 且 \(\|f\|_{p,\infty}\leq C\)。
\end{proposition}
\begin{proof}
第一项由 Chebyshev 得到。记 \(m=\mu(E)>0\)。若 \(A=0\)，则 \(f=0\) 几乎处处，结论直接成立。否则限制在 \(E\) 上的分布函数不超过 \(\min\{m,A^pt^{-p}\}\)。这两项在 \(a=Am^{-1/p}\) 处相等，故应在这个高度分开积分。层蛋糕给出
\[
 \begin{aligned}
 \int_E|f|^r\,\mathrm d\mu
 &\leq r\int_0^a t^{r-1}m\,\mathrm dt
       +rA^p\int_a^\infty t^{r-p-1}\,\mathrm dt\\
 &=m a^r+\frac{rA^p}{p-r}a^{r-p}
  =\frac p{p-r}A^rm^{1-r/p}.
 \end{aligned}
\]
这里 \(r<p\) 正好保证第二项在无穷端收敛。取 \(r\) 次方根即得；\(r=1<p\) 给出所列集合积分界。零测度集合的积分另为零。

证明反向，固定 \(t>0\)，取 \(F\subset\{|f|>t\}\) 且 \(0<\mu(F)<\infty\)。假设给出 \(t\mu(F)\leq C\mu(F)^{1/p'}\)，从而 \(\mu(F)\leq(C/t)^p\)。半有限性保证任意可测集合的测度等于其中有限测度子集测度的上确界，故 \(d_f(t)\leq(C/t)^p\)。这也解释了为什么反向需要该测度条件；Lebesgue 测度满足它。
\end{proof}
在 Lebesgue 空间上，\(p>1\) 的弱 \(L^p\) 函数由此局部可积。\(p=1\) 时不能照搬：例如 \(|x|^{-d}\cdot\mathbf1_{\{0<|x|<1\}}\) 属于弱 \(L^1\)，却在原点附近不可积。另一方面，降低到任意 \(0<r<p\) 的结论仍有效；严格不等号来自上面看得见的高度积分。

\ProseParagraphBreak

设 \((X,\mu),(Y,\nu)\) 为测度空间，\(1\leq p,q\leq\infty\)，\(T\) 把 \(L^p(X,\mu)\) 输入映为 \(Y\) 上的可测函数。若 \(\|Tf\|_{q,\infty}\leq A\|f\|_p\)，称它具有\term[ruoxing-guji]{弱型估计}[weak type estimate] \((p,q)\)，强型则把目标换成 \(L^q\)。当 \(q=\infty\) 时直接使用本性上确界。例如上一节的中心极大算子具有弱型 \((1,1)\)，但单位球输入的远场计算否定了强型 \((1,1)\)。

\ProseParagraphBreak

对整数 \(d\geq1\)、\(1<p<\infty\)，在 \(\mathbb R^d\) 上定义
\[N_p(f)\coloneqq\sup_{0<|E|<\infty}|E|^{-1/p'}\int_E|f|,\]
其中 \(E\) 可测。上一命题的两个方向给出
\[\|f\|_{p,\infty}\leq N_p(f)\leq p'\|f\|_{p,\infty}.\]
\(N_p\) 的三角不等式直接来自集合积分的三角不等式，零范数也只对应零等价类。因此它提供一个等价的真正范数。上确界可以只取有界集合，因为 \(E\cap B(0,n)\uparrow E\)，积分和测度都按单调收敛逼近。这种集合积分写法还适合卷积：它不要求原输入在全空间可积。

\Needspace{6\baselineskip}
\begin{proposition}[弱空间的完备性与卷积]\label{b2:prop:17-weak-complete}
设整数 \(d\geq1\)、\(1<p<\infty\)，\(p'=p/(p-1)\)。实或复函数空间 \(L^{p,\infty}(\mathbb R^d)\) 关于范数
\[N_p(f)=\sup_{\substack{E\subset\mathbb R^d\ \text{可测}\\0<|E|<\infty}}
 |E|^{-1/p'}\int_E|f|\]
完备。若 \(k\in L^1(\mathbb R^d)\)、\(f\in L^{p,\infty}(\mathbb R^d)\) 均为实或复函数，则 \(k*f\) 由几乎处处绝对收敛的积分定义，且
\[N_p(k*f)\leq\|k\|_1\cdot N_p(f).\]
\end{proposition}
\begin{proof}
设 \(f_n\) 在 \(N_p\) 中 Cauchy。对每个有限球 \(B\)，有
\[\|f_n-f_m\|_{L^1(B)}\leq |B|^{1/p'}\cdot N_p(f_n-f_m).\]
\(L^1(B)\) 的完备性给出局部极限；不同球上的极限因几乎处处唯一而相容，故拼成局部可积函数 \(f\)。沿整数半径球作对角抽取，可得到一个在全空间几乎处处趋于 \(f\) 的子列。

给定 \(\varepsilon>0\)，先取 \(N\)，使 \(N_p(f_n-f_m)\leq\varepsilon\) 对 \(n,m\geq N\) 成立。固定 \(n\geq N\)，对任意有限测度集合 \(E\) 沿上述子列使用 Fatou，便得
\[\int_E|f_n-f|\leq\varepsilon|E|^{1/p'}.\]
此界对全部 \(E\) 使用同一个 \(N\)，所以 \(N_p(f_n-f)\leq\varepsilon\)。再由 \(N_p\) 的三角不等式知 \(N_p(f)<\infty\)，完备性成立。

对卷积，欧氏空间上的 Lebesgue 积测度空间是 \(\sigma\)-有限的。Tonelli 和平移不变性给出
\[
 \begin{aligned}
 \int_E\int |k(y)|\cdot|f(x-y)|\,\mathrm dy\,\mathrm dx
 &=\int |k(y)|\int_{E-y}|f(z)|\,\mathrm dz\,\mathrm dy\\
 &\leq |E|^{1/p'}\|k\|_1\cdot N_p(f).
 \end{aligned}
\]
取整数半径球 \(E\)，有限右端说明内积分在全空间几乎处处有限，故定义了绝对收敛的卷积。它控制 \(|k*f|\)，再对全部有限测度 \(E\) 取上确界即得所需范数界。
\end{proof}
完备性不意味着有界或光滑紧支集函数在弱拟范数中稠密。设 \(d\geq1\)、\(0<p<\infty\)，令 \(f(x)=|x|^{-d/p}\cdot\mathbf1_{\{0<|x|<1\}}\)，并记 \(\omega_d=|B(0,1)|\)。当 \(t\geq1\) 时，\(d_f(t)=\omega_dt^{-p}\)，故 \(f\in L^{p,\infty}\)。对任意有界 \(g\)，在 \(t>\max\{1,\|g\|_\infty\}\) 时，\(\{|f|>2t\}\subset\{|f-g|>t\}\)，所以
\[\|f-g\|_{p,\infty}\geq\frac12\omega_d^{1/p}.\]
无论有界截断的高度多大，这个误差都不趋零。处理弱空间上的算子延拓时，必须使用已证明的积分定义或另选收敛方式，不能直接套用强 \(L^p\) 的稠密性。

\subsection{Lorentz 空间初步}
\label{b2:subsec:170203}

弱拟范数只取 \(t\cdot d_f(t)^{1/p}\) 的最高值，没有累加不同高度的贡献。为了区分同样临界的幂次尾部中是否还带有可积的对数改善，可以保留整个分布，而不只取一个上确界。先把函数值按大小排列：设 \((X,\mathcal A,\mu)\) 为测度空间、\(f:X\to\mathbb C\) 可测，定义其\term[dijian-chongpai]{递减重排}[decreasing rearrangement]
\[f^*(s)\coloneqq\inf\{\,t\geq0:d_f(t)\leq s\,\},\quad s>0,\]
空集的下确界取无穷。变量 \(s\) 表示测度大小，不是原空间的位置；这里的 \(f^*\) 也不表示微分理论中的点态代表。它是分布函数的广义逆：寻找使超水平集合缩到不超过 \(s\) 的最低高度。

\Needspace{6\baselineskip}
\begin{proposition}[重排与等分布]\label{b2:prop:17-rearrangement}
设 \((X,\mathcal A,\mu)\) 为测度空间，\(f:X\to\mathbb C\) 可测，\(d_f(t)=\mu\{|f|>t\}\)、\(t\geq0\)，并令 \(f^*(s)=\inf\{t\geq0:d_f(t)\leq s\}\)、\(s>0\)。则 \(f^*\) 单调不增，且对每个 \(t\geq0\)，
\[|\{\,s>0:f^*(s)>t\,\}|=d_f(t),\]
左侧取半直线的 Lebesgue 测度。因而对 \(0<p<\infty\)，
\[\int_X|f|^p\,\mathrm d\mu=\int_0^\infty f^*(s)^p\,\mathrm ds,\quad
 \|f\|_{p,\infty}=\sup_{s>0}s^{1/p}f^*(s),\]
各式允许无穷值。
\end{proposition}
\begin{proof}
增大 \(s\) 会扩大取下确界的集合，故 \(f^*\) 单调不增。若 \(d_f(t)\leq s\)，则 \(t\) 已在该集合中，\(f^*(s)\leq t\)。反之，若 \(d_f(t)>s\)，分布函数的右连续性保证存在 \(\delta>0\)，使 \(d_f(t+\delta)>s\)，于是 \(f^*(s)\geq t+\delta>t\)。因此
\[f^*(s)>t\quad\Longleftrightarrow\quad s<d_f(t).\]
这给出等分布恒等式；分别对 \(|f|\) 和 \(f^*\) 使用层蛋糕即得积分等式。

对弱拟范数，固定 \(s\)，让 \(0<t<f^*(s)\) 趋向 \(f^*(s)\)，上述等价关系给出 \(s^{1/p}f^*(s)\leq\|f\|_{p,\infty}\)。反向固定 \(t>0\)，让 \(0<s<d_f(t)\) 趋向 \(d_f(t)\)，得到反向上界。无穷值时按任意大的有限 \(s,t\) 解释同一论证。
\end{proof}
即使原空间有原子，这个等分布结论也成立；重排是在另一个连续测度坐标上表示原分布，并未要求把原子拆开。它也说明，两函数若分布相同，其强、弱范数相同，但这并不表示两函数在原空间上接近。

\ProseParagraphBreak

例如在 \(\mathbb R^d\) 取两个不交可测集 \(E,F\)，\(|E|=|F|=1\)，令 \(f=2\cdot\mathbf1_E+\mathbf1_F\)。\cref{b2:tab:17-distribution-rearrangement}将先前的两层模型按高度和测度两种坐标写出。原来的形状可以很不规则，重排却只留下两段水平阶梯。
\begin{table}[htbp]
\centering
\caption[两层函数的分布与重排]{两层函数的分布与重排。\(E,F\subset\mathbb R^d\) 不交且测度均为一，\(f=2\cdot\mathbf1_E+\mathbf1_F\)；\(t\) 是高度，\(s\) 是重排后的测度坐标。}
\label{b2:tab:17-distribution-rearrangement}
\begin{tabularx}{\textwidth}{@{}*{4}{>{\centering\arraybackslash}X}@{}}
\toprule
高度区间 & \(d_f(t)\) & 测度坐标区间 & \(f^*(s)\)\\
\midrule
\(0\leq t<1\) & \(2\) & \(0<s<1\) & \(2\)\\
\(1\leq t<2\) & \(1\) & \(1\leq s<2\) & \(1\)\\
\(t\geq2\) & \(0\) & \(s\geq2\) & \(0\)\\
\bottomrule
\end{tabularx}
\end{table}

对 \(0<p<\infty\)、\(0<q\leq\infty\)，\term[lorentz-kongjian]{Lorentz 空间}[Lorentz space]由下列拟范数定义：
\[
 \|f\|_{L^{p,q}}\coloneqq
 \left(\int_0^\infty[s^{1/p}f^*(s)]^q\frac{\mathrm ds}{s}\right)^{1/q}
 \quad(0<q<\infty),
\]
\(q=\infty\) 时取 \(\sup_{s>0}s^{1/p}f^*(s)\)。主指数 \(p\) 固定幂次尺度，第二指数 \(q\) 记录各测度尺度的累加方式。\(\mathrm ds/s\) 对伸缩 \(s\mapsto as\) 不变，因此每一倍增尺度受到同样对待。由重排命题，\(L^{p,p}=L^p\) 且范数相等，\(L^{p,\infty}\) 也与前面定义完全一致。
\symidx{\(f^*(s)\)：递减重排}
对表中的两层函数，可直接算得
\[
 \|f\|_{L^{p,q}}^q=\frac pq(2^q+2^{q/p}-1)\quad(q<\infty),\quad
 \|f\|_{p,\infty}=\max\{2,2^{1/p}\}.
\]
令 \(q=p\) 便恢复 \(\|f\|_p^p=2^p+1\)。这些等式展示的是同一个输入的不同大小信息，而非几种互不相干的空间定义。

\ProseParagraphBreak

固定 \(0<p<\infty\)，较小的第二指数确实给出较强的条件。对 \(0<q_0<\infty\)，单调性给出
\[
 \|f\|_{L^{p,q_0}}^{q_0}
 \geq f^*(s)^{q_0}\int_0^s r^{q_0/p-1}\,\mathrm dr
 =\frac p{q_0}[s^{1/p}f^*(s)]^{q_0}.
\]
故归一化重排 \(H(s)=s^{1/p}f^*(s)\) 的上确界受该拟范数控制。若 \(q_0<q_1<\infty\)，再用
\[\int_0^\infty H^{q_1}\frac{\mathrm ds}{s}
 \leq\|H\|_\infty^{q_1-q_0}\int_0^\infty H^{q_0}\frac{\mathrm ds}{s},\]
便得 \(L^{p,q_0}\subset L^{p,q_1}\)；\(q_1=\infty\) 已由上确界界涵盖。不能仅把它视为无穷测度空间上普通 \(L^q\) 的包含关系，关键额外信息是 \(f^*\) 的单调性。

\ProseParagraphBreak

临界幂次和对数使第二指数的作用更清楚。先设 \(a>0\)、\(0<p<\infty\)。径向积分给出
\[
 \begin{aligned}
 |x|^{-a}\cdot\mathbf1_{\{0<|x|<1\}}\in L^p
 &\quad\Longleftrightarrow\quad ap<d,\\
 |x|^{-a}\cdot\mathbf1_{\{|x|>1\}}\in L^p
 &\quad\Longleftrightarrow\quad ap>d.
 \end{aligned}
\]
在 \(ap=d\) 时，两者都只有临界弱界；相加所得纯幂次 \(|x|^{-a}\) 属于 \(L^{d/a,\infty}\)，却不属于任何全空间强 \(L^p\)。前一种失败来自高水平的原点奇性，后一种来自低水平的远尾。

\ProseParagraphBreak

再固定 \(0<p<\infty\)、\(\beta\in\mathbb R\)，在足够小的 \(0<s<s_0\) 上取
\[F(s)=s^{-1/p}(\log(e/s))^{-\beta},\]
并在 \(s\geq s_0\) 置零。选择 \(s_0\) 足够小使此函数在 \((0,s_0)\) 单调不增；例如 \(\log(e/s_0)>p\max\{\beta,0\}+1\) 即可。因此这个半直线上的模型本身就是自己的重排。对 \(q<\infty\)，代入定义并置 \(v=\log(e/s)\)，得到
\[\|F\|_{L^{p,q}}^q=\int_{\log(e/s_0)}^\infty v^{-\beta q}\,\mathrm dv.\]
所以 \(F\in L^{p,q}\) 当且仅当 \(\beta q>1\)，而 \(F\in L^{p,\infty}\) 当且仅当 \(\beta\geq0\)。特别地，强 \(L^p\) 要求 \(\beta p>1\)。纯临界幂次对应 \(\beta=0\)，只留下上确界控制；加入足够的对数衰减，才让尺度积分收敛。

\ProseParagraphBreak

这些更细的输入条件能够与弱端点输出配对，恢复原本可能发散的积分。
\Needspace{6\baselineskip}
\begin{proposition}[Lorentz 端点配对]\label{b2:prop:17-lorentz-pairing}
设 \((X,\mathcal A,\mu)\) 为测度空间，\(1<p<\infty\)、\(p'=p/(p-1)\)，实或复可测函数 \(f\in L^{p,1}(X,\mu)\)、\(g\in L^{p',\infty}(X,\mu)\)。Lorentz 拟范数采用递减重排及测度 \(\mathrm ds/s\) 的上述定义。则乘积绝对可积，且
\[\int_X|fg|\,\mathrm d\mu
 \leq\|f\|_{L^{p,1}}\cdot\|g\|_{p',\infty}.\]
\end{proposition}
\begin{proof}
对可测集 \(E\)，将层蛋糕用于 \(|g|\cdot\mathbf1_E\)，并用等分布恒等式，得
\[
 \int_E|g|\,\mathrm d\mu
 \leq\int_0^\infty\min\{\mu(E),d_g(u)\}\,\mathrm du
 =\int_0^{\mu(E)}g^*(s)\,\mathrm ds.
\]
最后的等式是在两个 Lebesgue 半直线之间使用非负积分换序，允许 \(\mu(E)=\infty\)。接着对测度 \(\nu(E)=\int_E|g|\,\mathrm d\mu\) 应用层蛋糕于 \(|f|\)，得到
\[
 \begin{aligned}
 \int_X|fg|\,\mathrm d\mu
 &=\int_0^\infty\int_{\{|f|>t\}}|g|\,\mathrm d\mu\,\mathrm dt\\
 &\leq\int_0^\infty\int_0^{d_f(t)}g^*(s)\,\mathrm ds\,\mathrm dt
 =\int_0^\infty f^*(s)g^*(s)\,\mathrm ds.
 \end{aligned}
\]
最后一步再用 \(s<d_f(t)\) 当且仅当 \(t<f^*(s)\)，只交换两个非负 Lebesgue 积分。代入 \(g^*(s)\leq\|g\|_{p',\infty}s^{-1/p'}\)，右端即为不超过 \(\|g\|_{p',\infty}\cdot\|f\|_{L^{p,1}}\) 的有限数。
\end{proof}
分层时需要注意恢复的对象：一般实、复函数的层蛋糕恢复的是 \(|f|\)，并非带符号或相位的 \(f\)。上述证明先控制绝对乘积，所以没有丢失相位条件。

\ProseParagraphBreak

这里一般不能把 \(L^{p,1}\) 放宽为 \(L^p\)。在 \((0,s_0)\) 上取同位置的两个函数
\[
 f(s)=s^{-1/p}(\log(e/s))^{-\beta},\quad g(s)=s^{-1/p'},
 \quad 1/p<\beta\leq1.
\]
前者是上面的对数模型，属于 \(L^p\)；后者属于弱 \(L^{p'}\)，但
\[\int_0^{s_0}|fg|\,\mathrm ds
 =\int_0^{s_0}s^{-1}(\log(e/s))^{-\beta}\,\mathrm ds=\infty.\]
若改为 \(\beta>1\)，则 \(f\in L^{p,1}\)，配对命题保证积分收敛。这一差别正好来自高度尺度上多一次可积求和。

\ProseParagraphBreak

最后考虑一种常见的输入端点：只先对集合示性函数检验算子。这样的\term[xianzhi-ruoxing-guji]{限制弱型估计}[restricted weak type estimate]比全部输入的弱型估计少一层要求，却自然产生 Lorentz 输入。
\Needspace{6\baselineskip}
\begin{proposition}[限制弱型与 Lorentz 输入]\label{b2:prop:17-restricted-weak}
设整数 \(d,n\geq1\)，\(1\leq p<\infty\)、\(1<q<\infty\)、\(q'=q/(q-1)\)。设复线性算子 \(T\) 定义在 \(\mathbb R^d\) 上有限测度支集的简单函数等价类中，输出 \(\mathbb R^n\) 上可测函数等价类。若存在 \(A<\infty\)，使每个有限测度可测集 \(E\subset\mathbb R^d\) 满足
\[\|T\mathbf1_E\|_{q,\infty}\leq A|E|^{1/p},\]
则 \(T\) 有唯一连续线性延拓 \(L^{p,1}(\mathbb R^d)\to L^{q,\infty}(\mathbb R^n)\)，且
\[N_q(Tf)\leq C_{p,q}A\|f\|_{L^{p,1}},\quad
 N_q(h)=\sup_{0<|B|<\infty}|B|^{-1/q'}\int_B|h|.\]
所有集合和积分取相应维数的 Lebesgue 测度。目标连续性可等价地用 \(N_q\) 范数或弱 \(L^q\) 拟范数表示。
\end{proposition}
\begin{proof}
若 \(f\) 是非负简单函数，\(f=\int_0^\infty\mathbf1_{\{f>t\}}\,\mathrm dt\) 实际只是有限个示性函数的正线性组合。\(N_q\) 的三角不等式与 \(N_q\leq q'\|\cdot\|_{q,\infty}\) 给出
\[N_q(Tf)\leq q'A\int_0^\infty d_f(t)^{1/p}\,\mathrm dt
 =\frac{q'A}{p}\|f\|_{L^{p,1}}.\]
最后的等式来自等分布关系和两个非负 Lebesgue 积分的换序。复简单函数分成实部、虚部的正负四部分，每部分的 Lorentz 量不超过原量，所以常数至多增加四倍。

为使用连续延拓，还须证明这类简单函数在输入空间稠密。先记一个拟三角界：超水平集合包含关系给出
\[(f+g)^*(s)\leq f^*(s/2)+g^*(s/2),\]
因而换元 \(r=s/2\) 得
\[\|f+g\|_{L^{p,1}}\leq2^{1/p}
 (\|f\|_{L^{p,1}}+\|g\|_{L^{p,1}}).\]
这里的重排界可以先取略高于两个广义逆的高度，使两个超水平集合的测度各不超过 \(s/2\)，再令高度下降取得。

给定 \(f\in L^{p,1}(\mathbb R^d)\)，保留 \(|x|\leq N\)、\(1/N\leq|f|\leq N\) 的部分，记为 \(f_N\)，令 \(e_N=|f-f_N|\)。则 \(e_N\downarrow0\) 几乎处处，且 \(e_N^*\leq f^*\)。由前面第二指数的上确界控制，\(d_f(t)<\infty\) 对每个 \(t>0\) 成立。因此 \(d_{e_N}(t)\downarrow0\)，进而 \(e_N^*(s)\to0\) 对每个 \(s>0\) 成立。以可积函数 \(s^{1/p-1}f^*(s)\) 支配，得到 \(\|f-f_N\|_{L^{p,1}}\to0\)。每个 \(f_N\) 有界且支集测度有限，可在该集合上作一致的简单函数逼近；若误差至多 \(\varepsilon\cdot\mathbf1_E\)，其 Lorentz 量至多 \(p\varepsilon|E|^{1/p}\)。结合拟三角界，即得所需稠密性。

最后对简单函数逼近列应用已证估计于两项之差，输出在完备的 \(N_q\) 范数中 Cauchy，因而有极限。差估计保证极限不依赖逼近列，并保留线性和同一上界；稠密性又保证延拓唯一。
\end{proof}
这个命题解释了为什么仅检验集合时，输入自然需要 \(L^{p,1}\) 的高度累加，而不是只用 \(L^p\) 的单一幂积分。目标条件 \(q>1\) 用于前面建立的等价 Banach 范数及完备性；在 \(q=1\) 时，这个证明不能原样使用。下一节的插值则从全部输入的两个弱端点出发，利用随输出高度改变的大小分解。

% !TeX root = ../../Book2.tex
\section{Marcinkiewicz 插值}
\label{b2:sec:1703}

弱型估计只控制输出超过某一高度的集合。若直接把它代入层蛋糕公式，端点处常会出现发散的高度积分。可以先把输入按大小分成两部分：大值部分交给较低的可积指数，小值部分交给较高的可积指数。分割高度随正在考察的输出高度改变，两个端点就分别控制同一个高度积分的两端。

\subsection{次线性算子}
\label{b2:subsec:170301}

设 \((X,\mathcal A,\mu)\)、\((Y,\mathcal B,\nu)\) 为测度空间，\(\mathbb K=\mathbb R\) 或 \(\mathbb C\)，\(\mathcal D\) 是 \(X\) 上有限值可测 \(\mathbb K\)-值函数的几乎处处等价类组成的向量空间。设 \(T\) 把 \(\mathcal D\) 映为 \(Y\) 上有限值可测函数的几乎处处等价类。若对所有 \(f,g\in\mathcal D\)、\(c\in\mathbb K\)，有
\[
 |T(f+g)|\leq |Tf|+|Tg|,
 \quad |T(cf)|=|c|\cdot|Tf|
 \quad\text{几乎处处},
\]
则称 \(T\) 为\term[cixianxing-suanzi]{次线性算子}[sublinear operator]。极大算子在 \(L^1+L^\infty\) 等已保证输出几乎处处有限的定义域上满足这些条件；线性算子以及线性算子的模也满足这些条件。这里保留的是输出大小的三角不等式，不要求输出本身能够线性叠加。

\ProseParagraphBreak

为什么把大值交给较小的指数？设 \(1\leq p_0<p<p_1<\infty\)，\(f\in L^p(X,\mu)\)，并取 \(a>0\)。令
\[
 f^{>a}=f\cdot\mathbf1_{\{|f|>a\}},
 \quad f^{\leq a}=f\cdot\mathbf1_{\{|f|\leq a\}}.
\]
在大值集合上，\(p_0-p<0\) 给出 \(|f|^{p_0-p}\leq a^{p_0-p}\)；小值集合上则有 \(|f|^{p_1-p}\leq a^{p_1-p}\)。所以
\[
 \int_{|f|>a}|f|^{p_0}\,\mathrm d\mu
 \leq a^{p_0-p}\|f\|_p^p,
 \quad
 \int_{|f|\leq a}|f|^{p_1}\,\mathrm d\mu
 \leq a^{p_1-p}\|f\|_p^p.
\]
若 \(p_1=\infty\)，后一条换成 \(\|f^{\leq a}\|_\infty\leq a\)。这也证明了每个中间 \(L^p\) 输入都属于和空间 \(L^{p_0}+L^{p_1}\)：两个截断确实落在算子的端点定义域中。

\ProseParagraphBreak

以上粗界适合检查截断的可积性，暂时不要把它们代入整个插值证明。这样代入后，两项都只剩下 \(t^{-p}\|f\|_p^p\) 一类的分布上界，层蛋糕积分又会遇到 \(\int\mathrm dt/t\)。要保留比弱型更多的信息，必须让集合 \(\{|f|>a\}\) 和 \(\{|f|\leq a\}\) 留在高度积分里面。

\subsection{弱型估计}
\label{b2:subsec:170302}

回顾\secref{b2:sec:1702}的记号，\(d_h(t)=\nu(\{|h|>t\})\)、\(t>0\)，表示目标函数的分布函数。若 \(T\) 在 \(L^{p_i}(X)\) 上满足弱型 \((p_i,p_i)\)，其含义是存在 \(A_i\geq0\)，使
\[
 d_{Th}(t)\leq
 \left(\dfrac{A_i\|h\|_{L^{p_i}(X)}}{t}\right)^{p_i}
 \quad(h\in L^{p_i}(X),\ t>0).
\]
当 \(p_i=\infty\) 时，使用的是强型条件
\(\|Th\|_{L^\infty(Y)}\leq A_i\|h\|_{L^\infty(X)}\)。

\ProseParagraphBreak

对定义域中的分解 \(f=f_0+f_1\)，如果 \(|Tf_0|\leq t/2\) 且 \(|Tf_1|\leq t/2\)，次线性便给出 \(|Tf|\leq t\)。因此，除一个零测集外，
\[
 \{|Tf|>t\}\subseteq
 \{|Tf_0|>t/2\}\cup\{|Tf_1|>t/2\}.
\]
由此可以分别使用两个端点，而不需要 \(Tf=Tf_0+Tf_1\)。若较高端点是 \(L^\infty\)，且其常数 \(A_1>0\)，取 \(a=t/(2A_1)\) 就使小值部分的输出不超过 \(t/2\)；此时只须估计大值部分的超水平集合。

\subsection{Marcinkiewicz 插值}
\label{b2:subsec:170303}

\term[marcinkiewicz-chazhi]{Marcinkiewicz 插值定理}[Marcinkiewicz interpolation theorem]先把两个同指数的弱型端点连接起来。这是后面从弱 \((1,1)\) 与强 \((2,2)\) 推出奇异积分中间指数估计时使用的版本。

\Needspace{6\baselineskip}
\begin{theorem}[Marcinkiewicz]\label{b2:thm:17-marcinkiewicz}
设 \((X,\mathcal A,\mu)\)、\((Y,\mathcal B,\nu)\) 为测度空间，\(\mathbb K=\mathbb R\) 或 \(\mathbb C\)，\(1\leq p_0<p_1\leq\infty\)。设 \(T\) 定义在 \(L^{p_0}(X;\mathbb K)+L^{p_1}(X;\mathbb K)\) 上，输出 \(Y\) 上可测函数的几乎处处等价类，且满足 \(|T(f+g)|\leq|Tf|+|Tg|\)、\(|T(cf)|=|c|\cdot|Tf|\) 几乎处处。设 \(A_0,A_1\geq0\)，并假设对每个有限端点 \(p_i\)，有
\[
 \nu(\{|Th|>t\})\leq
 \left(\dfrac{A_i\|h\|_{L^{p_i}(X)}}{t}\right)^{p_i}
 \quad(h\in L^{p_i}(X),\ t>0);
\]
若 \(p_1=\infty\)，该端点改为
\(\|Th\|_{L^\infty(Y)}\leq A_1\|h\|_{L^\infty(X)}\)。
对 \(p_0<p<p_1\)，以 \(1/\infty=0\) 约定唯一的 \(0<\theta<1\)，使
\[
 \frac1p=\frac{1-\theta}{p_0}+\frac\theta{p_1}.
\]
则对所有 \(f\in L^p(X;\mathbb K)\)，有
\[
 \|Tf\|_{L^p(Y)}
 \leq C_{p,p_0,p_1}A_0^{1-\theta}A_1^\theta\|f\|_{L^p(X)},
\]
其中常数只依赖所列指数。源和目标空间不要求无原子或 \(\sigma\)-有限。
\end{theorem}
\begin{proof}
先设 \(A_0,A_1>0\)，固定 \(f\in L^p(X)\)。这里 \(p<\infty\)，所以
\[
 E=\{|f|>0\}=\bigcup_{n\geq1}\{|f|>1/n\},
 \quad \mu(\{|f|>1/n\})\leq n^p\|f\|_p^p<\infty.
\]
因此 \(\mu\) 在 \(E\) 上的限制是 \(\sigma\)-有限的。以下与输入有关的非负重积分都支撑在 \(E\times(0,\infty)\) 上，Tonelli 定理可在那里使用；目标上的层蛋糕公式则用\cref{b2:thm:17-layer-cake}的一般测度版本。

先讨论 \(p_1<\infty\)。取一个待选择的 \(\kappa>0\)，在输出高度 \(t\) 处按输入高度 \(a=\kappa t\) 分割。前面已证明两个截断分别属于 \(L^{p_0}\)、\(L^{p_1}\)，故超水平集合估计给出
\begin{equation}\label{b2:eq:17-height-split}
 \begin{aligned}
 d_{Tf}(t)\leq{}&\left(\frac{2A_0}{t}\right)^{p_0}
       \int_{|f|>\kappa t}|f|^{p_0}\,\mathrm d\mu\\
 &+\left(\frac{2A_1}{t}\right)^{p_1}
       \int_{|f|\leq\kappa t}|f|^{p_1}\,\mathrm d\mu.
 \end{aligned}
\end{equation}
选择 \(a\) 与 \(t\) 成正比，使固定输入值 \(s=|f(x)|\) 对应的高度区间恰为 \(0<t<s/\kappa\) 或 \(t\geq s/\kappa\)。两个端点因而分别控制高度零点和无穷远处。

将\eqref{b2:eq:17-height-split}乘以 \(pt^{p-1}\)，对 \(t>0\) 积分。层蛋糕公式与 Tonelli 给出
\[
 \begin{aligned}
 \|Tf\|_p^p\leq{}&p(2A_0)^{p_0}
   \int_E |f|^{p_0}
      \left(\int_0^{|f|/\kappa}t^{p-p_0-1}\,\mathrm dt\right)\mathrm d\mu\\
 &+p(2A_1)^{p_1}
   \int_E |f|^{p_1}
      \left(\int_{|f|/\kappa}^\infty t^{p-p_1-1}\,\mathrm dt\right)\mathrm d\mu\\
 ={}&p\left(
   \frac{(2A_0)^{p_0}\kappa^{p_0-p}}{p-p_0}
  +\frac{(2A_1)^{p_1}\kappa^{p_1-p}}{p_1-p}\right)\|f\|_p^p.
 \end{aligned}
\]
第一段内积分有限恰好用到 \(p>p_0\)，第二段有限用到 \(p<p_1\)。左端在计算开始时允许为无穷，右端有限才证明 \(Tf\in L^p\)。若先取 \(\kappa=1\)，已经得到所求强型估计。

为了同时看清端点常数的作用，令两个分子中的常数因子相等，取
\[
 \kappa=\left(\frac{(2A_0)^{p_0}}{(2A_1)^{p_1}}\right)^{1/(p_1-p_0)}.
\]
由倒数指数关系，
\[
 p(1-\theta)=\frac{p_0(p_1-p)}{p_1-p_0},
 \quad p\theta=\frac{p_1(p-p_0)}{p_1-p_0}.
\]
于是两个分子均为
\(2^pA_0^{p(1-\theta)}A_1^{p\theta}\)，并得到显式的上界
\[
 \|Tf\|_p\leq
 2\left[p\left(\frac1{p-p_0}+\frac1{p_1-p}\right)\right]^{1/p}
 A_0^{1-\theta}A_1^\theta\|f\|_p.
\]
这项选择使两个分子相等，从而显示两个端点常数在上界中的幂次。

若 \(p_1=\infty\)，取 \(a=t/(2A_1)\)。小值部分的输出不超过 \(t/2\)，因此
\[
 d_{Tf}(t)\leq\left(\frac{2A_0}{t}\right)^{p_0}
       \int_{|f|>t/(2A_1)}|f|^{p_0}\,\mathrm d\mu.
\]
换序后，固定 \(x\) 的高度积分为
\(\int_0^{2A_1|f(x)|}t^{p-p_0-1}\,\mathrm dt\)。故
\[
 \|Tf\|_p^p\leq
 \frac{p(2A_0)^{p_0}(2A_1)^{p-p_0}}{p-p_0}\|f\|_p^p,
\]
即
\[
 \|Tf\|_p\leq
 2\left(\frac p{p-p_0}\right)^{1/p}
 A_0^{p_0/p}A_1^{1-p_0/p}\|f\|_p.
\]
此时 \(\theta=1-p_0/p\)，仍是同一个几何幂。最后，若某个 \(A_i=0\)，先用正的较大常数 \(A_i+\varepsilon\) 作上述估计，再令 \(\varepsilon\downarrow0\)，便得到相应结论。
\end{proof}

两个高度积分解释了为什么结论只在开区间内成立：当 \(p=p_0\) 或有限的 \(p=p_1\) 时，其中一个内积分成为对数发散。改变分割常数可以调整算子范数的上界，不能消除这一端点障碍。

\ProseParagraphBreak

定理假设算子已经在和空间上定义，因而截断可以直接送入算子。另一种常见情形是线性算子先只定义在有限测度支集的简单函数上，并在那里满足同样的两个端点界。这个类对高度截断封闭，证明先给出该类的中间强型界；它在有限指数 \(L^p\) 中稠密，且
\(\|Tf_n-Tf_m\|_p=\|T(f_n-f_m)\|_p\leq C\|f_n-f_m\|_p\)，
所以目标 \(L^p\) 的完备性产生唯一连续延拓。对一般次线性算子，若另外具有
\(|Tf-Tg|\leq|T(f-g)|\)，也可作同一构造。非负值次线性算子具有这项差值控制：分别把 \(f=g+(f-g)\)、\(g=f+(g-f)\) 代入次可加性即可。若只知道本节定义中的模不等式，输出的相位变化并不受控，不能由此断言算子本身连续或有唯一延拓。

\ProseParagraphBreak

输入和输出的指数也可以不同。下面保留一个适用于局部势核的版本；除了层蛋糕公式，还需要用到指数至少为一时的积分三角不等式。

\Needspace{6\baselineskip}
\begin{proposition}\label{b2:prop:17-off-diagonal}
设 \((X,\mathcal A,\mu)\)、\((Y,\mathcal B,\nu)\) 为测度空间，\(\mathbb K=\mathbb R\) 或 \(\mathbb C\)，\(1\leq r<\infty\)。设 \(T\) 定义在 \(L^1(X;\mathbb K)+L^\infty(X;\mathbb K)\) 上，输出 \(Y\) 上可测函数的几乎处处等价类，满足 \(|T(f+g)|\leq|Tf|+|Tg|\)、\(|T(cf)|=|c|\cdot|Tf|\) 几乎处处。假设存在 \(A\geq0\)、\(B>0\)，使
\[
 \begin{aligned}
 \nu(\{|Th|>t\})&\leq
 \left(\frac{A\|h\|_{L^1(X)}}t\right)^r
 &&(h\in L^1(X),\ t>0),\\
 \|Th\|_{L^\infty(Y)}&\leq B\|h\|_{L^\infty(X)}
 &&(h\in L^\infty(X)).
 \end{aligned}
\]
则对 \(1<p<\infty\)、\(f\in L^p(X;\mathbb K)\)，有
\[
 \|Tf\|_{L^{pr}(Y)}
 \leq 2\left(\frac p{p-1}\right)^{1/(pr)}
 A^{1/p}B^{1-1/p}\|f\|_{L^p(X)}.
\]
这里不要求两个测度空间 \(\sigma\)-有限。
\end{proposition}
\begin{proof}
固定 \(f\in L^p(X)\)，其非零支集 \(E\) 仍为 \(\sigma\)-有限。对输出高度 \(t>0\)，取
\[
 f_0=f\cdot\mathbf1_{\{|f|>t/(2B)\}},
 \quad f_1=f-f_0.
\]
有 \(\|f_0\|_1\leq(t/(2B))^{1-p}\|f\|_p^p\)、\(\|f_1\|_\infty\leq t/(2B)\)，所以两个端点都能使用，且小值输出不超过 \(t/2\)。于是
\[
 d_{Tf}(t)\leq\left(\frac{2A}t\right)^r
 \left(\int_E|f(x)|\cdot\mathbf1_{\{t<2B|f(x)|\}}\,\mathrm d\mu(x)\right)^r.
\]
置 \(q=pr\)，按层蛋糕公式，
\[
 \|Tf\|_q^q\leq q(2A)^r
 \int_0^\infty t^{q-r-1}
 \left(\int_E |f(x)|\cdot\mathbf1_{\{t<2B|f(x)|\}}\,\mathrm d\mu(x)\right)^r\mathrm dt.
\]
与同指数证明相比，输入积分现在还带有 \(r\) 次幂，不能直接把这次幂移进 \(x\) 积分。

记 \(\mathrm d\rho(t)=t^{q-r-1}\,\mathrm dt\)，并将括号内的非负函数记为 \(H(t)\)。在这个高度测度上使用积分 Minkowski 不等式，得到
\[
 \begin{aligned}
 \|H\|_{L^r(\rho)}
 &\leq\int_E |f(x)|
       \left(\int_0^{2B|f(x)|}t^{q-r-1}\,\mathrm dt\right)^{1/r}\mathrm d\mu(x)\\
 &=\frac{(2B)^{p-1}}{(q-r)^{1/r}}\|f\|_p^p.
 \end{aligned}
\]
这里 \(q-r=r(p-1)>0\) 保证零点处可积。所用的积分不等式也可以直接核验：\(r=1\) 时它就是 Tonelli；\(r>1\) 时，令 \(r'=r/(r-1)\)，对任意非负 \(h\in L^{r'}(\rho)\)、\(\|h\|_{r'}\leq1\)，先由 Tonelli 交换 \(x,t\) 积分，再对每个 \(x\) 用 Hölder，便有
\[
 \int_0^\infty H(t)\cdot h(t)\,\mathrm d\rho(t)
 \leq\int_E |f(x)|
       \left(\int_0^{2B|f(x)|}\mathrm d\rho(t)\right)^{1/r}\mathrm d\mu(x).
\]
为避免预设 \(H\in L^r(\rho)\)，先取
\(H_n=\min\{H,n\}\cdot\mathbf1_{[1/n,n]}\)。它有有限 \(L^r(\rho)\) 范数；非零时用
\(h=H_n^{r-1}/\|H_n\|_r^{r-1}\) 测试，得到同一个范数上界，再由单调收敛令 \(n\to\infty\)。所有重积分都在 \(E\times(0,\infty)\) 这个 \(\sigma\)-有限乘积空间上进行。

将 \(H\) 的估计代回，取 \(1/r\) 次方，得到
\[
 \|Tf\|_q^p
 \leq q^{1/r}(2A)\cdot\frac{(2B)^{p-1}}{(q-r)^{1/r}}\|f\|_p^p
 =2^p\left(\frac p{p-1}\right)^{1/r}AB^{p-1}\|f\|_p^p.
\]
再取 \(1/p\) 次方即得结论。若 \(A=0\)，可先用正的较大弱型常数，再取极限。
\end{proof}

这个指数可以在参数线上读出：对端点 \((1,r)\) 与 \((\infty,\infty)\)，取 \(\theta=1-1/p\)，则
\[
 \frac1q=\frac{1-\theta}r=\frac1{pr},
 \quad A^{1-\theta}B^\theta=A^{1/p}B^{1-1/p}.
\]
高度积分中的 Minkowski 论证使这个参数计算成为强型估计；假如 \(r<1\)，该积分三角不等式不再成立，以上证明便不能照用。

\ProseParagraphBreak

例如在 \(\mathbb R^d\) 上取 \(0<\beta<d\)，令
\[
 k(x)=|x|^{-\beta}\cdot\mathbf1_{\{0<|x|<1\}},
 \quad r=\frac d\beta>1,
 \quad \omega_d=|B(0,1)|.
\]
球坐标积分与超水平集合的半径给出
\[
 \|k\|_1=\frac{d\omega_d}{d-\beta},
 \quad
 d_k(s)=\begin{cases}\omega_d,&0<s<1,\\
                         \omega_ds^{-r},&s\geq1,
          \end{cases}
 \quad \|k\|_{r,\infty}=\omega_d^{1/r}.
\]
由\secref{b2:sec:1702}的集合积分范数
\(N_r(h)=\sup_{0<|F|<\infty}|F|^{-1/r'}\int_F|h|\)，有
\(N_r(k)\leq r'\omega_d^{1/r}\)。对 \(h\in L^1\)，把\cref{b2:prop:17-weak-complete}用于可积因子 \(h\) 与弱空间因子 \(k\)，得
\[
 \|k*h\|_{r,\infty}\leq N_r(k*h)
 \leq N_r(k)\cdot\|h\|_1.
\]
另一方面，\(k\in L^1\) 给出
\(\|k*h\|_\infty\leq\|k\|_1\cdot\|h\|_\infty\)。卷积积分在 \(L^1+L^\infty\) 上具有兼容的定义：可积部分的卷积几乎处处绝对收敛，有界部分的卷积处处绝对收敛；其和等于原输入的卷积积分几乎处处。于是命题给出
\[
 \|k*f\|_{pr}
 \leq 2\left(\frac p{p-1}\right)^{1/(pr)}
       (r'\omega_d^{1/r})^{1/p}
       \left(\frac{d\omega_d}{d-\beta}\right)^{1-1/p}\|f\|_p.
\]
取 \(d=3\)、\(\beta=1\)、\(p=2\)，就得到这个局部核的 \(L^2\to L^6\) 估计，两个端点常数可分别取 \(\tfrac32(4\pi/3)^{1/3}\) 与 \(2\pi\)。这是本命题给出的一种充分指数范围。若删除核的空间截断，输入 \(h=1\) 已使远场积分发散，所用的 \(L^\infty\) 端点便不成立；全空间分数积分需要另行估计远场。

\subsection{极大算子应用}
\label{b2:subsec:170304}

\Needspace{6\baselineskip}
\begin{corollary}\label{b2:thm:17-maximal-strong}
设 \(d\) 为正整数，\(1<p\leq\infty\)。在 \(\mathbb R^d\) 上取 Lebesgue 测度，\(M\)、\(M_u\) 分别为中心球和包含观察点的球上绝对平均的上确界；\(M_{\mathcal D}\) 为标准半开二进网格
\(\mathcal D=\{\,2^{-j}(k+[0,1)^d):j\in\mathbb Z,\ k\in\mathbb Z^d\,\}\)
上包含观察点的立方体绝对平均的上确界。则三个算子都把 \(L^p(\mathbb R^d;\mathbb C)\) 映入自身，并满足
\[
 \|Mf\|_p+\|M_uf\|_p+\|M_{\mathcal D}f\|_p
 \leq C_{d,p}\|f\|_p.
\]
各算子在 \(L^\infty\) 上的范数均不超过一。
\end{corollary}
\begin{proof}
\cref{b2:thm:17-maximal-weak}及 \(M_u\leq2^dM\) 给出前两个算子的弱 \((1,1)\) 界，\cref{b2:thm:17-dyadic-level}给出二进版本。每个绝对平均都不超过输入的本性上确界，故三个算子均为 \(L^\infty\) 压缩。在 \(L^1+L^\infty\) 上，可积部分的极大函数几乎处处有限，有界部分的极大函数处处有界；次可加性因而使三个算子都具有所需定义域。应用\cref{b2:thm:17-marcinkiewicz}的 \(p_0=1,p_1=\infty\) 情形，即得有限 \(p\) 的结论。
\end{proof}

对中心极大算子，还可保留一个以后用于端点问题的精细分布界。设 \(d\geq1\)、\(1\leq p<\infty\)、\(f\in L^p(\mathbb R^d)\)、\(t>0\)。分解
\(f=f\cdot\mathbf1_{\{|f|>t/2\}}+f\cdot\mathbf1_{\{|f|\leq t/2\}}\)，
大值部分可积，小值部分的极大平均不超过 \(t/2\)。对大值部分用弱型估计，得到
\begin{equation}\label{b2:eq:17-maximal-tail-refined}
 |\{Mf>t\}|\leq\frac{2\cdot5^d}{t}
                   \int_{|f|>t/2}|f|\,\mathrm dx.
\end{equation}
当 \(p>1\) 时，保留截断再作高度积分，便有
\[
 \|Mf\|_p^p\leq \frac{5^dp\,2^p}{p-1}\|f\|_p^p.
\]
分母 \(p-1\) 来自 \(\int_0^{2|f(x)|}t^{p-2}\,\mathrm dt\)。

\ProseParagraphBreak

这一常数在 \(p\downarrow1\) 时的增长，可以用具体输入检验。在实直线上取 \(f=\mathbf1_{[0,1]}\)。当 \(x>2\) 时，半径为 \(x\) 的中心区间给出 \(Mf(x)\geq1/(2x)\)，所以
\[
 \|Mf\|_p^p\geq2^{-p}\int_2^\infty x^{-p}\,\mathrm dx
 =\frac{2^{1-2p}}{p-1},
 \quad \|f\|_p=1.
\]
因此算子范数至少为一个固定常数乘以 \((p-1)^{-1/p}\)。在 \(1<p\leq2\) 上，这与 \((p-1)^{-1}\) 可比：若 \(\delta=p-1\in(0,1]\)，则
\[
 \frac{\delta^{-1/p}}{\delta^{-1}}
 =\delta^{\delta/(1+\delta)}\in[e^{-1/e},1].
\]
结合上界，实直线上的中心极大算子范数在 \(p\downarrow1\) 时确有 \(1/(p-1)\) 的量级。精确最佳常数还依赖所选的极大算子，不由这两个界确定。

\ProseParagraphBreak

极大函数还能同时控制一整族正则化算子，进而说明它们在 Lebesgue 点恢复原函数。下面的核可取复值，不要求非负。

\Needspace{6\baselineskip}
\begin{proposition}\label{b2:prop:17-approximate-identity}
设 \(d\) 为正整数，\(\psi\in L^1(\mathbb R^d;\mathbb C)\)、\(\int_{\mathbb R^d}\psi=1\)。假设存在 \(\varepsilon>0\)、\(0<C<\infty\)，使
\[
 |\psi(x)|\leq C(1+|x|)^{-d-\varepsilon}
 \quad\text{几乎处处}.
\]
对 \(t>0\)，置 \(\psi_t(x)=t^{-d}\cdot\psi(x/t)\)。若 \(1\leq p<\infty\)、\(f\in L^p(\mathbb R^d;\mathbb C)\)，则 \(\psi_t*f\to f\) 于 \(L^p\)。对每个满足 \(Mf(x)<\infty\) 的点，所有 \(t>0\) 的卷积积分均绝对收敛，且
\[
 \sup_{t>0}|\psi_t*f(x)|\leq C'\cdot Mf(x),
\]
其中 \(M\) 为中心 Hardy--Littlewood 极大算子，\(C'\) 只依赖 \(d,\varepsilon,C\)。对任一有限值可测代表 \(f\)，在它的每个 Lebesgue 点 \(x\) 都有 \(\psi_t*f(x)\to f(x)\) 当 \(t\downarrow0\)。
\end{proposition}
\begin{proof}
先说明共同的极大函数控制。把积分区域分为 \(|y|<t\) 及环带 \(2^{j-1}t\leq|y|<2^jt\)、\(j\geq1\)。内球上 \(|\psi_t|\leq Ct^{-d}\)，相应积分不超过 \(C\omega_d\cdot Mf(x)\)。第 \(j\) 个环带上，核的上界为
\(C2^{d+\varepsilon}t^{-d}2^{-j(d+\varepsilon)}\)，而 \(|f(x-y)|\) 的积分不超过
\(\omega_d(2^jt)^d\cdot Mf(x)\)。所以
\[
 \int_{\mathbb R^d}|\psi_t(y)|\cdot|f(x-y)|\,\mathrm dy
 \leq C_{d,\varepsilon}C
       \left(1+\sum_{j\geq1}2^{-j\varepsilon}\right)\cdot Mf(x).
\]
右端与 \(t\) 无关，且几何级数收敛。这给出所述上确界估计及全部尺度的绝对收敛；由弱型或强型极大估计，满足 \(Mf(x)<\infty\) 的点组成满测集。

范数收敛可以回用\cref{b2:thm:03-approximate-identity}：缩放保持 \(\int\psi_t=1\)、\(\|\psi_t\|_1=\|\psi\|_1\)，且对每个 \(a>0\)，
\(\int_{|y|\geq a}|\psi_t(y)|\,\mathrm dy
=\int_{|z|\geq a/t}|\psi(z)|\,\mathrm dz\to0\)。
其证明在本处也可直接写成
\[
 \|\psi_t*f-f\|_p
 \leq\int_{\mathbb R^d}|\psi(z)|\cdot
             \|f(\,\cdot-tz)-f\|_p\,\mathrm dz.
\]
有限指数的平移连续性使每个固定 \(z\) 的范数差趋零，而共同可积控制为
\(2\|f\|_p\cdot|\psi(z)|\)；支配收敛给出结论。

固定 \(f\) 的一个 Lebesgue 点 \(x\)，令
\[
 A_x(r)=\frac1{|B(x,r)|}\int_{B(x,r)}|f(y)-f(x)|\,\mathrm dy.
\]
有 \(A_x(r)\to0\) 当 \(r\downarrow0\)。此外 \(S_x=\sup_{r>0}A_x(r)<\infty\)：先选 \(r_0>0\) 控制 \(r<r_0\) 的平均；对 \(r\geq r_0\)，Hölder 给出
\[
 A_x(r)\leq (\omega_dr^d)^{-1/p}\|f\|_p+|f(x)|
 \leq(\omega_dr_0^d)^{-1/p}\|f\|_p+|f(x)|.
\]
这也说明 \(Mf(x)\leq S_x+|f(x)|<\infty\)，从而该点的卷积积分已有定义。由 \(\int\psi_t=1\)，将前面的环带估计用于 \(|f(x-y)-f(x)|\)，得到
\[
 |\psi_t*f(x)-f(x)|
 \leq C_{d,\varepsilon}C
       \left(A_x(t)+\sum_{j\geq1}2^{-j\varepsilon}A_x(2^jt)\right).
\]
给定误差要求，先利用 \(A_x\leq S_x\) 与几何级数选出一个固定的尾部起点，使尾和对所有 \(t\) 都小；剩下只有有限个半径，每个都随 \(t\downarrow0\) 趋零。因此右端趋零。
\end{proof}

这里恢复的是所选代表在 Lebesgue 点的取值。\cref{b2:thm:17-differentiation}保证这些点几乎处处存在；一个 \(L^p\) 等价类本身并不指定任意一点的函数值。

% !TeX root = ../../Book2.tex
\section{Riesz–Thorin 插值}
\label{b2:sec:1704}

上一节把输入按大小分开，再对输出的高度积分。若算子本身是线性的，而且两个端点已经给出强范数估计，还可以让输入随一个复参数变化：在条带两侧分别适合两个端点，在条带内部指定的位置回到原输入。输出再与一族测试函数配对，就把算子估计变成一个复函数的边界估计。这样得到的常数是两个端点常数的几何平均。

\ProseParagraphBreak

这条路线需要保留复线性结构。极大算子的绝对值和上确界一般不随复参数全纯，不能直接放入本节的证明；Fourier 变换、固定核卷积等线性算子则可以。先证明实际需要的三线估计，再从竖边的范数要求推导测试函数的选择。

\ProseParagraphBreak

这里所用的最大模原理只涉及有界矩形 \(Q\)：若 \(H\) 在矩形内部全纯并连续到闭矩形，则内部的 \(|H|\) 不超过边界上的最大值。所需的基本复分析事实是全纯函数在内部光滑，且实部、虚部调和；写 \(H=u+\mathrm iv\)，Cauchy--Riemann 方程 \(u_x=v_y\)、\(u_y=-v_x\) 给出 \(\Delta u=v_{yx}-v_{xy}=0\)，虚部同理。因此对每个单位复数 \(c\)，可把本卷\cref{b2:cor:13-boundary-maximum}用于实值调和函数 \(\operatorname{Re}(cH)\)。给定内点 \(z\)，选择 \(c\) 使 \(cH(z)=|H(z)|\)，便得到
\[
 |H(z)|=\operatorname{Re}(cH(z))
 \leq\max_{\partial Q}\operatorname{Re}(cH)
 \leq\max_{\partial Q}|H|.
\]
若 \(H(z)=0\)，最后的模界直接成立。这里全纯指在复开域中处处存在复导数；光滑性与 Cauchy--Riemann 方程是本处使用的局部复分析基础。相同的调和最大值论证也适用于有界扇形等有界区域，只要函数在内部全纯并连续到整个闭域。闭矩形的紧性保证边界最大值存在；下面仍须先控制每个有限矩形的横边，才能令矩形高度趋于无穷，不能直接把这条有界域结论套到无限条带。

\subsection{三线引理与解析族}
\label{b2:subsec:170401}

下面采用\term[hadamard-sanxian-yinli]{Hadamard 三线引理}[Hadamard three-lines lemma]的下述版本。

\Needspace{6\baselineskip}
\begin{lemma}[Hadamard 三线估计]\label{b2:thm:17-three-lines}
设 \(M_0,M_1\geq0\)。复值函数 \(F\) 在开条带 \(0<\operatorname{Re}z<1\) 全纯，在闭条带 \(0\leq\operatorname{Re}z\leq1\) 连续且有界，并对每个 \(t\in\mathbb R\) 满足 \( |F(\mathrm it)|\leq M_0\)、\( |F(1+\mathrm it)|\leq M_1\)。则对每个 \(0<\theta<1\)，有
\[
 |F(\theta)|\leq M_0^{1-\theta}M_1^\theta.
\]
\end{lemma}
\begin{proof}
先设 \(M_0,M_1>0\)，令 \(G(z)=F(z)M_0^{z-1}M_1^{-z}\)，幂以正实数的实对数定义。若 \(z=s+\mathrm it\)，则归一化因子的模为 \(M_0^{s-1}M_1^{-s}\)，与 \(t\) 无关；所以 \(G\) 在闭条带有界，两条竖边的模不超过一。

竖边界尚未控制有限矩形的横边。为此固定 \(\varepsilon>0\)，乘上 \(e^{\varepsilon(z^2-1)}\)。它在 \(s=0,1\) 时的模不超过一，在 \(t=\pm R\)、\(0\leq s\leq1\) 时的模不超过 \(e^{-\varepsilon R^2}\)。若 \(|G|\leq B\)，只须选 \(R\) 足够大使 \(B e^{-\varepsilon R^2}\leq1\)，四条边就都受一控制。矩形最大模原理给出
\[
 |G(\theta)|e^{\varepsilon(\theta^2-1)}\leq1,
 \quad
 |F(\theta)|\leq M_0^{1-\theta}M_1^\theta
 e^{\varepsilon(1-\theta^2)}.
\]
令 \(\varepsilon\downarrow0\) 得到结论。若某个 \(M_i=0\)，先以正数 \(M_i+\delta\) 作为两侧较大的界，再令 \(\delta\downarrow0\)。这同时处理了一侧或两侧恒为零的情形。
\end{proof}

三线估计要求的对象是一个在整个条带中有界的标量全纯函数，不只是两条竖边上的范数公式。下面取有限个简单函数的组合，既能逐项核对全纯性，也能直接证明内部有界。

\ProseParagraphBreak

设 \(U\subset\mathbb C\) 开，\(T_z\) 是定义在同一复测试向量空间上、\(z\in U\) 的复线性算子族，并固定可配对的测试输出空间。一个\term[jiexi-zu]{解析族}[analytic family]是复参数下的算子族，其对测试输入、测试输出的矩阵元为全纯函数。使用三线估计时必须检查边界范数和条带内增长；“公式含有 \(z\)”不是解析族的充分条件。下面先用固定算子和变化的测试函数，避免讨论一般算子的定义域变化。

\subsection{Riesz–Thorin 定理}
\label{b2:subsec:170402}

两个端点必须作用于同一个输入，并在交空间上给出同一输出。例如，一个端点从积分定义算子，另一个从完备化定义算子，就要先核对这两种定义在交空间上一致。此时可令 \(T(f_0+f_1)=T_0f_0+T_1f_1\)：交空间上的一致性保证结果与分解无关。下面直接把这项相容性写成和空间上的一个复线性算子。

\ProseParagraphBreak

本节把\cref{b2:subsec:04-complex-interpolation}中出现的 Riesz--Thorin 方法推广到一般线性算子，允许输入、输出的两个端点指数分别变化。约定 \(1/\infty=0\)，\(r'\) 表示满足 \(1/r+1/r'=1\) 的 Hölder 共轭指数。

\Needspace{6\baselineskip}
\begin{theorem}[Riesz--Thorin]\label{b2:thm:17-riesz-thorin}
设 \((X,\mathcal A,\mu)\)、\((Y,\mathcal B,\nu)\) 为 \(\sigma\)-有限测度空间，所有 \(L^r\) 空间均取复值函数的几乎处处等价类。设 \(1\leq p_i,q_i\leq\infty\)、\(A_i\geq0\)、\(i=0,1\)，且 \(T\) 是从 \(L^{p_0}(X,\mu)+L^{p_1}(X,\mu)\) 到 \(Y\) 上几乎处处有限的可测函数等价类空间的复线性算子，满足
\[
 \|Tf\|_{L^{q_i}(Y,\nu)}\leq A_i\|f\|_{L^{p_i}(X,\mu)}
 \quad\bigl(f\in L^{p_i}(X,\mu),\ i=0,1\bigr).
\]
给定 \(0<\theta<1\)，以 \(1/\infty=0\) 定义 \(p,q\)：
\[
 \dfrac1p=\dfrac{1-\theta}{p_0}+\dfrac\theta{p_1},\quad
 \dfrac1q=\dfrac{1-\theta}{q_0}+\dfrac\theta{q_1},
\]
则每个 \(f\in L^p(X,\mu)\) 都属于上述和空间，且
\[
 \|Tf\|_{L^q(Y,\nu)}
 \leq A_0^{1-\theta}A_1^\theta\|f\|_{L^p(X,\mu)}.
\]
\end{theorem}
\begin{proof}
先证明有限测度支集上的简单输入的估计。零输入无需处理；将输入 \(f\) 与有限测度支集上的简单测试函数 \(g\) 归一化为
\(\|f\|_p=\|g\|_{q'}=1\)。先设 \(p,q'<\infty\)。希望构造 \(f_z\)，使它在 \(z=\theta\) 等于 \(f\)，在 \(\operatorname{Re}z=i\)、\(i=0,1\) 时的 \(L^{p_i}\) 范数等于一。保留 \(f\) 的复相位，试取
\[
 f_z=\frac f{|f|}\cdot|f|^{\alpha(z)}.
\]
若 \(p_i<\infty\)，竖边的归一化要求
\(p_i\operatorname{Re}\alpha(i+\mathrm it)=p\)，这样积分恰好仍是 \(\int|f|^p=1\)。若 \(p_i=\infty\)，则要求实部为零，使支集上模恒为一。这两种要求统一为 \(\alpha(i)=p/p_i\)。取连接这两个值的仿射函数，就得到
\[
 \alpha(z)=p\left(\frac{1-z}{p_0}+\frac z{p_1}\right),
 \quad \alpha(\theta)=1.
\]
输出使用 Hölder 配对，所以应对 \(g\) 插值的是共轭指数。由于
\[
 \frac1{q'}=\frac{1-\theta}{q_0'}+\frac\theta{q_1'},
\]
相同的选择给出
\[
 \beta(z)=q'\left(\frac{1-z}{q_0'}+\frac z{q_1'}\right),
 \quad
 g_z=\frac g{|g|}\cdot|g|^{\beta(z)},
 \quad g_\theta=g.
\]
所有复幂都以非零函数值的模的实对数定义；原函数为零的位置统一赋零。即使某侧实部为零，也不把支集外的零改成一。在竖边 \(i=0,1\) 上，以上计算给出
\[
 \|f_{i+\mathrm it}\|_{p_i}
 =\|g_{i+\mathrm it}\|_{q_i'}=1.
\]
这里含无穷端点的等式，是相应函数在其非零支集上模恒为一，而非使用任何无穷次幂。

还要补足中间指数本身为无穷的构造。若 \(p=\infty\)，倒数插值式迫使 \(p_0=p_1=\infty\)。此时对任意归一化的 \(f\in L^\infty(X)\) 直接取 \(f_z=f\)，不要求它是简单函数或具有有限测度支集。若 \(q'=\infty\)，则 \(q_0'=q_1'=\infty\)，即两个输出端点均为 \(L^1\)；这时取 \(g_z=g\)，其中 \(g\) 仍为归一化的有限测度支集简单函数。其余一侧仍按上面的有限指数公式构造。这样所有情形都满足竖边范数不超过一，并在 \(\theta\) 处恢复原函数。

定义双线性配对
\[
 F(z)=\int_Y (Tf_z)(y)\cdot g_z(y)\,\mathrm d\nu(y).
\]
这里不对 \(g_z\) 取复共轭，否则参数一般会变成反全纯。对偶测试所需的共轭相位可以预先放进 \(g\) 中。若 \(f=\sum_j a_j\cdot\mathbf1_{E_j}\)、\(g=\sum_k b_k\cdot\mathbf1_{H_k}\)，则 \(F\) 是有限个整函数系数的乘积与固定数
\(\int_{H_k}T\mathbf1_{E_j}\,\mathrm d\nu\) 的线性组合。\(E_j,H_k\) 测度有限，端点估计及 Hölder 保证这些积分绝对收敛。若 \(p=\infty\)，只把其中的 \(T\mathbf1_{E_j}\) 换为固定的 \(Tf\)，同样由一个输出端点保证局部可积。因此 \(F\) 全纯并连续到闭条带，无需未经控制地交换求导和积分。

在闭条带上，复幂系数的模只依赖 \(\operatorname{Re}z\in[0,1]\)，因而各自有界；有限和使 \(F\) 在整个闭条带有界。两条竖边则满足
\[
 |F(i+\mathrm it)|
 \leq\|Tf_{i+\mathrm it}\|_{q_i}
 \cdot\|g_{i+\mathrm it}\|_{q_i'}
 \leq A_i.
\]
由\cref{b2:thm:17-three-lines}，在 \(z=\theta\) 得到
\[
 \left|\int_Y (Tf)\cdot g\,\mathrm d\nu\right|
 \leq A_0^{1-\theta}A_1^\theta.
\]
恢复两个归一化因子后，右侧再乘 \(\|f\|_p\cdot\|g\|_{q'}\)。

现在对所有这些 \(g\) 取上确界。对于在每个有限测度集上可积的可测函数 \(h\)，所用的范数刻画是
\[
 \|h\|_q
 =\sup_{\substack{g\text{ 为有限测度支集简单函数}\\\|g\|_{q'}\leq1}}
 \left|\int_Y h\cdot g\,\mathrm d\nu\right|.
\]
这个等式允许两边取无穷。对 \(1<q<\infty\)，先在有限测度集及 \(|h|\) 有界的部分上，取共轭相位 \(\overline h/|h|\) 乘 \(|h|^{q-1}\)，再归一化、简单逼近并穷竭；此过程也检验了原先未知的 \(L^q\) 可积性。对 \(q=1\)，测试函数改为该共轭相位在有限测度集上的截断；对 \(q=\infty\)，在 \(|h|\) 接近本性上确界的有限正测度子集上，使用该共轭相位除以集合测度的测试函数。所有相位在 \(h=0\) 处赋零。\(\sigma\)-有限性保证最后这样的子集能够选出，也允许用有限测度集穷竭 \(Y\)。当前的 \(Tf\) 在两个输出端点中，Hölder 保证它在每个有限测度集上可积，故上述测试适用。由此得到简单输入的估计；\(p=\infty\) 时，前面已经对任意有界输入证明了结论，无需再用稠密性。

最后设 \(p<\infty\)。\(\sigma\)-有限性允许先截取有限测度支集，再截断幅度并作简单逼近，得到 \(f_n\rightarrow f\) 于 \(L^p(X)\)。已证估计使 \(Tf_n\) 在 \(L^q(Y)\) 中 Cauchy，即使 \(q=\infty\) 也可用该空间的完备性，记极限为 \(v\)。这定义了唯一的有界线性延拓，但仍须核对 \(v\) 等于题设和空间中原有的 \(Tf\)。

为此将两个源端点按大小记为 \(r_0\leq p\leq r_1\)。对非零 \(h\in L^p\)，令 \(a=\|h\|_p\)，以 \(|h|>a\) 与 \(|h|\leq a\) 分解为 \(h_0+h_1\)。幂次比较给出
\[
 \|h_0\|_{r_0}\leq\|h\|_p,
 \quad \|h_1\|_{r_1}\leq\|h\|_p;
\]
若 \(r_1=\infty\)，第二项直接由 \(|h_1|\leq a\) 得到。于是 \(L^p\) 连续嵌入原来的端点和空间。对 \(h=f_n-f\) 使用这个分解，端点估计使 \(T h_0,T h_1\) 在各自的输出端点范数中趋零。Chebyshev 不等式说明它们在每个有限测度目标集上依测度趋零；无穷输出端点则直接给一致的本性上界趋零。因此 \(Tf_n\rightarrow Tf\) 在这些集合上依测度成立。另一方面 \(Tf_n\rightarrow v\) 的 \(L^q\) 收敛也蕴含同样的依测度收敛。极限唯一性和 \(Y\) 的可数穷竭给出 \(v=Tf\) 几乎处处，从而原算子本身满足所述估计。
\end{proof}

\begin{table}[htbp]
\centering
\caption[解析测试族的指数选择]{解析测试族的指数选择。有限中间指数时，复幂的实部使竖边范数归一，插值点的幂则为一；约定 \(1/\infty=0\)。中间指数 \(p\) 或 \(q'\) 本身为无穷时，相应的两个测试指数 \(p_i\) 或 \(q_i'\) 也均为无穷，应固定原函数，不能计算含 \(\infty\) 的复幂。}
\label{b2:tab:17-analytic-test-exponents}
\begin{tabularx}{\linewidth}{@{}p{9.4em}XXX@{}}
\toprule
测试族 & 左竖边的实部 & 右竖边的实部 & 插值点的取值\\
\midrule
\(p<\infty\) 的 \(f_z\) & \(\operatorname{Re}\alpha=p/p_0\) & \(\operatorname{Re}\alpha=p/p_1\) & \(\alpha(\theta)=1\)\\
\(q'<\infty\) 的 \(g_z\) & \(\operatorname{Re}\beta=q'/q_0'\) & \(\operatorname{Re}\beta=q'/q_1'\) & \(\beta(\theta)=1\)\\
\(p=\infty\) 的 \(f_z\) & 固定 \(f\) & 固定 \(f\) & \(f_\theta=f\)\\
\(q'=\infty\) 的 \(g_z\) & 固定 \(g\) & 固定 \(g\) & \(g_\theta=g\)\\
\bottomrule
\end{tabularx}
\end{table}
\cref{b2:tab:17-analytic-test-exponents}把两个作用分开：边界实部决定可用的端点范数，虚部只改变相位；插值点的幂等于一，才使最终估计针对原来的输入与测试函数。

\ProseParagraphBreak

解析算子族 \(T_z\) 的使用也应归结为同样的标量检验：对所选的变化测试函数，\(F(z)=\int(T_zf_z)\cdot g_z\) 必须在条带内部全纯、连续到边界。单独知道固定测试函数的矩阵元全纯时，还须说明变化测试族是有限解析组合，或给出相应的积分控制。若 \(F\) 在整个条带至多按 \(C_{f,g}e^{c|\operatorname{Im}z|}\) 增长，而两条竖边仍有常数界 \(A_i\)，对任意 \(\delta>0\) 乘上 \(e^{\delta(z-\theta)^2}\) 后便可用有界版本；其竖边界增加的因子至多是 \(e^{\delta(i-\theta)^2}\)。令 \(\delta\downarrow0\)，仍得精确的几何平均常数。

\ProseParagraphBreak

后文的某些族在竖边上也会增长。若对所有归一化测试函数有共同的 \(c\geq0\)，并且
\[
 |F(s+\mathrm it)|\leq C_{f,g}e^{ct^2}
 \quad(0\leq s\leq1),
 \quad |F(i+\mathrm it)|\leq A_i e^{ct^2}\quad(i=0,1),
\]
则取 \(\delta>c\) 后，同一乘子抑制了内部增长。应用三线估计得到
\[
 |F(\theta)|
 \leq e^{\delta\theta(1-\theta)}A_0^{1-\theta}A_1^\theta.
\]
因为 \(c\) 对测试函数一致，附加常数也一致；这里不能再令 \(\delta\) 趋零。曲面限制中的解析族允许这种 Gaussian 型增长。增长条件及其常数依赖必须在具体应用中验证，不能只引用有界版本后略去它们。

\ProseParagraphBreak

先看一个能直接算出尺度的应用。设 \(d\) 为正整数、\(\kappa,t>0\)，取热核
\[
 H_{\kappa,t}(x)=(4\pi\kappa t)^{-d/2}
 e^{-|x|^2/(4\kappa t)},
 \quad T_t f=H_{\kappa,t}*f.
\]
固定 \(1\leq p<\infty\)。核的积分为一，所以 Young 不等式给出 \(L^p\rightarrow L^p\) 的常数一；Hölder 给出 \(L^p\rightarrow L^\infty\) 的常数 \(\|H_{\kappa,t}\|_{p'}\)。这两个端点使用同一个积分卷积，故相容。对 \(p<q<\infty\)，输入指数保持为 \(p\)，目标指数满足 \(1/q=(1-\theta)/p\)，即 \(\theta=1-p/q\)。Gaussian 积分算出
\[
 \|H_{\kappa,t}\|_r
 =r^{-d/(2r)}(4\pi\kappa t)^{-\frac d2(1-1/r)}
 \quad(1\leq r<\infty),
 \quad \|H_{\kappa,t}\|_\infty=(4\pi\kappa t)^{-d/2}.
\]
因此
\[
 \|T_t f\|_q
 \leq C_{d,p,q}(4\pi\kappa t)^{-\frac d2(1/p-1/q)}\|f\|_p,
 \quad
 C_{d,p,q}=(p')^{-\frac{d(1-p/q)}{2p'}}.
\]
当 \(p=1\) 时该常数按 \(p'=\infty\) 的极限取一。\(q=p\) 与 \(q=\infty\) 分别由刚才两个端点直接得到；\(p=q=\infty\) 也由核的积分为一得到压缩。譬如 \(p=2,q=4\) 时，
\[
 \|T_t f\|_4\leq 2^{-d/8}(4\pi\kappa t)^{-d/8}\|f\|_2.
\]
正时间带来更高的可积性，而常数在 \(t\downarrow0\) 时按明确的幂发散。插值没有丢失尺度：两个端点中 \(\kappa t\) 的幂，按同一个 \(\theta\) 组合。这也解释了为什么应先记录尺度依赖，再对许多尺度求和。

\subsection{Hausdorff–Young 再论}
\label{b2:subsec:170403}

\Needspace{6\baselineskip}
\begin{corollary}[Hausdorff--Young]\label{b2:thm:17-hausdorff-young}
设 \(d\) 为正整数，\(1\leq p\leq2\)，\(p'\) 为 Hölder 共轭指数（\(p=1\) 时 \(p'=\infty\)）。对 \(f\in L^p(\mathbb R^d;\mathbb C)\)，令 \(\widehat f\) 为使用核 \(e^{-2\pi\mathrm ix\cdot\xi}\) 的可积 Fourier 变换或其从有限测度支集简单函数延拓所得的变换。则
\[
 \|\widehat f\|_{p'}\leq\|f\|_p.
\]
\end{corollary}
\begin{proof}
可积 Fourier 变换满足 \(L^1\rightarrow L^\infty\) 的范数界一；\cref{b2:thm:04-plancherel-extension}给出 \(L^2\rightarrow L^2\) 的等距延拓，并明确说明它在 \(L^1\cap L^2\) 上与可积变换一致。因此二者给出同一个和空间算子。对 \(1<p<2\)，取
\[
 \theta=2-\frac2p,
 \quad \frac1p=1-\frac\theta2,
 \quad \frac1{p'}=\frac\theta2.
\]
将端点 \((p_0,q_0)=(1,\infty)\)、\((p_1,q_1)=(2,2)\) 及 \(A_0=A_1=1\) 代入\cref{b2:thm:17-riesz-thorin}，即得所需估计。两个端点由原有结论直接处理。对有限测度支集简单函数的 \(L^p\) 逼近产生唯一连续延拓；对一般 \(L^p\) 输入，这个结论不要求普通 Fourier 积分逐点绝对收敛。
\end{proof}
常数一来自两个端点的相同上界。Gaussian 的可计算范数已经提示最优常数问题比有界性更细。逆变换具有相同的两个端点，因而同样得到 \(L^p\) 到 \(L^{p'}\) 的估计。

\ProseParagraphBreak

指数上限也能用本章可完成的有限计算检验。设 \(d\) 为正整数、\(2<p\leq\infty\)，置 \(r=p'\in[1,2)\)。选非负、非零的 \(\psi\in C_c^\infty((0,1/4)^d)\)，并对符号 \(\varepsilon_k\in\{-1,1\}\) 令
\[
 f_N(x)=\sum_{k=1}^N\varepsilon_k\cdot\psi(x-ke_1),
 \quad
 \widehat f_N(\xi)=\widehat\psi(\xi)\cdot
 \sum_{k=1}^N\varepsilon_k e^{-2\pi\mathrm ik\xi_1}.
\]
平移后的支集不交，所以
\(\|f_N\|_p=N^{1/p}\|\psi\|_p\)，包括 \(p=\infty\) 时的约定 \(N^{1/p}=1\)。若符号全部为正，只在 \(|\xi_1|\lesssim1/N\) 的窄峰上计算，所得 \(L^r\) 下界是 \(N^{1-1/r}=N^{1/p}\)，与输入同阶，不能据此排除有界性。需要让变换在固定大小的区域上保留约 \(\sqrt N\) 的平均量级。

\ProseParagraphBreak

这里的符号平均只是在全部 \(2^N\) 种选择上取算术平均，记作 \(\mathbb E\)；事件的概率 \(\mathbb P\) 就是相应选择所占的比例。对固定 \(\xi\)，令 \(S_N=\sum_k\varepsilon_k e^{-2\pi\mathrm ik\xi_1}\)。展开乘积时，含有某个符号的奇数次幂的项平均为零；在四次乘积中，留下的下标至多有三种两两配对方式。因此
\[
 \mathbb E|S_N|^2=N,
 \quad \mathbb E|S_N|^4\leq3N^2.
\]
设 \(B=\{\,|S_N|^2\geq N/2\,\}\)。其补集上二次量小于 \(N/2\)，故 Cauchy--Schwarz 给出
\[
 \frac N2\leq\mathbb E\bigl(|S_N|^2\cdot\mathbf1_B\bigr)
 \leq\bigl(\mathbb E|S_N|^4\bigr)^{1/2}\cdot\mathbb P(B)^{1/2}.
\]
于是 \(\mathbb P(B)\geq1/12\)，并有
\[
 \mathbb E|S_N|^r\geq\frac1{12}(N/2)^{r/2}.
\]
由于 \(\widehat\psi(0)=\int\psi>0\) 且变换连续，存在固定的有限正测度频率盒 \(J\) 与 \(a>0\)，使 \(|\widehat\psi|\geq a\) 于 \(J\)。对 \(\xi\in J\) 积分并与有限符号平均交换，得到
\[
 \mathbb E\int_J|\widehat f_N(\xi)|^r\,\mathrm d\xi
 \geq\frac{a^r\cdot|J|}{12}(N/2)^{r/2}.
\]
至少有一组符号使 \(\|\widehat f_N\|_r\geq c_{\psi,r}\sqrt N\)。因此
\[
 \frac{\|\widehat f_N\|_{p'}}{\|f_N\|_p}
 \geq c_{\psi,p}N^{1/2-1/p}\longrightarrow\infty.
\]
这些输入本身光滑且紧支撑，已经排除了 \(p>2\) 时任何与输入无关的 \(L^p\rightarrow L^{p'}\) 常数。符号的作用是让固定频率区域上的平均大小保持为平方根量级，而不是只制造越来越窄的同号峰。

\subsection{插值方法与 \texorpdfstring{\(K\)}{K} 泛函}
\label{b2:subsec:170404}

Marcinkiewicz 方法使用大小分解和分布函数，能处理极大算子这样的次线性对象；Riesz--Thorin 方法保留复线性结构，从两个强型端点得到几何平均常数。对奇异积分，常见路线是先用弱 \((1,1)\) 与强 \((2,2)\) 得到 \(1<p<2\)，再由伴随算子的估计处理另一半指数。选用哪种方法，取决于已知端点的实际强弱，不能把弱型范数直接当作解析插值中的强型范数。

\ProseParagraphBreak

大小分解也提示另一种问题：不预先规定截断高度，而把一个输入分配给两个空间，能达到的最小总代价是什么？设 \(X_0,X_1\) 是同一实或复标量域上的 Banach 空间，连续嵌入同一个 Hausdorff 线性拓扑空间，使 \(f_0+f_1\) 有共同的意义。对 \(f\in X_0+X_1\) 及 \(t>0\)，定义其 \term[k-fanhan]{\(K\) 泛函}[K-functional]
\[
 K(t,f;X_0,X_1)\coloneqq
 \inf_{\substack{f=f_0+f_1\\f_i\in X_i}}
 \bigl(\|f_0\|_{X_0}+t\|f_1\|_{X_1}\bigr).
\]
参数 \(t\) 调整第二个空间的代价。\(t\) 小时可以把更多部分交给 \(X_1\)，\(t\) 大时则倾向于减少 \(X_1\) 的范数。因此下确界中的分解必须允许随 \(t\) 改变。对 \(L^1,L^\infty\)，有限的 \(L^\infty\) 代价对应削平函数的大值，剩余高度由 \(L^1\) 支付；这个分配可以逐点调整幅度，不要求把某个水平集切成任意指定测度的子集。

\Needspace{6\baselineskip}
\begin{proposition}\label{b2:prop:17-k-functional}
设 \((X,\mathcal A,\mu)\) 为 \(\sigma\)-有限测度空间，\(f\in L^1(X,\mu)+L^\infty(X,\mu)\) 为实或复值函数。令 \(d_f(\lambda)=\mu\{|f|>\lambda\}\)、\(\lambda\geq0\)，并定义递减重排 \(f^*(s)=\inf\{\lambda\geq0:d_f(\lambda)\leq s\}\)、\(s>0\)，空集的下确界取无穷。对每个 \(t>0\)，令 \(K(t,f;L^1,L^\infty)=\inf_{f=u+v}(\|u\|_1+t\|v\|_\infty)\)，其中 \(u\in L^1(X,\mu)\)、\(v\in L^\infty(X,\mu)\)。则
\[
 K(t,f;L^1,L^\infty)=\int_0^t f^*(s)\,\mathrm ds.
\]
另设 \((X_0,X_1)\)、\((Y_0,Y_1)\) 为同一实或复标量域上的相容 Banach 空间对（每对连续嵌入同一个 Hausdorff 线性拓扑空间），\(T:X_0+X_1\rightarrow Y_0+Y_1\) 是线性算子。若 \(T:X_i\rightarrow Y_i\) 的范数分别至多 \(A_i>0\)、\(i=0,1\)，则对每个 \(f\in X_0+X_1\) 和 \(t>0\)，有
\[
 K(t,Tf;Y_0,Y_1)\leq A_0\cdot K(tA_1/A_0,f;X_0,X_1).
\]
\end{proposition}
\begin{proof}
先把重排积分改写为分布函数积分。对 \(s>0\)、\(\lambda\geq0\)，若 \(d_f(\lambda)\leq s\)，则 \(\lambda\) 已属于定义下确界的集合，故 \(f^*(s)\leq\lambda\)。若 \(d_f(\lambda)>s\)，右连续性保证某个 \(\delta>0\) 满足 \(d_f(\lambda+\delta)>s\)，于是 \(f^*(s)\geq\lambda+\delta>\lambda\)。因此 \(f^*(s)>\lambda\) 当且仅当 \(s<d_f(\lambda)\)。对 \(s\in(0,t)\) 使用层蛋糕公式及 Tonelli 定理，得到
\[
 J_t\coloneqq\int_0^t f^*(s)\,\mathrm ds
 =\int_0^\infty\min\{t,d_f(\lambda)\}\,\mathrm d\lambda.
\]
这个恒等式暂时允许两边为无穷；当某个 \(d_f(\lambda)=\infty\) 时，被积函数的相应值就是有限数 \(t\)。因此不必假定各个超水平集合都有有限测度。

取任意分解 \(f=u+v\)，其中 \(u\in L^1\)、\(v\in L^\infty\)，并令 \(b=\|v\|_\infty\)。在 \(\lambda>b\) 时，\(|f|\leq|u|+b\) 几乎处处蕴含
\(d_f(\lambda)\leq d_u(\lambda-b)\)。将高度积分在 \(b\) 处分开，得到
\[
 J_t\leq tb+\int_b^\infty d_u(\lambda-b)\,\mathrm d\lambda
 =t\|v\|_\infty+\|u\|_1.
\]
最后一步是 \(u\) 的层蛋糕公式。至少存在一个这样的分解，所以 \(J_t<\infty\)；再对所有分解取下确界，得 \(J_t\leq K(t,f;L^1,L^\infty)\)。整个比较通过高度分布完成，因此也适用于带原子的测度空间。

反向令 \(a=f^*(t)\)。重排单调不增，所以 \(ta\leq J_t<\infty\)，从而 \(a\) 是有限非负数。由广义逆定义及分布函数的右连续性，\(d_f(\lambda)>t\) 对 \(0\leq\lambda<a\) 成立，而 \(d_f(\lambda)\leq t\) 对 \(\lambda\geq a\) 成立。因此
\[
 J_t=ta+\int_a^\infty d_f(\lambda)\,\mathrm d\lambda.
\]
当 \(a=0\) 时，第一段高度区间为空，等式仍成立。现在以保持相位的截断定义
\[
 f_1(x)=
 \begin{cases}f(x),&|f(x)|\leq a,\\a f(x)/|f(x)|,&|f(x)|>a,\end{cases}
 \quad f_0=f-f_1.
\]
这里 \(f_1\) 满足 \(\|f_1\|_\infty\leq a\)，且 \(|f_0|=(|f|-a)_+\)。再次由层蛋糕公式，
\[
 \|f_0\|_1=\int_a^\infty d_f(\lambda)\,\mathrm d\lambda
 =J_t-ta<\infty.
\]
因此这是一个合法的 \(L^1+L^\infty\) 分解，其代价至多为 \(J_t-ta+ta=J_t\)。两边合起来即得所需等式。水平集即使是不可分割的原子，仍可对该原子上的函数值作幅度截断，因而不影响这个分解。

对于算子估计，每个分解产生 \(Tf=Tf_0+Tf_1\)，代价至多 \(A_0\|f_0\|_{X_0}+tA_1\|f_1\|_{X_1}\)。对全部分解取下确界即可。
\end{proof}

例如，取单点空间 \(X=\{x_0\}\)，令 \(\mu(\{x_0\})=1\)、\(f(x_0)=1\)、\(t=1/2\)。它的递减重排在 \(0<s<1\) 上恒为一，故 \(K(1/2,f;L^1,L^\infty)=1/2\)。直接取 \(f_1=f\)、\(f_0=0\) 就达到这个代价。另一方面，满足 \(\mu(E)\leq1/2\) 的集合只有空集，因此
\[
 \sup_{\mu(E)\leq1/2}\int_E|f|\,\mathrm d\mu=0.
\]
原子不能切成半个集合，但可以把原子上的函数值截断。这个例子说明，上述集合积分上确界不能在一般测度空间中代替 \(K\) 泛函，重排公式仍然成立则来自幅度截断。

\ProseParagraphBreak

对相容 Banach 空间对 \((X_0,X_1)\)、\(0<\theta<1\)、\(1\leq q\leq\infty\)，将所有尺度的分解代价合起来，可定义实插值量
\[
 \|f\|_{\theta,q;K}\coloneqq
 \left(\int_0^\infty
 \bigl[t^{-\theta}\cdot K(t,f;X_0,X_1)\bigr]^q
 \frac{\mathrm dt}{t}\right)^{1/q}
 \quad(f\in X_0+X_1),
\]
当 \(q=\infty\) 时取 \(\sup_{t>0}t^{-\theta}\cdot K(t,f;X_0,X_1)\)。乘法测度 \(\mathrm dt/t\) 使不同数量级的尺度具有同样权重；因子 \(t^{-\theta}\) 则决定在两个端点之间的位置。

\ProseParagraphBreak

对同一 \(\sigma\)-有限测度空间上的 \(L^1,L^\infty\)，令 \(1/p=1-\theta\)，故 \(1<p<\infty\)、\(\theta=1/p'\)。对 \(f\in L^1+L^\infty\)，上一命题给出的公式使实插值量与\secref{b2:sec:1702}中定义的 Lorentz 量满足
\[
 \|f\|_{L^{p,q}}\leq\|f\|_{\theta,q;K}
 \leq p'\|f\|_{L^{p,q}}.
\]
先看左侧。由于 \(f^*\) 单调不增，
\[
 K(t,f)=\int_0^t f^*(s)\,\mathrm ds
 \geq t\cdot f^*(t),
 \quad t^{-1/p'}\cdot K(t,f)\geq t^{1/p}f^*(t).
\]
对 \(t\) 取 \(L^q(\mathrm dt/t)\) 范数便得所需下界。反向令 \(h(s)=s^{1/p}f^*(s)\)，换元 \(s=tu\) 得到
\[
 t^{-1/p'}\cdot K(t,f)
 =\int_0^1u^{-1/p}\cdot h(tu)\,\mathrm du.
\]
在 \(\mathrm dt/t\) 中，固定 \(u>0\) 的缩放 \(t\mapsto tu\) 保持范数。因此有限 \(q\geq1\) 时，Minkowski 给出
\[
 \|f\|_{\theta,q;K}
 \leq\int_0^1u^{-1/p}\,\mathrm du\cdot
 \|h\|_{L^q(\mathrm dt/t)}
 =p'\|f\|_{L^{p,q}}.
\]
\(q=\infty\) 时把上确界移入积分得到同一界。这里 \(p>1\) 保证 \(\int_0^1u^{-1/p}\,\mathrm du<\infty\)，\(q\geq1\) 则保证所用 Minkowski 有效。不能只保留形式相同的重排表达而删除这两个条件。

\ProseParagraphBreak

上一命题的算子估计又给出相容 Banach 空间对之间的结论。若复线性或实线性算子 \(T:X_0+X_1\rightarrow Y_0+Y_1\) 的两个端点范数分别至多 \(A_0,A_1>0\)，则
\[
 K(t,Tf;Y_0,Y_1)
 \leq A_0\cdot K(tA_1/A_0,f;X_0,X_1).
\]
在实插值量中令 \(u=tA_1/A_0\)，有
\[
 t^{-\theta}A_0\cdot K(tA_1/A_0,f)
 =A_0^{1-\theta}A_1^\theta u^{-\theta}\cdot K(u,f).
\]
由于 \(\mathrm dt/t=\mathrm du/u\)，对有限 \(q\) 积分，或在 \(q=\infty\) 时取上确界，得到
\[
 \|Tf\|_{\theta,q;K(Y_0,Y_1)}
 \leq A_0^{1-\theta}A_1^\theta
 \|f\|_{\theta,q;K(X_0,X_1)}.
\]
这里的下标仅说明分别用哪一对空间计算 \(K\) 泛函。端点常数的总次数为一，与算子缩放 \(T\mapsto cT\) 时范数乘 \(|c|\) 相符。这一计算解释了几何平均常数为什么在两种插值方法中都会出现：复方法来自三线估计，实方法来自分解代价的尺度换元。

\ProseParagraphBreak

实方法还记录 Lorentz 的第二指数，但上述 Banach 空间对估计并不自动给出任意两个弱型端点之间的最佳 Lorentz 范围。弱端点首先提供的是分布函数控制，能否使用等价范数、输入如何分解以及目标是否完备，都须按具体指数检查。固定的可积核卷积可以使用强型端点和解析插值；临界弱核常更适合大小分解及集合积分；需要区分整个高度分布时，\(K\) 泛函又提供了相应的量。

\ProseParagraphBreak

这些方法也都不能仅靠开放区间的结论恢复端点。极大函数在每个 \(p>1\) 上有界，但单位球示性函数的像有不可积尾部；当估计常数在 \(p\downarrow1\) 时发散，便没有可用于取极限的统一控制。类似地，Fourier 的两个端点给出中间强型界，并不决定最优常数、核的抵消或主值存在性。使用插值之后，仍应按当前问题分别证明这些额外性质。


\begin{chaptersummary}
极大函数的弱型估计控制的是“大于某个高度的点有多少”。单独一个弱端点仍可能产生对数发散；Marcinkiewicz 方法让输入截断随输出高度改变，再交换积分次序，才得到严格位于两个端点之间的强型估计。二进分解还提供坏块的平均上下界，为下一章利用零平均抵消准备了空间分块。

\ProseParagraphBreak

Riesz--Thorin 方法则把端点指数编码在解析测试函数中，由三线引理得到范数常数的几何平均。使用时先核实两端是同一个兼容线性算子，并区分有限指数的稠密延拓与无穷端点的直接论证。Lorentz 空间进一步区分临界幂次下的对数可积性，\(K\) 泛函把分解代价写成重排积分。后面选择哪种方法，取决于手中掌握的是弱型还是强型、算子是否线性，以及是否需要保留尺度参数和第二指数。
\end{chaptersummary}

\BookExercises

% 自拟：围绕本章已经建立的工具，改变条件或组合方法。
\Needspace{6\baselineskip}
\begin{exercise}\label{b2:ex:17-new-1}
在带 Lebesgue 测度的 \(\mathbb R\) 上令 \(f=\mathbf1_{(0,1)}\)。求 \(x>1\) 时的中心和非中心极大函数，并由此说明二者不相等。
\end{exercise}
\begin{solution}
中心区间半径 \(r\leq x-1\) 时积分为零；\(x-1<r<x\) 时平均为 \((r-x+1)/(2r)\)，随 \(r\) 增大；\(r\geq x\) 时平均为 \(1/(2r)\)，随 \(r\) 减小。所以 \(Mf(x)=1/(2x)\)。包含 \(x\) 的非中心区间若左端为 \(a\in[0,1]\)，平均至多 \((1-a)/(x-a)\leq1/x\)；左端小于零或右端超过 \(x\) 只增大分母。端点极限给出 \(M_uf(x)=1/x\)。
\end{solution}

% 自拟：组合两个分布尾部得到弱型乘积界，并检验强型结论的失败。
\Needspace{6\baselineskip}
\begin{exercise}\label{b2:ex:17-new-2}
设 \((X,\mathcal A,\mu)\) 为测度空间，\(0<p,q<\infty\)，并以 \(1/r=1/p+1/q\) 定义 \(r\)。对 \(f\in L^{p,\infty}(X)\)、\(g\in L^{q,\infty}(X)\)，证明
\[
 \|fg\|_{r,\infty}\leq 2^{1/r}\|f\|_{p,\infty}\cdot\|g\|_{q,\infty}.
\]
再在 \((0,1)\) 上构造这样的 \(f,g\)，使 \(fg\notin L^r(0,1)\)。不要求求出弱型乘积界的最佳常数。
\end{exercise}
\begin{solution}
记 \(A=\|f\|_{p,\infty}\)、\(B=\|g\|_{q,\infty}\)。若任一常数为零，相应函数几乎处处为零，结论成立。设 \(A,B>0\)。对任意 \(a,t>0\)，有
\[
 \{|fg|>t\}\subset\{|f|>a\}\cup\{|g|>t/a\},
 \quad d_{fg}(t)\leq A^pa^{-p}+B^qa^qt^{-q}.
\]
选择 \(a\) 使右边两项相等，即 \(a^{p+q}=A^pt^q/B^q\)。利用 \(r=pq/(p+q)\)，两项都等于 \((AB/t)^r\)，故 \(t\cdot d_{fg}(t)^{1/r}\leq2^{1/r}AB\)。对高度取上确界便得估计。

在 \((0,1)\) 上取 \(f(x)=x^{-1/p}\)、\(g(x)=x^{-1/q}\)。两者的弱拟范数均为一，而 \(|f(x)g(x)|^r=x^{-1}\)，其积分发散。因此两个弱型因子的乘积一般仍只能得到弱型控制。
\end{solution}

% 自拟：用统一弱型尾部与依测度收敛取得严格低于端点的范数收敛。
\Needspace{6\baselineskip}
\begin{exercise}\label{b2:ex:17-new-3}
设 \((X,\mathcal A,\mu)\) 为有限测度空间，\(0<p<\infty\)，可测函数列 \(f_n\) 满足 \(\sup_n\|f_n\|_{p,\infty}\leq A<\infty\)，且对每个 \(t>0\) 都有 \(\mu\{|f_n|>t\}\to0\)。证明对每个 \(0<r<p\)，都有 \(\|f_n\|_r\to0\)。在 \((0,1)\) 上给出反例，说明一般不能取 \(r=p\)。
\end{exercise}
\begin{solution}
记 \(m=\mu(X)\)，固定 \(R>0\)。层蛋糕公式把范数分成
\[
 \|f_n\|_r^r
 =r\int_0^R t^{r-1}d_{f_n}(t)\,\mathrm dt
  +r\int_R^\infty t^{r-1}d_{f_n}(t)\,\mathrm dt.
\]
第一段的被积函数逐点趋零，由可积函数 \(rm t^{r-1}\) 控制，因而趋零。第二段不超过
\[
 rA^p\int_R^\infty t^{r-p-1}\,\mathrm dt
 =\frac{rA^p}{p-r}R^{r-p}.
\]
先令 \(n\to\infty\)，再令 \(R\to\infty\)，得到结论。两种假设各有作用：依测度收敛处理有限高度，统一弱型界处理高值尾部。

反例取 \(f_n=n^{1/p}\cdot\mathbf1_{(0,1/n)}\)。它逐点趋零，弱 \(L^p\) 拟范数与强 \(L^p\) 范数都恒为一，而 \(\|f_n\|_r^r=n^{r/p-1}\to0\)。严格不等式 \(r<p\) 不能删除。
\end{solution}

% 自拟：精算幅度截断在弱Lp与L1中的不同误差。
\Needspace{6\baselineskip}
\begin{exercise}\label{b2:ex:17-new-4}
在 Lebesgue 测度空间 \((0,1)\) 上，设 \(f(x)=x^{-1/p}\)、\(1<p<\infty\)，并令 \(f_N(x)=\min\{f(x),N\}\)、\(N\geq1\)。求 \(f-f_N\) 的弱 \(L^p\) 拟范数；再求 \(\|f-f_N\|_1\) 并判断其极限。解释为什么这两个结论并不矛盾。
\end{exercise}
\begin{solution}
令 \(h_N=f-f_N=(f-N)_+\)。对 \(t>0\)，
\[
 d_{h_N}(t)=(N+t)^{-p},\quad
 \|h_N\|_{p,\infty}=\sup_{t>0}\frac{t}{N+t}=1.
\]
另一方面，层蛋糕公式给出
\[
 \|h_N\|_1=\int_0^\infty(N+t)^{-p}\,\mathrm dt
 =\frac{N^{1-p}}{p-1}\longrightarrow0.
\]
幅度截断消除了越来越少的 \(L^1\) 质量，但在任意高的水平上仍留下同一临界幂次尾部。弱拟范数取遍所有高度的上确界，因此不会随截断误差的积分一起趋零。此题也给出了上一题严格降低目标指数时的一个可算模型。
\end{solution}

% 自拟：将精细的极大函数尾部界用于LlogL输入的局部可积性。
\Needspace{6\baselineskip}
\begin{exercise}\label{b2:ex:17-new-5}
设 \(d\geq1\)，\(f:\mathbb R^d\to\mathbb C\) 可测且
\(\int_{\mathbb R^d}|f(x)|\log(e+|f(x)|)\,\mathrm dx<\infty\)。对任意有限 Lebesgue 测度集合 \(E\subset\mathbb R^d\)，证明中心极大函数满足
\[
 \int_E Mf\leq |E|+C_d\int_{\mathbb R^d}|f|\log(e+|f|).
\]
解释这一结论为什么没有推出 \(M:L^1\to L^1\)。
\end{exercise}
\begin{solution}
假设先保证 \(f\in L^1\)。对每个 \(t>0\)，把 \(f\) 分为大于 \(t/2\) 的部分和其余部分。后者的极大函数至多为 \(t/2\)，对前者用\cref{b2:thm:17-maximal-weak}，得
\[
 |\{Mf>t\}|\leq\frac{2\cdot5^d}{t}\int_{|f|>t/2}|f|.
\]
对 \(Mf\) 在 \(E\) 上用层蛋糕公式，高度 \(0<t<1\) 的积分至多为 \(|E|\)。其余部分用上式及非负积分换序，得到
\[
 \begin{aligned}
 \int_1^\infty |E\cap\{Mf>t\}|\,\mathrm dt
 &\leq2\cdot5^d\int_{|f|>1/2}|f(x)|
       \int_1^{2|f(x)|}\frac{\mathrm dt}{t}\,\mathrm dx\\
 &=2\cdot5^d\int_{|f|>1/2}|f|\log(2|f|).
 \end{aligned}
\]
当 \(s>1/2\) 时，\(\log(2s)\leq2\log(e+s)\)，故取 \(C_d=4\cdot5^d\) 即可。这里一方面对输入增加了对数可积性，另一方面只在有限测度集合上积分输出。即使 \(f\) 是单位球的示性函数，\(Mf\) 的全空间积分仍因远场尾部而发散。
\end{solution}

% 自拟：围绕本章已经建立的工具，改变条件或组合方法。
\Needspace{6\baselineskip}
\begin{exercise}\label{b2:ex:17-new-6}
设 \(d\geq1\)，\(K\in L^1(\mathbb R^d)\cap L^\infty(\mathbb R^d)\) 是复值核，\(Tf=K*f\)，Fourier 变换采用 \(e^{-2\pi\mathrm ix\cdot\xi}\) 的约定。用两个端点证明，对 \(1\leq p\leq2\) 及 \(f\in L^p(\mathbb R^d)\)，有
\[
 \|Tf\|_{p'}\leq\|K\|_\infty^{2/p-1}\cdot
 \|\widehat K\|_\infty^{2-2/p}\cdot\|f\|_p.
\]
其中 \(1/p+1/p'=1\)；\(K=0\) 时单独理解结论。
\end{exercise}
\begin{solution}
绝对积分给出 \(L^1\to L^\infty\) 常数 \(\|K\|_\infty\)。对 \(f\in L^1\cap L^2\)，卷积公式与 Plancherel 等式给出
\[
 \|K*f\|_2=\|\widehat K\,\widehat f\|_2
 \leq\|\widehat K\|_\infty\cdot\|f\|_2.
\]
由 \(L^2\) 稠密性得到平方端点；它与积分卷积在交空间上一致。对这两个兼容端点应用\cref{b2:thm:17-riesz-thorin}，源指数要求 \(1/p=1-\theta/2\)，故 \(\theta=2-2/p\)。目标倒数为 \(\theta/2=1/p'\)，恰得所需公式。一般输入的卷积还由 \(K\in L^1\) 的 Young 估计定义，它与上述延拓相同。若 \(\|\widehat K\|_\infty=0\)，Fourier 唯一性给出 \(K=0\)；若 \(\|K\|_\infty=0\)，也有同一结论，因而端点没有零的零次幂问题。
\end{solution}

% 自拟：围绕本章已经建立的工具，改变条件或组合方法。
\Needspace{6\baselineskip}
\begin{exercise}\label{b2:ex:17-new-7}
在测度空间 \((X,\mathcal A,\mu)\) 中，设 \(N\geq1\)，\(Q_1,\ldots,Q_N\) 是两两不交、有限正测度的可测集合，并记 \(U=\bigcup_jQ_j\)。对 \(f\in L^1(U;\mathbb C)\) 定义
\[
 Pf=\sum_{j=1}^N\left(\frac1{\mu(Q_j)}\int_{Q_j}f\,\mathrm d\mu\right)\cdot\mathbf1_{Q_j}.
\]
证明 \(P^2=P\)，且 \(P\) 在每个 \(L^p(U)\)、\(1\leq p\leq\infty\) 上的算子范数恰为一。进一步，若 \(f\in L^2(U)\)，证明 \(Pf\) 是所有在各 \(Q_j\) 上分别为常数的函数中唯一的最佳 \(L^2\) 逼近。
\end{exercise}
\begin{solution}
二次平均不改变常数，所以 \(P^2=P\)。积分三角不等式给出 \(\|Pf\|_1\leq\|f\|_1\)，每个平均不超过本性上确界又给出 \(\|Pf\|_\infty\leq\|f\|_\infty\)。在有限测度空间 \(U\) 上应用\cref{b2:thm:17-riesz-thorin}，所有中间指数的范数也至多为一。实函数可先把同一积分公式作用在复函数上，再限制回实函数。取非零函数 \(\mathbf1_U\)，它被 \(P\) 保持，故范数恰为一。

若 \(v\) 在各块上为常数，则每个块上 \(\int_{Q_j}(f-Pf)=0\)，从而
\(\int_U(f-Pf)\overline{(Pf-v)}\,\mathrm d\mu=0\)。展开平方模的积分便得
\[
 \|f-v\|_2^2=\|f-Pf\|_2^2+\|Pf-v\|_2^2.
\]
最小值在且仅在 \(v=Pf\) 几乎处处时取得。这同时说明为何平方误差下应取算术平均，而不必另猜每个块上的最优常数。
\end{solution}

% 自拟：用临界幂函数检验K泛函与Lorentz弱范数等价中的上界常数。
\Needspace{6\baselineskip}
\begin{exercise}\label{b2:ex:17-new-8}
在 Lebesgue 测度空间 \((0,1)\) 上，设 \(1<p<\infty\)、\(p'=p/(p-1)\)、\(f(x)=x^{-1/p}\)。计算 \(K(t,f;L^1,L^\infty)\)、\(t>0\)，并求
\[
 \sup_{t>0}t^{-1/p'}K(t,f;L^1,L^\infty).
\]
由此判断正文中 \(\|f\|_{1/p',\infty;K}\leq p'\|f\|_{p,\infty}\) 的系数 \(p'\) 能否普遍减小。进一步说明该函数为何不属于任何 \(L^{p,q}(0,1)\)、\(0<q<\infty\)。
\end{exercise}
\begin{solution}
此函数可积，重排在 \(0<s<1\) 时为 \(s^{-1/p}\)，在 \(s\geq1\) 时为零。由\cref{b2:prop:17-k-functional}，
\[
 K(t,f;L^1,L^\infty)=
 \begin{cases}
 p't^{1/p'},&0<t\leq1,\\
 p',&t>1.
 \end{cases}
\]
因此所求上确界为 \(p'\)，在每个 \(0<t\leq1\) 处取到，而 \(\|f\|_{p,\infty}=1\)。正文已经对任意输入证明了系数 \(p'\) 的上界，本例又达到该上界，所以这个等价估计的系数确实不能减小。

若 \(q<\infty\)，则 \(s^{1/p}f^*(s)=1\) 在 \((0,1)\) 上成立，Lorentz 量的 \(q\) 次幂包含 \(\int_0^1\mathrm ds/s=\infty\)。\(K\) 量也表现为所有小尺度上同一个常数：它的上确界有限，但对乘法测度 \(\mathrm dt/t\) 积分便发散。
\end{solution}

% !TeX root = ../../Book2.tex
\chapter{奇异积分与 Calderón–Zygmund 理论}
\label{b2:ch:18}
\BookChapterNumber{b2:ch:18}
\begin{chapterabstract}
若卷积核在原点附近像距离的负维数次幂那样增长，取绝对值后的积分通常发散，Young 不等式便不能直接使用。然而，核的正负部分可能相互抵消，使对称截断仍有极限。本章从 Hilbert 变换的具体计算出发，说明怎样把这种抵消转化为范数估计：先将输入分成有界的好部与零平均的坏块，再把坏块的零平均移到核的差分上。由此得到弱端点、中间指数的有界性与主值实现，最后用于共轭调和函数和椭圆方程的二阶导数估计。
\end{chapterabstract}

\begin{learninggoals}\begin{itemize}
\item\label{b2:item:18-goal-1}计算 Hilbert 变换的主值与 Fourier 符号，区分平方范数有界和点态极限。
\item\label{b2:item:18-goal-2}构造 Calderón--Zygmund 分解，说明坏块的零平均怎样改善远场衰减。
\item\label{b2:item:18-goal-3}核对标准核的大小、光滑性及离支集表示，并辨认平方有界性这一独立假设。
\item\label{b2:item:18-goal-4}完成弱型、插值、伴随与极大截断的论证，把稠密类上的主值传给一般输入。
\item\label{b2:item:18-goal-5}用 Riesz 变换描述 Poisson 共轭结构、频率投影与椭圆二阶估计。
\end{itemize}\end{learninggoals}

对区间示性函数作 Hilbert 变换，答案是一个对数。输入有界且可积，输出却在区间端点无界，并在无穷远只按倒数衰减。另一方面，它的 Fourier 符号模长为一，平方范数完全保持。同一个算子为什么能够同时具有这些性质？要回答这个问题，必须分清两个作用：绝对值估计控制大小，正负抵消则能消去绝对值估计中最危险的部分。

\ProseParagraphBreak

\chapref{b2:ch:17}已经准备了三件工具。二进微分定理让局部平均反映函数值，极大立方体分解把大平均的位置收集起来，Marcinkiewicz 插值则把弱端点与平方有界性接成中间的强型估计。Fourier 部分沿用\secref{b2:sec:0401}的 Plancherel 延拓，归一化固定为 \(e^{-2\pi\mathrm i x\cdot\xi}\)；涉及分布时，只使用\chapref{b2:ch:05}中的测试函数配对、分布导数和 Fourier 变换。除\subsecref{b2:subsec:180502}的周期比较外，底空间均为带 Lebesgue 测度的 \(\mathbb R^d\)。

\ProseParagraphBreak

首次阅读可先沿 Hilbert 变换、好坏分解、标准核、弱型证明这一顺序走到\subsecref{b2:subsec:180403}。这样先看清一个具体核怎样进入一般定理。随后读极大截断与主值传递，再到\secref{b2:sec:1805}看 Riesz 变换的实现和应用。向量值估计与 BMO 端点解释的是同一方法怎样服务于后续章节，可以在掌握标量证明后回读。

\ProseParagraphBreak

进一步阅读可参考 Elias M. Stein 的《Singular Integrals and Differentiability Properties of Functions》\cite{Stein1970Singular}与 Loukas Grafakos 的《Classical Fourier Analysis》\cite{Grafakos2014Classical}。前者适合继续理解核抵消与微分性质的联系，后者可用于扩展奇异积分、乘子与函数空间的背景。较弱的核条件还可产生更广的奇异积分理论，阅读专著时可比较其光滑性假设与这里的标准核条件。

% !TeX root = ../../Book2.tex
\section{Hilbert 变换}
\label{b2:sec:1801}

卷积核可积时，可以先估计核的绝对值，再用 Young 不等式控制卷积。在实直线上，核 \(1/x\) 的绝对积分在原点和无穷远都发散，这条路线无法使用。不过，\(1/x\) 是奇函数，原点两侧的贡献可能抵消。Hilbert 变换提供了一个完整模型：我们先说明抵消怎样定义积分，再用 Fourier 变换证明有界性，最后辨认哪些结论还需要一般奇异积分理论。

\subsection{主值}
\label{b2:subsec:180101}

对 \(f\in\mathcal S(\mathbb R)\)，定义\term[hilbert-bianhuan]{Hilbert 变换}[Hilbert transform]
\[
 Hf(x)\coloneqq\lim_{\varepsilon\downarrow0}
 \frac1\pi\int_{|x-y|>\varepsilon}\frac{f(y)}{x-y}\,\mathrm dy.
\]
\symidx{\(Hf\)：Hilbert 变换}
这里的极限称为\term[zhuzhi-jifen]{主值积分}[principal value integral]。截断围绕观察点 \(x\) 对称，而不是围绕固定原点对称。换元 \(t=x-y\) 后，奇性移到了 \(t=0\)。在对称区间中常数项的积分为零，因而
\[
 \int_{\varepsilon<|t|<1}\frac{f(x-t)}t\,\mathrm dt
 =\int_{\varepsilon<|t|<1}\frac{f(x-t)-f(x)}t\,\mathrm dt.
\]
均值定理给出
\[
 \left|\frac{f(x-t)-f(x)}t\right|
 \leq\sup_{|z-x|\leq1}|f'(z)|\quad(0<|t|<1),
\]
所以右端在 \(\varepsilon\downarrow0\) 时有绝对收敛的极限。远处的积分则由 Schwartz 衰减控制。也可以先把正、负两侧配成一对：
\[
 Hf(x)=\frac1\pi\int_0^\infty
 \frac{f(x-t)-f(x+t)}t\,\mathrm dt.
\]
这个配对后的积分绝对收敛；原来未配对的积分在 \(t=0\) 附近一般并不绝对收敛。减去 \(f(x)\) 只在有限对称区间上进行，不能把它写成两个各自发散的全直线积分之差。

\ProseParagraphBreak

先看一个可以逐项计算的例子。取 \(f=\mathbf1_{(0,1)}\)。若 \(x\notin[0,1]\)，普通积分直接给出对数；若 \(0<x<1\)，取 \(0<\varepsilon<\min\{x,1-x\}\)，则
\[
 \begin{aligned}
 \int_0^{x-\varepsilon}\frac{\mathrm dy}{x-y}
 +\int_{x+\varepsilon}^1\frac{\mathrm dy}{x-y}
 &=\log x-\log\varepsilon+\log\varepsilon-\log(1-x)\\
 &=\log\frac{x}{1-x}.
 \end{aligned}
\]
因此，在 \(x\ne0,1\) 时，
\[
 H\mathbf1_{(0,1)}(x)=\frac1\pi\cdot\log\left|\frac{x}{x-1}\right|.
\]
在 \(x=0\) 只有右侧输入，在 \(x=1\) 只有左侧输入，两端分别发散到 \(-\infty\) 和 \(+\infty\)。这并不与区间内部的抵消冲突：主值存在是关于观察点的性质。改变输入在单个点上的取值不影响这些积分。

\subsection{Fourier 乘子}
\label{b2:subsec:180102}

为了合法使用卷积的 Fourier 公式，先同时截去原点和无穷远，令
\[
 h_{\varepsilon,R}(t)=\frac{\mathbf1_{\{\varepsilon<|t|<R\}}}{\pi t},
 \quad 0<\varepsilon<R<\infty.
\]
这个核属于 \(L^1(\mathbb R)\)。其奇性消去了 Fourier 积分中的余弦部分，故
\[
 \widehat h_{\varepsilon,R}(\xi)
 =-\frac{2\mathrm i}{\pi}\int_\varepsilon^R
       \frac{\sin(2\pi\xi t)}t\,\mathrm dt.
\]
负号来自本卷 Fourier 变换的指数 \(e^{-2\pi\mathrm i x\xi}\)。换元会把 \(2\pi|\xi|\) 同时从正弦和分母中消去，因此它不在最终乘子的幅值中留下额外因子。

\Needspace{6\baselineskip}
\begin{lemma}\label{b2:thm:18-dirichlet-integral}
对 \(0<\varepsilon<R<\infty\) 和 \(\xi\in\mathbb R\)，令
\[
 m_{\varepsilon,R}(\xi)
 =-\frac{2\mathrm i}{\pi}\int_\varepsilon^R
                 \frac{\sin(2\pi\xi t)}t\,\mathrm dt.
\]
存在与 \(\varepsilon,R,\xi\) 均无关的有限常数 \(C\)，使
\(|m_{\varepsilon,R}(\xi)|\leq C\)。对每个 \(\xi\ne0\)，当 \(\varepsilon\downarrow0\)、\(R\uparrow\infty\) 时，
\[
 m_{\varepsilon,R}(\xi)\longrightarrow
       -\mathrm i\cdot\operatorname{sgn}\xi.
\]
\end{lemma}
\begin{proof}
当 \(\xi=0\) 时乘子为零。其余情形换元 \(u=2\pi|\xi|t\)，只需考察 \(\int_a^b\sin u/u\,\mathrm du\)。在 \((0,1)\) 上，被积函数的绝对值不超过一；在 \(1\leq a<b\) 时，分部积分给出
\[
 \left|\int_a^b\frac{\sin u}u\,\mathrm du\right|
 \leq\frac1a+\frac1b+\int_a^b\frac{\mathrm du}{u^2}
 =\frac2a.
\]
这既证明了一致有界性，也证明了无穷端的收敛。

为计算全积分，固定 \(\eta>0\)，加入阻尼 \(e^{-\eta u}\)。对实参数 \(s\) 求导是合法的，因为求导后的被积函数被 \(e^{-\eta u}\) 控制，且
\[
 \frac{\mathrm d}{\mathrm ds}\int_0^\infty
 e^{-\eta u}\cdot\frac{\sin(su)}u\,\mathrm du
 =\int_0^\infty e^{-\eta u}\cdot\cos(su)\,\mathrm du
 =\frac{\eta}{\eta^2+s^2}.
\]
零参数处积分为零，所以取 \(s=1\) 得到 \(\arctan(1/\eta)\)。还需说明去掉阻尼后得到的是原来的广义积分。函数
\(w_\eta(u)=e^{-\eta u}/u\) 非负且递减，对 \(B>A>0\)，同一次分部积分给出
\[
 \left|\int_A^B w_\eta(u)\cdot\sin u\,\mathrm du\right|
 \leq w_\eta(A)+w_\eta(B)
          +\int_A^B|w_\eta'(u)|\,\mathrm du
 =2w_\eta(A)\leq\frac2A.
\]
这个尾界对 \(\eta>0\) 一致，无阻尼时也成立。先在固定的 \([0,A]\) 上令 \(\eta\downarrow0\)，再令 \(A\uparrow\infty\)，便得到
\(\int_0^\infty\sin u/u\,\mathrm du=\pi/2\)。正频率给出所述极限，负频率由正弦的奇性给出相反符号。
\end{proof}
乘子在 \(\xi=0\) 的取值不影响 \(L^2\) 算子。这是频率空间中单点的零测修改，与空间积分如何避开核的奇点是两件不同的事。

\subsection{\texorpdfstring{\(L^2\)}{L2} 有界性}
\label{b2:subsec:180103}

有限截断的核范数 \(\|h_{\varepsilon,R}\|_1\) 随截断范围扩大而增长，Young 不等式不能给出统一界。上一引理却控制了其 Fourier 乘子。Plancherel 定理恰好能利用这种控制，而且不损失抵消。

\Needspace{6\baselineskip}
\begin{theorem}\label{b2:thm:18-hilbert-l2}
在 \(\mathcal S(\mathbb R)\) 上以对称主值定义的 Hilbert 变换，具有唯一的有界线性延拓
\(H:L^2(\mathbb R;\mathbb C)\to L^2(\mathbb R;\mathbb C)\)。对每个 \(f\in L^2(\mathbb R)\)，
\[
 \widehat{Hf}(\xi)=-\mathrm i\cdot\operatorname{sgn}(\xi)\cdot\widehat f(\xi)
 \quad\text{几乎处处},\quad
 \|Hf\|_2=\|f\|_2.
\]
令 \(h_{\varepsilon,R}(t)=\mathbf1_{\{\varepsilon<|t|<R\}}/(\pi t)\)。当 \(\varepsilon\downarrow0\)、\(R\uparrow\infty\) 时，\(h_{\varepsilon,R}*f\to Hf\) 于 \(L^2\)。此外，对每个固定 \(\varepsilon>0\)，积分
\begin{equation}\label{b2:eq:18-hilbert-hard-truncation}
 H_\varepsilon f(x)\coloneqq
 \frac1\pi\int_{|t|>\varepsilon}\frac{f(x-t)}t\,\mathrm dt
\end{equation}
对每个 \(x\in\mathbb R\) 绝对收敛，并定义一个 \(L^2\) 函数；当 \(\varepsilon\downarrow0\) 时，\(H_\varepsilon f\to Hf\) 于 \(L^2\)。
\end{theorem}
\begin{proof}
先对任意 \(f\in L^2\)，用乘子
\(-\mathrm i\cdot\operatorname{sgn}\xi\) 定义一个 \(L^2\) 算子 \(H\)。由 Plancherel，\(\|Hf\|_2=\|f\|_2\)。有限截断满足
\[
 \widehat{h_{\varepsilon,R}*f}=m_{\varepsilon,R}\cdot\widehat f.
\]
上一引理的一致界和逐点极限，使
\[
 \int_{\mathbb R}
 |m_{\varepsilon,R}(\xi)+\mathrm i\cdot\operatorname{sgn}\xi|^2
 \cdot|\widehat f(\xi)|^2\,\mathrm d\xi\longrightarrow0.
\]
这里使用了以常数倍 \(|\widehat f|^2\) 为控制函数的支配收敛。因而有限截断在 \(L^2\) 中趋于 \(Hf\)。若 \(f\) 是 Schwartz 函数，前面已经证明有限截断逐点趋于主值；从 \(L^2\) 收敛抽出几乎处处收敛子列，就把两个定义识别为同一个等价类。Schwartz 函数在 \(L^2\) 中稠密，所以有界延拓唯一。

接着固定 \(\varepsilon>0\)。Cauchy--Schwarz 不等式给出
\[
 \int_{|t|>\varepsilon}\frac{|f(x-t)|}{|t|}\,\mathrm dt
 \leq\|f\|_2\cdot
       \left(\int_{|t|>\varepsilon}\frac{\mathrm dt}{t^2}\right)^{1/2}
 =\left(\frac2\varepsilon\right)^{1/2}\|f\|_2.
\]
所以硬截断积分绝对收敛。对固定的 \(\varepsilon\)，令 \(R\uparrow\infty\)，有限截断的乘子趋于
\[
 m_\varepsilon(\xi)
 =-\frac{2\mathrm i}{\pi}\int_\varepsilon^\infty
                  \frac{\sin(2\pi\xi t)}t\,\mathrm dt,
\]
并仍有共同上界。因此有限截断也在 \(L^2\) 中趋于乘子为 \(m_\varepsilon\) 的函数。另一方面，空间积分的尾部满足
\[
 |H_\varepsilon f(x)-(h_{\varepsilon,R}*f)(x)|
 \leq\frac1\pi\left(\frac2R\right)^{1/2}\|f\|_2
 \quad(R>\varepsilon),
\]
所以它逐点趋于 \(H_\varepsilon f\)。再次用几乎处处收敛子列识别极限，便知 \(H_\varepsilon f\in L^2\)，其乘子正是 \(m_\varepsilon\)。最后令 \(\varepsilon\downarrow0\)，乘子几乎处处趋于 \(-\mathrm i\cdot\operatorname{sgn}\xi\)，同样的支配收敛给出 \(L^2\) 收敛。
\end{proof}
上述结论区分了三个层次：Schwartz 输入的逐点主值、任意 \(L^2\) 输入的乘子定义，以及硬截断到该定义的 \(L^2\) 收敛。最后一项只保证存在几乎处处收敛的截断子列；要证明所有 \(\varepsilon\downarrow0\) 的逐点极限存在，还需本章后面的极大截断估计。

\ProseParagraphBreak

从乘子还能立即读出
\[
 H^2f=-f,\quad H^*=-H\quad(f\in L^2(\mathbb R)).
\]
第一式来自 \((-\mathrm i\cdot\operatorname{sgn}\xi)^2=-1\) 几乎处处，第二式来自乘子的复共轭等于其负值。乘子是纯虚数并不表示实输入变成虚函数：空间截断的核是实的，对实 Schwartz 函数输出为实，再由 \(L^2\) 逼近可知实 \(L^2\) 输入的输出也为实。Hilbert 变换不是保正算子，其有界性来自抵消，不能靠正核估计解释。

\ProseParagraphBreak

另一个可算模型是 \(P(x)=1/(1+x^2)\)。它属于 \(L^1\cap L^2\)，并且
\[
 HP(x)=\frac{x}{1+x^2}\quad\text{几乎处处}.
\]
可以不用预先知道 \(P\) 的 Fourier 变换来验证。对 \(y\ne x\) 作部分分式分解：
\[
 \frac1{(1+y^2)(x-y)}
 =\frac1{1+x^2}
   \left(\frac1{x-y}+\frac{y+x}{1+y^2}\right).
\]
先在有限区间 \([-R,R]\) 内排去 \((x-\varepsilon,x+\varepsilon)\)，其中 \(R>|x|+1\)，然后令 \(\varepsilon\downarrow0\)。第一项的主值为
\(\log|x+R|-\log|x-R|\)，在 \(R\uparrow\infty\) 时趋于零；第二项中的奇函数 \(y/(1+y^2)\) 积分为零，而 \(x/(1+y^2)\) 的积分趋于 \(\pi x\)。乘回 \(1/\bigl(\pi(1+x^2)\bigr)\) 就得到公式。这里先在有限区间上拆分，避免把两个不绝对可积的全直线积分擅自相加。原积分在无穷远绝对收敛，所得逐点主值再由定理中的 \(L^2\) 收敛识别为 \(HP\)。与示性函数相比，这个平滑输入没有产生对数尖峰，输出却仍有 \(1/x\) 型远尾。

\subsection{\texorpdfstring{\(L^p\)}{Lp} 有界性}
\label{b2:subsec:180104}

核 \(K(x,y)=1/(\pi(x-y))\) 在 \(x\ne y\) 时满足大小条件，而且当 \(|y-y'|\leq|x-y|/2\) 时，
\[
 \begin{aligned}
 |K(x,y)|&=\frac1{\pi|x-y|},\\
 |K(x,y)-K(x,y')|
 &=\frac{|y-y'|}{\pi|x-y|\cdot|x-y'|}
 \leq\frac2\pi\frac{|y-y'|}{|x-y|^2}.
 \end{aligned}
\]
交换变量可得同样的第一变量差分界。若 \(f\in L^2(\mathbb R)\) 有紧支集 \(F\)，则 \(f\in L^1\)，且在 \(x\notin F\) 时普通核积分绝对收敛。对这样的 \(x\)，当 \(\varepsilon<\operatorname{dist}(x,F)\) 时，\(H_\varepsilon f(x)\) 已等于该普通积分。将硬截断的 \(L^2\) 收敛与一个几乎处处收敛子列结合，便得到
\[
 Hf(x)=\frac1\pi\int_F\frac{f(y)}{x-y}\,\mathrm dy
 \quad\text{对几乎处处的 }x\notin F.
\]
因此我们已经具备一般 Calderón--Zygmund 定理所需的三项输入：\(L^2\) 有界性、远离支集的核表达，以及两个变量的核差分控制。

\ProseParagraphBreak

\Cref{b2:thm:18-cz-lp}将在\secref{b2:sec:1804}由这些条件推出
\[
 \|Hf\|_p\leq C_p\|f\|_p\quad(1<p<\infty).
\]
这个结论先在 \(L^p\cap L^2\) 上成立，再由稠密性定义 \(L^p\) 上的有界延拓；在交集上，它与已有的 \(L^2\) 算子相同。对一般 \(L^p\) 输入，等价类意义的延拓和逐点主值表示仍需区分。该节的极大截断估计与\eqref{b2:eq:18-principal-value-transfer}的逼近论证，会把它们识别起来；本节证明的硬截断兼容性提供了其中所需的 \(L^2\) 极限。

\ProseParagraphBreak

两个端点的障碍可以从 \(\mathbf1_{(0,1)}\) 同时看出。在 \(x=0,1\) 附近，已算得的对数函数本性无界，故不存在 \(L^\infty\to L^\infty\) 的一致有界估计。在 \(|x|\to\infty\) 时，
\[
 H\mathbf1_{(0,1)}(x)=\frac1{\pi x}+O(|x|^{-2}),
\]
所以输出不属于 \(L^1\)，强型 \((1,1)\) 也失败。两种障碍的位置不同：前者发生在输入跳跃的边缘，后者发生在远场。输入同时可积且有界并不能消除它们。对于这个模型，\(1<p<\infty\) 恰好使对数奇性和 \(1/x\) 尾部的 \(p\) 次方都可积；弱型 \((1,1)\) 则允许 \(1/x\) 这样的尾部。

\ProseParagraphBreak

零积分会改善远场，却未必改善局部奇性。具体地，若 \(f\in L^1(\mathbb R)\) 支撑于 \([-R_0,R_0]\)，其中 \(R_0>0\)，且 \(\int f=0\)，则在 \(|x|>2R_0\) 的普通核积分中可以减去常数项：
\[
 \begin{aligned}
 \frac1\pi\int\frac{f(y)}{x-y}\,\mathrm dy
 &=\frac1\pi\int f(y)\cdot
       \left(\frac1{x-y}-\frac1x\right)\mathrm dy,\\
 \left|\frac1\pi\int\frac{f(y)}{x-y}\,\mathrm dy\right|
 &\leq\frac2{\pi|x|^2}\int|y|\cdot|f(y)|\,\mathrm dy.
 \end{aligned}
\]
若还要求 \(f\in L^2\)，这就是 \(Hf\) 在支集外的估计；对于一般可积输入，此处只讨论普通核积分。例如
\(f=\mathbf1_{(0,1)}-\mathbf1_{(1,2)}\) 的 Hilbert 变换在远处是 \(O(|x|^{-2})\)，但区间端点的对数奇性依然存在。下一节将把这种零平均带来的尾部改善推广到多维坏块。局部平均振荡所对应的无穷端点空间，则见\appref{b2:app:C}中的 BMO 理论。

% !TeX root = ../../Book2.tex
\section{Calderón–Zygmund 分解}
\label{b2:sec:1802}

一个 \(L^1\) 函数的总质量虽有限，高度却可能任意大。若直接用已有的 \(L^2\) 算子界处理它，高尖峰就会造成困难。Calderón--Zygmund 分解的做法是先选出局部平均过大的区域，把其中的函数压到受控高度；被削下来的部分仍保留在原区域内，但总积分被安排为零。这两部分分别适合不同的估计：高度控制使好部分进入 \(L^2\)，零平均使坏部分经过奇异积分后在远处衰减得更快。

\subsection{好/坏分解}
\label{b2:subsec:180201}

沿用上一章的标准半开二进立方体网格 \(\mathcal D\)。对局部可积函数 \(f\)，记
\[
 f_Q=\frac1{|Q|}\int_Q f(y)\,\mathrm dy.
\]
选择坏区域时要看 \(|f|\) 的平均，避免正负抵消掩盖尖峰；替换函数值时则取 \(f\) 本身的平均。两种平均的作用不同。下面给出\term[calderon-zygmund-fenjie]{Calderón--Zygmund 分解}[Calderón--Zygmund decomposition]的精确形式。

\Needspace{6\baselineskip}
\begin{theorem}[Calderón--Zygmund 分解]\label{b2:thm:18-cz-decomposition}
设 \(d\geq1\) 为整数，\(f\in L^1(\mathbb R^d;\mathbb C)\)，\(\alpha>0\)。存在至多可数个互不相交的半开二进立方体 \(Q_j\)，令 \(E=\bigcup_jQ_j\)，并定义
\[
 g(x)=\begin{cases}
 f(x),&x\notin E,\\
 f_{Q_j},&x\in Q_j,
 \end{cases}
 \quad b_j=(f-f_{Q_j})\cdot\mathbf1_{Q_j}.
\]
则 \(f=g+\sum_jb_j\) 几乎处处，级数在 \(L^1\) 中收敛，并有
\[
 \begin{gathered}
 \alpha<\frac1{|Q_j|}\int_{Q_j}|f|\leq2^d\alpha,
 \quad \sum_j|Q_j|\leq\alpha^{-1}\|f\|_1,\\
 \|g\|_\infty\leq2^d\alpha,
 \quad\|g\|_1\leq\|f\|_1,
 \quad\int b_j=0,
 \quad\sum_j\|b_j\|_1\leq2\|f\|_1.
 \end{gathered}
\]
若另外 \(f\in L^2(\mathbb R^d)\)，则 \(g\in L^2\)，\(\sum_jb_j\) 也在 \(L^2\) 中收敛，并且
\[
 \|g\|_2^2\leq2^d\alpha\|f\|_1,
 \quad \|f\|_2^2=\|g\|_2^2+\sum_j\|b_j\|_2^2.
\]
\end{theorem}
\begin{proof}
用\cref{b2:thm:17-dyadic-level}选择满足
\(|Q|^{-1}\int_Q|f|>\alpha\) 的极大二进立方体。回忆该选择的理由：二进立方体相交时必有一个包含另一个；沿祖先向上，平均值不超过 \(\|f\|_1/|Q|\)，所有足够大的祖先都不再超水平。在仍超水平的有限多个祖先中取最大的一个，就得到所需的极大者。这些极大者互不相交，其二进父立方体 \(\widetilde Q_j\) 的平均不超过 \(\alpha\)。因此
\[
 \int_{Q_j}|f|>\alpha|Q_j|,
 \quad \int_{Q_j}|f|\leq\int_{\widetilde Q_j}|f|
 \leq\alpha|\widetilde Q_j|=2^d\alpha|Q_j|.
\]
对第一式求和得到坏区域的测度界。对 \(E\) 外的点，所有二进平均都不超过 \(\alpha\)；由上一章已经证明的二进微分结论，\(|f|\leq\alpha\) 在该处几乎处处成立。

按陈述定义 \(g\)。在 \(E\) 内，
\[
 |f_{Q_j}|\leq\frac1{|Q_j|}\int_{Q_j}|f|\leq2^d\alpha,
 \quad \int_{Q_j}|g|=|Q_j|\cdot|f_{Q_j}|
                          \leq\int_{Q_j}|f|.
\]
结合 \(E\) 外的结论，得到 \(g\) 的两个范数界。每个坏块满足
\[
 \int b_j=\int_{Q_j}f-|Q_j|\cdot f_{Q_j}=0,
 \quad \|b_j\|_1\leq\int_{Q_j}|f|+|Q_j|\cdot|f_{Q_j}|
                         \leq2\int_{Q_j}|f|.
\]
故 \(\sum_j\|b_j\|_1\leq2\|f\|_1\)，级数在 \(L^1\) 中绝对收敛。不交性还说明每个点至多落在一个坏块中，因而逐点表达与 \(L^1\) 极限相同，给出分解等式。绝对值估计和平均值的线性使这个分解同样适用于复值函数。

若 \(f\in L^2\)，则
\[
 \|g\|_2^2\leq\|g\|_\infty\cdot\|g\|_1
                         \leq2^d\alpha\|f\|_1.
\]
又 \(b=f-g\in L^2\)，且 \(b_j=b\cdot\mathbf1_{Q_j}\)。不交性给出
\(\sum_j\|b_j\|_2^2=\|b\|_2^2\)，所以坏块级数的尾部 \(L^2\) 范数趋于零。更精确地，在每个 \(Q_j\) 上，常数函数 \(f_{Q_j}\) 与 \(b_j\) 正交，因为
\[
 \int_{Q_j}b_j\overline{f_{Q_j}}
 =\overline{f_{Q_j}}\int_{Q_j}b_j=0.
\]
在各块上使用平方范数的勾股恒等式，再与 \(E\) 外的积分合并，便得到陈述中的精确 \(L^2\) 恒等式。
\end{proof}
为什么必须用平均值替换坏块，而不是简单地将它置零？若在 \(Q\) 上用常数 \(a\) 替换，剩余部分的积分为
\(\int_Qf-a|Q|\)。使其为零的唯一选择正是 \(a=f_Q\)。因此，好部分保留了每块的总质量，坏部分只保留块内相对于平均值的变化。若改用 \((|f|)_Q\)，一般就失去了这一抵消。好函数的名称仅指高度受控，并不表示它连续或光滑。

\ProseParagraphBreak

以一维函数 \(f=4\cdot\mathbf1_{[0,1/4)}\) 和高度 \(\alpha=1\) 为例。区间 \([0,1/2)\) 上的平均为 \(2\)，其二进父区间 \([0,1)\) 上的平均恰为 \(1\)；所有更大的祖先平均更小。因此唯一的极大坏区间是 \(Q=[0,1/2)\)，分解为
\[
 g=2\cdot\mathbf1_{[0,1/2)},
 \quad b=2\cdot\mathbf1_{[0,1/4)}
             -2\cdot\mathbf1_{[1/4,1/2)}.
\]
\Cref{b2:fig:18-cz-step-model}显示，削去的正面积由右半块的负面积抵消。这里
\(\|f\|_1=\|g\|_1=\|b\|_1=1\)，而平方范数分别为 \(4,2,2\)，可直接验证定理中的恒等式。坏块并没有变得处处小；它获得的是零平均。

\begin{figure}[htbp]
 \centering
 \begin{tikzpicture}[x=3.6cm,y=0.38cm,font=\small]
  \begin{scope}
   \fill[AnalysisBlueLight] (0,0) rectangle (0.25,4);
   \draw[AnalysisBlue,thick] (0,0)--(0,4)--(0.25,4)--(0.25,0)--(0.72,0);
   \draw[->,AnalysisGray] (-0.03,0)--(0.77,0);
   \draw[->,AnalysisGray] (0,-2.3)--(0,4.7);
   \node[above] at (0.35,5) {原函数 \(f\)};
   \node[left] at (0,4) {\(4\)};
   \foreach \x/\lab in {0/0,0.25/{1/4},0.5/{1/2}} {
    \draw (\x,0.12)--(\x,-0.12);
    \node at (\x,-2.8) {\(\lab\)};
   }
  \end{scope}
  \begin{scope}[xshift=3.8cm]
   \fill[AnalysisBlueLight] (0,0) rectangle (0.5,2);
   \draw[AnalysisBlue,thick] (0,0)--(0,2)--(0.5,2)--(0.5,0)--(0.72,0);
   \draw[->,AnalysisGray] (-0.03,0)--(0.77,0);
   \draw[->,AnalysisGray] (0,-2.3)--(0,4.7);
   \node[above] at (0.35,5) {好部分 \(g\)};
   \node[left] at (0,2) {\(2\)};
   \foreach \x/\lab in {0/0,0.25/{1/4},0.5/{1/2}} {
    \draw (\x,0.12)--(\x,-0.12);
    \node at (\x,-2.8) {\(\lab\)};
   }
  \end{scope}
  \begin{scope}[xshift=7.6cm]
   \fill[AnalysisBlueLight] (0,0) rectangle (0.25,2);
   \fill[AnalysisGrayLight] (0.25,0) rectangle (0.5,-2);
   \draw[AnalysisBlue,thick] (0,0)--(0,2)--(0.25,2)--(0.25,-2)--(0.5,-2)--(0.5,0)--(0.72,0);
   \draw[->,AnalysisGray] (-0.03,0)--(0.77,0);
   \draw[->,AnalysisGray] (0,-2.3)--(0,4.7);
   \node[above] at (0.35,5) {坏部分 \(b\)};
   \node[left] at (0,2) {\(2\)};
   \node[left] at (0,-2) {\(-2\)};
   \foreach \x/\lab in {0/0,0.25/{1/4},0.5/{1/2}} {
    \draw (\x,0.12)--(\x,-0.12);
    \node at (\x,-2.8) {\(\lab\)};
   }
  \end{scope}
 \end{tikzpicture}
 \caption[阶梯函数的好坏分解]{在标准二进网格中，对 \(f=4\cdot\mathbf1_{[0,1/4)}\) 取 \(\alpha=1\)。坏区间为 \([0,1/2)\)，以平均值 \(2\) 替换后得到 \(g\)；剩余的 \(b\) 有等量的正、负面积，故 \(\int b=0\)。图中只画出非零部分附近，三者在所示区间外均为零。}
 \label{b2:fig:18-cz-step-model}
\end{figure}

\subsection{零平均与远场估计}
\label{b2:subsec:180202}

设观察点 \(x\) 离坏块很远。此时 \(K(x,y)\) 随块内变量 \(y\) 的变化较小，可以先把它近似为块中心处的常数 \(K(x,c)\)。这个常数乘上坏块的总积分，恰好为零；真正需要估计的仅是核的变化量。下面的引理把这一思路量化，并保留核差分常数的依赖。

\Needspace{6\baselineskip}
\begin{lemma}\label{b2:thm:18-cancellation-tail}
设 \(d\geq1\) 为整数，\(Q\subset\mathbb R^d\) 是边长 \(\ell>0\)、中心 \(c\) 的轴平行立方体。设 \(b\in L^1(\mathbb R^d;\mathbb C)\) 在 \(Q\) 外几乎处处为零，且 \(\int b=0\)。设复值 Borel 可测核 \(K\) 在 \(x\ne y\) 上取有限值，并存在 \(\delta>0\) 和 \(A_K\geq0\)，使
\[
 |K(x,y)-K(x,c)|
 \leq A_K\frac{|y-c|^\delta}{|x-c|^{d+\delta}}
 \quad\text{当 }x\ne c,\ |x-c|\geq2|y-c|.
\]
令 \(Q^*=8\sqrt d\,Q\) 表示同中心、边长为 \(8\sqrt d\,\ell\) 的立方体。则对 \(x\notin Q^*\)，普通积分 \(\int_QK(x,y)\cdot b(y)\,\mathrm dy\) 绝对收敛，而且
\[
 \int_{\mathbb R^d\setminus Q^*}
 \left|\int_QK(x,y)\cdot b(y)\,\mathrm dy\right|\,\mathrm dx
 \leq C_{d,\delta}A_K\|b\|_1,
\]
其中 \(C_{d,\delta}\) 只依赖 \(d,\delta\)，与 \(Q,b,K\) 无关。
\end{lemma}
\begin{proof}
若 \(x\notin Q^*\)、\(y\in Q\)，则
\[
 |x-c|\geq4\sqrt d\,\ell,
 \quad |y-c|\leq\frac{\sqrt d\,\ell}{2}.
\]
因此差分假设适用。固定这样的 \(x\)，\(K(x,c)\) 有限，核差分在 \(Q\) 上有界，故 \(K(x,\cdot)b\) 绝对可积。零平均给出合法的恒等式
\[
 \int_QK(x,y)\cdot b(y)\,\mathrm dy
 =\int_Q\bigl(K(x,y)-K(x,c)\bigr)\cdot b(y)\,\mathrm dy.
\]
取绝对值，并对非负积分使用 Tonelli，得到所求左端不超过
\[
 A_K\left(\frac{\sqrt d\,\ell}{2}\right)^\delta\|b\|_1
 \int_{|x-c|\geq4\sqrt d\,\ell}|x-c|^{-d-\delta}\,\mathrm dx.
\]
记单位球面的面积为 \(\sigma_{d-1}\)；在 \(d=1\) 时取 \(\sigma_0=2\)。极坐标计算最后一项为
\[
 \sigma_{d-1}\int_{4\sqrt d\,\ell}^\infty r^{-1-\delta}\,\mathrm dr
 =\frac{\sigma_{d-1}}\delta(4\sqrt d\,\ell)^{-\delta}.
\]
这正好消去前面的 \(\ell^\delta\)，给出常数
\(C_{d,\delta}=\sigma_{d-1}/(\delta8^\delta)\) 的一个可用选择。
\end{proof}
这里有两个值得分开的量：坏块的 \(L^1\) 大小决定估计的规模，核变化的快慢决定远场能否积分。若只用大小界 \(|K(x,y)|\lesssim|x-c|^{-d}\)，径向积分会出现 \(\int^\infty\mathrm dr/r\)，无法得到有限上界。减去常数项后，多出的 \(|x-c|^{-\delta}\) 才使它收敛。\(\delta>0\) 正是这一计算所需的条件；尾界对 \(A_K\) 的线性依赖，也使后文能够分别跟踪算子的 \(L^2\) 范数和核的光滑常数。

\ProseParagraphBreak

扩大立方体并不改变坏函数，也不使它产生额外抵消。它只将靠近坏块、不能使用核差分的观察点暂时排除。扩大倍数的具体数值并不重要，只要保证观察距离比块的直径大一个固定倍数；相应的测度代价只依赖维数。

\subsection{弱型估计的分解策略}
\label{b2:subsec:180203}

现在把两个估计合在一起。为清楚看出每项假设的用途，暂设 \(T\) 是 \(L^2(\mathbb R^d)\) 上的有界复线性算子，记 \(B=\|T\|_{2\to2}\)。还假设它在每个紧支集 \(L^2\) 输入的支集外，几乎处处由核 \(K\) 的普通积分表示；核对每个立方体中心都满足上一引理的差分条件，且使用共同的 \(\delta,A_K\)。下一节会把这些条件组织成标准算子定义。

\ProseParagraphBreak

给定 \(f\in L^1\cap L^2\) 和输出阈值 \(t>0\)，在高度 \(\alpha=t\) 分解 \(f=g+b\)。好部分使用 \(L^2\) Chebyshev 不等式：
\[
 \begin{aligned}
 |\{|Tg|>t/2\}|
 &\leq\frac4{t^2}\|Tg\|_2^2
 \leq\frac{4B^2}{t^2}\|g\|_2^2\\
 &\leq\frac{2^{d+2}B^2}{t}\|f\|_1.
 \end{aligned}
\]
好部分的有界性把它的平方积分控制在 \(t\|f\|_1\) 的量级，这正是分解带来的收益。

\ProseParagraphBreak

对坏部分，先丢弃扩大坏区域
\[
 E^*=\bigcup_jQ_j^*,
 \quad |E^*|\leq(8\sqrt d)^d\sum_j|Q_j|
          \leq\frac{(8\sqrt d)^d}{t}\|f\|_1.
\]
原来的 \(Q_j\) 互不相交，扩大的 \(Q_j^*\) 可能重叠；这里用的是并集测度不超过测度之和，不需要扩大后仍然不交。

\ProseParagraphBreak

在 \((E^*)^c\) 上使用坏块的零平均。先取有限和 \(b^{(N)}=\sum_{j=1}^Nb_j\)。每个 \(b_j\) 属于 \(L^2\) 且支撑于一个有界立方体，因而支集外核表示与远场引理给出
\[
 \begin{aligned}
 \int_{(E^*)^c}|Tb^{(N)}(x)|\,\mathrm dx
 &\leq\sum_{j=1}^N\int_{(Q_j^*)^c}|Tb_j(x)|\,\mathrm dx\\
 &\leq C_{d,\delta}A_K\sum_{j=1}^N\|b_j\|_1
 \leq2C_{d,\delta}A_K\|f\|_1.
 \end{aligned}
\]
若坏块有无穷多个，定理保证 \(b^{(N)}\to b\) 于 \(L^2\)，于是 \(Tb^{(N)}\to Tb\) 于 \(L^2\)。抽出几乎处处收敛的子列，再对 \((E^*)^c\) 使用 Fatou 引理，得到
\[
 \int_{(E^*)^c}|Tb(x)|\,\mathrm dx
 \leq2C_{d,\delta}A_K\|f\|_1.
\]
这不要求 \(\sum_jTb_j\) 在全空间的 \(L^1\) 中收敛。我们得到的仅是远离坏区域后的积分控制，正好足够用于 Chebyshev 不等式：
\[
 |\{x\notin E^*:|Tb(x)|>t/2\}|
 \leq\frac{4C_{d,\delta}A_K}{t}\|f\|_1.
\]

由 \(Tf=Tg+Tb\)，其 \(t\) 超水平集包含在好部分的 \(t/2\) 超水平集、扩大坏区域，以及坏部分在其外的 \(t/2\) 超水平集三者之并中。合并以上估计即得
\[
 |\{|Tf|>t\}|
 \leq\frac{C_{d,\delta}(1+B^2+A_K)}t\|f\|_1
 \quad(f\in L^1\cap L^2,\ t>0),
\]
其中最后的 \(C_{d,\delta}\) 可以增大以吸收各项维数常数。好部分靠 \(L^2\) 有界性，坏区域靠测度控制，坏部分的远场靠零平均与核差分；三条估计承担着不同任务。由稠密性把算子延拓到全部 \(L^1\)，以及由弱端点推出 \(L^p\) 强型的步骤，将在\secref{b2:sec:1804}完成。

% !TeX root = ../../Book2.tex
\section{Calderón–Zygmund 核}
\label{b2:sec:1803}

Hilbert 核与 Riesz 核都有两种不同的性质：核本身在对角线附近不可积，移开一个变量后得到的核之差却衰减得更快。前一种性质决定积分为什么需要截断，后一种性质使坏块的零平均能够发挥作用。本节把这些条件分开写出，再说明怎样把核与一个已经有平方估计的算子联系起来。

\subsection{大小条件}
\label{b2:subsec:180301}

设整数 \(d\geq1\)，\(K:(\mathbb R^d\times\mathbb R^d)\setminus\{x=y\}\to\mathbb C\) 可测。若存在有限常数 \(C_0\geq0\)，使对每个 \(x\ne y\) 都有
\[
 |K(x,y)|\leq C_0|x-y|^{-d},
\]
则称 \(K\) 满足标准大小条件。半径为 \(r\) 的环带有 \(r^d\) 量级的体积，所以这个上界在每一倍增环带上只给出常数量级的绝对积分；对角线附近有无穷多个这样的环带，不能直接求和。

\ProseParagraphBreak

若 \(f\in L^2\) 支撑在紧集 \(F\) 内，则 \(f\in L^1\)。在与 \(F\) 有正距离的点 \(x\) 处，大小条件给出
\[
 \int_F|K(x,y)|\cdot|f(y)|\,\mathrm dy
 \leq C_0\cdot\operatorname{dist}(x,F)^{-d}\cdot\|f\|_1<\infty.
\]
因此支集外的核积分具有普通绝对积分的意义。这里同时用了输入的有限支集和观察点到支集的距离；支集内部的积分尚未由此定义。

\ProseParagraphBreak

卷积型核写成 \(K(x,y)=k(x-y)\)。例如取有界球面函数 \(\Omega\)，令 \(k(z)=\Omega(z/|z|)\cdot|z|^{-d}\)。若光滑输入在 \(x\) 附近恒等于一，则局部截断积分含有
\[
 \int_{\epsilon<|z|<1}k(z)\,\mathrm dz
 =\left(\int_{S^{d-1}}\Omega\,\mathrm d\sigma\right)\cdot
    \log\frac1\epsilon.
\]
球面平均不为零时，这一项发散。齐次次数 \(-d\) 给出大小条件，却没有提供抵消；球面上的奇函数平均为零，Hilbert 核就是一维的例子。

\subsection{光滑性条件}
\label{b2:subsec:180302}

坏块的零平均使 \(K(x,y)\) 被替换为 \(K(x,y)-K(x,c_Q)\)。因此需要控制的是移动输入变量 \(y\) 后的变化。对上述离开对角线的核，固定 \(0<\delta\leq1\)。标准光滑条件要求存在有限常数 \(C_1\geq0\)，使
\[
 \begin{aligned}
 |K(x,y)-K(x,y')|&\leq C_1\frac{|y-y'|^\delta}{|x-y|^{d+\delta}}
 &&\left(|y-y'|\leq\frac{|x-y|}2\right),\\
 |K(x,y)-K(x',y)|&\leq C_1\frac{|x-x'|^\delta}{|x-y|^{d+\delta}}
 &&\left(|x-x'|\leq\frac{|x-y|}2\right).
 \end{aligned}
\]
所有出现的原始点均满足 \(x\ne y\)。距离限制使新点仍离开对角线，且新旧距离可互相比较。第一行控制原算子的坏块；取伴随后输入、输出互换，第二行才成为伴随坏块所需的条件。

\ProseParagraphBreak

核的梯度界是一个常用的检查方法。若 \(K\) 在对角线外为 \(C^1\)，且
\[
 |\nabla_xK(x,y)|+|\nabla_yK(x,y)|\leq C|x-y|^{-d-1},
\]
则沿连接新旧点的线段积分导数即可取 \(\delta=1\)。例如 \(|y-y'|\leq|x-y|/2\) 时，线段上每点到 \(x\) 的距离至少为 \(|x-y|/2\)，所以积分中的梯度有同一个上界。

\ProseParagraphBreak

弱型证明实际上只用到了输入变量差的积分控制。设 \(h=|y-y'|>0\)，上述 Hölder 界给出
\[
 \begin{aligned}
 \int_{|x-y|>2h}|K(x,y)-K(x,y')|\,\mathrm dx
 &\leq C_1|S^{d-1}|h^\delta
            \int_{2h}^\infty r^{-1-\delta}\,\mathrm dr\\
 &=\frac{C_1|S^{d-1}|}{\delta\,2^\delta}.
 \end{aligned}
\]
\(h=0\) 时差为零。因而标准光滑性蕴含积分 Hörmander 条件
\[
 \sup_{y,y'}\int_{|x-y|>2|y-y'|}
       |K(x,y)-K(x,y')|\,\mathrm dx<\infty.
\]
它直接对应上一节对坏块远场的积分；点态条件还将用于截断比较和平均振荡估计。

\ProseParagraphBreak

也可以把幂次换成较慢的连续模量。设 \(\omega:[0,1]\to[0,\infty)\) 单调不减、\(\omega(0)=0\)，且
\[
 |K(x,y)-K(x,y')|\leq C|x-y|^{-d}\cdot
 \omega\!\left(\frac{|y-y'|}{|x-y|}\right)
 \quad\left(|y-y'|\leq\frac{|x-y|}2\right).
\]
径向积分换元 \(s=h/r\)，得到
\[
 \int_{|x-y|>2h}|K(x,y)-K(x,y')|\,\mathrm dx
 \leq C|S^{d-1}|\int_0^{1/2}\omega(s)\frac{\mathrm ds}s.
\]
因此 \(\int_0^1\omega(s)\,\mathrm ds/s<\infty\) 也足以控制坏块尾部。若已经给定有界 \(L^2\) 算子及支集外表示，这个积分条件可代替弱型证明中的输入变量 Hölder 条件；伴随也满足对应条件时，再由插值和对偶得到全部有限的中间指数。

\ProseParagraphBreak

取 \(\omega(s)=s^\delta\) 就回到幂次情形。取 \(\omega(s)=\bigl(\log(e/s)\bigr)^{-2}\) 时，换元 \(v=\log(e/s)\) 得到可积尾部 \(\int^\infty v^{-2}\,\mathrm dv\)；改成 \(\bigl(\log(e/s)\bigr)^{-1}\) 则发散。这些计算检验了模量所控制的尾积分；具体算子的附加抵消可能进一步改善估计。按二进环带估计，同一个条件写成
\(\sum_{j\geq1}\omega(2^{-j})<\infty\)：每个环带的体积消去核的 \(|x-y|^{-d}\)，留下的正是模量。单调性给出
\[
 (\log2)\cdot\omega(2^{-j-1})
 \leq\int_{2^{-j-1}}^{2^{-j}}\omega(s)\frac{\mathrm ds}s
 \leq(\log2)\cdot\omega(2^{-j}),
\]
从而级数与积分具有相同的收敛性。

\subsection{奇异积分算子}
\label{b2:subsec:180303}

核的条件只描述离开对角线的作用。还须指定一个算子，并独立证明其平方有界性。
\Needspace{6\baselineskip}
\begin{definition}\label{b2:thm:18-cz-operator-definition}
设整数 \(d\geq1\)，\(T:L^2(\mathbb R^d;\mathbb C)\to L^2(\mathbb R^d;\mathbb C)\) 为有界复线性算子。设可测核 \(K:(\mathbb R^d\times\mathbb R^d)\setminus\{x=y\}\to\mathbb C\)，并存在有限常数 \(C_0,C_1\geq0\)、\(0<\delta\leq1\)，使
\[
 |K(x,y)|\leq C_0|x-y|^{-d}
 \quad(x\ne y),
\]
\[
 \begin{aligned}
 &|K(x,y)-K(x',y)|+|K(y,x)-K(y,x')|\\
 &\quad\leq C_1\frac{|x-x'|^\delta}{|x-y|^{d+\delta}}
 \quad\left(x\ne y,\ |x-x'|\leq\frac{|x-y|}2\right).
 \end{aligned}
\]
若对每个支撑在某个紧集 \(F\) 内的 \(f\in L^2(\mathbb R^d)\)，在 \(\mathbb R^d\setminus F\) 上几乎处处有
\[
 Tf(x)=\int_FK(x,y)\cdot f(y)\,\mathrm dy,
\]
则称 \(T\) 为具有核 \(K\) 的\term[calderon-zygmund-suanzi]{Calderón--Zygmund 算子}[Calderón--Zygmund operator]。上述核条件称为标准核条件；第二个显示式同时包含两个变量的光滑性。
\end{definition}

平方有界性没有被核条件包含。例如零核既对应零算子，也对应任意有界乘法算子 \(M_af=af\)、\(a\in L^\infty\)：输入支集外两种输出都为零。后面的有界性理论适用于这些算子，但若只知道支集外的核，就不能判定给定算子是否等于硬截断的主值。

\ProseParagraphBreak

伴随所需的核可以直接从内积求出。取紧支集 \(f,g\in L^2\)，并假设两个支集有正距离。它们均可积，且
\[
 \int\!\int |K(x,y)|\cdot|f(y)|\cdot|g(x)|\,\mathrm dy\,\mathrm dx
 \leq C\|f\|_1\cdot\|g\|_1<\infty.
\]
这里 \(C\) 还依赖两个支集的距离。Fubini 与 \(\langle Tf,g\rangle=\langle f,T^*g\rangle\) 给出
\[
 \int f(y)\cdot\overline{T^*g(y)}\,\mathrm dy
 =\int f(y)\cdot
       \overline{\int\overline{K(x,y)}\cdot g(x)\,\mathrm dx}\,\mathrm dy.
\]
固定 \(g\)，让 \(f\) 遍历任一离开 \(\operatorname{supp}g\) 的有界紧区域上的 \(L^2\) 函数，即可识别两边的局部 \(L^2\) 函数。再以可数个这样的区域穷竭支集外，得到伴随的核
\[
 K^*(x,y)=\overline{K(y,x)}.
\]
它的大小与两个变量光滑常数仍可取 \(C_0,C_1\)，且 \(\|T^*\|_{2\to2}=\|T\|_{2\to2}\)。输入变量的光滑性在伴随后变成输出变量的光滑性，两者不能只保留其中一条。

\ProseParagraphBreak

下面从一个具体核构造算子。这里的目标是先取得独立的 \(L^2\) 界，再使用下一节的一般有界性定理。
\Needspace{6\baselineskip}
\begin{theorem}\label{b2:thm:18-homogeneous-kernel}
设整数 \(d\geq2\)，\(\Omega\in C^1(S^{d-1};\mathbb C)\)，且 \(\int_{S^{d-1}}\Omega\,\mathrm d\sigma=0\)，其中 \(\sigma\) 为球面面积测度。则对每个 \(f\in\mathcal S(\mathbb R^d)\)、每个 \(x\in\mathbb R^d\)，主值
\[
 T_\Omega f(x)=\lim_{\epsilon\downarrow0}
 \int_{|y|>\epsilon}\frac{\Omega(y/|y|)}{|y|^d}\cdot f(x-y)\,\mathrm dy
\]
存在。它给出兼容的有界线性算子 \(T_\Omega:L^p(\mathbb R^d)\to L^p(\mathbb R^d)\)、\(1<p<\infty\)，且
\[
 \|T_\Omega f\|_p\leq C_{d,p}\|\Omega\|_{C^1(S^{d-1})}\cdot\|f\|_p.
\]
其 \(L^2\) Fourier 乘子在 \(\xi\ne0\) 时为
\[
 m(\xi)=-\int_{S^{d-1}}\Omega(\theta)\cdot
 \left(\log|2\pi\xi\cdot\theta|
       +\frac{\pi\mathrm i}{2}\cdot\operatorname{sgn}(\xi\cdot\theta)\right)\mathrm d\sigma(\theta).
\]
对数奇性按球面上的可积函数解释，\(m\) 为零次齐次；其在 \(\xi=0\) 的取值可任意指定。
\end{theorem}
\begin{proof}
先固定 \(f\in\mathcal S\)、\(x\)。在 \(0<|y|<1\) 的积分中减去 \(f(x)\)，球面平均零使常数部分消失。由微积分基本定理，
\[
 |f(x-y)-f(x)|\leq |y|\sup_{|z-x|\leq1}|\nabla f(z)|,
\]
所以近原点的绝对积分由 \(C_x\int_0^1\mathrm dr\) 控制。远处 \(|x-y|\) 随 \(|y|\) 增大，Schwartz 衰减使径向尾部可积。因此主值在每个点存在。

为了求平方界，先在径向乘上 \(e^{-\eta r}\)、\(\eta>0\)，得到同样按球面平均零定义的算子 \(T_\eta\)。对 \(\operatorname{Re}z>0\)，有
\[
 \int_0^\infty\frac{e^{-zr}-e^{-r}}r\,\mathrm dr=-\log z.
\]
右半平面取连续对数：在其每个紧子集上，对 \(z\) 求导的积分由一个指数衰减函数控制，导数为 \(-1/z\)；再由 \(z=1\) 的值为零确定常数。减项 \(e^{-r}\) 与角变量无关，乘 \(\Omega\) 后的积分为零。由此得到
\[
 m_\eta(\xi)=-\int_{S^{d-1}}\Omega(\theta)\cdot
                  \log(\eta+2\pi\mathrm i\xi\cdot\theta)\,\mathrm d\sigma(\theta).
\]
具体地，先在两个有限径向截断之间计算，再与 Schwartz 频率测试函数配对。近 \(r=0\) 时，差
\(e^{-(\eta+2\pi\mathrm i\xi\cdot\theta)r}-e^{-r}\)
除以 \(r\) 后由 \(C_\eta(1+|\xi|)\) 控制；乘测试函数后可积。远处则由 \(e^{-\eta r}+e^{-r}\) 控制。因而可以在这些配对中换序并放开截断，得到 \(\widehat{T_\eta f}=m_\eta\cdot\widehat f\) 的温和分布恒等式。

需要的估计是对 \(\eta\) 与频率同时一致的。对 \(\xi\ne0\)，记 \(\omega=\xi/|\xi|\)、\(a=\eta/|\xi|\)。平均零先去掉 \(\log|\xi|\)。若 \(a\geq1\)，再去掉 \(\log a\)，余下的 \(\log(1+2\pi\mathrm i\omega\cdot\theta/a)\) 一致有界。若 \(0<a\leq1\)，对数虚部的绝对值不超过 \(\pi/2\)，实部的正部一致有界，负部则由
\(C+\bigl|\log|\omega\cdot\theta|\bigr|\) 控制。旋转不变性给出
\[
 \sup_{\omega\in S^{d-1}}
 \int_{S^{d-1}}\bigl|\log|\omega\cdot\theta|\bigr|\,\mathrm d\sigma(\theta)<\infty:
\]
旋转到一个坐标方向后，在赤道附近令 \(s=\omega\cdot\theta\)，面积密度有界，该部分归结为
\(\int_{-1/2}^{1/2}\bigl|\log|s|\bigr|\,\mathrm ds<\infty\)；其余球面上对数有界。因此
\[
 \sup_{\eta>0}\|m_\eta\|_\infty\leq C_d\|\Omega\|_\infty.
\]
对几乎每个球面点令 \(\eta\downarrow0\)，复对数趋向
\(\log|2\pi\xi\cdot\theta|+\mathrm i\pi\cdot\operatorname{sgn}(\xi\cdot\theta)/2\)。上述控制给出 \(m_\eta(\xi)\to m(\xi)\)。Plancherel 与频率上的支配收敛于是给出有界 \(L^2\) 乘子算子 \(T\)，且 \(T_\eta f\to Tf\) 于 \(L^2\)。

在物理空间中，减去 \(f(x)\) 后的近原点积分和 Schwartz 远尾，都有与 \(0<\eta\leq1\) 无关的可积控制，所以 \(T_\eta f(x)\to T_\Omega f(x)\) 对每个 \(x\) 成立。取 \(L^2\) 收敛的几乎处处子列，便把这个主值识别为 \(Tf\)。平均零也从乘子公式中消去 \(\log|\xi|\)，给出零次齐次性。

核的大小与梯度常数由 \(C_d\|\Omega\|_{C^1}\) 控制。它对光滑紧支集输入的支集外表示已由主值公式给出；对一般紧支集 \(L^2\) 输入，在任意到其支集有正距离的有界区域上，用共用一个稍大紧支集的光滑 \(L^2\) 逼近。该区域到共同支集有正距离，核积分的差由 \(L^1\) 误差一致控制；\(L^2\) 输出收敛再抽子列，即得所需支集外表示。以可数个这样的区域穷竭支集外，说明 \(T\) 满足标准算子定义。\cref{b2:thm:18-cz-lp}因此给出全部 \(1<p<\infty\) 的界。先将非零的 \(\Omega\) 除以其 \(C^1\) 范数应用该结论，再还原尺度，就得到陈述中的线性范数因子；\(\Omega=0\) 时直接成立。
\end{proof}

乘子公式的两部分来自不同的对称性。若 \(\Omega\) 为奇函数，对数项的球面积分为零；若 \(\Omega\) 为偶函数，符号函数项的积分为零。再假设 \(\Omega\) 实值时，前一种乘子为纯虚函数，后一种为实偶函数。二维的 \(\Omega(\theta)=\theta_1\theta_2\) 给出偶核 \(x_1x_2/|x|^4\)，而 \(\Omega(\theta)=\theta_j\) 给出 Riesz 型奇核。相反，\(\Omega=1\) 产生前面的径向对数发散；人为扣除发散常数属于另外一种归一化，不能据此声称原硬截断主值已经存在。

\subsection{截断}
\label{b2:subsec:180304}

取整数 \(d\geq1\)，设核在对角线外为 Borel 可测，且满足 \(|K(x,y)|\leq C_0|x-y|^{-d}\)。前面的标准核满足这一可测性：在远离对角线的小邻域内，先移动 \(x\) 再移动 \(y\)，两个 Hölder 界共同给出连续性。对 \(\epsilon>0\) 定义硬截断
\[
 T_\epsilon f(x)=\int_{|x-y|>\epsilon}K(x,y)\cdot f(y)\,\mathrm dy.
\]
若 \(f\in L^p(\mathbb R^d)\)、\(1<p<\infty\)，Hölder 给出
\[
 \int_{|x-y|>\epsilon}|K(x,y)|\cdot|f(y)|\,\mathrm dy
 \leq C_{d,p}C_0\epsilon^{-d/p}\|f\|_p.
\]
这里径向控制核属于 \(L^{p'}\)，因为 \(dp'>d\)。\(p=1\) 时改用 \(|K|\leq C_0\epsilon^{-d}\)，仍有绝对收敛。故每个固定截断在这些输入上有普通积分的意义；一般有界输入却可能在无穷远处失去可积性。

\ProseParagraphBreak

把截断半径移向零，还需另外的论证。给定的标准 \(L^2\) 算子可以含有一个支集外核看不见的乘法项，所以截断的极限不必等于它。常见的识别办法，是先对 \(C_c^\infty\) 输入验证截断的逐点极限，再证明这些截断在 \(L^2\) 中趋于给定算子；平方收敛的子列才能把两个极限识别起来。下一节先证明极大截断的有界性，再把这项识别从光滑类传给一般 \(L^p\) 输入。

\ProseParagraphBreak

不同截断约定的比较也应写出实际差值。取实函数 \(\chi\in C^\infty([0,\infty))\)，在 \([0,1]\) 为零、在 \([2,\infty)\) 为一，且 \(0\leq\chi\leq1\)，令
\[
 T^\chi_\epsilon f(x)=\int K(x,y)\cdot
                 \chi(|x-y|/\epsilon)\cdot f(y)\,\mathrm dy.
\]
与硬截断的差仅在 \(\epsilon<|x-y|<2\epsilon\)，所以对有限指数的 \(L^p\) 输入，
\[
 |T^\chi_\epsilon f(x)-T_\epsilon f(x)|
 \leq C_dC_0\epsilon^{-d}\int_{B(x,2\epsilon)}|f|
 \leq C_d'C_0\cdot Mf(x).
\]
此界能传递 \(1<p<\infty\) 的极大范数估计，但右端不随 \(\epsilon\) 消失，不能单凭这一界识别极限。

\ProseParagraphBreak

若硬截断在某点已有有限极限，则两种截断其实有相同极限，而不必再要求核的抵消。由 \(\int_1^2\chi'(s)\,\mathrm ds=1\)，对上述绝对收敛积分换序，得
\[
 T^\chi_\epsilon f(x)=\int_1^2\chi'(s)\cdot T_{s\epsilon}f(x)\,\mathrm ds.
\]
因此它与硬截断极限 \(L\) 的差不超过
\[
 \left(\int_1^2|\chi'(s)|\,\mathrm ds\right)
 \sup_{\epsilon\leq r\leq2\epsilon}|T_rf(x)-L|\longrightarrow0.
\]
对球面平均零的齐次核，还能在光滑输入上直接比较：环带差核的积分为零，减去 \(f(x)\) 后，差不超过
\(C_d\|\Omega\|_\infty\epsilon\sup_{B(x,2\epsilon)}|\nabla f|\)。这给出共同的光滑类极限；向一般输入推广仍须控制极大算子的逼近误差。这里比较的是两种截断；它们与预先给定的算子 \(T\) 是否相等，仍是前一段所说的识别问题。

\ProseParagraphBreak

截断还可以先保留一个有限的外半径，再逐次放开。对 \(L^p\)、\(1\leq p<\infty\)，固定内半径后的远场已有上述绝对积分控制，外半径趋于无穷不会改变结果。对常函数或一般有界输入，远场可能发散或仅有条件收敛，需要另行指定归一化；有限指数上建立的定义不能不加说明地用于这些输入。

% !TeX root = ../../Book2.tex
\section{\texorpdfstring{\(L^p\)}{Lᵖ} 有界性}
\label{b2:sec:1804}

上一节把核条件与平方有界性分开。本节先把两者用于好坏分解：好函数的高度受控，可以使用平方估计；坏函数的零平均则使支集外的尾部可积。弱端点成立后，大小分解插值给出 \(1<p<2\)，伴随再给出 \(p>2\)。最后才讨论截断是否恢复这个有界算子。

\subsection{\texorpdfstring{\(L^2\)}{L2} 假设}
\label{b2:subsec:180401}

设 \(T\) 是标准 Calderón--Zygmund 算子，记其平方算子范数为 \(N\)，大小与光滑常数为 \(C_0,C_1\)。对卷积型例子，\(N\) 可以从有界 Fourier 乘子得到；非平移不变算子则需要另外的平方估计。以下证明不从核条件推导 \(N\)，而是说明已经取得的平方界怎样用于其他指数。

\ProseParagraphBreak

若 \(f\in L^1\cap L^2\)，在高度 \(\alpha>0\) 作上一节的分解 \(f=g+\sum_jb_j\)。好函数满足
\[
 \|g\|_2^2\leq\|g\|_\infty\cdot\|g\|_1
 \leq2^d\alpha\|f\|_1.
\]
坏块虽然可能有很大的平方范数，仍各自属于算子的原定义域。事实上，在不交的 \(Q_j\) 上有
\[
 \|b_j\|_2^2
 =\int_{Q_j}|f|^2-|Q_j|\cdot|f_{Q_j}|^2
 \leq\int_{Q_j}|f|^2.
\]
因此坏块级数在 \(L^2\) 中收敛于 \(b=f-g\)，有限和的算子像也在 \(L^2\) 中趋于 \(Tb\)。这项 \(L^2\) 收敛使下一段可以先估计有限个坏块，再取极限。

\subsection{弱型 \texorpdfstring{\((1,1)\)}{(1,1)}}
\label{b2:subsec:180402}

\Needspace{6\baselineskip}
\begin{theorem}[Calderón--Zygmund 弱型]\label{b2:thm:18-cz-weak}
设整数 \(d\geq1\)、\(0<\delta\leq1\)，\(C_0,C_1,N\geq0\) 为有限常数。设 \(T:L^2(\mathbb R^d;\mathbb C)\to L^2(\mathbb R^d;\mathbb C)\) 为复线性算子，\(\|T\|_{2\to2}\leq N\)，其离开对角线的可测核 \(K\) 满足
\[
 |K(x,y)|\leq C_0|x-y|^{-d}\quad(x\ne y),
\]
\[
 \begin{aligned}
 &|K(x,y)-K(x',y)|+|K(y,x)-K(y,x')|\\
 &\quad\leq C_1\frac{|x-x'|^\delta}{|x-y|^{d+\delta}}
 \quad\left(x\ne y,\ |x-x'|\leq\frac{|x-y|}2\right).
 \end{aligned}
\]
假设对每个紧支集 \(h\in L^2\)，在其支集外几乎处处有
\(Th(x)=\int K(x,y)\cdot h(y)\,\mathrm dy\)。
则对 \(f\in L^1(\mathbb R^d)\cap L^2(\mathbb R^d)\)、\(t>0\)，有
\[
 \bigl|\{\,|Tf|>t\,\}\bigr|\leq
 C_{d,\delta}(N+C_0+C_1)t^{-1}\|f\|_1.
\]
算子 \(T\) 有唯一的连续线性延拓
\(T_1:L^1(\mathbb R^d)\to L^{1,\infty}(\mathbb R^d)\)，目标连续性使用弱 \(L^1\) 拟范数；延拓满足同一个弱型界，并在 \(L^1\cap L^2\) 上等于原算子。
\end{theorem}
\begin{proof}
记 \(B=N+C_0+C_1\)。\(B=0\) 时 \(T=0\)，以下设 \(B>0\)，并取分解高度 \(\alpha=t/B\)。由\cref{b2:thm:18-cz-decomposition}，好部分给出
\[
 \bigl|\{\,|Tg|>t/2\,\}\bigr|
 \leq\frac{4N^2}{t^2}\|g\|_2^2
 \leq C_d\frac{N^2}{Bt}\|f\|_1.
\]
取 \(Q_j^*=8\sqrt d\,Q_j\)，记 \(E^*=\bigcup_jQ_j^*\)、\(G=\mathbb R^d\setminus E^*\)。扩大坏区域的测度满足
\[
 |E^*|\leq(8\sqrt d)^d\sum_j|Q_j|
 \leq C_d\frac B t\|f\|_1.
\]
对每个坏块，在 \(Q_j^*\) 外用支集外表示和零平均减去 \(K(x,c_{Q_j})\)。两个相关距离可比，输入变量光滑条件因而满足\cref{b2:thm:18-cancellation-tail}的假设，给出
\[
 \int_{\mathbb R^d\setminus Q_j^*}|Tb_j(x)|\,\mathrm dx
 \leq C_{d,\delta}C_1\|b_j\|_1.
\]
对有限和先积分取三角不等式。坏块在 \(L^2\) 中收敛，故其算子像的某个子列几乎处处趋于 \(Tb\)；在固定的 \(G\) 上用 Fatou，得到
\[
 \int_G|Tb|
 \leq C_{d,\delta}C_1\sum_j\|b_j\|_1
 \leq2C_{d,\delta}C_1\|f\|_1.
\]
这里的可积性只在扩大坏区域之外成立。于是
\[
 \bigl|\{\,|Tf|>t\,\}\bigr|
 \leq \bigl|\{\,|Tg|>t/2\,\}\bigr|+|E^*|
       +\frac2t\int_G|Tb|
 \leq C_{d,\delta}\frac B t\|f\|_1,
\]
其中使用 \(N^2/B\leq B\)、\(C_1\leq B\)。选择高度 \(t/B\) 也使这个弱型常数随算子大小线性缩放。

现在构造 \(L^1\) 延拓，记刚才的弱型常数为 \(W=C_{d,\delta}B\)。对 \(f\in L^1\)，取 \(f_n\in L^1\cap L^2\) 使 \(f_n\to f\) 于 \(L^1\)。抽成子列 \(g_n\)，使 \(\|g_{n+1}-g_n\|_1\leq2^{-3n}\)。线性与弱界给出
\[
 \bigl|\{\,|Tg_{n+1}-Tg_n|>2^{-n}\,\}\bigr|\leq W2^{-2n}.
\]
这些测度的和有限，故相应集合的并集尾部测度趋零。除一个零测集外，输出的相邻差最终不超过 \(2^{-n}\)，所以 \(Tg_n\) 逐点为 Cauchy 列并有有限极限，记为 \(T_1f\)。

对 \(0<s<t\)，有
\(\mathbf1_{\{|T_1f|>t\}}\leq\liminf_n\mathbf1_{\{|Tg_n|>s\}}\)。
Fatou 给出 \(\bigl|\{\,|T_1f|>t\,\}\bigr|\leq W\|f\|_1/s\)，令 \(s\uparrow t\) 即传递弱界。更重要的是，对任一固定的原逼近项 \(f_n\)，对 \(T(f_n-g_m)\) 用同一论证，得
\[
 \bigl|\{\,|Tf_n-T_1f|>t\,\}\bigr|
 \leq \frac Wt\|f_n-f\|_1.
\]
因而原列整体在弱拟范数中趋于 \(T_1f\)，另取任何 \(L^1\cap L^2\) 逼近列也给出同一极限。和与数乘可在公共逼近列上取局部依测度极限，故 \(T_1\) 线性，并有
\[
 \|T_1f-T_1h\|_{1,\infty}\leq W\|f-h\|_1.
\]
对原定义域中的输入可取常值逼近列，说明两种定义一致。稠密性与上述超水平集合的差估计也给出连续延拓的唯一性。
\end{proof}

弱端点保留的是超水平集合估计。Hilbert 变换的区间输入已有不可积远尾，所以不能把这个结论改成强 \(L^1\) 有界。下面只对严格的中间指数插值。

\subsection{\texorpdfstring{\(1<p<\infty\)}{1<p<∞}}
\label{b2:subsec:180403}

\Needspace{6\baselineskip}
\begin{theorem}[Calderón--Zygmund]\label{b2:thm:18-cz-lp}
设整数 \(d\geq1\)、\(0<\delta\leq1\)，\(C_0,C_1,N\geq0\) 为有限常数。设 \(T:L^2(\mathbb R^d;\mathbb C)\to L^2(\mathbb R^d;\mathbb C)\) 为复线性算子，\(\|T\|_{2\to2}\leq N\)，其离开对角线的可测核 \(K\) 满足
\[
 |K(x,y)|\leq C_0|x-y|^{-d}\quad(x\ne y),
\]
\[
 \begin{aligned}
 &|K(x,y)-K(x',y)|+|K(y,x)-K(y,x')|\\
 &\quad\leq C_1\frac{|x-x'|^\delta}{|x-y|^{d+\delta}}
 \quad\left(x\ne y,\ |x-x'|\leq\frac{|x-y|}2\right).
 \end{aligned}
\]
假设对每个紧支集 \(h\in L^2\)，在其支集外几乎处处有
\(Th(x)=\int K(x,y)\cdot h(y)\,\mathrm dy\)。
则对每个 \(1<p<\infty\)，原算子在 \(L^2\cap L^p\) 上有唯一有界延拓 \(T_p:L^p(\mathbb R^d)\to L^p(\mathbb R^d)\)，且
\[
 \|T_pf\|_p\leq C_{d,\delta,p}(N+C_0+C_1)\|f\|_p.
\]
这些延拓在任意两个输入空间的交集上一致，并与上述核条件产生的 \(L^1\to L^{1,\infty}\) 连续延拓在 \(L^1\cap L^p\) 上一致。
\end{theorem}
\begin{proof}
上一条弱型延拓 \(T_1\) 与原算子在 \(L^1\cap L^2\) 上一致，所以
\(T(f_1+f_2)=T_1f_1+Tf_2\)
定义了和空间 \(L^1+L^2\) 上的同一个复线性算子，结果与分解无关。弱 \((1,1)\) 与强 \((2,2)\) 于是满足\cref{b2:thm:17-marcinkiewicz}的定义域及端点条件，得到 \(1<p<2\) 的强型界。若 \(B=N+C_0+C_1>0\)，先对 \(T/B\) 插值，两端常数均由 \(d,\delta\) 控制，再还原 \(B\)，便得到陈述中的线性因子；\(B=0\) 时直接成立。

当 \(p>2\) 时，\(p'=p/(p-1)\in(1,2)\)。上一节已经证明 \(T^*\) 的核为 \(\overline{K(y,x)}\)，具有相同的大小、光滑常数和平方范数界。因此刚才的结论给出它的 \(L^{p'}\) 界。对 \(f\in L^p\cap L^2\)、\(g\in L^{p'}\cap L^2\)，有
\[
 \left|\int Tf\cdot\overline g\right|
 =\left|\int f\cdot\overline{T^*g}\right|
 \leq C_{d,\delta,p'}B\|f\|_p\cdot\|g\|_{p'}.
\]
取有限测度支集的简单 \(g\) 便已足以恢复 \(L^p\) 范数：在有限球内再截断 \(|Tf|\)，沿 \(Tf/|Tf|\) 的相位乘以 \(|Tf|^{p-1}\) 作测试，并归一化、简单逼近，最后穷竭；零点处取零。配对中已有 \(\overline g\)，因此这样选择才产生 \(|Tf|^p\)。这个范数刻画允许原先未知的 \(\|Tf\|_p\) 为无穷；上式的有限右端证明它实际有限。\(L^p\cap L^2\) 的稠密性随后给出唯一有界延拓。

还须识别不同指数上的输出。若 \(f\in L^p\cap L^r\)、\(1<p,r<\infty\)，取
\(f_n=f\cdot\mathbf1_{\{|x|\leq n,\ |f|\leq n\}}\)。
它们属于 \(L^2\cap L^p\cap L^r\)，并在后两个空间同时趋于 \(f\)。它们的输出在两个指数上都是原来的 \(Tf_n\)，各自的范数极限都给出每个有限球上的依测度极限，唯一性使 \(T_pf=T_rf\) 几乎处处。若 \(f\in L^1\cap L^p\)，同样的截断还在 \(L^1\) 收敛；弱型差估计也给出局部依测度收敛，故 \(T_1f=T_pf\)。
\end{proof}

这项证明中，原算子坏块消去常数用的是输入变量的光滑性；得到另一半指数时，伴随坏块需要原核输出变量的光滑性。它不要求核齐次，也不要求算子平移不变。若把点态 Hölder 条件换成积分 Hörmander 条件，相同的弱型、插值及对偶路线仍适用，但必须同时给出核及伴随核的相应积分界。下面的截断和振荡估计继续使用点态条件。

\subsection{极大截断、向量值估计与 \texorpdfstring{\(\mathrm{BMO}\)}{BMO} 端点}
\label{b2:subsec:180404}

先把截断本身定义清楚。对标准核及 \(f\in L^p\)、\(1\leq p<\infty\)，上一节已经证明每个固定 \(\epsilon>0\) 的积分
\[
 T_\epsilon f(x)=\int_{|x-y|>\epsilon}K(x,y)\cdot f(y)\,\mathrm dy
\]
绝对收敛。固定 \(x\)，在 \(\epsilon_0>0\) 附近可用 \(|x-y|>\epsilon_0/2\) 的绝对积分控制，球面又为零测集，所以支配收敛使 \(\epsilon\mapsto T_\epsilon f(x)\) 连续。因此其\term[jida-jieduan]{极大截断}[maximal truncation]
\[
 T_*f(x)\coloneqq\sup_{\epsilon>0}|T_\epsilon f(x)|
       =\sup_{\epsilon\in\mathbb Q,\ \epsilon>0}|T_\epsilon f(x)|
\]
可测，暂时允许为无穷。\(T^*\) 始终表示线性伴随，与 \(T_*\) 不同。

\Needspace{6\baselineskip}
\begin{proposition}\label{b2:thm:18-maximal-truncation}
设整数 \(d\geq1\)、\(0<\delta\leq1\)，\(C_0,C_1,N\geq0\) 为有限常数。设 \(T:L^2(\mathbb R^d;\mathbb C)\to L^2(\mathbb R^d;\mathbb C)\) 为复线性算子，\(\|T\|_{2\to2}\leq N\)，其离开对角线的可测核 \(K\) 满足
\[
 |K(x,y)|\leq C_0|x-y|^{-d}\quad(x\ne y),
\]
\[
 \begin{aligned}
 &|K(x,y)-K(x',y)|+|K(y,x)-K(y,x')|\\
 &\quad\leq C_1\frac{|x-x'|^\delta}{|x-y|^{d+\delta}}
 \quad\left(x\ne y,\ |x-x'|\leq\frac{|x-y|}2\right).
 \end{aligned}
\]
假设对每个紧支集 \(h\in L^2\)，在其支集外几乎处处有
\(Th(x)=\int K(x,y)\cdot h(y)\,\mathrm dy\)。
令 \(B=N+C_0+C_1\)，对 \(1\leq p<\infty\) 的输入以绝对积分定义
\(T_\epsilon f(x)=\int_{|x-y|>\epsilon}K(x,y)\cdot f(y)\,\mathrm dy\)，
并置 \(T_*f=\sup_{\epsilon>0}|T_\epsilon f|\)。若 \(1<q<\infty\)、\(f\in C_c^\infty(\mathbb R^d)\)，则
\[
 T_*f(x)\leq C_{d,\delta,q}
 \left[M(|Tf|^q)(x)^{1/q}
       +B\,M(|f|^q)(x)^{1/q}\right]
\]
几乎处处成立，其中 \(M\) 为中心 Hardy--Littlewood 极大算子。因此对每个 \(1<p<\infty\)、\(f\in L^p(\mathbb R^d)\)，有
\[
 \|T_*f\|_p\leq C_{d,\delta,p}B\|f\|_p.
\]
这两个估计不要求预先知道 \(T_\epsilon\to T\) 或主值存在。
\end{proposition}
\begin{proof}
先取光滑紧支集 \(f\)，固定 \(x,\epsilon\)，分解
\(f_1=f\cdot\mathbf1_{B(x,2\epsilon)}\)、\(f_2=f-f_1\)。
对几乎处处 \(z\in B(x,\epsilon/4)\)，\(Tf_2(z)\) 有支集外核表示。比较它与 \(T_\epsilon f(x)\)，在 \(\epsilon<|x-y|<2\epsilon\) 上，用大小界和球平均得到至多 \(C_dC_0\cdot Mf(x)\)。在 \(|x-y|>2\epsilon\) 上，用输出变量光滑性；分成半径约 \(2^j\epsilon\)、\(j\geq1\) 的环带后，每一带至多为
\[
 C_dC_1\epsilon^\delta(2^j\epsilon)^{-d-\delta}
          \int_{B(x,2^{j+1}\epsilon)}|f|
 \leq C_d'C_1 2^{-j\delta}\cdot Mf(x).
\]
故以 \(A_{d,\delta}=C_d(C_0+C_1\sum_{j\geq1}2^{-j\delta})\) 记这部分常数，有
\[
 |T_\epsilon f(x)|\leq |Tf(z)|+|Tf_1(z)|
                       +A_{d,\delta}\cdot Mf(x)
\]
对球内几乎处处 \(z\) 成立。取 \(z\) 上的归一化 \(L^q\) 范数，第一项不超过 \(M(|Tf|^q)(x)^{1/q}\)。第二项由已证的 \(L^q\) 有界性控制为
\[
 |B(x,\epsilon/4)|^{-1/q}\cdot\|Tf_1\|_q
 \leq C_{d,\delta,q}B\epsilon^{-d/q}
                   \|f\cdot\mathbf1_{B(x,2\epsilon)}\|_q
 \leq C'_{d,\delta,q}B\,M(|f|^q)(x)^{1/q}.
\]
Hölder 还给出 \(Mf\leq M(|f|^q)^{1/q}\)。对 \(\epsilon\) 取上确界，即得点态估计。对指定的 \(p>1\)，选 \(1<q<p\)，极大算子在 \(L^{p/q}\) 上有界，\(T\) 在 \(L^p\) 上有界，所以
\(\|T_*f\|_p\leq C_{d,\delta,p}B\|f\|_p\)。

最后将这个界用于积分定义的任意 \(f\in L^p\)。取 \(g_n\in C_c^\infty\) 趋于 \(f\) 于 \(L^p\)。对每个固定 \(\epsilon>0\)，Hölder 给出
\[
 |T_\epsilon(g_n-f)(x)|
 \leq C_{d,p}C_0\epsilon^{-d/p}\|g_n-f\|_p\longrightarrow0.
\]
因此 \(T_*f\leq\liminf_nT_*g_n\)；可先对有限个有理半径取最大值，再取其上确界得到这个不等式。Fatou 把 \(L^p\) 界传给 \(f\)。于是 \(T_*f\) 几乎处处有限；再由 \(|T_*f-T_*h|\leq T_*(f-h)\)，这个积分定义与从光滑类取得的连续延拓相同。
\end{proof}

现在可以证明主值的传递。固定 \(1<p<\infty\)，额外假设对每个 \(g\in C_c^\infty\)，硬截断几乎处处存在全参数有限极限，且这个极限等于给定的 \(Tg\)。一种足够的核验方式，是同时证明光滑类的逐点收敛及 \(T_\epsilon g\to Tg\) 于 \(L^2\)：后者的几乎处处子列识别前者的极限。单独的 \(L^2\) 收敛不保证全参数逐点收敛；单独的逐点收敛也可能得到含乘法差异的另一个算子。

\ProseParagraphBreak

对一般 \(f\in L^p\)，取光滑 \(g\)，在光滑类已知极限的满测集上有
\begin{equation}\label{b2:eq:18-principal-value-transfer}
 \limsup_{\epsilon\downarrow0}|T_\epsilon f-T_pf|
 \leq T_*(f-g)+|T_p(f-g)|.
\end{equation}
左侧超过任意 \(a>0\) 的集合 \(E_a\) 因而满足
\[
 |E_a|\leq\left(\frac2a\right)^p
        \left(C_*^p+\|T_p\|_{p\to p}^p\right)\cdot\|f-g\|_p^p,
\]
其中 \(C_*\) 是刚才的极大算子界。集合 \(E_a\) 只由 \(f,a\) 决定，令光滑逼近误差趋零，就得 \(|E_a|=0\)。对 \(a=1/n\) 取可数并，得到 \(T_\epsilon f\to T_pf\) 几乎处处。再由
\(|T_\epsilon f-T_pf|\leq T_*f+|T_pf|\in L^p\)
的支配收敛，得到同一截断的 \(L^p\) 收敛。

\ProseParagraphBreak

\cref{b2:thm:18-homogeneous-kernel}已在 Schwartz 输入上把主值与所构造的平方算子识别，因而满足上述光滑类条件；Hilbert 与 Riesz 变换也可这样使用。反例 \(T=M_a\)、\(K=0\) 则有全部截断为零，却未必 \(Tf=0\)，说明识别条件不能从极大范数界中删去。\appref{b2:app:C}将继续讨论这些算子的端点空间。

\ProseParagraphBreak

应区分三个弱端点：\(M\) 的弱 \((1,1)\) 来自覆盖，\(T\) 的弱 \((1,1)\) 来自好坏分解，\(T_*\) 的弱 \((1,1)\) 还需要额外的端点论证。上面的 \(q>1\) 估计不能直接用于一般 \(L^1\) 输入，因为 \(|f|^q\) 未必可积。进一步的结果可参看 Loukas Grafakos 的《Estimates for Maximal Singular Integrals》\cite[Theorem 1]{Grafakos2003Maximal}：该结果在核及转置核的统一环带大小、积分 Hörmander 条件以及关联 \(L^2\) 算子的支集外表示下，给出双截断极大算子的弱 \((1,1)\) 界。当前标准核满足这些条件；对 \(f\in L^1\)，固定内半径后积分绝对收敛，令外半径趋于无穷即可包含这里的单截断。

\ProseParagraphBreak

平方函数应用还需要同时处理许多输出。若逐坐标把估计相加，常数可能随坐标个数增大；适合无限频带的估计应只依赖核的算子范数和总平方界。
\Needspace{6\baselineskip}
\begin{proposition}\label{b2:prop:18-hilbert-valued-cz}
设整数 \(d\geq1\)，\(\mathcal H_1,\mathcal H_2\) 为同取实数域或同取复数域的有限维 Hilbert 空间。设线性算子
\(T:L^2(\mathbb R^d;\mathcal H_1)\to L^2(\mathbb R^d;\mathcal H_2)\)
满足 \(\|T\|_{2\to2}\leq N<\infty\)。设可测核
\(K(x,y)\in\mathcal B(\mathcal H_1,\mathcal H_2)\)、\(x\ne y\)，存在有限常数 \(C_0,C_1\geq0\)、\(0<\delta\leq1\)，使
\[
 \|K(x,y)\|\leq C_0|x-y|^{-d},
\]
\[
 \begin{aligned}
 &\|K(x,y)-K(x',y)\|+\|K(y,x)-K(y,x')\|\\
 &\quad\leq C_1\frac{|x-x'|^\delta}{|x-y|^{d+\delta}}
 \quad\left(x\ne y,\ |x-x'|\leq\frac{|x-y|}2\right),
 \end{aligned}
\]
这里均为算子范数。假设对每个紧支集 \(h\in L^2(\mathbb R^d;\mathcal H_1)\)，在其支集外几乎处处有
\(Th(x)=\int K(x,y)h(y)\,\mathrm dy\)，积分按有限维向量积分解释。则 \(T\) 有兼容的弱型 \((1,1)\) 延拓及强型 \((p,p)\) 延拓、\(1<p<\infty\)，且
\[
 \begin{aligned}
 \bigl|\{\,x:\|Tf(x)\|_{\mathcal H_2}>t\,\}\bigr|
 &\leq C_{d,\delta}(N+C_0+C_1)t^{-1}
                   \|f\|_{L^1(\mathcal H_1)},\\
 \|Tf\|_{L^p(\mathcal H_2)}
 &\leq C_{d,\delta,p}(N+C_0+C_1)\|f\|_{L^p(\mathcal H_1)}.
 \end{aligned}
\]
第一式对 \(f\in L^1(\mathbb R^d;\mathcal H_1)\)、\(t>0\) 成立，第二式对 \(f\in L^p(\mathbb R^d;\mathcal H_1)\)、\(1<p<\infty\) 成立。常数与两个 Hilbert 空间的维数无关。
\end{proposition}
\begin{proof}
先取 \(f\in L^1(\mathcal H_1)\cap L^2(\mathcal H_1)\)。对标量函数 \(\|f(x)\|_{\mathcal H_1}\) 在高度 \(\alpha\) 选择二进坏块，再在每块内取向量平均 \(f_Q=|Q|^{-1}\int_Qf\)。三角不等式给出
\[
 \|g\|_\infty\leq2^d\alpha,
 \quad\|g\|_1\leq\|f\|_1,
 \quad\|g\|_2^2\leq2^d\alpha\|f\|_1.
\]
令 \(b_Q=(f-f_Q)\cdot\mathbf1_Q\)。它满足向量等式 \(\int b_Q=0\)，且
\[
 \sum_Q\|b_Q\|_1\leq2\|f\|_1,
 \quad\sum_Q|Q|\leq\alpha^{-1}\|f\|_1,
 \quad\sum_Q\|b_Q\|_2^2\leq\|f\|_2^2.
\]
最后一式仍由平均的平方恒等式得到，因此坏块的向量 \(L^2\) 级数收敛。在 \(Q^*=8\sqrt d\,Q\) 外，
\[
 \|Tb_Q(x)\|_{\mathcal H_2}
 \leq\int_Q\|K(x,y)-K(x,c_Q)\|
                       \cdot\|b_Q(y)\|_{\mathcal H_1}\,\mathrm dy.
\]
径向尾积分给出 \(\int_{(Q^*)^c}\|Tb_Q\|\leq C_{d,\delta}C_1\|b_Q\|_1\)。有限和、\(L^2\) 子列及 Fatou 因而给出与标量证明相同的坏区域外积分界；取 \(\alpha=t/(N+C_0+C_1)\)，好部分 Chebyshev 与坏部分积分合起来给出所述弱型界。若分母为零，算子为零。

\(L^1\) 延拓也可按前面的可求和差值构造：几乎每点的向量相邻差之范数最终不超过 \(2^{-n}\)，Hilbert 空间的完备性保证点态极限存在。弱界对范数的超水平集合传递，差估计给出连续性和唯一性。

对 \(1<p<2\)，把输入分解为
\(f\cdot\mathbf1_{\{\|f\|>a\}}+f\cdot\mathbf1_{\{\|f\|\leq a\}}\)。
第一部分在 \(L^1(\mathcal H_1)\)，第二部分在 \(L^2(\mathcal H_1)\)，其幂次估计与标量证明一致，只需把 \(|f|\) 换成 \(\|f\|_{\mathcal H_1}\)。输出范数满足次可加性，\cref{b2:thm:17-marcinkiewicz}证明中的高度积分因而给出同一中间强型界，不出现坐标求和常数。

对 \(p>2\)，伴随核为 \(K(y,x)^*\)，两个 Hilbert 空间互换，核范数及差分界不变。\(L^2\) 内积关系及
\[
 |\langle u,v\rangle|\leq\|u\|_{\mathcal H_2}\cdot\|v\|_{\mathcal H_2}
\]
给出与前面相同的 \(L^p\)--\(L^{p'}\) 对偶估计。有限测度支集上沿输出向量方向归一化作测试，就能恢复其 \(L^p\) 范数，所用常数仍是一；稠密性与共同截断再识别各个延拓。
\end{proof}

例如设 \(T_j\) 是一族标量 Calderón--Zygmund 算子，其 \(N,C_0,C_1,\delta\) 对 \(j\) 统一。对有限组输入，取对角算子 \((f_j)\mapsto(T_jf_j)\)。它的核范数为 \(\sup_j|K_j(x,y)|\)，平方界由各坐标平方积分相加得到，范数至多为 \(N\)。故
\[
 \left\|\left(\sum_j|T_jf_j|^2\right)^{1/2}\right\|_p
 \leq C_{d,\delta,p}(N+C_0+C_1)
       \left\|\left(\sum_j|f_j|^2\right)^{1/2}\right\|_p.
\]
对可数族，先用有限部分和，再在两侧用单调收敛即可。一般矩阵核则必须检查算子范数；逐个矩阵元的统一上界不能代替它。

\ProseParagraphBreak

无穷输入端点适合用平均振荡衡量。对 \(h\in L^1_{\mathrm{loc}}(\mathbb R^d)\)，记
\[
 \|h\|_{\mathrm{BMO}}\coloneqq
 \sup_Q\frac1{|Q|}\int_Q|h-h_Q|,
 \quad h_Q=\frac1{|Q|}\int_Qh,
\]
其中 \(Q\) 遍历轴平行立方体。有限者构成\term[youjie-pingjun-zhendang]{有界平均振荡空间}[bounded mean oscillation]，简称 \(\mathrm{BMO}\)。加上常数不改变这个半范数，故把相差常数的函数视为同一类。下面先对实际属于 \(L^2\) 定义域的有界紧支集输入证明估计，再解释一般有界输入的模常数意义。

\Needspace{6\baselineskip}
\begin{proposition}\label{b2:prop:18-infinity-bmo}
设整数 \(d\geq1\)、\(0<\delta\leq1\)，\(C_0,C_1,N\geq0\) 为有限常数。设 \(T:L^2(\mathbb R^d;\mathbb C)\to L^2(\mathbb R^d;\mathbb C)\) 为复线性算子，\(\|T\|_{2\to2}\leq N\)，其离开对角线的可测核 \(K\) 满足
\[
 |K(x,y)|\leq C_0|x-y|^{-d}\quad(x\ne y),
\]
\[
 \begin{aligned}
 &|K(x,y)-K(x',y)|+|K(y,x)-K(y,x')|\\
 &\quad\leq C_1\frac{|x-x'|^\delta}{|x-y|^{d+\delta}}
 \quad\left(x\ne y,\ |x-x'|\leq\frac{|x-y|}2\right).
 \end{aligned}
\]
假设对每个紧支集 \(h\in L^2\)，在其支集外几乎处处有
\(Th(x)=\int K(x,y)\cdot h(y)\,\mathrm dy\)。
对局部可积函数采用轴平行立方体的半范数
\(\|h\|_{\mathrm{BMO}}=\sup_Q|Q|^{-1}\int_Q|h-h_Q|\)、\(h_Q=|Q|^{-1}\int_Qh\)。则每个有界紧支集复函数 \(f\) 的原 \(L^2\) 输出满足
\[
 \|Tf\|_{\mathrm{BMO}}\leq C_{d,\delta}(N+C_1)\|f\|_\infty.
\]
常数与输入支集的位置、大小无关。
\end{proposition}
\begin{proof}
固定立方体 \(Q\)，记中心为 \(c_Q\)、边长为 \(\ell\)，取 \(Q^*=8\sqrt d\,Q\)。令
\(f_1=f\cdot\mathbf1_{Q^*}\)、\(f_2=f-f_1\)。局部部分用平方界：
\[
 \frac1{|Q|}\int_Q|Tf_1|
 \leq |Q|^{-1/2}\cdot N\cdot\|f_1\|_2
 \leq(8\sqrt d)^{d/2}N\|f\|_\infty.
\]
远场在 \(Q\) 内有核表示。取比较常数
\[
 c_Q(f)=\int_{(Q^*)^c}K(c_Q,y)\cdot f(y)\,\mathrm dy.
\]
输入紧支集且积分离开对角线，所以它绝对收敛。对 \(x\in Q\)，输出变量光滑性与距离可比性给出
\[
 \begin{aligned}
 |Tf_2(x)-c_Q(f)|
 &\leq C_dC_1\|f\|_\infty\ell^\delta
       \int_{|y-c_Q|\geq4\sqrt d\,\ell}
             |y-c_Q|^{-d-\delta}\,\mathrm dy\\
 &\leq C_{d,\delta}C_1\|f\|_\infty
 \end{aligned}
\]
几乎处处成立。故 \(Q\) 上相对于 \(c_Q(f)\) 的平均偏差受控。再用
\[
 |(Tf)_Q-c_Q(f)|\leq\frac1{|Q|}\int_Q|Tf-c_Q(f)|
\]
把比较常数换成真正平均，至多增加因子二。对全部 \(Q\) 取上确界即得结论。
\end{proof}

对一般 \(f\in L^\infty\)，\(c_Q(f)\) 的远场积分未必收敛，但下面的局部归一化有意义：在每个 \(Q\) 内置
\[
 F_Q(x)=T(f\cdot\mathbf1_{Q^*})(x)
       +\int_{(Q^*)^c}[K(x,y)-K(c_Q,y)]\cdot f(y)\,\mathrm dy.
\]
第一项的输入属于 \(L^2\)，第二项由刚才的衰减估计绝对收敛；\(F_Q\in L^1(Q)\)，且相对于零的平均绝对值受 \(C_{d,\delta}(N+C_1)\|f\|_\infty\) 控制。

\ProseParagraphBreak

这些局部输出在重叠处只相差常数。为验证这一点，令 \(f_R=f\cdot\mathbf1_{B(0,R)}\)，并取 \(R\) 足够大使 \(Q^*\subset B(0,R)\)。在 \(Q\) 内有
\[
 Tf_R-c_{R,Q}=T(f\cdot\mathbf1_{Q^*})
  +\int_{B(0,R)\setminus Q^*}[K(x,y)-K(c_Q,y)]\cdot f(y)\,\mathrm dy,
\]
其中 \(c_{R,Q}=\int_{B(0,R)\setminus Q^*}K(c_Q,y)\cdot f(y)\,\mathrm dy\)。
右端在 \(L^1(Q)\) 中趋于 \(F_Q\)：尾部的核差积分对 \(x\in Q\) 一致趋零。对两个有正测度重叠的立方体 \(Q,P\)，相减后左端为常数 \(c_{R,P}-c_{R,Q}\)，故其 \(L^1(Q\cap P)\) 极限 \(F_Q-F_P\) 也为常数。

\ProseParagraphBreak

取一列递增并穷竭全空间的同心立方体，逐次给 \(F_Q\) 加常数，使相邻两个在较小立方体上相等，即可拼成全空间局部可积函数。对任意其他立方体，它与相应的 \(F_Q\) 仍只相差常数，因而保留上述平均振荡界。不同拼接归一化给出的结果相差一个全局常数；这便定义了 \(L^\infty\) 输入的模常数 \(\mathrm{BMO}\) 输出，并与紧支集输入的原算子相容。弱积分端点使用输入变量的差分，这里的平均振荡与远场归一化则使用输出变量的差分。

% !TeX root = ../../Book2.tex
\section{典型算子}
\label{b2:sec:1805}

Hilbert 变换用符号的正负区分一维频率。在高维中，频率不再只有两个方向；把向量 \(\xi\) 归一化为 \(\xi/|\xi|\)，就得到一组只记录方向、不改变频率大小的零阶算子。这组 Riesz 变换既能组成共轭调和函数，也能把 Poisson 方程的右端转为解的二阶导数。本节逐一说明这些联系中的核、常数和函数空间。

\subsection{Riesz 变换}
\label{b2:subsec:180501}

\Needspace{6\baselineskip}
\begin{definition}\label{b2:thm:18-riesz-definition}
设 \(d\geq1\) 为整数，\(1\leq j\leq d\)。对 \(f\in L^2(\mathbb R^d;\mathbb C)\)，以本卷 \(e^{-2\pi\mathrm ix\cdot\xi}\) 的 Fourier 约定定义\term[riesz-bianhuan]{Riesz 变换}[Riesz transform] \(R_j\)：
\[
 \widehat{R_jf}(\xi)\coloneqq
 -\mathrm i\frac{\xi_j}{|\xi|}\cdot\widehat f(\xi)
 \quad(\xi\ne0).
\]
零频处的乘子可任意补值；变换及输出均按 \(L^2\) 等价类理解。
\end{definition}
\symidx{\(R_jf\)：第 \(j\) 个 Riesz 变换}

由\cref{b2:thm:04-plancherel-extension}的等距性得到 \(\|R_jf\|_2\leq\|f\|_2\)，但单个方向通常不保持全部能量。把所有方向合在一起，才有
\[
 \sum_{j=1}^d\|R_jf\|_2^2=\|f\|_2^2,
 \quad \sum_{j=1}^dR_j^2f=-f.
\]
第一式来自 \(\sum_j\xi_j^2/|\xi|^2=1\)，第二式还使用两个 \(-\mathrm i\) 相乘得到的负号。若形式上把 \((-\Delta)^{1/2}\) 的乘子记为 \(2\pi|\xi|\)，则 \(R_j=-\partial_j(-\Delta)^{-1/2}\)：负号由本卷约定决定。这个记法只解释方向与微分的关系，实际定义仍是上面的有界乘子，不需先在零频率除法定义负次幂。

\ProseParagraphBreak

为了找到空间核，先保留正高度 \(t\)。Poisson 核在内部没有奇点，其共轭核将随 \(t\downarrow0\) 产生主值；这一过程还能确定核的归一化常数。

\Needspace{6\baselineskip}
\begin{proposition}[Riesz 核的正则化与主值实现]\label{b2:prop:18-riesz-realization}
设 \(d\geq1\) 为整数、\(1\leq j\leq d\)，并令
\[
 \begin{gathered}
 c_d=\frac{\Gamma\left(\frac{d+1}{2}\right)}{\pi^{(d+1)/2}},\\
 P_t(x)=c_d\frac{t}{(t^2+|x|^2)^{(d+1)/2}},
 \quad Q_{j,t}(x)=c_d\frac{x_j}{(t^2+|x|^2)^{(d+1)/2}}
 \quad(t>0).
 \end{gathered}
\]
则 \(P_t\) 的积分为一，且
\[
 \widehat P_t(\xi)=e^{-2\pi t|\xi|},
 \quad
 \widehat Q_{j,t}(\xi)=-\mathrm i\frac{\xi_j}{|\xi|}e^{-2\pi t|\xi|}
 \quad(\xi\ne0),
\]
第二式为温和分布恒等式，右端在原点任意补有限值。定义
\[
 k_j(x)=c_d\frac{x_j}{|x|^{d+1}}\quad(x\ne0),
 \quad
 R_{j,\varepsilon}f(x)=\int_{|y|>\varepsilon}k_j(y)\cdot f(x-y)\,\mathrm dy.
\]
对每个 \(f\in L^2(\mathbb R^d;\mathbb C)\)，硬截断 \(R_{j,\varepsilon}f\) 在 \(\varepsilon\downarrow0\) 时按 \(L^2\) 范数趋于乘子定义的 \(R_jf\)。\(R_j\) 具有兼容的 \(L^p\) 有界延拓，\(1<p<\infty\)，并对这些输入满足几乎处处的主值表示
\[
 R_jf(x)=\lim_{\varepsilon\downarrow0}R_{j,\varepsilon}f(x),
 \quad \|R_jf\|_p\leq C_{d,p}\|f\|_p.
\]
\end{proposition}
\begin{proof}
先核对常数。记 \(a=(d+1)/2\)，使用正实数上的 Euler 积分
\(\Gamma(a)=\int_0^\infty s^{a-1}e^{-s}\,\mathrm ds\)。代换给出
\((1+|x|^2)^{-a}=\Gamma(a)^{-1}\int_0^\infty s^{a-1}e^{-s(1+|x|^2)}\,\mathrm ds\)。Tonelli 与 Gaussian 积分于是给出
\[
 \int_{\mathbb R^d}(1+|x|^2)^{-a}\,\mathrm dx
 =\frac{\pi^{d/2}}{\Gamma(a)}
   \int_0^\infty s^{-1/2}e^{-s}\,\mathrm ds
 =\frac{\pi^{(d+1)/2}}{\Gamma(a)}.
\]
这里 \(\Gamma(1/2)=\sqrt\pi\) 也是一维 Gaussian 积分的代换形式。由 \(x=tz\) 得 \(\int P_t=1\)，并恢复\subsecref{b2:subsec:030404}的 Poisson 核。

还须在一般维数证明其 Fourier 公式；\chapref{b2:ch:03}已算出的一维公式不能直接当作高维径向公式。对 \(t>0\)，核满足 \((\partial_t^2+\Delta_x)P_t=0\)。例如令 \(D=t^2+|x|^2\)，将两部分导数相加后得到
\[
 (\partial_t^2+\Delta_x)\bigl(tD^{-a}\bigr)
 =2atD^{-a-1}\bigl(2(a+1)-d-3\bigr)=0.
\]
各阶核导数在固定的正高度区间上有共同的可积控制，故对固定 \(\xi\)，函数 \(v(t)=\widehat P_t(\xi)\) 满足
\(v''=4\pi^2|\xi|^2v\)。又有 \(|v(t)|\leq\|P_t\|_1=1\)，以及由缩放和支配收敛得到的 \(v(t)\to1\)（\(t\downarrow0\)）。当 \(\xi\ne0\) 时，有界性排除增长的指数解，初值确定剩余系数，故 \(v(t)=e^{-2\pi t|\xi|}\)；在 \(\xi=0\) 时则由总质量直接得到 \(v(t)=1\)。

共轭核满足 \(\partial_tQ_{j,t}=\partial_jP_t\)。为由这条微分关系确定 \(Q_{j,t}\)，必须同时给出无穷高度的条件。核的缩放给出
\[
 \|Q_{j,t}\|_\infty\leq C_dt^{-d},
 \quad\|\partial_jP_t\|_\infty\leq C_dt^{-d-1}.
\]
因此对每个 \(\varphi\in\mathcal S\)，有 \(\langle Q_{j,t},\varphi\rangle\to0\)（\(t\rightarrow\infty\)），并且
\[
 Q_{j,t}=-\int_t^\infty\partial_jP_s\,\mathrm ds
 \quad\text{在 }\mathcal S'\text{ 中}.
\]
时间积分对每个测试配对都绝对收敛。取 Fourier 变换后可以继续换序，因为
\[
 \int_t^\infty2\pi|\xi_j|e^{-2\pi s|\xi|}\,\mathrm ds
 =\frac{|\xi_j|}{|\xi|}e^{-2\pi t|\xi|}\leq1
 \quad(\xi\ne0),
\]
再乘任一 Schwartz 测试函数的绝对值即可积分。故整个温和分布正是声明中的有界函数乘子，没有遗漏支撑于 \(\xi=0\) 的项。只在非零频率上解微分关系而不使用无穷高度条件，不能作出这一识别。

对 \(f\in\mathcal S\)，\(Q_{j,t}*f\) 的变换因而为
\(-\mathrm i\xi_j|\xi|^{-1}e^{-2\pi t|\xi|}\cdot\widehat f\)。这里卷积与变换的分布规则来自\cref{b2:prop:05-tempered-convolution}，其输入一方为 Schwartz 函数，满足该命题条件。Plancherel 与支配收敛给出 \(Q_{j,t}*f\rightarrow R_jf\) 于 \(L^2\)。空间中，在 \(|y|<1\) 内利用奇性减去 \(f(x)\)：
\[
 Q_{j,t}*f(x)
 =\int_{|y|<1}Q_{j,t}(y)\cdot\bigl(f(x-y)-f(x)\bigr)\,\mathrm dy
  +\int_{|y|\geq1}Q_{j,t}(y)\cdot f(x-y)\,\mathrm dy.
\]
近处用 \(|Q_{j,t}(y)|\leq c_d|y|^{-d}\) 与一阶差 \(O_x(|y|)\)，远处用输入的快速衰减，均可支配取极限。所得正是 \(k_j\) 的对称球形主值。这也表明 \(Q_{j,t}\rightarrow\operatorname{pv}k_j\) 在 \(\mathcal S'\) 中成立。

上述平滑极限还不是硬截断的范数收敛。记
\[
 \begin{aligned}
 D_\varepsilon(y)&=k_j(y)\cdot\mathbf1_{\{|y|>\varepsilon\}}-Q_{j,\varepsilon}(y)\\
 &=\varepsilon^{-d}\cdot D_1(y/\varepsilon).
 \end{aligned}
\]
在 \(|y|<1\) 时，\(D_1=-Q_{j,1}\) 局部可积；在 \(|y|>1\) 时，利用 \((1+s)^{-a}=1+O(s)\) 得
\(|D_1(y)|\leq C_d|y|^{-d-2}\)。故 \(D_1\in L^1\)，并因奇性有 \(\int D_1=0\)。因此
\[
 \|D_\varepsilon*f\|_2
 \leq\int|D_1(z)|\cdot
 \|f(\,\cdot-\varepsilon z)-f\|_2\,\mathrm dz
 \longrightarrow0.
\]
平移的 \(L^2\) 连续性及控制量 \(2\|f\|_2\cdot|D_1(z)|\) 证明最后的极限。结合 \(R_{j,\varepsilon}f=Q_{j,\varepsilon}*f+D_\varepsilon*f\)，得到 Schwartz 输入的硬截断 \(L^2\) 收敛。

推广到一般 \(L^2\) 时，需要分清两种卷积界。\(Q_{j,t}\) 的乘子模至多为一，故其卷积在 Schwartz 类上满足平方界，并连续延拓到 \(L^2\)；这里使用 Plancherel，而不要求 \(Q_{j,t}\in L^1\)。对差核则可用 Young 不等式，因为 \(\|D_\varepsilon\|_1=\|D_1\|_1\)。上述分解因此给出与 \(\varepsilon\) 无关的界
\[
 \|R_{j,\varepsilon}f\|_2
 \leq\bigl(1+\|D_1\|_1\bigr)\cdot\|f\|_2.
\]
结合 Schwartz 稠密性，就将硬截断的范数收敛传给全部 \(L^2\) 输入。还须识别这些延拓与普通积分。固定 \(t,\varepsilon>0\) 时，\(Q_{j,t}\) 与 \(k_j\cdot\mathbf1_{\{|y|>\varepsilon\}}\) 都属于 \(L^2\)。若 \(f_n\to f\) 于 \(L^2\)，对这两个固定核中的任一个 \(A\)，Hölder 给出
\[
 \bigl|A*(f_n-f)(x)\bigr|\leq\|A\|_2\cdot\|f_n-f\|_2\longrightarrow0
 \quad\text{对每个 }x.
\]
取 Schwartz 逼近并与算子延拓的 \(L^2\) 收敛子列比较，便知积分与延拓几乎处处相同。

最后，核满足
\[
 |k_j(y)|\leq C_d|y|^{-d},
 \quad
 \partial_\ell k_j(y)=c_d\left(
 \frac{\delta_{j\ell}}{|y|^{d+1}}
 -(d+1)\frac{y_jy_\ell}{|y|^{d+3}}\right).
\]
这给出两个变量的标准差分条件。对紧支集 \(L^2\) 输入，在其支集外硬截断最终就是普通核积分；从已证的 \(L^2\) 收敛取几乎处处子列，便识别了支集外的 \(R_jf\)。故可用\cref{b2:thm:18-cz-lp}得到兼容的 \(L^p\) 有界性，并用\cref{b2:thm:18-maximal-truncation}控制极大截断。对光滑紧支集输入，全参数主值已逐点存在，硬截断的 \(L^2\) 收敛又将它识别为 \(R_jf\)。于是可用\subsecref{b2:subsec:180404}中\eqref{b2:eq:18-principal-value-transfer}的稠密性论证，得到所有 \(1<p<\infty\) 输入的几乎处处主值及与连续延拓的一致性。
\end{proof}

一维时 \(c_1=1/\pi\)，\(R_1\) 就是\cref{b2:thm:18-hilbert-l2}定义的 Hilbert 变换。高维核也保留同样的奇性抵消；但 \(Q_{j,t}\) 在远处的绝对积分具有 \(\int_1^\infty\mathrm dr/r\) 的发散，因而不属于 \(L^1\)。上面的分布识别与截断比较，分别说明了为何仍能定义它，以及为何得到的正是指定的硬主值。

\subsection{共轭函数}
\label{b2:subsec:180502}

先在一维取实值 \(f\in\mathcal S(\mathbb R)\)，令 \(u(x,t)=P_t*f(x)\)、\(v(x,t)=Q_{1,t}*f(x)\)。核满足
\[
 \partial_tQ_{1,t}=\partial_xP_t,
 \quad \partial_xQ_{1,t}=-\partial_tP_t,
\]
因而 \(v_t=u_x\)、\(v_x=-u_t\)。这正是 \(u+\mathrm iv\) 关于 \(z=x+\mathrm it\) 的 Cauchy--Riemann 方程，此时称 \(v\) 为 \(u\) 的\term[tiaohe-gonge]{调和共轭}[harmonic conjugate]。若改用核 \(-Q_{1,t}\)，则应同时把解析函数改为 \(u-\mathrm iv\)，不能只改其中一个符号。

\ProseParagraphBreak

对实值 \(f\in L^p(\mathbb R)\)、\(1<p<\infty\)，应把第二个分量写为 \(v=P_t*Hf\)。在平方交空间上，其乘子与 \(Q_{1,t}*f\) 一致；再用 \(L^p\) 逼近得到一般情形的一致性。固定 \(t>0\) 时两个核的导数提供充分的可积控制，故上述微分关系仍成立，且 \(u,v\) 分别在边界按 \(L^p\) 范数趋于 \(f,Hf\)。例如取 \(a>0\)、\(f=P_a\)，可以完全算出
\[
 \begin{gathered}
 u(x,t)=\frac{a+t}{\pi\bigl((a+t)^2+x^2\bigr)},
 \quad v(x,t)=\frac{x}{\pi\bigl((a+t)^2+x^2\bigr)},\\
 u+\mathrm iv=\frac{\mathrm i}{\pi(z+\mathrm ia)}.
 \end{gathered}
\]
前两式由 Poisson 乘子的乘法和共轭核公式得到。最后一个函数的唯一极点位于下半平面，直接显示上半平面中的全纯性；正的 Poisson 边界数据产生偶的实部与奇的共轭分量。

\ProseParagraphBreak

Hilbert 变换还把任意复边界数据分成正负频率。对 \(f\in L^2(\mathbb R;\mathbb C)\)，令
\[
 P_+=\frac{I+\mathrm iH}{2},\quad P_-=\frac{I-\mathrm iH}{2}.
\]
它们的乘子分别为 \(\mathbf1_{\{\xi>0\}}\)、\(\mathbf1_{\{\xi<0\}}\)，零频单点不影响这个平方空间定义。因此
\[
 P_++P_-=I,\quad P_\pm^2=P_\pm,\quad P_+P_-=P_-P_+=0.
\]
在 \(L^p(\mathbb R;\mathbb C)\)、\(1<p<\infty\) 中，同样的公式先由 \(H\) 的有界性定义有界算子，再在 \(L^2\cap L^p\) 上核对这些恒等式，并由其稠密性推广。这里的“投影”指有界线性幂等算子；平方空间中它还正交，其他指数下不使用正交性。

\ProseParagraphBreak

若 \(f\in L^2\) 的变换只支于正半轴，则 \(Hf=-\mathrm if\)，故 \(P_t*Hf=-\mathrm iP_t*f\)。\(P_t*f\) 只由正频率的 \(e^{2\pi\mathrm i\xi(x+\mathrm it)}\) 组成；对 \(t>0\)，Cauchy--Schwarz 和指数衰减使这个积分及其各阶导数绝对收敛，从而给出全纯函数。一般实输入满足 \(f+\mathrm iHf=2P_+f\)，它的 Poisson 延拓才成为上述解析函数。原来的全部数据 \(f\) 一般并非这个解析函数本身的边界值。

\ProseParagraphBreak

同一投影还能移动频率阈值。对实数 \(a\)，调制 \(\mathcal M_a f(x)=e^{2\pi\mathrm iax}\cdot f(x)\) 在每个 \(L^p\) 上等距，而 \(\mathcal M_aP_+\mathcal M_{-a}\) 在 \(L^2\) 上的乘子为 \(\mathbf1_{\{\xi>a\}}\)。因此 \(a<b\) 时两者的差给出区间 \((a,b)\) 的 Fourier 投影，并有与位置、长度无关的 \(L^p\) 界。对 \(\mathbb R^d\) 上的矩形 \(\prod_j(a_j,b_j)\)，逐坐标应用这个算子，Fubini 给出至多 \((2\|P_+\|_{p\rightarrow p})^d\) 的范数界。矩形的频率条件可以逐坐标分离；球面的边界方向随位置变化，不能用同一分解自动得到球频率投影的有界性。

\ProseParagraphBreak

在圆周上，正负整数之间还多一个真正的常数模态，不能照搬实直线的零频处理。为避免 \(2\pi\) 的歧义，以下用角变量 \(\theta\) 表示 \(\mathbb R/(2\pi\mathbb Z)\)，测度取 \(\mathrm d\theta/(2\pi)\)。这只是\chapref{b2:ch:02} \(\mathbb T=\mathbb R/\mathbb Z\) 的变量替换。

\Needspace{6\baselineskip}
\begin{proposition}[周期共轭算子]\label{b2:prop:18-periodic-hilbert}
在圆周 \(\mathbb R/(2\pi\mathbb Z)\) 上，令
\(\widehat f(n)=(2\pi)^{-1}\int_{-\pi}^\pi f(\theta)e^{-\mathrm in\theta}\,\mathrm d\theta\)。以乘子 \(-\mathrm i\operatorname{sgn}n\)（\(n=0\) 时取零）定义 \(L^2\) 算子 \(H_{\mathbb T}\)。它在每个 \(L^p\)、\(1<p<\infty\) 上具有兼容的有界延拓，且对这些输入几乎处处有
\[
 H_{\mathbb T}f(\theta)=\lim_{\varepsilon\downarrow0}
 \frac1{2\pi}\int_{\varepsilon<|\alpha|<\pi}
 \cot(\alpha/2)\cdot f(\theta-\alpha)\,\mathrm d\alpha.
\]
若 \(\Pi_0f=\widehat f(0)\) 是常数投影，则
\[
 H_{\mathbb T}^2=-I+\Pi_0,
 \quad
 \Pi_+=\frac{I+\mathrm iH_{\mathbb T}-\Pi_0}{2},
 \quad
 \Pi_-=\frac{I-\mathrm iH_{\mathbb T}-\Pi_0}{2}
\]
给出正、负整数频率投影，并满足 \(\Pi_++\Pi_-+\Pi_0=I\)。
\end{proposition}
\begin{proof}
对 \(0<r<1\)，几何级数给出一对正则核
\[
 \begin{aligned}
 P_r(\alpha)&=1+2\sum_{n\geq1}r^n\cos(n\alpha)
 =\frac{1-r^2}{1-2r\cos\alpha+r^2},\\
 Q_r(\alpha)&=2\sum_{n\geq1}r^n\sin(n\alpha)
 =\frac{2r\sin\alpha}{1-2r\cos\alpha+r^2}.
 \end{aligned}
\]
其级数绝对一致收敛，故卷积的 Fourier 系数分别乘 \(r^{|n|}\) 和 \(-\mathrm i\operatorname{sgn}(n)r^{|n|}\)。Parseval 与支配收敛说明 \(Q_r*f\rightarrow H_{\mathbb T}f\) 于 \(L^2\)。对光滑输入，\(Q_r\) 的奇性允许在积分中减去 \(f(\theta)\)；又有 \(|Q_r(\alpha)|\leq C/|\alpha|\)、\(0<|\alpha|\leq\pi\)，而函数差为 \(O(|\alpha|)\)。于是支配收敛识别极限核为 \(\cot(\alpha/2)\)，并识别为所写的对称主值。

为证明周期 \(L^p\) 界，不需另行假定圆周上的奇异积分定理。核差
\[
 b(\alpha)=\frac1{2\pi}\cot(\alpha/2)-\frac1{\pi\alpha}
 \quad(-\pi<\alpha<\pi, \alpha\ne0)
\]
在原点可连续补值，且可积。将 \(f\) 作周期延拓，再在 \((-2\pi,4\pi)\) 之外置零，得到实直线函数 \(g\)。对 \(0<\theta<2\pi\)，\(Hg(\theta)\) 中 \(|\alpha|<\pi\) 的部分恰为周期主值的 \(1/(\pi\alpha)\) 部分。其余积分的 \(|\alpha|\geq\pi\)，核有界，总量受有限个周期上的 \(L^1\) 范数控制。核差 \(b\) 的卷积则由 Minkowski 控制，因为圆周平移保持 \(L^p\) 范数。因此实直线的 \(L^p\) 有界性给出
\(\|H_{\mathbb T}f\|_p\leq C_p\|f\|_p\)，先对光滑输入成立，再由三角多项式稠密性延拓。对一般 \(L^p\) 输入，同一分解中的实直线主值几乎处处存在，两个余项绝对可积，故周期主值也存在，并与延拓一致。

平方关系和三个投影的公式先逐个整数频率核对，再由稠密性推广。特别地，\(H_{\mathbb T}1=0\)，而对正整数 \(n\)，
\(H_{\mathbb T}\cos(n\theta)=\sin(n\theta)\)、\(H_{\mathbb T}\sin(n\theta)=-\cos(n\theta)\)。常数模态使平方关系多出 \(\Pi_0\)，并非核符号的错误。
\end{proof}

若周期数据 \(f\in L^p(\mathbb R/(2\pi\mathbb Z);\mathbb R)\)、\(1<p<\infty\)，则其 Poisson 延拓与共轭延拓满足
\[
 (P_r*f)(\theta)+\mathrm i(Q_r*f)(\theta)
 =\widehat f(0)+2\sum_{n>0}\widehat f(n)r^ne^{\mathrm in\theta}.
\]
右侧是圆盘变量 \(w=re^{\mathrm i\theta}\) 的幂级数，在每个较小圆盘绝对收敛，因而全纯。常数保留在解析函数中，但其共轭分量为零；这把核计算、频率投影与圆盘解析函数联系起来。

\ProseParagraphBreak

正频率在乘法下封闭，还产生一个具体的乘积恒等式。

\Needspace{6\baselineskip}
\begin{proposition}[Hilbert 变换的乘积恒等式]\label{b2:prop:18-hilbert-product}
设 \(f,g\in L^4(\mathbb R;\mathbb C)\)，\(H\) 为本章在 \(L^4\) 与 \(L^2\) 上兼容的 Hilbert 变换。则各乘积属于 \(L^2\)，并在该空间中有
\[
 H(f\cdot Hg+g\cdot Hf)=Hf\cdot Hg-fg.
\]
\end{proposition}
\begin{proof}
先取实值 Schwartz 输入，且 Fourier 变换光滑、紧支撑并避开零频。此时 \(F=f+\mathrm iHf\)、\(G=g+\mathrm iHg\) 的变换均支于正半轴上的紧集。乘积 \(FG\) 的变换是两者的卷积，仍支于正半轴，故 \(H(FG)=-\mathrm iFG\)。写成
\[
 FG=(fg-Hf\cdot Hg)+\mathrm i(f\cdot Hg+g\cdot Hf),
\]
并利用 \(H\) 保持实值，比较虚部即得所需等式。

下面证明所用输入确实稠密，不能把低频删除当作任意 \(L^4\) 函数已经有的性质。先用实值 \(h\in\mathcal S\) 逼近实值 \(f\in L^4\)。选实偶函数 \(\chi\in C_c^\infty(\mathbb R)\)，在零点某邻域为一。对 \(R>1\)、\(0<\varepsilon<1\)，定义
\[
 \widehat h_{R,\varepsilon}(\xi)
 =\bigl(\chi(\xi/R)-\chi(\xi/\varepsilon)\bigr)\cdot\widehat h(\xi).
\]
它光滑、紧支撑且在零频附近为零，逆变换是实值 Schwartz 函数。高频截断 \(\chi(\xi/R)\) 的逆核为 \(R\check\chi(Rx)\)，积分为一且 \(L^1\) 范数一致有界，故由平移连续性得到其卷积趋于 \(h\) 于 \(L^4\)。低频项的逆核则为 \(\varepsilon\check\chi(\varepsilon x)\)，由 Young 有
\[
 \|\varepsilon\check\chi(\varepsilon\,\cdot)*h\|_4
 \leq\varepsilon^{3/4}\|\check\chi\|_4\cdot\|h\|_1
 \longrightarrow0.
\]
这里先选 \(h\in\mathcal S\subset L^1\)，才可使用这个衰减估计。依次选择初始逼近、足够大的 \(R\) 和足够小的 \(\varepsilon\)，就得到所需的实值逼近列 \(f_n,g_n\)。

Hilbert 的 \(L^4\) 连续性使 \(Hf_n\rightarrow Hf\)、\(Hg_n\rightarrow Hg\) 于 \(L^4\)。例如
\[
 \|f_n\cdot Hg_n-f\cdot Hg\|_2
 \leq\|f_n-f\|_4\cdot\|Hg_n\|_4
  +\|f\|_4\cdot\|Hg_n-Hg\|_4\longrightarrow0.
\]
其余乘积同样在 \(L^2\) 中收敛，外层 \(H\) 再用 \(L^2\) 连续性取极限，完成实输入证明。最后将复输入分成实部、虚部；等式两边关于 \(f,g\) 均为复双线性，逐项应用实输入结论即得复值版本。
\end{proof}

令 \(f=g\) 得 \(2H(f\cdot Hf)=(Hf)^2-f^2\)。它来自正频率的乘法结构，不是把 \(H\) 当作普通求导来使用 Leibniz 法则。\(L^4\) 条件使乘积落在已有强型界的 \(L^2\) 中；仅有两个 \(L^2\) 输入时，乘积只保证 \(L^1\)，不能直接沿用这个范数极限证明。

\subsection{Poisson 延拓}
\label{b2:subsec:180503}

一维共轭函数由两个分量组成。高维中，一个标量的各个切向导数需要由一个向量来配合，所以自然取 \(u=P_t*f\) 与 \(v_j=P_t*R_jf\)。Poisson 核负责内部光滑化，Riesz 变换则选择各方向的边界共轭数据。

\Needspace{6\baselineskip}
\begin{proposition}[高维共轭 Poisson 延拓]\label{b2:prop:18-poisson-conjugacy}
设整数 \(d\geq1\)、\(1<p<\infty\)，\(f\in L^p(\mathbb R^d;\mathbb C)\)，\(R_j\) 是本节的兼容 \(L^p\) Riesz 变换。对 \(t>0\)，以
\(P_t(x)=c_dt(t^2+|x|^2)^{-(d+1)/2}\)、\(c_d=\Gamma\bigl((d+1)/2\bigr)/\pi^{(d+1)/2}\) 定义
\[
 u(x,t)=P_t*f(x),\quad v_j(x,t)=P_t*R_jf(x)\quad(1\leq j\leq d).
\]
这些函数在上半空间光滑、调和，并满足
\[
 \partial_tv_j=\partial_ju,\quad
 \sum_{j=1}^d\partial_jv_j=-\partial_tu,\quad
 \partial_iv_j=\partial_jv_i\quad(1\leq i,j\leq d).
\]
此外 \(\|u(\cdot,t)\|_p\leq\|f\|_p\)、\(\|v_j(\cdot,t)\|_p\leq C_{d,p}\|f\|_p\)，并在 \(t\downarrow0\) 时分别按 \(L^p\) 范数趋于 \(f\)、\(R_jf\)。
\end{proposition}
\begin{proof}
先取 Schwartz 输入。\cref{b2:prop:18-riesz-realization}给出的乘子使三个微分关系分别化为
\[
 \begin{aligned}
 (-2\pi|\xi|)\frac{-\mathrm i\xi_j}{|\xi|}
 &=2\pi\mathrm i\xi_j,\\
 \sum_j(2\pi\mathrm i\xi_j)\frac{-\mathrm i\xi_j}{|\xi|}
 &=2\pi|\xi|,\\
 (2\pi\mathrm i\xi_i)\frac{-\mathrm i\xi_j}{|\xi|}
 &=(2\pi\mathrm i\xi_j)\frac{-\mathrm i\xi_i}{|\xi|}.
 \end{aligned}
\]
每项都还乘有 \(e^{-2\pi t|\xi|}\cdot\widehat f\)。第一式是切向与法向导数的配合，第二式把向量散度与标量法向变化配对，第三式表示向量场无旋。核满足的调和方程同时给出所有分量的调和性。

对一般 \(L^p\) 输入，固定 \(0<a\leq t\leq b<\infty\)。Poisson 核的任意时空导数均有共同的 \(L^1\) 及 \(L^{p'}\) 控制，并在这些范数中随参数连续；对卷积用 Hölder 与差商极限，可逐阶求导，得到内部光滑性。用 \(L^p\) 中的 Schwartz 逼近以及 \(R_j\) 的 \(L^p\) 连续性，将刚才的关系传给极限。范数界由 \(\|P_t\|_1=1\) 及 Young 得到。边界范数极限则对 \(f\) 和 \(R_jf\) 分别使用\cref{b2:thm:03-approximate-identity}；两者都属于有限指数 \(L^p\)，满足平移连续性这一关键条件。
\end{proof}

这些关系没有把 \(R_j\) 本身变成平滑算子。它的乘子只改方向和相位，并不压低高频；正高度的 \(e^{-2\pi t|\xi|}\) 才压低高频，产生内部光滑化。边界共轭数据与内部正则性由不同因素控制，因而不能从 \(R_j\) 有界性单独推断输入具有经典导数。

\subsection{椭圆型偏微分方程应用}
\label{b2:subsec:180504}

Poisson 方程中的另一个问题是：右端只有 \(L^p\) 大小，解的二阶导数能否仍由同一指数控制？\chapref{b2:ch:13}已用密度的 Hölder 增量消去 Newton 二阶核的奇性；现在使用本章的奇异积分界，得到不要求右端连续的范数估计。首先对光滑紧支集源识别算子，再用磨光处理弱解。

\Needspace{6\baselineskip}
\begin{proposition}[Newton 势的二阶导数]\label{b2:prop:18-newton-second-derivatives}
设整数 \(d\geq1\)，\(f\in C_c^\infty(\mathbb R^d;\mathbb C)\)。取\cref{b2:thm:05-laplace-fundamental}的标准径向基本解
\[
 \Phi_d(x)=
 \begin{cases}
 \dfrac{|x|^{2-d}}{(d-2)\sigma_{d-1}},&d\geq3,\\
 -(2\pi)^{-1}\log|x|,&d=2,\\
 -|x|/2,&d=1,
 \end{cases}
 \quad \sigma_{d-1}=|S^{d-1}|,
\]
并令 \(u=\Phi_d*f\)。则对 \(1\leq i,j\leq d\)，在温和分布以及每个 \(L^p\)、\(1<p<\infty\) 中有
\[
 \partial_i\partial_j u=R_iR_jf,
 \quad \|D^2u\|_p\leq C_{d,p}\|f\|_p.
\]
这里 \(|D^2u|=(\sum_{i,j}|\partial_i\partial_ju|^2)^{1/2}\)。在平方指数处还有等式
\[
 \|D^2u\|_2^2=\sum_{i,j}\|\partial_i\partial_ju\|_2^2=\|f\|_2^2.
\]
特别地，对任意 \(w\in C_c^\infty(\mathbb R^d;\mathbb C)\)，有
\(\|D^2w\|_p\leq C_{d,p}\|\Delta w\|_p\)。
\end{proposition}
\begin{proof}
\Cref{b2:prop:05-fundamental-convolution}保证 \(u\) 光滑且 \(-\Delta u=f\)，\cref{b2:prop:05-tempered-convolution}的增长估计又保证 \(u\) 温和。不能直接把它的变换写成普通函数 \(\widehat f/(4\pi^2|\xi|^2)\)：二维和一维的 \(|\xi|^{-2}\) 在零点不可积，而分布方程还可能允许支撑于零频率的附加项。

先在空间一侧证明 \(w_{ij}=\partial_i\partial_ju\in L^2\)。在任一有界区域上，由光滑性得到局部有界；选 \(A>0\) 使 \(\operatorname{supp}f\subset B(0,A)\)。当 \(d\geq2\)、\(|x|>2A\) 时可在普通积分中求导，且标准基本解满足 \(|D^2\Phi_d(x-y)|\leq C_d|x|^{-d}\)，所以
\[
 |w_{ij}(x)|\leq C_d|x|^{-d}\cdot\|f\|_1.
\]
其平方在远处可积。一维则直接由方程得到 \(u''=-f\in L^2\)。因此所有维数下 \(w_{ij}\) 都属于 \(L^2\)，这一步没有使用待证的奇异积分估计。

对 \(-\Delta w_{ij}=\partial_i\partial_jf\) 作温和分布 Fourier 变换，使用\eqref{b2:eq:05-distribution-operations}，得到
\[
 4\pi^2|\xi|^2\widehat w_{ij}
 =-4\pi^2\xi_i\xi_j\widehat f.
\]
只在 \(\xi\ne0\) 的开集上除法，得到
\(\widehat w_{ij}=-\xi_i\xi_j|\xi|^{-2}\cdot\widehat f\)。两边都已经是 \(L^2\) 函数，所以它们的差即使支撑于零点，也只能为零：一个零测集合不支持非零 \(L^2\) 等价类。由此排除了所有可能的零频分布歧义，得到整个空间中的等式。两个 Riesz 符号相乘正是 \(-\xi_i\xi_j/|\xi|^2\)，故 \(w_{ij}=R_iR_jf\)。

各 \(R_j\) 的兼容 \(L^p\) 有界性给出每个二阶导数的估计，有限求和得到 Hesse 矩阵的范数界。在 \(p=2\) 时直接计算
\[
 \sum_{i,j}\|w_{ij}\|_2^2
 =\int_{\mathbb R^d}
 \frac{\sum_{i,j}\xi_i^2\xi_j^2}{|\xi|^4}
 \cdot|\widehat f(\xi)|^2\,\mathrm d\xi
 =\|f\|_2^2.
\]
最后若 \(w\in C_c^\infty\)，取 \(f=-\Delta w\)。紧支集允许把导数在分布卷积中移到基本解上，故
\(\Phi_d*(-\Delta w)=(-\Delta\Phi_d)*w=\delta_0*w=w\)。对这个 Newton 势使用刚才的估计，得到最后一项，而无需声称全空间 Poisson 方程在所有光滑函数中唯一。
\end{proof}

这个识别还恢复了空间核中的局部项。对 \(d\geq2\)，\cref{b2:lem:13-hessian-kernel}已经证明
\[
 \partial_i\partial_j\Phi_d
 =\operatorname{pv}K_{ij}-\frac{\delta_{ij}}d\delta_0,
 \quad
 K_{ij}(x)=\frac1{\sigma_{d-1}}\left(
 \frac{d x_ix_j}{|x|^{d+2}}-\frac{\delta_{ij}}{|x|^d}\right).
\]
因此对上述光滑输入，\(R_iR_jf=\operatorname{pv}K_{ij}*f-(\delta_{ij}/d)f\)。核 \(K_{ij}\) 的球面平均为零；对 \(i=j\) 求和，主值核的迹消失，留下 \(\Delta u=-f\)。这个空间点质量不能删除，也不是 Fourier 乘子在零频率任意补一个数的问题。它对应算子中真实的局部乘法项。

\ProseParagraphBreak

估计控制的是二阶导数，不是整个非齐次 \(W^{2,p}\) 范数。二维若 \(\int f\ne0\)，\chapref{b2:ch:05}的远场计算给出 \(u(x)=-(\int f)\log|x|/(2\pi)+O(|x|^{-1})\)，所以 \(u\) 本身不在任何有限指数的全空间 \(L^p\) 中。基本解的标准归一化也不能任意替换：加调和多项式虽不改变方程，却可能改变二阶导数及其全空间可积性。

\ProseParagraphBreak

对局部弱解，必须先把全空间估计应用于真正满足其条件的紧支集函数。截断产生的一阶交换项也要保留。

\Needspace{6\baselineskip}
\begin{proposition}[内部二阶估计]\label{b2:prop:18-interior-w2p}
设整数 \(d\geq1\)、\(1<p<\infty\)、\(r>0\)、\(x_0\in\mathbb R^d\)，并记 \(B_s=B(x_0,s)\)。若 \(u\in W^{1,p}(B_{2r};\mathbb C)\)、\(f\in L^p(B_{2r};\mathbb C)\)，且 \(-\Delta u=f\) 在 \(B_{2r}\) 中按分布成立，则 \(u\in W^{2,p}(B_r)\)，并有
\[
 \|D^2u\|_{L^p(B_r)}
 \leq C_{d,p}\left(
 \|f\|_{L^p(B_{2r})}
 +r^{-1}\|\nabla u\|_{L^p(B_{2r})}
 +r^{-2}\|u\|_{L^p(B_{2r})}\right).
\]
常数不依赖 \(u,f,x_0,r\)。
\end{proposition}
\begin{proof}
取实值 \(\eta\in C_c^\infty(B_{2r})\)，在 \(B_r\) 为一，且 \(|\nabla\eta|\leq C_d/r\)、\(|D^2\eta|\leq C_d/r^2\)。将 \(v=\eta\cdot u\) 在球外延为零，得到紧支集的全空间 \(W^{1,p}\) 函数。分布乘积法则给出
\[
 -\Delta v=\eta\cdot f-2\nabla\eta\cdot\nabla u-u\cdot\Delta\eta
 \eqqcolon F.
\]
由于截断的支集紧含于球内，延零不会产生边界分布。右端是紧支集的 \(L^p\) 函数，并满足
\[
 \|F\|_p\leq C_d\left(
 \|f\|_{L^p(B_{2r})}
 +r^{-1}\|\nabla u\|_{L^p(B_{2r})}
 +r^{-2}\|u\|_{L^p(B_{2r})}\right).
\]

此时尚未知道 \(D^2v\in L^p\)，不能直接把 \(v\) 当作前一命题中的光滑输入。取非负 \(\rho\in C_c^\infty\)、\(\int\rho=1\)，令 \(\rho_\varepsilon(x)=\varepsilon^{-d}\rho(x/\varepsilon)\)，并置 \(v_\varepsilon=\rho_\varepsilon*v\)、\(F_\varepsilon=\rho_\varepsilon*F\)。它们均光滑紧支撑，满足 \(-\Delta v_\varepsilon=F_\varepsilon\)。由前一命题，
\[
 \|D^2v_\varepsilon\|_p\leq C_{d,p}\|F_\varepsilon\|_p
 \leq C_{d,p}\|F\|_p.
\]
还可以对 \(v_\varepsilon-v_\delta\) 使用同一估计，得到
\[
 \|D^2v_\varepsilon-D^2v_\delta\|_p
 \leq C_{d,p}\|F_\varepsilon-F_\delta\|_p\longrightarrow0.
\]
逼近恒等给出 \(F_\varepsilon\rightarrow F\) 于 \(L^p\)，所以 Hesse 矩阵在 \(L^p\) 中有强极限。同时 \(v_\varepsilon\rightarrow v\) 于 \(L^p\)，故对紧支集测试函数分部积分后取极限，该矩阵极限正是 \(D^2v\) 的弱导数。于是 \(v\in W^{2,p}\)，保留同一个范数界；在 \(B_r\) 中 \(v=u\)，得到结论。
\end{proof}

这一结论的起点是指定的 \(W^{1,p}\) 解。若最初只有 \(u\in H^1\)、\(f\in L^2\)，可以直接取 \(p=2\)；要得到更高指数，仍须先改善右端和一阶导数的可积性。截断公式本身不提供任意 \(p\) 的自举起点。以上也都是恒系数算子的内部估计，边界正则性、变系数及端点 \(p=1,\infty\) 还需各自的条件与论证。


\begin{chaptersummary}
Calderón--Zygmund 分解在每个坏立方体上保留平均值，把剩余部分变成零平均函数。远离这个立方体时，核可以减去中心值；核差分比核本身多衰减一个正幂次，于是坏部的远场可以绝对积分。好部由平方估计控制，坏区由体积估计控制，这三项合起来得到弱 \((1,1)\)。再经插值与伴随，标准 Calderón--Zygmund 算子在全部 \(1<p<\infty\) 上有兼容的延拓。

\ProseParagraphBreak

使用这一结论时，还要分别核对核条件、\(L^2\) 有界性和截断实现。离支集核并不决定算子在对角线上的作用，\(L^p\) 范数延拓也不自动等于逐点主值。极大截断提供统一的误差控制，具体核在光滑稠密类上的收敛提供起点，两者结合才把主值传给一般输入。Hilbert 与 Riesz 变换满足这些条件，它们的零阶 Fourier 符号又把共轭调和结构和二阶椭圆估计连接起来。区域问题仍须经过截断和局部化，所得估计中的低阶项不能在没有额外条件时删去。
\end{chaptersummary}

\BookExercises


% 自拟：由区间模型精算分布函数与弱范数，训练端点的定量判断。
\Needspace{6\baselineskip}
\begin{exercise}\label{b2:ex:18-new-1}
在带 Lebesgue 测度的 \(\mathbb R\) 上，令
\[
 h=H\mathbf1_{(0,1)},\quad
 h(x)=\frac1\pi\log\left|\frac{x}{x-1}\right|\quad(x\ne0,1),
\]
并在两个端点任意补有限值。对 \(t>0\)，求出 \(\{|h|>t\}\) 的区间组成，证明其分布函数与弱 \(L^1\) 准范数分别为
\[
 d_h(t)=\frac2{\sinh(\pi t)},
 \quad \|h\|_{L^{1,\infty}}
       \coloneqq\sup_{t>0}t\cdot d_h(t)=\frac2\pi.
\]
分别解释 \(t\downarrow0\) 和 \(t\uparrow\infty\) 的分布行为，并据此判断 \(h\) 的有限指数可积性。
\end{exercise}
\begin{solution}
固定 \(t>0\)，记 \(q=e^{\pi t}>1\)。条件 \(|h(x)|>t\) 等价于
\[
 \left|\frac{x}{x-1}\right|<q^{-1}
 \quad\text{或}\quad
 \left|\frac{x}{x-1}\right|>q.
\]
在 \(x<0\) 上，比值的绝对值为 \(-x/(1-x)<1\)，第一项不等式给出
\(-1/(q-1)<x<0\)。在 \(0<x<1\) 上，比值为 \(x/(1-x)\)，两项分别给出
\(0<x<1/(q+1)\) 和 \(q/(q+1)<x<1\)。在 \(x>1\) 上，比值为 \(x/(x-1)>1\)，第二项给出
\(1<x<q/(q-1)\)。因此，忽略有限多个点后，超水平集就是
\[
 \left(-\frac1{q-1},\frac1{q+1}\right)
 \ \cup\
 \left(\frac q{q+1},\frac q{q-1}\right).
\]
两段长度之和为
\[
 d_h(t)=2\left(\frac1{q-1}+\frac1{q+1}\right)
       =\frac{4q}{q^2-1}=\frac2{\sinh(\pi t)}.
\]
由 \(\sinh s\geq s\)（\(s>0\)），有 \(t\cdot d_h(t)\leq2/\pi\)；又由
\(\sinh s/s\to1\)（\(s\downarrow0\)），这个上界在极限中达到。所以弱范数等于 \(2/\pi\)，其上确界无需在某个正高度取到。

更具体地，低高度时 \(d_h(t)\sim2/(\pi t)\)，记录的是远处 \(1/x\) 尾部占据的长区间；高高度时 \(d_h(t)\sim4e^{-\pi t}\)，记录的是两个对数奇点附近迅速缩小的区域。用分布函数恒等式
\[
 \|h\|_p^p=p\int_0^\infty t^{p-1}\cdot d_h(t)\,\mathrm dt
 \quad(0<p<\infty),
\]
无穷端因指数衰减总可积，零端则等价于 \(\int_0^1t^{p-2}\,\mathrm dt\)。故恰在 \(p>1\) 时 \(h\in L^p\)，在 \(0<p\leq1\) 时不属于 \(L^p\)。此外高高度的分布函数始终为正，故 \(h\notin L^\infty\)。
\end{solution}

% 自拟：将 Hilbert 变换的代数恒等式与正负频率结合，判断可逆性及精确平方范数。
\Needspace{6\baselineskip}
\begin{exercise}\label{b2:ex:18-new-2}
设 \(1<p<\infty\)，\(a,b\in\mathbb C\)，\(H\) 取本章兼容的 \(L^p(\mathbb R;\mathbb C)\) 延拓。证明 \(A=aI+bH\) 为可逆有界算子，当且仅当 \(a^2+b^2\ne0\)，并在可逆时求其逆。再在 \(L^2\) 上证明
\[
 \|A\|_{2\to2}=\max\{|a-\mathrm ib|,|a+\mathrm ib|\},
 \quad
 \|A^{-1}\|_{2\to2}
 =\frac1{\min\{|a-\mathrm ib|,|a+\mathrm ib|\}}.
\]
必要性与范数下界须给出实际输入，不能只在一个频率点上检验符号。
\end{exercise}
\begin{solution}
先将正文的平方恒等式传到所用空间。在 \(L^2\) 中已有 \(H^2=-I\)。对 \(\varphi\in\mathcal S\)，其 \(L^2\) 与 \(L^p\) 实现相容，且 \(H\varphi\in L^2\cap L^p\)，所以该恒等式也在 \(L^p\) 中成立。对一般 \(f\in L^p\)，取 \(\varphi_n\in\mathcal S\) 在 \(L^p\) 中逼近；两次使用 \(H\) 的有界性，就得到 \(H^2f=-f\)。因此
\[
 (aI+bH)(aI-bH)=(aI-bH)(aI+bH)=(a^2+b^2)I.
\]
若 \(a^2+b^2\ne0\)，则
\[
 A^{-1}=\frac{aI-bH}{a^2+b^2}
\]
是所需的有界双侧逆。

由于 \(a^2+b^2=(a-\mathrm ib)(a+\mathrm ib)\)，退化时至少有一个因子为零。取非零 \(\psi_+\in C_c^\infty\bigl((1,2)\bigr)\)，令 \(f_+=\mathcal F^{-1}\psi_+\in\mathcal S\)。其频率全为正，故
\(Hf_+=-\mathrm if_+\)，从而 \(Af_+=(a-\mathrm ib)f_+\)。同样取非零频率测试函数支撑于 \((-2,-1)\)，得到 \(f_-\in\mathcal S\) 且
\(Af_-=(a+\mathrm ib)f_-\)。每个 Schwartz 函数都属于当前的 \(L^p\)，所以任一因子为零都使 \(A\) 有非零核，不能可逆。

在 \(L^2\) 中，Plancherel 给出
\[
 \|Af\|_2^2
 =|a-\mathrm ib|^2\int_{\xi>0}|\widehat f(\xi)|^2\,\mathrm d\xi
  +|a+\mathrm ib|^2\int_{\xi<0}|\widehat f(\xi)|^2\,\mathrm d\xi.
\]
这给出最大值所对应的上界，而选取 \(f_+\) 或 \(f_-\) 就得到相应下界。可逆时，逆算子在两个半轴上的符号分别为
\((a-\mathrm ib)^{-1}\) 和 \((a+\mathrm ib)^{-1}\)，用同样的输入得到其精确范数。正、负半轴都是正测度集，所以这些范数由真正的 \(L^2\) 函数实现。
\end{solution}

% 自拟：区分算子范数最优与非零函数能否取等，训练频率集中构造。
\Needspace{6\baselineskip}
\begin{exercise}\label{b2:ex:18-new-3}
设整数 \(d\geq2\)、\(1\leq j\leq d\)，\(R_j\) 为乘子
\(-\mathrm i\xi_j/|\xi|\) 定义的 \(L^2(\mathbb R^d;\mathbb C)\) Riesz 变换。证明
\(\|R_j\|_{2\to2}=1\)，但每个非零 \(f\in L^2\) 都满足
\(\|R_jf\|_2<\|f\|_2\)。构造范数为一的 Schwartz 函数族，使其输出范数趋于一。
\end{exercise}
\begin{solution}
由 Plancherel，
\[
 \|f\|_2^2-\|R_jf\|_2^2
 =\int_{\mathbb R^d}
   \frac{\sum_{k\ne j}\xi_k^2}{|\xi|^2}
       \cdot|\widehat f(\xi)|^2\,\mathrm d\xi\geq0.
\]
权重只在第 \(j\) 条坐标轴上为零，原点处的补值不影响积分。由于 \(d\geq2\)，这条轴的 \(d\) 维 Lebesgue 测度为零。若 \(f\ne0\)，则 \(|\widehat f|>0\) 的集合有正测度，故上式严格为正。特别地，范数等式不能由某个非零输入取到。

为逼近上界，记 \(e_j\) 为第 \(j\) 个单位坐标向量，取非零
\(\psi\in C_c^\infty\bigl(B(0,1)\bigr)\)。对 \(0<\varepsilon<1/2\)，定义
\[
 \widehat f_\varepsilon(\xi)
 =\frac{\varepsilon^{-d/2}}{\|\psi\|_2}
       \cdot\psi\left(\frac{\xi-e_j}{\varepsilon}\right).
\]
该函数光滑且紧支撑，所以 \(f_\varepsilon\in\mathcal S\)；换元给出
\(\|f_\varepsilon\|_2=1\)。在其频率支集上，\(\xi_j\geq1-\varepsilon\)、\(|\xi|\leq1+\varepsilon\)，所以
\[
 \left(\frac{1-\varepsilon}{1+\varepsilon}\right)^2
 \leq\|R_jf_\varepsilon\|_2^2\leq1.
\]
令 \(\varepsilon\downarrow0\)，输出范数便趋于一。这说明最优算子范数可以由集中族逼近，即使没有任何非零取等函数。维数条件也有实质作用：在一维，\(R_1=H\) 在全部 \(L^2\) 上等距。
\end{solution}

% 自拟：将坏块零平均推广到有限阶矩消失，并用有限差分构造精确阶数。
\Needspace{6\baselineskip}
\begin{exercise}\label{b2:ex:18-new-4}
设 \(m\geq0\) 为整数、\(R>0\)，\(b\in L^1(\mathbb R;\mathbb C)\) 在 \([-R,R]\) 外几乎处处为零。定义
\[
 \mu_\ell\coloneqq\int_{\mathbb R}y^\ell\cdot b(y)\,\mathrm dy
 \quad(0\leq\ell\leq m),
\]
并假设 \(\mu_\ell=0\) 对 \(0\leq\ell<m\) 成立。当 \(m=0\) 时这一组消失条件为空。此题只在支集外定义普通绝对积分
\[
 Hb(x)\coloneqq\frac1\pi\int_{\mathbb R}\frac{b(y)}{x-y}\,\mathrm dy
 \quad(|x|>R),
\]
不预设一般 \(L^1\) 输入在支集内的主值。证明当 \(|x|>2R\) 时，
\[
 \left|Hb(x)-\frac{\mu_m}{\pi x^{m+1}}\right|
 \leq\frac{2R^{m+1}}{\pi|x|^{m+2}}\|b\|_1.
\]
再对
\[
 b_m(y)=\sum_{k=0}^m(-1)^k\binom mk\cdot\mathbf1_{(k,k+1)}(y)
\]
证明前 \(m\) 个矩为零，而 \(\mu_m=(-1)^m m!\)，从而该尾部阶数可以实际出现。
\end{exercise}
\begin{solution}
紧支集保证每个所用的矩都绝对收敛。在 \(|x|>2R\)、\(|y|\leq R\) 时，有限几何展开给出恒等式
\[
 \frac1{x-y}=\sum_{\ell=0}^m\frac{y^\ell}{x^{\ell+1}}
                  +\frac{y^{m+1}}{x^{m+1}(x-y)}.
\]
乘 \(b(y)\) 积分，前 \(m\) 项因矩条件消失。余项中
\(|x-y|\geq|x|/2\)，故
\[
 \begin{aligned}
 \left|Hb(x)-\frac{\mu_m}{\pi x^{m+1}}\right|
 &\leq\frac2{\pi|x|^{m+2}}
             \int|y|^{m+1}\cdot|b(y)|\,\mathrm dy\\
 &\leq\frac{2R^{m+1}}{\pi|x|^{m+2}}\|b\|_1.
 \end{aligned}
\]
估计对正、负两侧的 \(x\) 都成立，首项的符号保留在 \(x^{m+1}\) 中。

对给定的 \(b_m\)，它支撑于 \([0,m+1]\)，因而属于 \(L^1\)。将每个区间换元为 \(y=z+k\)，得到
\[
 \int y^\ell\cdot b_m(y)\,\mathrm dy
 =\int_0^1\sum_{k=0}^m(-1)^k\binom mk(z+k)^\ell\,\mathrm dz.
\]
记前向差分为 \(\Delta P(z)=P(z+1)-P(z)\)。反复展开可得
\[
 \Delta^mP(z)=\sum_{k=0}^m(-1)^{m-k}\binom mk\cdot P(z+k).
\]
差分每次将多项式次数降低一，且把最高次系数乘以原来的次数。因此对 \(P(z)=z^\ell\)，\(\ell<m\) 时 \(\Delta^mP=0\)，\(\ell=m\) 时 \(\Delta^mP=m!\)。上面的矩积分于是分别为零和 \((-1)^m m!\)；\(m=0\) 时同样得到 \(\mu_0=1\)。由于第 \(m\) 个矩非零，所得尾部的首项确为 \(x^{-(m+1)}\) 阶，而非仅有一个可能不尖锐的上界。
\end{solution}

% 自拟：从离支集局部性识别对角线自由度，补足零核乘法例子的结构解释。
\Needspace{6\baselineskip}
\begin{exercise}\label{b2:ex:18-new-5}
设整数 \(d\geq1\)，\(T_1,T_2\) 是 \(L^2(\mathbb R^d;\mathbb C)\) 上的有界复线性算子。设可测核 \(K\) 在 \(x\ne y\) 上满足
\(|K(x,y)|\leq A|x-y|^{-d}\)，其中 \(A<\infty\)；并且对每个紧集 \(F\) 与每个在 \(F\) 外几乎处处为零的 \(f\in L^2\)，有
\[
 T_1f(x)=T_2f(x)=\int_FK(x,y)\cdot f(y)\,\mathrm dy
 \quad\text{对几乎处处的 }x\notin F.
\]
证明差算子 \(S=T_1-T_2\) 必为乘法算子：存在唯一的 \(a\in L^\infty(\mathbb R^d)\) 等价类，使
\[
 Sf=a\cdot f\quad(f\in L^2),
 \quad\|S\|_{2\to2}=\|a\|_\infty.
\]
证明中须从紧支集输入出发，解释为什么 \(S\) 对任意可测集合都保持支撑，而不能只对开集作形式判断。
\end{exercise}
\begin{solution}
若 \(f\) 支撑于紧集 \(F\)，共同核表示相减便知 \(Sf=0\) 在 \(F\) 外几乎处处成立。先把这个局部性推广到任意 Lebesgue 可测集 \(E\)。设 \(f\in L^2\) 在 \(E\) 外几乎处处为零，选定有限值可测代表，并令
\[
 E_n=E\cap B(0,n)\cap\{|f|\leq n\}.
\]
这些集合递增且覆盖 \(f\) 的非零部分，至多相差零测集。由有限测度集合的内正则性，选紧集 \(K_n\subset E_n\)，使
\(|E_n\setminus K_n|<n^{-4}\)。令 \(f_n=f\cdot\mathbf1_{K_n}\)，则
\[
 \|f-f_n\|_2^2
 \leq\int_{E\setminus E_n}|f|^2+n^2|E_n\setminus K_n|
 \leq\int_{E\setminus E_n}|f|^2+n^{-2}\longrightarrow0.
\]
每个 \(Sf_n\) 都在 \(K_n\subset E\) 外为零，且 \(Sf_n\to Sf\) 于 \(L^2\)。乘以 \(\mathbf1_{E^c}\) 是有界算子，故极限 \(Sf\) 也在 \(E\) 外为零。这里不要求 \(K_n\) 递增。

对任意 \(f\in L^2\)，将它分成 \(\mathbf1_E\cdot f\) 与 \(\mathbf1_{E^c}\cdot f\)，两部分的输出分别留在 \(E\) 与 \(E^c\)。因此
\[
 S(\mathbf1_E\cdot f)=\mathbf1_E\cdot Sf.
\]
现在令 \(D_n=(-n,n)^d\)、\(a_n=S\mathbf1_{D_n}\)。当 \(m<n\) 时，
\[
 \mathbf1_{D_m}\cdot a_n
 =S(\mathbf1_{D_m}\cdot\mathbf1_{D_n})
 =S\mathbf1_{D_m}=a_m.
\]
取去掉可数个零测例外集的共同满测集，这些函数就在各 \(D_n\) 上相容，可以拼成一个可测函数 \(a\)，满足 \(a=a_n\) 于 \(D_n\)。对每个可测 \(E\subset D_n\)，又有
\[
 S\mathbf1_E=\mathbf1_E\cdot a,
 \quad\|a\cdot\mathbf1_E\|_2
      \leq\|S\|_{2\to2}\cdot|E|^{1/2}.
\]
若 \(\{|a|>\|S\|_{2\to2}+\varepsilon\}\cap D_n\) 有正测度，以上式中的 \(E\) 取为此集合，平方后就与 \(L^2\) 积分下界矛盾。因此对每个 \(n\) 和 \(\varepsilon>0\)，该集合测度为零，得到 \(\|a\|_\infty\leq\|S\|_{2\to2}\)。

由线性，\(Sf=a\cdot f\) 对支撑在某个 \(D_n\) 内的简单函数成立。这样的简单函数在 \(L^2\) 中稠密：先截去空间远尾和过高幅度，再将有界的复值范围作有限网格近似。两边均为有界算子，故等式推广到全部 \(L^2\)。反向不等式
\(\|S\|_{2\to2}=\|M_a\|_{2\to2}\leq\|a\|_\infty\)
来自直接积分，于是范数相等。若另一个 \(\widetilde a\) 给出相同乘法算子，作用于各 \(\mathbf1_{D_n}\) 即得 \(a=\widetilde a\) 几乎处处，证明唯一性。可见离支集核所遗漏的，恰好是这样的有界乘法项。
\end{solution}

% 自拟：组合方向乘子得到向量投影，连接奇异积分与无散条件。
\Needspace{6\baselineskip}
\begin{exercise}\label{b2:ex:18-new-6}
设整数 \(d\geq2\)、\(1<p<\infty\)，\(R_j\) 取本章兼容的 \(L^p(\mathbb R^d;\mathbb C)\) Riesz 变换。对向量函数 \(F=(F_1,\ldots,F_d)\in L^p(\mathbb R^d;\mathbb C^d)\)，定义
\[
 (\mathbb PF)_i=F_i+\sum_{j=1}^dR_iR_jF_j
 \quad(1\leq i\leq d),
\]
其中向量范数采用逐点欧氏范数。证明 \(\mathbb P\) 在该向量值 \(L^p\) 空间上有界，\(\mathbb P^2=\mathbb P\)，且
\(\operatorname{div}(\mathbb PF)=0\) 于分布意义；对每个 \(\varphi\in\mathcal S(\mathbb R^d)\)，还有 \(\mathbb P\nabla\varphi=0\)。再证明在 \(p=2\) 时，\(\mathbb P\) 是到闭子空间
\[
 \mathcal V=\{F\in L^2(\mathbb R^d;\mathbb C^d):
                         \operatorname{div}F=0\text{ 于分布意义}\}
\]
的正交投影，且 \(\|\mathbb P\|_{2\to2}=1\)。
\end{exercise}
\begin{solution}
先看有界性。每个分量的 \(L^p\) 范数不超过 \(\|F\|_{L^p(\mathbb C^d)}\)，且 \(R_i,R_j\) 均为有界 \(L^p\) 算子。有限次使用三角不等式，例如
\[
 \begin{aligned}
 \|\mathbb PF\|_{L^p(\mathbb C^d)}
 &\leq\sum_{i=1}^d\|(\mathbb PF)_i\|_p\\
 &\leq\sum_i\|F_i\|_p
   +\sum_{i,j}\|R_i\|_{p\to p}\cdot\|R_j\|_{p\to p}\cdot\|F_j\|_p\\
 &\leq C_{d,p}\|F\|_{L^p(\mathbb C^d)},
 \end{aligned}
\]
就给出足够的界，不要求这个常数最优。

在 \(L^2\) 中，\(R_iR_j\) 的符号为 \(-\xi_i\xi_j/|\xi|^2\)，所以 \(\mathbb P\) 的矩阵符号是
\[
 P(\xi)=I-\frac{\xi\xi^{\mathsf T}}{|\xi|^2}
 \quad(\xi\ne0).
\]
它是到 \(\xi^\perp\subset\mathbb C^d\) 的欧氏正交投影，因此
\(P(\xi)^2=P(\xi)=P(\xi)^*\)、\(\xi^{\mathsf T}P(\xi)=0\)，且矩阵范数不超过一。Plancherel 把这些等式传给 \(\mathbb P\)，得到幂等性、自伴性、平方范数上界及无散性。

还须识别投影的像。对 \(F\in L^2\)，分布恒等式 \(\operatorname{div}F=0\) 等价于
\(\xi\cdot\widehat F(\xi)=0\) 几乎处处：Fourier 变换将散度化为 \(2\pi\mathrm i\xi\cdot\widehat F\)，后者局部可积，一个局部可积函数给出零分布当且仅当它几乎处处为零。这样的 \(F\) 满足 \(P(\xi)\widehat F=\widehat F\)，即 \(\mathbb PF=F\)。而任意输出已经无散，所以像恰为 \(\mathcal V\)。若一列无散向量在 \(L^2\) 中收敛，与每个紧支集光滑测试函数的梯度配对都收敛，极限仍无散，故 \(\mathcal V\) 闭。结合自伴与幂等，得到所述正交投影。

范数确实为一，而不只是至多一。取 \(\psi\in C_c^\infty(\mathbb R^d)\)，使
\((\partial_1\psi,\partial_2\psi)\) 不恒为零，并令
\[
 G=(\partial_2\psi,-\partial_1\psi,0,\ldots,0).
\]
混合偏导相等使 \(\operatorname{div}G=0\)，且 \(G\ne0\)，故 \(\mathbb PG=G\) 给出下界一。另一方面，\(\widehat{\nabla\varphi}=2\pi\mathrm i\xi\cdot\widehat\varphi\) 沿 \(\xi\) 方向，矩阵 \(P(\xi)\) 将其消去，所以 \(\mathbb P\nabla\varphi=0\) 在 \(L^2\) 中成立。

最后处理一般 \(p\)。对 Schwartz 向量，上述等式先在 \(L^2\) 成立；它和其中的 \(\mathbb P\) 输出均属于 \(L^2\cap L^p\)，各个 Riesz 变换的相容性使幂等式也在 \(L^p\) 中成立。取 Schwartz 向量在 \(L^p\) 中逼近任意 \(F\)，利用有界性即可传递 \(\mathbb P^2=\mathbb P\)。对每个紧支集光滑测试函数，其梯度属于 \(L^{p'}\)，所以输出的无散分布等式也能在该极限中传递。梯度消去公式同样由相容性成立。这里的正交性与范数一专指 \(p=2\)，未对其他指数作相同断言。
\end{solution}

% 自拟：把热核积分、预解方程与二阶 Riesz 乘子连接，并区分有质量与零质量估计。
\Needspace{6\baselineskip}
\begin{exercise}\label{b2:ex:18-new-7}
设整数 \(d\geq1\)、\(1<p<\infty\)、\(\lambda>0\)，\(f\in L^p(\mathbb R^d;\mathbb C)\)。令
\[
 H_t(x)=(4\pi t)^{-d/2}e^{-|x|^2/(4t)},
 \quad u_\lambda=\int_0^\infty e^{-\lambda t}\cdot(H_t*f)\,\mathrm dt,
\]
时间积分按 \(L^p\) 范数中的积分理解。证明它收敛，且
\[
 (\lambda-\Delta)u_\lambda=f\text{ 于分布意义},
 \quad\|u_\lambda\|_p\leq\lambda^{-1}\|f\|_p.
\]
先对 \(f\in\mathcal S\) 求出 \(\widehat u_\lambda\) 并证明 \(u_\lambda\in\mathcal S\)，再由
\[
 \partial_i\partial_ju_\lambda
 =R_iR_j(f-\lambda u_\lambda)
\]
证明全部二阶分布导数属于 \(L^p\)，且
\[
 \sum_{i,j=1}^d\|\partial_i\partial_ju_\lambda\|_p
 \leq C_{d,p}\|f\|_p,
\]
其中常数与 \(\lambda>0\) 无关。最后用缩放说明不存在对所有 \(u\in C_c^\infty(\mathbb R^d)\) 都成立的统一估计
\(\|u\|_p\leq C\|\Delta u\|_p\)。
\end{exercise}
\begin{solution}
先核对时间积分的连续性。固定 \(t_0>0\)，当 \(t_0/2\leq t\leq2t_0\) 时，热核的显式公式给出共同可积控制 \(C_{d,t_0}e^{-|x|^2/(8t_0)}\)。支配收敛说明 \(\|H_t-H_{t_0}\|_1\to0\)，Young 不等式因而给出
\[
 \|H_t*f-H_{t_0}*f\|_p
 \leq\|H_t-H_{t_0}\|_1\cdot\|f\|_p\longrightarrow0.
\]
零时的 \(L^p\) 连续性由\cref{b2:thm:12-heat-convolution}给出。热核非负且积分为一，Young 不等式还给出
\[
 \|H_t*f\|_p\leq\|f\|_p,
 \quad\int_0^\infty e^{-\lambda t}\|H_t*f\|_p\,\mathrm dt
                  \leq\lambda^{-1}\|f\|_p.
\]
因而可先在紧时间区间内积分，再放开两端，得到收敛的 \(L^p\) 值积分及所求零阶界。特别地，\(f\mapsto u_\lambda\) 是范数不超过 \(\lambda^{-1}\) 的线性算子。

若 \(f\in\mathcal S\)，还可在 \(L^1\) 中作同一积分，因为
\(\int_0^\infty e^{-\lambda t}\|H_t*f\|_1\,\mathrm dt<\infty\)。Fubini 因而允许先取 Fourier 变换，得到
\[
 \widehat u_\lambda(\xi)
 =\int_0^\infty e^{-t(\lambda+4\pi^2|\xi|^2)}
                 \cdot\widehat f(\xi)\,\mathrm dt
 =\frac{\widehat f(\xi)}{\lambda+4\pi^2|\xi|^2}.
\]
分母严格为正，其倒数是光滑函数，各阶导数至多作多项式增长，所以乘上 Schwartz 函数仍为 Schwartz 函数。Fourier 反变换说明 \(u_\lambda\in\mathcal S\)，并立即给出
\((\lambda-\Delta)u_\lambda=f\)。

记 \(g=f-\lambda u_\lambda\)。由刚求出的公式，
\[
 \widehat g(\xi)=\frac{4\pi^2|\xi|^2}{\lambda+4\pi^2|\xi|^2}
                       \cdot\widehat f(\xi).
\]
两个 Riesz 符号相乘为 \(-\xi_i\xi_j/|\xi|^2\)，故
\[
 \widehat{R_iR_jg}(\xi)
 =-\frac{4\pi^2\xi_i\xi_j}{\lambda+4\pi^2|\xi|^2}
                       \cdot\widehat f(\xi)
 =\widehat{\partial_i\partial_ju_\lambda}(\xi).
\]
零频处任意补值不影响等式。使用 \(R_iR_j\) 的 \(L^p\) 有界性与零阶界，有
\[
 \begin{aligned}
 \sum_{i,j}\|\partial_i\partial_ju_\lambda\|_p
 &\leq C_{d,p}\|f-\lambda u_\lambda\|_p\\
 &\leq C_{d,p}(\|f\|_p+\lambda\|u_\lambda\|_p)
 \leq2C_{d,p}\|f\|_p.
 \end{aligned}
\]
吸收因子二即可得到与 \(\lambda\) 无关的常数。

对一般 \(f\in L^p\)，取 \(f_n\in\mathcal S\) 在 \(L^p\) 中逼近。零阶界使对应的 \(u_{\lambda,n}\to u_\lambda\) 于 \(L^p\)，而
\(R_iR_j(f_n-\lambda u_{\lambda,n})\)
也在 \(L^p\) 中收敛到 \(R_iR_j(f-\lambda u_\lambda)\)。对任意 \(\varphi\in C_c^\infty\)，Hölder 不等式保证
\[
 \int u_{\lambda,n}\cdot\partial_i\partial_j\varphi
 \longrightarrow\int u_\lambda\cdot\partial_i\partial_j\varphi.
\]
所以二阶分布导数的极限正是 \(\partial_i\partial_ju_\lambda\)，并由上述 \(L^p\) 极限表示。对方程作同样的测试配对，得到其分布版本；范数估计也随 \(L^p\) 收敛传递。

零质量参数不能提供同样的零阶控制。固定非零 \(\phi\in C_c^\infty\)，使 \(\Delta\phi\ne0\)，令 \(v_R(x)=\phi(x/R)\)、\(R>0\)。换元给出
\[
 \|v_R\|_p=R^{d/p}\|\phi\|_p,
 \quad\|\Delta v_R\|_p=R^{d/p-2}\|\Delta\phi\|_p.
\]
若存在所述统一常数，则 \(R^2\|\phi\|_p\leq C\|\Delta\phi\|_p\) 对所有 \(R\) 成立，与 \(R\uparrow\infty\) 矛盾。二阶估计中相同的缩放因子可以相消，零阶估计却不能，这解释了 \(\lambda>0\) 在前一估计中的作用。
\end{solution}

% 自拟：用光滑乘法系数的差消去对角线奇性，比较差核的逐点极限与平方范数极限。
\Needspace{6\baselineskip}
\begin{exercise}\label{b2:ex:18-new-8}
设 \(d\geq1\) 为整数，\(T\) 是 \(\mathbb R^d\) 上\cref{b2:thm:18-cz-operator-definition}的标准 Calder\'on--Zygmund 算子，核为 \(K\)。另设硬截断
\[
 T_\epsilon g(x)=\int_{|x-y|>\epsilon}K(x,y)\cdot g(y)\,\mathrm dy
\]
对每个 \(g\in C_c^\infty(\mathbb R^d;\mathbb C)\)，当 \(\epsilon\downarrow0\) 时在 \(L^2\) 中趋于 \(Tg\)。给定 \(a,f\in C_c^\infty(\mathbb R^d;\mathbb C)\)，令 \(C_af=a\cdot Tf-T(a\cdot f)\)。证明
\[
 C_af(x)=\int_{\mathbb R^d}\bigl(a(x)-a(y)\bigr)\cdot K(x,y)\cdot f(y)\,\mathrm dy
\]
几乎处处成立，且右端对每个 \(x\) 都绝对收敛；解释乘法系数的差如何改善对角线附近的奇性。再证明对每个 \(1<p<\infty\)，有
\[
 \|C_af\|_p\leq2\|a\|_\infty\cdot\|T\|_{L^p\to L^p}\cdot\|f\|_p.
\]
这里 \(T\) 的 \(L^p\) 实现取\cref{b2:thm:18-cz-lp}的兼容延拓，不预先假定 \(T_\epsilon f\) 单独几乎处处收敛。
\end{exercise}
\begin{solution}
设核的大小常数为 \(A\)，\(a\) 的全局 Lipschitz 常数为 \(L_a\)。当 \(|x-y|<1\) 时，
\[
 \Bigl|\bigl(a(x)-a(y)\bigr)\cdot K(x,y)\cdot f(y)\Bigr|
 \leq A L_a\|f\|_\infty\cdot|x-y|^{1-d}.
\]
这在单位球上可积，径向积分为 \(\int_0^1\mathrm dr\)；系数之差提供的 \(|x-y|\) 消去了一阶奇性。在 \(|x-y|\geq1\) 时，绝对积分不超过 \(2A\|a\|_\infty\cdot\|f\|_1\)。因此对每个 \(x\)，差核的全积分绝对收敛，并且
\[
 a(x)\cdot T_\epsilon f(x)-T_\epsilon(a\cdot f)(x)
 =\int_{|x-y|>\epsilon}\bigl(a(x)-a(y)\bigr)\cdot K(x,y)\cdot f(y)\,\mathrm dy
\]
随 \(\epsilon\downarrow0\) 趋于该积分。另一方面，\(f,a\cdot f\in C_c^\infty\)，题设的 \(L^2\) 收敛及乘法算子的有界性使左端在 \(L^2\) 中趋于 \(a\cdot Tf-T(a\cdot f)\)。从任意趋零截断序列中抽取一个 \(L^2\) 收敛的逐点子列，就把后者与已存在的全参数差核极限识别，得到所求几乎处处等式。两项之差中的核奇性已被消去，因而其逐点极限可以直接由绝对收敛识别。

最后利用 \(T\) 的 \(L^p\) 有界性及 \(L^2\)、\(L^p\) 实现的相容性，
\[
 \|a\cdot Tf-T(a\cdot f)\|_p
 \leq\|a\|_\infty\cdot\|Tf\|_p+\|T(a\cdot f)\|_p
 \leq2\|a\|_\infty\cdot\|T\|_{L^p\to L^p}\cdot\|f\|_p.
\]
绝对积分表示使用了 \(a\) 的 Lipschitz 正则性，而这个范数估计只用其有界性，两项结论的假设作用不同。
\end{solution}

% !TeX root = ../../Book2.tex
\chapter{振荡积分}
\label{b2:ch:21}
\BookChapterNumber{b2:ch:21}
\begin{chapterabstract}
振荡积分的衰减来自相位变化造成的抵消。本章证明一维 van der Corput 引理、多维非驻相法、非退化驻相展开和混合 Hesse 矩阵 可逆时的振荡算子 \(L^2\) 估计，并解释曲率、参数退化及 Fourier 积分算子的初步几何。
\end{chapterabstract}

\begin{learninggoals}\begin{itemize}
\item\label{b2:item:21-goal-1}能区分边界项、驻点贡献与非驻相余项，并跟踪估计常数。
\item\label{b2:item:21-goal-2}能用有限阶导数下界证明一维衰减，并判断适用假设。
\item\label{b2:item:21-goal-3}能计算 Hesse 矩阵 行列式与符号，得到驻相领先项及下一阶。
\item\label{b2:item:21-goal-4}能使用 \(TT^*\) 和 Schur 估计证明局部振荡算子的 \(L^2\) 界。
\end{itemize}\end{learninggoals}

在前几章，Fourier 估计主要把振荡放进乘子或核的结构中。本章直接研究积分里的相位，说明曲率如何变成幂次衰减。所有高频参数均取 \(\lambda\ge1\)，所有相位均为实值。先处理标量积分，再进入算子估计；算子部分所需的非退化条件不同于标量驻相条件，不能互相替换。

% !TeX root = ../../Book2.tex
\section{一维振荡积分}
\label{b2:sec:2101}

正积分的大小由振幅总量决定，振荡积分还受到相位变化的制约。令
\[
 I_\lambda(a,\phi)=\int_\alpha^\beta e^{2\pi\mathrm i\lambda\phi(t)}a(t)\,\mathrm dt,
 \qquad \lambda\ge1.
\]
这里相位为实函数，振幅允许取复值。衰减估计中的常数必须说明依赖哪些导数与非退化下界；若这些量随参数恶化，单写一个与参数无关的 \(O(\lambda^{-1})\) 就失去了内容。

\subsection{分部积分}
\label{b2:subsec:210101}

假设 \(\phi'\ne0\)，并令 \(a\) 绝对连续。由
\[
 e^{2\pi\mathrm i\lambda\phi}
 =\frac1{2\pi\mathrm i\lambda\phi'}\frac{\mathrm d}{\mathrm dt}e^{2\pi\mathrm i\lambda\phi}
\]
得到边界项与余项的准确分离：
\[
 I_\lambda=\left[\frac{e^{2\pi\mathrm i\lambda\phi}a}{2\pi\mathrm i\lambda\phi'}\right]_\alpha^\beta
 -\frac1{2\pi\mathrm i\lambda}\int_\alpha^\beta e^{2\pi\mathrm i\lambda\phi}
       \left(\frac a{\phi'}\right)'\,\mathrm dt.
\]
若振幅在两端消失，第一项为零；若不消失，就必须保留它。对于线性相位与常数振幅，全部积分恰好由边界给出。这也说明增加振幅的光滑性而不处理边界，通常不能得到任意阶衰减。

\ProseParagraphBreak
在 \(|\phi'|\ge c>0\) 时，上式给出的上界为
\[
 |I_\lambda|\le\frac1{2\pi\lambda}
 \left[c^{-1}\bigl(|a(\alpha)|+|a(\beta)|+\|a'\|_1\bigr)
       +c^{-2}\int |a\phi''|\right].
\]
分母中出现的是相位导数；大的相位本身并不保证振荡，例如加上任意常数只改变整体复相位。

\subsection{非驻相}
\label{b2:subsec:210102}

\begin{proposition}[非驻相衰减]\label{b2:prop:21-nonstationary-one}
若 \(a\in C_c^\infty((\alpha,\beta))\)、\(\phi\in C^\infty\)，且在振幅支撑的一个邻域内 \(|\phi'|\ge c>0\)，则对每个整数 \(N\ge1\)，
\[
 |I_\lambda|\le C_N\lambda^{-N}.
\]
常数只依赖该邻域上 \(c^{-1}\)、相位的有限阶导数及振幅的有限阶导数。
\end{proposition}
\begin{proof}
把上一个公式反复用于振幅，定义
\(a_0=a\)、\(a_{j+1}=-(a_j/\phi')'\)。
每个 \(a_j\) 仍紧支撑于同一个集合，因此没有边界项，且
\(I_\lambda=(2\pi\mathrm i\lambda)^{-N}\int e^{2\pi\mathrm i\lambda\phi}a_N\)。
对 \(a_N\) 取 \(L^1\) 范数即得结论。商的每次求导只增加有限阶相位导数及 \(\phi'\) 的负幂，说明了常数依赖。
\end{proof}
紧支撑应理解为支撑远离正在使用公式的积分域边界。把光滑函数硬截断在一个区间上，会产生新的边界项，不能直接套用此命题。

\subsection{van der Corput 引理}
\label{b2:subsec:210103}

下面采用\term[van-der-corput-yinli]{van der Corput 引理}[van der Corput lemma]的下述版本。

\begin{theorem}[van der Corput 引理]\label{b2:thm:21-vdc}
设 \(a\) 绝对连续，\(\phi\in C^k([\alpha,\beta])\)。若 \(k=1\)，假设 \(|\phi'|\ge1\) 且 \(\phi'\) 单调；若 \(k\ge2\)，假设 \(|\phi^{(k)}|\ge1\)。则
\[
 |I_\lambda|\le C_k\lambda^{-1/k}
 \bigl(\|a\|_\infty+\|a'\|_1\bigr).
\]
\end{theorem}
\begin{proof}
先取振幅为 \(1\)。\(k=1\) 时分部积分，单调性给出
\(\int|(\phi')^{-2}\phi''|\le2\)，故任意子区间上的积分不超过 \(C\lambda^{-1}\)。
下面证明所需的次水平集估计：若 \(|h^{(r)}|\ge1\)，则
\[
 |\{t:|h(t)|\le\delta\}|\le C_r\delta^{1/r}.
\]
记左边为 \(L>0\)。按该集合的累积测度选择 \(r+1\) 个点，使任意两相邻点的距离至少 \(L/(2r)\)；可先以略低于 \(L\) 的测度作选择，再取极限。插值多项式的最高次系数为
\(\sum_{j=0}^r h(t_j)/\prod_{\ell\ne j}(t_j-t_\ell)\)，绝对值至多 \(C_r\delta L^{-r}\)。
逐次应用 Rolle 定理，某点的 \(h^{(r)}/r!\) 等于这个系数，得到所需估计。

对 \(h=\phi'\)、\(r=k-1\) 应用它。集合 \(|\phi'|\le\delta\) 上的积分至多 \(C_k\delta^{1/(k-1)}\)。其补集可分成有限个 \(\phi'\) 单调且 \(|\phi'|>\delta\) 的区间：\(\phi^{(k)}\) 无零点，Rolle 定理说明 \(\phi''\) 至多有 \(k-2\) 个零点；每个单调区间再按 \(\pm\delta\) 切分。第一阶估计在这些区间给出总上界 \(C_k(\lambda\delta)^{-1}\)。
取 \(\delta=\lambda^{-(k-1)/k}\)，两项均为 \(C_k\lambda^{-1/k}\)。

最后设 \(F(t)=\int_\alpha^t e^{2\pi\mathrm i\lambda\phi(s)}\,\mathrm ds\)。
上面的证明适用于每个子区间，所以 \(\|F\|_\infty\le C_k\lambda^{-1/k}\)。
公式 \(I_\lambda=a(\beta)F(\beta)-\int a'F\) 给出一般振幅的结论。
\end{proof}
第一阶中单调性不能删除。相位导数始终非零但反复变化时，分部积分会积累导数的总变差；只有下界而没有该变差控制，便不能保证统一常数。

\subsection{衰减估计}
\label{b2:subsec:210104}

若相位的第 \(k\) 阶导数下界为 \(c\)，对 \(\lambda c\) 应用定理，可得到 \(C_k(\lambda c)^{-1/k}\) 的上界。因此相位趋于退化时，常数会变大。以
\(\phi(t)=t^k\)、振幅在原点等于 \(1\) 为例，令 \(s=\lambda^{1/k}t\)，便出现自然尺度 \(\lambda^{-1/k}\)。对于二次相位，该尺度确实达到；更高阶时适当选取单侧振幅也可避免两侧抵消。

\ProseParagraphBreak
常数下界与有限类型条件也可局部使用。若紧支撑上每一点至少有某个 \(1\le j\le k\) 满足 \(|\phi^{(j)}|\ge c\)，有限覆盖与光滑分割把积分分成若干项。第一阶部分还需细分成有受控变差的区间；其余部分可直接使用 van der Corput 引理。这里有限覆盖的大小、分割函数导数和非退化下界都进入常数。有限类型给出一个统一衰减率，并不说明每个点都以最坏的 \(k\) 阶速度贡献。

% 扩写定位：b2:subsec:210104
有限型估计在临界点碰撞时比逐点驻相更稳定。取
\[
 I(\lambda,s)=\int e^{2\pi\mathrm i\lambda(t^3-st)}a(t)\,\mathrm dt,
 \qquad a\in C_c^\infty(\mathbb R).
\]
三阶导数恒为六，因此 van der Corput 给出
\(|I(\lambda,s)|\leq C_a\lambda^{-1/3}\)，常数与 \(s\) 无关。若 \(s>0\)，两个临界点为 \(\pm\sqrt{s/3}\)，二阶导数大小为 \(2\sqrt{3s}\)；分别使用驻相会出现 \(s^{-1/4}\lambda^{-1/2}\) 的因子，不能在 \(s\downarrow0\) 时保持统一。

\ProseParagraphBreak

换元 \(t=\lambda^{-1/3}u\) 后，相位成为 \(u^3-(s\lambda^{2/3})u\)。因此碰撞区域的自然参数是 \(s\lambda^{2/3}\)，自然宽度是 \(s\asymp\lambda^{-2/3}\)。在这个区域，两点驻相的分离距离已经与波动尺度相当，应保持三次相位整体；当 \(s\lambda^{2/3}\) 很大时，才适合分别围绕两个临界点分析。这里的统一估计无需写出 Airy 函数，但缩放已经解释了为什么二次渐近在碰撞附近失效。

\ProseParagraphBreak

非驻相区域也需要与支撑相配合。若 \(s<0\)，则 \(\phi_s'(t)=3t^2+|s|\geq|s|\)，固定的负 \(s\) 给出任意阶衰减，但分部积分常数会随 \(|s|\downarrow0\) 增长；统一三阶界仍然有效。若振幅支撑远离零点，则即使 \(s\) 接近零，一阶导数也有统一下界，此时可以恢复快速衰减。由此可见，退化不是相位公式本身的孤立属性，必须连同振幅支撑和参数范围一起判断。

% !TeX root = ../../Book2.tex
\section{驻相法}
\label{b2:sec:2102}

一阶导数消失后，分部积分的分母发生奇异，主要贡献便集中在驻点附近。非退化驻点的尺度是 \(\lambda^{-1/2}\)，其领先系数由 Hesse 矩阵 决定。以下把 Fourier 指数中的 \(2\pi\) 保留在相位内，以便与全卷的变换约定一致。

\subsection{非退化临界点}
\label{b2:subsec:210201}

\begin{definition}\label{b2:def:21-critical}
若 \(\nabla\phi(x_0)=0\)，称 \(x_0\) 为驻点；若同时
\(\det D^2\phi(x_0)\ne0\)，称其为\term[fei-tuihua-zhudian]{非退化驻点}[nondegenerate critical point]。
\end{definition}
逆函数定理用于映射 \(\nabla\phi\)，说明非退化驻点在一个小邻域内是唯一的。反过来，孤立驻点未必非退化：\(t^4\) 的原点就是例子。第二种情形的衰减阶由更高阶项决定，不能套用二次驻相公式。

\ProseParagraphBreak
在一个紧集上，若所有驻点非退化，则驻点只有有限个。否则驻点列有聚点，连续性使聚点仍为驻点，与局部孤立性矛盾。于是可以分别计算各驻点贡献，再对剩余部分使用非驻相法；不同贡献的复相位可能抵消，因此各项上界之和只是上界。

\subsection{Hesse 矩阵}
\label{b2:subsec:210202}

对实对称可逆矩阵 \(H\)，记其正特征值数减负特征值数为 \(\operatorname{sgn}H\)。在非退化驻点附近存在光滑坐标 \(u\)，使
\[
 \phi(x)=\phi(x_0)+\tfrac12\sum_{j=1}^d\varepsilon_j u_j^2,
 \qquad \varepsilon_j\in\{1,-1\}.
\]
这是局部 Morse 归约。为说明它所用的事实，先平移驻点并作固定线性变换，使 Hesse 矩阵 对角化。Taylor 积分余项写成
\(\phi(x)-\phi(0)=\tfrac12x^{\mathsf T}H(x)x\)，其中
\(H(x)=2\int_0^1(1-t)D^2\phi(tx)\,\mathrm dt\) 光滑，且 \(H(0)\) 可逆。
在足够小的邻域内，按固定顺序作对称消元，各主元仍不为零。由消元得到光滑的 \(L(x)D(x)L(x)^{\mathsf T}\) 分解；各主元的符号不变，把其绝对值的平方根吸收到三角矩阵内，便有
\(H(x)=B(x)^{\mathsf T}S B(x)\)，\(S=\operatorname{diag}(\varepsilon_j)\)。
令 \(u=B(x)x\)，其在原点的导数可逆，逆函数定理完成归约。

\ProseParagraphBreak
坐标变换的 Jacobi 因子满足
\[
 \left|\det\frac{\partial x}{\partial u}(0)\right|
      =|\det D^2\phi(x_0)|^{-1/2}.
\]
这个因子不是可忽略的规范化。一个 Hesse 矩阵 特征值变小时，对应方向的驻点邻域变宽，领先系数也随之增大。

\subsection{驻相引理}
\label{b2:subsec:210203}

下面采用\term[zhuxiang-gongshi]{驻相公式}[stationary phase formula]的下述版本。

\begin{theorem}[驻相公式]\label{b2:thm:21-stationary-phase}
设 \(a\in C_c^\infty(\mathbb R^d)\)、\(\phi\) 在支撑附近光滑，支撑内只有一个驻点 \(x_0\)，且该点非退化。记 \(H=D^2\phi(x_0)\)，则
\[
 \int e^{2\pi\mathrm i\lambda\phi(x)}a(x)\,\mathrm dx
 =e^{2\pi\mathrm i\lambda\phi(x_0)}
   \lambda^{-d/2}e^{\mathrm i\pi\operatorname{sgn}H/4}
   \frac{a(x_0)}{|\det H|^{1/2}}
   +O(\lambda^{-d/2-1}).
\]
\end{theorem}
\begin{proof}
分割振幅使一部分支撑于 Morse 坐标邻域，剩余部分由非驻相法降到任意阶。换坐标后只需处理
\(\int e^{\pi\mathrm i\lambda u^{\mathsf T}Su}b(u)\,\mathrm du\)。
对角矩阵 \(S\) 的每个分量应用复 Gaussian 积分，先加因子 \(e^{-\epsilon|u|^2}\)，再令 \(\epsilon\downarrow0\)，得到分布恒等式
\[
 \int e^{\pi\mathrm i\lambda u^{\mathsf T}Su}b(u)\,\mathrm du
 =\lambda^{-d/2}e^{\mathrm i\pi\operatorname{sgn}S/4}
   \int e^{-\pi\mathrm i\eta^{\mathsf T}S\eta/\lambda}\widehat b(\eta)\,\mathrm d\eta.
\]
Gaussian 积分在 \(\epsilon>0\) 时绝对收敛；配方所得 Fourier 变换的模有统一上界，乘上 Schwartz 函数 \(\widehat b\) 后可用支配收敛。这说明取极限与配对都是合法的。
右侧指数等于 \(1+O(|\eta|^2/\lambda)\)，而
\(\int\widehat b=b(0)\)，\(\int|\eta|^2|\widehat b|<\infty\)。
再代入 \(b(0)=a(x_0)|\det H|^{-1/2}\)，结论成立。
\end{proof}
该公式是固定振幅、固定相位下的渐近式。若振幅支撑随 \(\lambda\) 缩小或导数增长，余项常数可能增长；此时应先把自然尺度规范化。

% 扩写定位：b2:subsec:210203
边界还会改变驻相的首项。考虑一个临界点恰好位于积分端点的模型
\[
 J(\lambda)=\int_0^\infty e^{\pi\mathrm i\lambda t^2}a(t)\,\mathrm dt,
 \qquad a\in C_c^\infty([0,\infty)).
\]
此处光滑是指可延拓到零点邻域的光滑函数。若 \(a(0)\ne0\)，则
\[
 J(\lambda)=\frac12e^{\pi\mathrm i/4}a(0)\lambda^{-1/2}+O(\lambda^{-1}).
\]
取在零点附近为一的偶截断 \(\chi\)，把 \(a=a(0)\chi+t b\) 分解。第一项是全直线二次 Gaussian 积分的一半，其远离零点部分可反复分部积分；第二项令 \(v=t^2\)，成为 \(\frac12\int_0^\infty e^{\pi\mathrm i\lambda v}b(\sqrt v)\,\mathrm dv\)。该振幅的导数在零点至多有 \(v^{-1/2}\) 奇性，仍可积，分部积分得到 \(O(\lambda^{-1})\)。这就证明了首项中的一半。

\ProseParagraphBreak

若临界点离开端点而进入内部，其首项重新是完整的二次驻相系数；若移到区间外，则固定参数下没有临界点而只剩边界项。随着临界点在距离 \(\lambda^{-1/2}\) 的区域内移动，应使用保留端点的 Fresnel 积分，不能把内部公式和非驻相公式硬拼成一个统一展开。这个模型同时说明，振幅紧支于开集和振幅在闭区间端点非零是不同假设。

\subsection{渐近展开}
\label{b2:subsec:210204}

上一个证明还给出任意阶展开。令 \(D_S=\sum_j\varepsilon_j\partial_{u_j}^2\)，则
\[
 \int e^{\pi\mathrm i\lambda u^{\mathsf T}Su}b(u)\,\mathrm du
 =\lambda^{-d/2}e^{\mathrm i\pi\operatorname{sgn}S/4}
 \left[\sum_{j=0}^{N-1}\frac{(\mathrm i/(4\pi\lambda))^j}{j!}
              D_S^j b(0)+O(\lambda^{-N})\right].
\]
证明是在指数的 Taylor 余项上使用
\(|e^{\mathrm it}-\sum_{j<N}(\mathrm it)^j/j!|\le C_N|t|^N\)，再利用 Fourier 求导恒等式
\(\int(\eta^{\mathsf T}S\eta)^j\widehat b
 =(-4\pi^2)^{-j}D_S^j b(0)\)。
因此第 \(j\) 个系数由驻点处至多 \(2j\) 阶的振幅导数及坐标变换导数决定。

\ProseParagraphBreak
领先系数为零并不意味着整个积分为零。例如一维二次相位配以偶振幅 \(a(t)=t^2\chi(t)\)，首项消失，第二项通常非零，衰减变成 \(\lambda^{-3/2}\)。若振幅为奇函数、截断也对称，积分则确实为零。前者是局部消失阶带来的增益，后者是全局对称性抵消；两种机制应分清。

% 扩写定位：b2:subsec:210204
驻相公式也适用于一个临界子流形，前提是法向二次型非退化。这里可以给出一个明确的局部版本，无需提前使用 Fourier 积分算子理论。
\begin{proposition}\label{b2:prop:21-clean-stationary}
设局部坐标为 \((y,z)\in\mathbb R^{d-r}\times\mathbb R^r\)，临界集为 \(z=0\)。假定在振幅的一个固定紧支撑附近有
\(\mathrm d\phi(y,0)=0\)，法向 Hesse 矩阵 \(H(y)=\partial_z^2\phi(y,0)\) 一致可逆，且无其他临界点。若该片临界集连通，则 \(\phi(y,0)=c\) 为常数。对 \(a\in C_c^\infty\)，
\[
 \int e^{2\pi\mathrm i\lambda\phi(y,z)}a(y,z)\,\mathrm dy\,\mathrm dz
 =e^{2\pi\mathrm i\lambda c}\lambda^{-r/2}
 \int e^{\pi\mathrm i\operatorname{sgn}H(y)/4}
 \frac{a(y,0)}{|\det H(y)|^{1/2}}\,\mathrm dy
 +O(\lambda^{-r/2-1}).
\]
\end{proposition}
\begin{proof}
临界条件使沿 \(y\) 方向的一阶导数为零，所以相位在连通临界片上常值。对每个固定 \(y\)，在法向 \(z\) 上作 Taylor 展开和前面构造的 Morse 坐标变换。Hesse 矩阵 的一致可逆性与紧支撑允许在有限个参数邻域中统一选择坐标；矩阵分解、逆函数定理和振幅导数都具有统一的有限阶界。在这些片上，法向驻相公式的余项一致为 \(O(\lambda^{-r/2-1})\)。对 \(y\) 积分并加上有限分割便得结论。远离 \(z=0\) 的部分梯度具有正下界，反复分部积分把它放入任意小的幂次余项中。
\end{proof}
衰减幂次由真正振荡的法向维数 \(r\) 决定，而不是积分空间的总维数。例如 \(\phi(y,z)=|z|^2/2\) 与乘积振幅 \(a=b(y)c(z)\)，对 \(y\) 的积分没有任何 \(\lambda\) 衰减；只有 \(z\) 变量贡献 \(\lambda^{-r/2}\)。圆柱曲面的 Fourier 变换也体现这种秩信息：弯曲方向提供驻相增益，平直方向是否衰减取决于振幅和 Fourier 方向，不能从总维数套用球面的公式。

\ProseParagraphBreak

这一局部公式不是任意退化临界点的公式。若 Hesse 矩阵 在法向也退化，Morse 坐标不再存在；若临界集有交叉或尖点，也不能选成单一光滑片。比如 \(t^4\) 在零点的衰减为 \(\lambda^{-1/4}\)，已与任何一维非退化二次型的 \(\lambda^{-1/2}\) 不同。所需的是有限型分析，而不是把行列式的零值放进分母。

\ProseParagraphBreak

非二次相位的下一项并不只取决于振幅的二阶导数。这个差别可以用一维局部计算明确看到。将驻点移到零，并归一化 \(\phi''(0)=1\)，写
\[
 \phi(t)=\frac{t^2}2+\frac c6t^3+\frac d{24}t^4+O(t^5).
\]
Morse 变量为 \(y=t\sqrt{2\phi(t)/t^2}\)。在零点取使导数为正的一支，逐项展开再反解得到
\[
 t=y-\frac c6y^2+\left(\frac{5c^2}{72}-\frac d{24}\right)y^3+O(y^4).
\]
新振幅是 \(b(y)=a(t(y))t'(y)\)，而不只是 \(a(t(y))\)。计算其二阶导数得到
\[
 b''(0)=a''(0)-c\,a'(0)
 +\left(\frac{5c^2}{12}-\frac d4\right)a(0).
\]
二次模型的展开于是给出
\[
 I_\lambda=e^{\pi\mathrm i/4}\lambda^{-1/2}
 \left[a(0)+\frac{\mathrm i}{4\pi\lambda}b''(0)
 +O(\lambda^{-2})\right],
\]
这里假设相位在该局部片只有这个驻点，远离它的部分已用非驻相法处理。

\ProseParagraphBreak

即使原振幅在驻点附近为常数，三次和四次相位也会改变次项，因为坐标变换的体积因子不是常数。若只把二次公式中的 \(a''(0)\) 原样用于任意相位，就会遗漏这些项。计算首项时只需变换在驻点的一阶微分；计算更高项则需要相位与变换的更高阶 Taylor 数据。这解释了为什么完整渐近展开的导数假设随所需阶数增加。

\ProseParagraphBreak

参数族的微分也应先去掉迅速振荡的临界值。设相位和振幅光滑依赖参数 \(v\)，唯一非退化临界点 \(x_0(v)\) 位于固定片内，并有统一下界与支撑。隐函数定理使 \(x_0(v)\) 光滑。把积分乘以
\(e^{-2\pi\mathrm i\lambda\phi(x_0(v),v)}\lambda^{d/2}\)，则该归一化量具有以 \(\lambda^{-1}\) 为幂的渐近展开，系数光滑依赖 \(v\)，并在紧参数集上连同每个固定阶参数导数一致成立。

\ProseParagraphBreak

证明仍用含参数的 Morse 坐标：矩阵分解和逆映射都光滑依赖参数，新振幅的所需导数因此一致有界。二次 Fourier 计算的 Taylor 余项可以在这些半范数下统一估计。对远离临界点的部分，参数微分可能带来若干 \(\lambda\) 因子，但预先增加相应次数的非驻相分部积分即可吸收。因此“对参数一致的函数值余项”与“对参数一致的全部所需导数余项”要分别核对；只有前者不足以支持后续对参数求导。

\ProseParagraphBreak

有限个驻点的贡献需要先相加，才能讨论积分的整体大小。假定支撑中只有非退化驻点 \(x_1,\ldots,x_\ell\)，选各点附近互不相交的截断，再把其余部分留给非驻相法，便有
\[
 I_\lambda=\lambda^{-d/2}\sum_{\nu=1}^{\ell}
 e^{2\pi\mathrm i\lambda\phi(x_\nu)}
 e^{\pi\mathrm i\operatorname{sgn}H_\nu/4}
 \frac{a(x_\nu)}{|\det H_\nu|^{1/2}}
 +O(\lambda^{-d/2-1}).
\]
这里每个局部首项都可能非零，但它们的和仍可能抵消。取一维相位 \(\phi(t)=t^3/3-t\)，振幅只支在 \(\pm1\) 的两个小邻域，并在两个驻点等于一。两点的相位值为 \(\mp2/3\)，二阶导数为 \(\pm2\)，故首项合成为
\[
 \sqrt2\,\lambda^{-1/2}
 \cos\left(\frac{4\pi\lambda}{3}-\frac\pi4\right).
\]
沿余弦消失的一列正参数，积分只剩 \(O(\lambda^{-3/2})\)。因此一般驻相定理提供上界与展开；若要求统一的下界，还必须排除不同驻点之间的干涉。反过来，只有一个驻点且该点振幅非零时，首项模有正下界，余项便保证充分大参数下确实具有 \(\lambda^{-d/2}\) 的大小。

\ProseParagraphBreak

首项中的行列式与体积因子也有坐标无关的含义。设 \(x=F(y)\) 是局部光滑换元，\(F(y_0)=x_0\)，且 \(B=DF(y_0)\) 可逆。在驻点处一阶相位导数为零，二阶链式法则中的变换二阶导数项消失，因而新 Hesse 矩阵 为 \(B^{\mathsf T}HB\)，新振幅为 \(a(x_0)|\det B|\)。于是
\[
 \frac{a(x_0)|\det B|}
 {|\det(B^{\mathsf T}HB)|^{1/2}}
 =\frac{a(x_0)}{|\det H|^{1/2}}.
\]
实可逆合同也保持正负惯性指数，所以相位因子不变。这说明驻相系数与选择哪一套 Morse 坐标无关。若只改变相位而忘记变换振幅中的 Jacobi 因子，就会破坏这种不变性；前面的一维次项计算正是这一原则在更高 Taylor 阶上的体现。

% !TeX root = ../../Book2.tex
\section{多维振荡积分}
\label{b2:sec:2103}

多维情况下既可能完全没有驻点，也可能只在某些方向振荡。估计的有效维数取决于可利用的非退化方向数，而不单由积分变量的个数决定。

\subsection{相位函数}
\label{b2:subsec:210301}

若 \(\nabla\phi\ne0\)，定义
\[
 L=\frac{\nabla\phi\cdot\nabla}{2\pi\mathrm i\lambda|\nabla\phi|^2},
 \qquad Le^{2\pi\mathrm i\lambda\phi}=e^{2\pi\mathrm i\lambda\phi}.
\]
对紧支撑振幅反复转移 \(L\) 的形式伴随，可得任意阶 \(\lambda^{-N}\) 衰减。
展开伴随时出现的每一项都是相位有限阶导数、\(|\nabla\phi|^{-1}\) 的幂与振幅导数的乘积。在 \(|\nabla\phi|\ge c\) 的紧集上这些系数有界，故证明和一维完全相同，只是导数变成多重指标。

\ProseParagraphBreak
若积分域有边界，散度定理同时产生边界积分。比如平直立方体上的线性相位，即使内域完全无驻点，某些面上的切向梯度仍可能为零。这解释了为什么光滑截断与特征函数截断的 Fourier 衰减差别很大。

\subsection{振幅}
\label{b2:subsec:210302}

振幅支撑的大小、光滑性和消失阶都参与估计。给定 \(a_\lambda(x)=a(\lambda^\theta x)\)，其第 \(j\) 阶导数有 \(\lambda^{j\theta}\) 的增长；直接套用固定振幅定理会把这些增长藏进余项。先换变量 \(u=\lambda^\theta x\)，新的相位参数变为 \(\lambda^{1-2\theta}\)（二次相位），才看得清有效振荡尺度。

\ProseParagraphBreak
振幅可以包含较慢的参数，但必须统一控制其支撑和导数。一个常用情形是 \(a(x,z)\) 支撑于固定紧集，\(z\) 属于紧参数集，并且所需 \(x\) 导数一致有界。若相位也具有统一非退化下界，所有前述估计都可对 \(z\) 一致成立。参数集不紧时应另行核对这些界。

\subsection{曲率}
\label{b2:subsec:210303}

设光滑超曲面在局部写成 \(x_d=h(x')\)。曲面测度的 Fourier 变换在方向 \(\omega\) 上具有相位
\(\phi_\omega(x')=-\omega'\cdot x'-\omega_d h(x')\)。
驻点条件说明 \(\omega\) 平行于曲面的法向；驻点处 Hesse 矩阵 为
\(-\omega_dD^2h\)。在相应坐标中 Gaussian 曲率非零，等价于这个 Hesse 矩阵 可逆。因此紧支撑光滑曲面密度满足
\[
 |\widehat{a\,\mathrm d\sigma}(\xi)|\le C(1+|\xi|)^{-(d-1)/2}
\]
（假设支撑上 Gaussian 曲率处处非零）。
对方向单位球面作有限覆盖，分出驻点邻域与非驻点区域，驻相公式和非驻相法分别控制它们，便得到方向一致的常数。

\ProseParagraphBreak
若曲面是平面，沿法向相位是常数，便没有这样的衰减。若曲率只有 \(r\) 个非零方向，局部二次模型只有 \(r\) 个振荡变量，通常只能期待 \(|\xi|^{-r/2}\)。柱面是一个直接例子：圆周方向提供曲率，轴向仍可保持相干。仅说“曲面光滑”不足以推出球面型衰减。

\subsection{局部化}
\label{b2:subsec:210304}

局部化需要同时记录所选覆盖和每片的非退化下界。若 \(\chi_j\) 构成振幅支撑附近的有限光滑分割，则
\[
 I_\lambda(a,\phi)=\sum_j I_\lambda(a\chi_j,\phi).
\]
在非驻相片使用梯度下界，在驻相片使用 Hesse 矩阵 下界；余项的阶数可以统一选取，再将有限项相加。
对于一个逼近退化曲面的参数族，片数或下界可能变化。此时统一估计需要新的尺度分解，不能仅把每个固定参数的驻相公式相加。

\ProseParagraphBreak
以 \(\phi_s(t)=t^3-st\) 为例，\(s>0\) 时有两个非退化驻点，\(s=0\) 时合并为三阶驻点，\(s<0\) 时没有驻点。两个驻点的 Hesse 矩阵 在 \(s\downarrow0\) 时趋于零，分别使用二次公式的常数会发散。令 \(t=\lambda^{-1/3}u\)，可见转变发生在 \(s\lambda^{2/3}\) 为常数量级时。这个例子说明统一渐近式为何需要将参数与空间尺度一起组织。

% !TeX root = ../../Book2.tex
\section{振荡积分算子}
\label{b2:sec:2104}

标量积分的衰减不自动给出算子范数衰减，因为输入函数也会随参数改变，并可抵消一部分振荡。算子理论通过比较两个输出点，使不能同时抵消的相位差显现出来。

\subsection{算子核}
\label{b2:subsec:210401}

考虑
\[
 T_\lambda f(x)=\int_{\mathbb R^d}e^{2\pi\mathrm i\lambda\Phi(x,y)}a(x,y)f(y)\,\mathrm dy,
 \qquad a\in C_c^\infty(\mathbb R^{2d}).
\]
固定 \(x\) 后对核应用驻相法，只适用于足够光滑且导数受控的 \(f\)。\(L^2\) 单位球中的函数不满足这种控制，所以需要另一个方法。
若 \(\Phi(x,y)=\psi(x)+\theta(y)\)，则输入可取
\(f(y)=e^{-2\pi\mathrm i\lambda\theta(y)}g(y)\)，完全消去 \(y\) 振荡；输出只是固定非振荡积分乘一个模为 \(1\) 的因子，范数不衰减。

\ProseParagraphBreak
混合 Hesse 矩阵 \(\partial_x\partial_y\Phi\) 正是排除这种分离相位的量。它描述改变输出点时，输入变量中的相位梯度如何改变。对于双线性相位 \(x\cdot y\)，算子等于缩放后的 Fourier 变换，Plancherel 给出准确的 \(\lambda^{-d/2}\) 范数。

\subsection{\texorpdfstring{\(TT^*\)}{TT*} 方法}
\label{b2:subsec:210402}

伴随为
\[
 T_\lambda^*g(y)=\int e^{-2\pi\mathrm i\lambda\Phi(x,y)}
                      \overline{a(x,y)}g(x)\,\mathrm dx,
\]
故 \(T_\lambda T_\lambda^*\) 的核是
\[
 K_\lambda(x,z)=\int e^{2\pi\mathrm i\lambda(\Phi(x,y)-\Phi(z,y))}
                          a(x,y)\overline{a(z,y)}\,\mathrm dy.
\]
Hilbert 空间恒等式 \(\|T_\lambda\|_{2\to2}^2
=\|T_\lambda T_\lambda^*\|_{2\to2}\) 可由
\(\|T^*g\|^2=\langle TT^*g,g\rangle\) 与 \(\|T\|=\|T^*\|\) 直接证明。
此时核的主对角线 \(x=z\) 不振荡，但离开对角线后相位差可有统一梯度下界。估计目标因此是一个宽度约 \(\lambda^{-1}\) 的对角邻域。

\ProseParagraphBreak
若 \(\sup_x\int|K(x,z)|\,\mathrm dz\le A\) 且
\(\sup_z\int|K(x,z)|\,\mathrm dx\le B\)，则 Schur 估计给出
\(\|K\|_{2\to2}\le\sqrt{AB}\)。证明是先对
\(\int |K(x,z)||g(z)|\,\mathrm dz\) 用带权 Cauchy–Schwarz，再交换非负积分。这里绝对值估计发生在已经利用振荡得到衰减之后；若一开始对原核取绝对值，振荡增益便消失了。

\subsection{\texorpdfstring{\(L^2\)}{L2} 估计}
\label{b2:subsec:210403}

\begin{theorem}[非退化振荡算子的局部 \(L^2\) 估计]\label{b2:thm:21-oscillatory-l2}
若 \(\det\partial_x\partial_y\Phi(x_0,y_0)\ne0\)，则存在该点的邻域，使支撑在此邻域内的光滑振幅满足
\[
 \|T_\lambda f\|_2\le C\lambda^{-d/2}\|f\|_2,\qquad \lambda\ge1.
\]
若混合 Hesse 矩阵 在一个固定紧支撑附近处处可逆，同样的结论成立。
\end{theorem}
\begin{proof}
先取支撑足够小，使混合 Hesse 矩阵 一致接近固定可逆矩阵 \(H_0\)。由均值公式，
\[
 \nabla_y\Phi(x,y)-\nabla_y\Phi(z,y)
 =\left[\int_0^1\partial_x\nabla_y\Phi(z+t(x-z),y)\,\mathrm dt\right](x-z),
\]
其长度至少 \(c|x-z|\)。把相位差除以 \(|x-z|\) 后，该相位的有限阶 \(y\) 导数一致有界，梯度一致有正下界。多维非驻相估计给出
\[
 |K_\lambda(x,z)|\le C_N(1+\lambda|x-z|)^{-N}.
\]
当 \(\lambda|x-z|\le1\) 时用振幅的积分上界补上这个式子即可。取 \(N>d\)，两侧积分均不超过 \(C\lambda^{-d}\)，Schur 估计和 \(TT^*\) 恒等式得出结论。
一般紧支撑由有限个小矩形覆盖，以光滑分割写成有限个局部算子。每个局部算子满足相同幂次的估计，三角不等式给出整体结论；只改变常数。
\end{proof}
这个局部化步骤很必要：处处可逆的矩阵沿一段线段积分后未必仍可逆。证明中的固定矩阵近似只在小片内成立，不能在大支撑上直接使用。

\subsection{参数依赖性}
\label{b2:subsec:210404}

在一个参数族中，若混合 Hesse 矩阵 最小奇异值一致大于 \(c>0\)，相位与振幅的所需导数一致有界，支撑处于同一紧集，则有限覆盖可统一选择，上述常数也统一。若非退化下界趋于零，必须显示其依赖。
例如 \(\Phi_\epsilon(x,y)=\epsilon x\cdot y\)，在有效振荡范围 \(\lambda\epsilon\ge1\) 中，缩放 Fourier 变换给出 \((\lambda\epsilon)^{-d/2}\)；当 \(\lambda\epsilon<1\) 时，振幅的普通积分上界更合适。

\ProseParagraphBreak
若两侧维数不同，混合 Hesse 矩阵 不再是方阵，可逆条件也失去意义。此时应考虑其秩以及余下变量中的几何关系，不能沿用 \(\lambda^{-d/2}\) 的公式。例如把上述算子乘一个无振荡的额外变量积分，新增维度并不会增加衰减。这种秩与曲率的配合将在限制问题和多线性问题中再次出现。

% 扩写定位：b2:subsec:210404
统一常数在使用 \(TT^*\) 时有两个来源。第一是局部相位梯度差的下界，需要混合矩阵与某个固定可逆矩阵足够接近；第二是分部积分中各个振幅和倒梯度系数的导数界。仅给出行列式不为零还不够得到一族算子的统一常数，例如矩阵 \(\operatorname{diag}(\epsilon,\epsilon^{-1})\) 的行列式恒为一，但导数大小和最小奇异值都失去控制。

\ProseParagraphBreak

在有限阶估计中，实际上只需相位和振幅的有限多个导数；取 \(TT^*\) 的非驻相次数 \(N>d\) 后，所需阶数由 \(N,d\) 决定。若要全部阶的渐近展开，则应有全部相应阶导数。两种命题虽然都常写为“光滑”，对参数族的要求并不相同。记录常数时应说明固定了哪些半范数，而不能把全部依赖藏在一个无法随参数核对的 \(C\) 中。

\ProseParagraphBreak

该 \(\lambda^{-d/2}\) 幂次在一般非退化算子中不能改善。取 \(\Phi(x,y)=x\cdot y\)、振幅 \(a(x,y)=\chi(x)\eta(y)\)，其中 \(\chi,\eta\) 在零点附近不为零。固定 \(g\in C_c^\infty\) 满足 \(\eta g\not\equiv0\)，令输入为 \(g\)。则
\[
 T_\lambda g(x)=\chi(x)\widehat{\eta g}(-\lambda x),\qquad
 \|T_\lambda g\|_2^2=
 \lambda^{-d}\int|\chi(v/\lambda)|^2|\widehat{\eta g}(-v)|^2\,\mathrm dv.
\]
支配收敛给出 \(\lambda^d\|T_\lambda g\|_2^2\to|\chi(0)|^2\|\eta g\|_2^2>0\)。因而算子范数至少是 \(c\lambda^{-d/2}\)。这个例子直接测试了上界的幂次，所用输入不必依赖于 \(\lambda\)。

% !TeX root = ../../Book2.tex
\OptionalSection{Besov 与 Triebel–Lizorkin 空间}{b2:sec:1905}

同样一组频率块，可以记录两种不同的信息。先量每一块在整个空间中的大小，再比较不同尺度，得到一个数列范数；先在每个位置比较各尺度，再对位置积分，得到另一种范数。\subsecref{b2:subsec:190404}的平方空间已经采用了后一种次序。本节允许尺度指数变化，并通过差分、乘积和平坦边界的迹，说明这个指数为什么不能省略。一般空间可以选读；\chapref{b2:ch:22}的普通乘子正文只需要上一节已经建立的平方情形。

\subsection{二进定义}
\label{b2:subsec:190501}

先把分割固定下来，使后面的每个范数都有确定含义。低频块的作用也必须保留：所有齐次环带都消去多项式，而非齐次分解还能记录零频附近的信息。

\Needspace{6\baselineskip}
\begin{definition}\label{b2:def:19-besov-triebel-lizorkin}
设整数 \(d\geq1\)，取实径向函数 \(\chi\in C_c^\infty(\mathbb R^d)\)，满足 \(0\leq\chi\leq1\)，在 \(|\xi|\leq1\) 为一、\(|\xi|\geq2\) 为零，并沿半径单调不增。令 \(\varphi=\chi-\chi(2\,\cdot)\)。对复值温和分布 \(f\in\mathcal S'(\mathbb R^d)\)，定义
\[
 P_0f=\mathcal F^{-1}(\chi\cdot\widehat f),
 \quad P_jf=\mathcal F^{-1}\bigl(\varphi(2^{-j}\,\cdot)\cdot\widehat f\bigr)
 \quad(j\geq1).
\]
设 \(s\in\mathbb R\)、\(1\leq p,q\leq\infty\)。满足
\[
 \|f\|_{B^s_{p,q}}
 \coloneqq\left\|\bigl\{2^{js}\|P_jf\|_p\bigr\}_{j\geq0}\right\|_{\ell^q}<\infty
\]
的温和分布组成\term[besov-kongjian]{Besov 空间}[Besov space]，记为 \(B^s_{p,q}(\mathbb R^d)\)。当 \(1\leq p<\infty\)、\(1\leq q\leq\infty\) 时，满足
\[
 \|f\|_{F^s_{p,q}}
 \coloneqq\left\|\left\|\bigl\{2^{js}P_jf(x)\bigr\}_{j\geq0}\right\|_{\ell^q}\right\|_{L^p_x}<\infty
\]
的温和分布记为 \(F^s_{p,q}(\mathbb R^d)\)，称为一般的 Triebel--Lizorkin 空间。序列指数 \(q=\infty\) 时，\(\ell^q\) 范数取上确界；Besov 定义中的 \(p=\infty\) 取空间本性上确界。
\end{definition}
\symidx{\(B^s_{p,q}\)：非齐次 Besov 空间}
\symidx{\(F^s_{p,q}\)：非齐次 Triebel--Lizorkin 空间}
这里的块正是\cref{b2:def:19-inhomogeneous-square}中的块。每个块都是光滑函数，但未必有有限的全空间范数；定义的有限性不是自动得到的。由于
\[
 P_0f+\sum_{j=1}^NP_jf
 =\mathcal F^{-1}\bigl(\chi(2^{-N}\,\cdot)\cdot\widehat f\bigr)
 \longrightarrow f\quad\text{于 }\mathcal S',
\]
这些范数为零时所有块为零，继而 \(f=0\)。三角不等式分别来自 \(L^p\) 与 \(\ell^q\) 的三角不等式，所以在本节的指数范围内它们确实是范数。

\ProseParagraphBreak

三个参数分别回答三个问题：\(s\) 规定高频块要乘多大的尺度权；\(p\) 规定怎样量空间大小；\(q\) 规定允许多少个尺度同时显著。只写光滑阶数而不写后两个指数，通常不能判断函数是否属于该空间。

\subsection{Sobolev 空间的重新解释}
\label{b2:subsec:190502}

\Cref{b2:prop:19-inhomogeneous-identification}已证明
\[
 F^0_{p,2}(\mathbb R^d)=L^p(\mathbb R^d)
 \quad(1<p<\infty),
 \quad F^s_{2,2}(\mathbb R^d)=H^s(\mathbb R^d)
 \quad(s\in\mathbb R),
\]
等号指空间相同、范数等价。前一识别包括反向的函数代表问题：若温和分布的非齐次平方函数属于 \(L^p\)，合成估计使其块和在 \(L^p\) 中收敛，分布重构再将极限识别为原分布。不能只对已知属于 \(L^p\) 的函数证明不等式，就宣布对所有温和分布完成了空间识别。

\ProseParagraphBreak

在 \(p=q=2\) 时，Tonelli 允许交换位置积分与尺度求和，直接给出 \(B^s_{2,2}=F^s_{2,2}=H^s\)。一般 \(p=q<\infty\) 也有定义范数完全相等的恒等式
\[
 \sum_{j\geq0}2^{jsp}\|P_jf\|_p^p
 =\int_{\mathbb R^d}\sum_{j\geq0}2^{jsp}|P_jf(x)|^p\,\mathrm dx.
\]
这一次是两个积分次序的交换；\(H^s\) 与频块范数的比较还涉及频率权重，所以通常只有等价而不是数值相等。

\ProseParagraphBreak

另一个常用描述从符号 \((1+4\pi^2|\xi|^2)^{s/2}\) 定义 \(J^s\)，再要求 \(J^sf\in L^p\)。该符号光滑、各阶导数至多作多项式增长，因而可作用于温和分布。对于 \(1<p<\infty\)，把这个条件与 \(F^s_{p,2}\) 识别，需要统一控制环带上的归一化符号及其逆；这一乘子论证见\cref{b2:prop:22-bessel-square}。本章的平方可积情形已由 Plancherel 独立完成。齐次势 \(I_\alpha\) 在零频奇异，而 \(J^{-\alpha}\) 的符号在那里有限；这种低频差别正是全空间同指数估计不同的原因，不能把二者混作同一个算子。

\subsection{Besov 空间}
\label{b2:subsec:190503}

先说明所选分割如何影响 Besov 范数。若 \((\widetilde P_k)\) 来自另一套满足定义条件的截断，则高频新块只与固定个数的旧块相交，低频处也只有固定个数的例外。分布重构因此给出有限和
\[
 \widetilde P_kf=\sum_{j\sim k}\widetilde P_kP_jf.
\]
这里 \(j\sim k\) 表示 \(|j-k|\leq L\)，可增大固定整数 \(L\) 以覆盖低频例外。新块卷积核的 \(L^1\) 范数一致有界，且在这些指标上 \(2^{ks}\) 与 \(2^{js}\) 可比，故
\[
 2^{ks}\|\widetilde P_kf\|_p
 \leq C_s\sum_{j\sim k}2^{js}\|P_jf\|_p.
\]
对序列取 \(\ell^q\) 范数，有限次平移不改变该范数，便得到一个方向的界；交换新旧分割即得反向界。这个证明覆盖 \(1\leq p,q\leq\infty\)，常数依赖两套固定截断。

\ProseParagraphBreak

为了具体看见第二指数，取 \(s>0\)，在彼此相隔固定个数二进尺度的频率小球内选 \(f_n\in\mathcal S\)，使 \(\|f_n\|_2=1\)。记相应尺度为 \(j_n=Ln\)，其中固定 \(L\) 足够大。令
\[
 f=\sum_{n\geq1}a_n2^{-j_ns}f_n,
\]
并先设 \((a_n)\) 有界。不同球互不相交，正 \(s\) 使级数在 \(L^2\) 中收敛。每个球只与固定个数的频块相交，不同球对应的指标组互不交叠；在每组中，分割和为一，至少有一个块记录了固定比例的平方大小。Plancherel 和有限个数项的范数等价因此给出
\[
 \|f\|_{B^s_{2,q}}\asymp\|(a_n)\|_{\ell^q}.
\]
于是 \(a_n=1/n\) 给出属于 \(B^s_{2,2}\) 而不属于 \(B^s_{2,1}\) 的函数；\(a_n=1\) 给出属于 \(B^s_{2,\infty}\) 而不属于 \(H^s\) 的函数。频率位置相同的尺度权并不能替代尺度之间的可和性。

\ProseParagraphBreak

当光滑阶数小于一时，还可以不用 Fourier 变换来辨认一种 Besov 范数。频率块在步长远小于其波长时变化较小，在步长较大时只能用两倍振幅控制；下面的两段几何级数正来自这一区别。

\Needspace{6\baselineskip}
\begin{proposition}\label{b2:prop:19-besov-difference}
设整数 \(d\geq1\)、\(0<s<1\)、\(1\leq p<\infty\)，采用\cref{b2:def:19-besov-triebel-lizorkin}的非齐次块。对每个 \(f\in L^p(\mathbb R^d;\mathbb C)\)，有
\[
 \|f\|_{B^s_{p,p}}\asymp
 \|f\|_p+
 \left(\int_{\mathbb R^d}
 \frac{\|f(\,\cdot+h)-f\|_p^p}{|h|^{d+sp}}\,\mathrm dh\right)^{1/p}.
\]
允许两边同时为无穷；比较常数只依赖 \(d,s,p\) 与固定截断。因此该 Besov 空间与按右侧差分范数定义的分数 Sobolev 空间相同。
\end{proposition}
\begin{proof}
先设左边有限，记 \(u_j=P_jf\)、\(a_j=2^{js}\|u_j\|_p\)。由于 \(s>0\)，Hölder 的序列形式（\(p=1\) 时直接求和）给出
\[
 \sum_{j\geq0}\|u_j\|_p
 =\sum_{j\geq0}2^{-js}a_j
 \leq C_{s,p}\|(a_j)\|_{\ell^p}.
\]
所以块和在 \(L^p\) 中绝对收敛；分布重构说明其和就是 \(f\)，并且已经得到所需的 \(\|f\|_p\) 界。

沿线段积分及 Bernstein 估计给出 \(\|u_j(\cdot+h)-u_j\|_p\leq C2^j|h|\cdot\|u_j\|_p\)；平移不变性另给出两倍函数范数的界。因此
\[
 \|u_j(\,\cdot+h)-u_j\|_p
 \leq C\min\{1,2^j|h|\}\cdot\|u_j\|_p.
\]
在 \(2^{-k-1}<|h|\leq2^{-k}\)、\(k\geq0\) 上，对绝对重构求差并相加，得到
\[
 2^{ks}\|f(\,\cdot+h)-f\|_p
 \leq C\sum_{j\leq k}2^{-(k-j)(1-s)}a_j
       +C\sum_{j>k}2^{-(j-k)s}a_j.
\]
右边是两段可和几何序列与 \((a_j)\) 的卷积，离散 Young 不等式控制其 \(\ell^p_k\) 范数。每个步长环带上的 \(\int|h|^{-d}\,\mathrm dh\) 相同，且 \(|h|^{-sp}\) 与 \(2^{ksp}\) 可比，所以这个序列界正给出 \(|h|\leq1\) 的差分积分。\(|h|>1\) 的部分由 \(2\|f\|_p\) 以及 \(\int_{|h|>1}|h|^{-d-sp}\,\mathrm dh<\infty\) 控制。

反向设右边有限。对 \(j\geq1\)，令 \(k_j(h)=2^{jd}\check\varphi(2^jh)\)。其积分为零，因此在 \(L^p\) 中有
\[
 P_jf=\int k_j(h)\cdot\bigl(f(\,\cdot-h)-f\bigr)\,\mathrm dh.
\]
Minkowski 后对测度 \(|k_j(h)|\,\mathrm dh\) 使用 Hölder，得到
\[
 \|P_jf\|_p^p
 \leq\|k_j\|_1^{p-1}
       \int|k_j(h)|\cdot\|f(\,\cdot-h)-f\|_p^p\,\mathrm dh.
\]
\(p=1\) 时同式直接成立。Schwartz 衰减又给出
\[
 \sum_{j\geq1}2^{jsp}|k_j(h)|\leq C|h|^{-d-sp}\quad(h\ne0).
\]
具体地，写 \(r=|h|\)，每项至多为 \(Cr^{-d-sp}(2^jr)^{d+sp}(1+2^jr)^{-N}\)。在 \(2^jr\leq1\) 与 \(>1\) 两侧分别求几何级数，取 \(N>d+sp\) 即可。乘以差分范数后用 Tonelli，便控制全部高频项；低频块用 \(\|P_0f\|_p\leq\|\check\chi\|_1\cdot\|f\|_p\)。两向估计由 \(L^p\) 差分与各频带的范数联系起来。
\end{proof}
\(s<1\) 保证小步长对低频块的几何衰减可和，\(s>0\) 保证高频重构与大步长积分收敛。若 \(s\geq1\)，一般必须换用更高阶差分；对非零光滑梯度，\(s=1\) 的一阶差分积分已经出现 \(\mathrm dr/r\) 的发散。

\ProseParagraphBreak

这个判别可以在一个跳跃函数上精算。取单位立方体 \(Q=(0,1)^d\)。当每个 \(|h_i|<1\) 时，交集体积为 \(\prod_i(1-|h_i|)\)，因而对任意有限 \(p\geq1\)，
\[
 \|\mathbf1_Q(\,\cdot+h)-\mathbf1_Q\|_p^p
 =|Q\mathbin{\triangle}(Q-h)|
 =2\left(1-\prod_{i=1}^d(1-|h_i|)\right).
\]
若 \(|h|<1/2\)，括号介于 \(\max_i|h_i|\) 与 \(\sum_i|h_i|\) 之间，故与 \(|h|\) 可比。小步长积分于是与 \(\int_0^{1/2}r^{-sp}\,\mathrm dr\) 同时收敛。得到精确条件
\[
 \mathbf1_Q\in B^s_{p,p}\quad\Longleftrightarrow\quad sp<1
 \quad(0<s<1,\ 1\leq p<\infty).
\]
区域内外的函数都很简单，限制光滑阶数的是跨越边界时保留的跳跃。

\ProseParagraphBreak

绝对尺度求和还控制函数乘积。由 Bernstein，
\[
 \sum_{j\geq0}\|P_jf\|_\infty
 \leq C\sum_{j\geq0}2^{jd/p}\|P_jf\|_p
 =C\|f\|_{B^{d/p}_{p,1}}.
\]
所以 \(B^{d/p}_{p,1}\) 有一个有界连续代表：每块连续，上确界范数绝对可和，故块和一致收敛；分布重构识别这个代表。这也说明为什么临界处的序列指数一有特殊作用。

\Needspace{6\baselineskip}
\begin{proposition}\label{b2:prop:19-besov-algebra}
设整数 \(d\geq1\)、\(1\leq p<\infty\)、\(s\geq d/p\)，采用\cref{b2:def:19-besov-triebel-lizorkin}的非齐次 Besov 空间。则它的元素具有有界连续代表，其逐点乘积满足
\[
 \|fg\|_{B^s_{p,1}}
 \leq C_{d,s,p,\chi}\|f\|_{B^s_{p,1}}\cdot\|g\|_{B^s_{p,1}}.
\]
特别地，\(B^s_{p,1}(\mathbb R^d)\) 在乘法下封闭。
\end{proposition}
\begin{proof}
令 \(f_k=P_kf\)、\(g_l=P_lg\)。因为 \(s\geq d/p>0\)，两组块都在 \(L^p\) 和 \(L^\infty\) 中绝对求和，前述嵌入给出有界连续代表。双级数 \(\sum_{k,l}f_kg_l\) 在 \(L^p\) 中绝对收敛，并等于 \(fg\)。

每项乘积的 Fourier 支集包含于两频率球之和，故在半径 \(C2^{\max\{k,l\}}\) 的球内。这一支集规则可以对紧支集频率分布作卷积来理解；也可先用相同频率截断将 \(L^p\) 光滑逼近保持在略大的固定频带内，再把 Schwartz 函数的卷积公式传到分布极限。因而存在固定整数 \(L\)，使 \(P_j(f_kg_l)=0\) 当 \(j>\max\{k,l\}+L\)。

若 \(k\geq l\)，由块核的统一 \(L^1\) 界及 \(s>0\)，
\[
 \begin{aligned}
 \sum_{j\geq0}2^{js}\|P_j(f_kg_l)\|_p
 &\leq C\sum_{0\leq j\leq k+L}2^{js}\|f_k\|_p\cdot\|g_l\|_\infty\\
 &\leq C_s2^{ks}\|f_k\|_p\cdot\|g_l\|_\infty.
 \end{aligned}
\]
\(l>k\) 时交换两因子。先对有限双和估计，再用 \(L^p\) 收敛及 Fatou，得到
\[
 \begin{aligned}
 \|fg\|_{B^s_{p,1}}
 \leq C_s&\left(\sum_k2^{ks}\|f_k\|_p\right)
          \left(\sum_l\|g_l\|_\infty\right)\\
 +C_s&\left(\sum_k\|f_k\|_\infty\right)
          \left(\sum_l2^{ls}\|g_l\|_p\right).
 \end{aligned}
\]
两组上确界范数和均由对应 Besov 范数控制，即得结论。证明只用了 \(s\geq d/p\) 与 \(s>0\)，所以包括 \(s=d/p\) 的临界情形。
\end{proof}
若两个相乘的块频率相差很大，输出仍以较高频为主；若二者频率接近，却可能相互抵消而产生低频。上面求和的是所有 \(j\leq\max\{k,l\}+L\)，已经把这些低频输出包括在内，不能只保留高低频相互作用。这项乘积结论依赖已得到的函数代表，不能推及任意两个温和分布。

\subsection{Triebel--Lizorkin 空间}
\label{b2:subsec:190504}

这类范数先在位置 \(x\) 处比较不同频率，再作空间积分。若许多块集中在同一小区域，点态序列范数会同时记录它们；若这些块分布在彼此远离的区域，同一位置只见到少数项。先把每块的全空间范数算成一个数，会丢掉这种位置关系。两种次序可以比较，却通常不能交换。

\Needspace{6\baselineskip}
\begin{proposition}\label{b2:prop:19-besov-triebel-embeddings}
设整数 \(d\geq1\)、\(s\in\mathbb R\)、\(1\leq p<\infty\)、\(1\leq q\leq\infty\)，所有空间使用\cref{b2:def:19-besov-triebel-lizorkin}中的同一分割。则有连续嵌入
\[
 B^s_{p,\min\{p,q\}}\subset F^s_{p,q}
 \subset B^s_{p,\max\{p,q\}}.
\]
此外，对 \(1\leq p\leq r\leq\infty\)、\(1\leq q\leq\infty\)、\(s\in\mathbb R\)，有
\[
 B^s_{p,q}\subset B^{s-d(1/p-1/r)}_{r,q},
\]
其中 \(1/\infty=0\)。
\end{proposition}
\begin{proof}
先对有限个块令 \(a_j(x)=2^{js}|P_jf(x)|\)。若 \(q\leq p\)，Minkowski 在 \(L^{p/q}\) 中给出
\[
 \left\|\sum_j a_j^q\right\|_{p/q}
 \leq\sum_j\|a_j^q\|_{p/q},
\]
即 \(\|f\|_{F^s_{p,q}}\leq\|f\|_{B^s_{p,q}}\)。逐点的 \(\ell^p\leq\ell^q\) 又给出 \(\|f\|_{B^s_{p,p}}\leq\|f\|_{F^s_{p,q}}\)。若 \(p\leq q<\infty\)，Minkowski 在序列空间 \(\ell^{q/p}\) 中给出
\[
 \left(\sum_j\left[\int a_j^p\right]^{q/p}\right)^{p/q}
 \leq\int\left(\sum_j a_j^q\right)^{p/q},
\]
从而 \(\|f\|_{B^s_{p,q}}\leq\|f\|_{F^s_{p,q}}\)；逐点的 \(\ell^q\leq\ell^p\) 给出另一边。\(q=\infty\) 时直接使用
\[
 \sup_j\|a_j\|_p\leq\|\sup_j a_j\|_p
 \leq\left(\sum_j\|a_j\|_p^p\right)^{1/p}.
\]
有限块数趋于无穷，单调收敛证明第一组嵌入。

第二组嵌入对每个块使用 Bernstein：
\[
 2^{j[s-d(1/p-1/r)]}\|P_jf\|_r
 \leq C2^{js}\|P_jf\|_p.
\]
低频块尺度为一，满足同样估计。取 \(\ell^q\) 范数即得结论，包括两个无穷端点。
\end{proof}
这些都是范数不等式；只有另有严格性例子时，才可把包含号改成真包含。\(p=q\) 时第一组嵌入就是前面由 Tonelli 得到的相等，因而不可能一概宣称严格。

\ProseParagraphBreak

分割独立性也要区分指数。对 \(q=2\)、\(1<p<\infty\)，写
\(\widetilde P_kf=\sum_{|j-k|\leq L}\widetilde P_kP_jf\)，并将有限个相邻偏移分别处理。每个固定偏移对应一族 \(\widetilde P_k\)：其缩放核有统一的标准核常数，平方范数一致有界，所以\cref{b2:prop:18-hilbert-valued-cz}控制 \(L^p(\ell^2)\) 范数。权重 \(2^{ks}\) 与 \(2^{js}\) 在相邻指标上可比，得到新旧 \(F^s_{p,2}\) 范数等价。\(p=1\) 或一般 \(q\) 不在这一 Hilbert 值强型论证的范围内。这里的定义仍使用固定分割，前面的差分、嵌入和例子也均已在声明的分割下证明；一般 Triebel--Lizorkin 分割独立性的进一步理论可参看 Hans Triebel 的《Theory of Function Spaces》\cite{Triebel1983Theory}，不能用标量 Young 不等式代替其中的向量估计。

\subsection{Besov 空间的平面迹与延拓}
\label{b2:subsec:190505}

最后考虑在 \(x_d=0\) 上取边界值。一个频率为 \(2^j\) 的块，在法向上可变化的长度约为 \(2^{-j}\)。把这一层的 \(d\) 维积分变成边界的 \((d-1)\) 维积分，函数范数会出现 \(2^{j/p}\)，因此 \(s\) 阶的自然边界阶数为 \(s-1/p\)。下面把这项尺度计算落实为一个有界迹与右逆，范围限定为平坦边界、序列指数一。

\Needspace{6\baselineskip}
\begin{proposition}[平坦边界的 Besov 迹与延拓]
\label{b2:prop:19-flat-besov-trace}
设整数 \(d\geq2\)、\(1\leq p<\infty\)、\(s>1/p\)。在 \(\mathbb R^d\) 与 \(\mathbb R^{d-1}\) 上各固定\cref{b2:def:19-besov-triebel-lizorkin}的一套非齐次分割。Schwartz 函数上的限制 \(u\mapsto u(\,\cdot,0)\) 有唯一有界线性延拓
\[
 \operatorname{Tr}:B^s_{p,1}(\mathbb R^d)
 \longrightarrow B^{s-1/p}_{p,1}(\mathbb R^{d-1}),
\]
并存在有界线性算子 \(E\) 从目标空间映回 \(B^s_{p,1}(\mathbb R^d)\)，满足 \(\operatorname{Tr}Eg=g\)。各常数只依赖固定参数和截断。
\end{proposition}
\begin{proof}
令 \(u_j=P_ju\)。其法向频率不超过 \(C2^j\)。取一维 Schwartz 函数 \(k\)，使 \(\widehat k=1\) 于 \([-C,C]\)，则由法向频率截断，
\[
 u_j(x',0)=\int_{\mathbb R}2^jk(-2^jt)\cdot u_j(x',t)\,\mathrm dt.
\]
对每个 \(x'\) 使用 Hölder，\(p'=p/(p-1)\) 给出核的 \(L^{p'}\) 范数 \(2^{j/p}\|k\|_{p'}\)；\(p=1\) 时使用 \(p'=\infty\)。再积分切向变量，得到
\[
 \|u_j(\,\cdot,0)\|_{L^p(\mathbb R^{d-1})}
 \leq C2^{j/p}\|u_j\|_{L^p(\mathbb R^d)}.
\]
卷积表示给出所用的法向连续代表；频率块本身光滑，因此这就是其普通限制。边界频率仍在半径 \(C2^j\) 的球内。记边界块为 \(P'_k\)，其核的 \(L^1\) 范数一致有界，而且 \(P'_k[u_j(\cdot,0)]=0\) 当 \(k>j+L\)。于是
\[
 \begin{aligned}
 \|u_j(\,\cdot,0)\|_{B^{s-1/p}_{p,1}}
 &\leq C\sum_{0\leq k\leq j+L}2^{k(s-1/p)}
              \cdot2^{j/p}\|u_j\|_p\\
 &\leq C_s2^{js}\|u_j\|_p.
 \end{aligned}
\]
最后的几何级数使用了 \(s-1/p>0\)。对 \(j\) 求和，边界块的限制之和在 \(L^p\) 中绝对收敛，记为 \(\operatorname{Tr}u\)；逐个边界块取极限并用 Fatou，得到目标 Besov 界。

补足延拓的唯一性。因为序列指数为一，有限频率和在 \(B^s_{p,1}\) 中稠密：对剩余块使用有限重叠及统一卷积界，其范数受原尺度序列的尾和控制。若一个有限频率和 \(v\in L^p\) 的频率支集在固定球内，选光滑紧频率乘子 \(\eta\)，在该球邻域内为一。取 \(h_n\in C_c^\infty\) 在 \(L^p\) 中逼近 \(v\)，则 \(\eta(D)h_n\in\mathcal S\)，且
\[
 \|\eta(D)h_n-v\|_p
 \leq\|\check\eta\|_1\cdot\|h_n-v\|_p\longrightarrow0.
\]
误差有共同有界频带，故只含有限个非齐次块，\(L^p\) 收敛也给出 \(B^s_{p,1}\) 收敛。由此证明 Schwartz 类稠密。对 Schwartz 函数，非齐次截断在 Schwartz 拓扑中趋于原函数，边界限制与所定义的块和相同，所以所得连续延拓唯一。

为构造右逆，将 \(g\) 分解为边界块 \(g_j=P'_jg\)。选一维 Schwartz 函数 \(\rho\)，使 \(\rho(0)=1\) 且 \(\widehat\rho\) 紧支撑，令
\[
 v_j(x',t)=g_j(x')\cdot\rho(2^jt),
 \quad Eg=\sum_{j\geq0}v_j.
\]
每项 Fourier 支集在半径 \(C2^j\) 的球内，且
\[
 \|v_j\|_{L^p(\mathbb R^d)}
 =2^{-j/p}\|g_j\|_{L^p(\mathbb R^{d-1})}\cdot\|\rho\|_{L^p(\mathbb R)}.
\]
对体内输出块再求和，由 \(s>0\) 得
\[
 \|v_j\|_{B^s_{p,1}(\mathbb R^d)}
 \leq C_s2^{js}\|v_j\|_p
 \leq C2^{j(s-1/p)}\|g_j\|_p.
\]
右侧可和，所以 \(\sum v_j\) 在 \(L^p\) 中收敛，并由块估计与 Fatou 得到 \(Eg\in B^s_{p,1}\) 及相应范数界；同一尾和估计也证明级数在 Besov 范数中收敛。每个 \(v_j\) 只有有限多个体内频率块，其块和就是本身，故由定义 \(\operatorname{Tr}v_j=v_j(\cdot,0)=g_j\)。连续性将有限和等式传给极限，得到 \(\operatorname{Tr}Eg=\sum_jg_j=g\)。
\end{proof}
临界 \(s=1/p\) 时，边界输出尺度的求和不再有几何衰减，不能由本证明得到同样的 Besov 目标；但是只求边界 \(L^p\) 范数时，绝对求和还能保留一个临界版本，见\cref{b2:ex:19-new-3}。曲边界则另需坐标变换的有界性，不能把平面上的公式直接搬到任意区域。迹的尺度损失、序列指数与边界几何，各自承担不同的条件。


\begin{chaptersummary}
非驻相区域可反复分部积分，非退化驻点提供每个方向半个幂次的衰减。van der Corput 引理把高阶退化与较慢衰减联系起来。对输入允许任意 \(L^2\) 振荡的算子，混合 Hesse 矩阵 与 \(TT^*\) 才能揭示统一增益。局部化、振幅导数、边界处理和参数下界共同决定常数；它们都是估计的一部分。
\end{chaptersummary}

\BookExercises

% 自拟：围绕本章已经建立的工具，改变条件或组合方法。
\begin{exercise}\label{b2:ex:21-new-1}
计算 \(\int_0^1e^{2\pi\mathrm i\lambda t}\,\mathrm dt\)。说明为什么这个例子不能仅凭相位无驻点就宣称任意阶衰减。
\end{exercise}
\begin{solution}
积分为 \((e^{2\pi\mathrm i\lambda}-1)/(2\pi\mathrm i\lambda)\)。沿半整数参数，其绝对值为 \(1/(\pi\lambda)\)。振幅在边界不消失，分部积分留下边界项，因此非驻相命题的紧支撑条件不满足。
\end{solution}

% 自拟：围绕本章已经建立的工具，改变条件或组合方法。
\begin{exercise}\label{b2:ex:21-new-2}
设 \(\phi''\ge c>0\)、\(a\) 绝对连续，证明 \(I_\lambda=O((\lambda c)^{-1/2})\)，并说明区间长度为什么无需进入主常数。
\end{exercise}
\begin{solution}
对相位 \(\phi/c\) 使用二阶 van der Corput 引理，参数为 \(\lambda c\)，得上界 \(C(\lambda c)^{-1/2}(\|a\|_\infty+\|a'\|_1)\)。次水平集 \(|\phi'|\le\delta\) 的长度由导数下界控制为 \(2\delta/c\)，其余区间使用单调第一阶估计，所以没有额外的区间长度项。
\end{solution}

% 自拟：围绕本章已经建立的工具，改变条件或组合方法。
\begin{exercise}\label{b2:ex:21-new-3}
对 \(a\in C_c^\infty(\mathbb R)\)，求 \(\int e^{\pi\mathrm i\lambda t^2}a(t)\,\mathrm dt\) 的前两项。
\end{exercise}
\begin{solution}
二次模型公式给出 \(\lambda^{-1/2}e^{\mathrm i\pi/4}[a(0)+\mathrm i a''(0)/(4\pi\lambda)]+O(\lambda^{-5/2})\)。若 \(a(0)=0\) 而 \(a''(0)\ne0\)，领先阶变为 \(\lambda^{-3/2}\)。
\end{solution}

% 自拟：围绕本章已经建立的工具，改变条件或组合方法。
\begin{exercise}\label{b2:ex:21-new-4}
设 \(H\) 是实对称可逆矩阵，计算相位 \(\phi(x)=x^{\mathsf T}Hx/2\) 的驻相主项，分别讨论正定与不定的情形。
\end{exercise}
\begin{solution}
主项为 \(\lambda^{-d/2}e^{\mathrm i\pi\operatorname{sgn}H/4}a(0)|\det H|^{-1/2}\)。正定时符号为 \(d\)；若有 \(r\) 个正、\(d-r\) 个负特征值，符号为 \(2r-d\)。绝对值系数相同，复相位不同。
\end{solution}

% 自拟：围绕本章已经建立的工具，改变条件或组合方法。
\begin{exercise}\label{b2:ex:21-new-5}
设 \(a\in C_c^\infty(\mathbb R^d)\)。证明相位 \(\xi\cdot x\) 的积分在 \(|\xi|\to\infty\) 时任意阶衰减，并把它与 Schwartz 函数的 Fourier 变换联系起来。
\end{exercise}
\begin{solution}
选择坐标 \(j\) 使 \(|\xi_j|\ge|\xi|/\sqrt d\)，沿该方向分部积分 \(N\) 次，得 \(|\int e^{2\pi\mathrm i\xi\cdot x}a|\le(2\pi|\xi_j|)^{-N}\|\partial_j^Na\|_1\)。对所有 \(j\) 取最大值即可。乘以任意坐标幂后仍适用同样估计，因此 Fourier 变换属于 Schwartz 空间。
\end{solution}

% 自拟：围绕本章已经建立的工具，改变条件或组合方法。
\begin{exercise}\label{b2:ex:21-new-6}
对 \(\Phi_\epsilon(x,y)=\epsilon x\cdot y\)，解释 \(\epsilon\downarrow0\) 时振荡算子估计为何不可能有统一的 \(\lambda^{-d/2}\) 常数。
\end{exercise}
\begin{solution}
取分离振幅 \(a(x,y)=\chi(x)\chi(y)\) 及非零固定非负输入 \(f\)。当 \(\lambda\epsilon\to0\)，支配收敛使输出趋于 \(\chi(x)\int\chi f\)，其范数为正。若取 \(\epsilon=\lambda^{-2}\)，统一的 \(\lambda^{-d/2}\) 界将使该范数趋零，矛盾。真正参数是 \(\lambda\epsilon\)。
\end{solution}

% 自拟：围绕本章已经建立的工具，改变条件或组合方法。
\begin{exercise}\label{b2:ex:21-new-7}
在小邻域内考虑 \(\Phi(x,y)=x\cdot y+|x|^2|y|^2\)。证明其满足\chapref{b2:ch:21}的振荡 \(L^2\) 估计。
\end{exercise}
\begin{solution}
混合 Hesse 矩阵 为 \(I+4xy^{\mathsf T}\)。当 \(|x||y|<1/8\) 时扰动算子范数小于 \(1/2\)，最小奇异值至少 \(1/2\)。振幅紧支撑于该邻域后，\cref{b2:thm:21-oscillatory-l2}给出 \(\|T_\lambda\|_{2\to2}\le C\lambda^{-d/2}\)。仅固定一个变量检查二阶 Hesse 矩阵 并不是这里所需的条件。
\end{solution}

% 自拟：围绕本章已经建立的工具，改变条件或组合方法。
\begin{optionalexercise}\label{b2:ex:21-new-8}
对半波相位 \(x\cdot\theta+t|\theta|-y\cdot\theta\)，计算典范关系，并说明 \(\theta=0\) 为什么应另行处理。
\end{optionalexercise}
\begin{solution}
驻点方程为 \(x+t\theta/|\theta|-y=0\)，两端余切向量均为 \(\theta\)，所以关系为 \((x,\theta;y,\theta)\)、\(y=x+t\theta/|\theta|\)、\(\theta\ne0\)。零点处 \(|\theta|\) 不光滑，齐次高频相位表示不适用；低频光滑截断产生的核可直接处理，不属于所讨论的高频典范关系。
\end{solution}

% !TeX root = ../../Book2.tex
\chapter{Fourier 乘子}
\label{b2:ch:22}
\BookChapterNumber{b2:ch:22}
\begin{chapterabstract}
本章把 Fourier 侧的有界符号转化为 \(L^p\) 算子。完整证明一个有限导数版本的 Mihlin 定理及 Sobolev 环带版本的 Hörmander 定理，并应用于 Riesz 变换、Bessel 位势、平方函数正则性和常系数椭圆估计。
\end{chapterabstract}

\begin{learninggoals}\begin{itemize}
\item\label{b2:item:22-goal-1}掌握 \(L^2\) 乘子的准确刻画及尺度、对偶和复合规则。
\item\label{b2:item:22-goal-2}能把尺度导数条件转换成频率环带核的统一估计。
\item\label{b2:item:22-goal-3}理解 Hörmander 积分差分条件如何取代核的逐点光滑条件。
\item\label{b2:item:22-goal-4}能区分同指数乘子界、分数积分界和带导数的椭圆先验估计。
\end{itemize}\end{learninggoals}

卷积算子在 Fourier 一侧成为乘法，但 \(L^p\) 范数并不被 Fourier 变换保持。\chapref{b2:ch:18}建立的奇异积分理论与\chapref{b2:ch:19}的频率分解共同弥补这一缺口：每个环带有快速衰减核，所有尺度的叠加则通过零阶奇异积分处理。以下定理中的导数数目是明确的充分条件，凡未证明的锐利改进都不纳入正文前置。

% !TeX root = ../../Book2.tex
\section{乘子算子}
\label{b2:sec:2201}

平移不变算子在 Fourier 一侧表现为逐点乘法。本章讨论怎样从符号的局部光滑性得到整个 \(L^p\) 范数界，并区别阶数为零的有界乘子、改变可积指数的分数积分和具有导数损失的微分算子。

\subsection{Fourier 乘子}
\label{b2:subsec:220101}

\begin{definition}\label{b2:def:22-multiplier}
对可测函数 \(m\)，先在使乘积有意义的测试函数上定义
\[
 T_m f=\mathcal F^{-1}(m\widehat f).
\]
若该算子延拓为 \(L^p(\mathbb R^d)\) 上的有界算子，称 \(m\) 为 \(L^p\) Fourier 乘子，并记其算子范数为 \(\|m\|_{M_p}\)。
\end{definition}
乘子在零测集上的值不影响算子。对于一般有多项式增长的光滑 \(m\)，乘法可在 Schwartz 空间与温和分布之间解释；对于仅可测的有界 \(m\)，最稳妥的起点是 \(L^2\)。
若 \(m=\widehat k\)、\(k\in L^1\)，则 \(T_mf=k*f\)，Young 不等式给出所有 \(1\le p\le\infty\) 上的有界性。

\ProseParagraphBreak
乘法和卷积的对应不意味着每个乘子都有普通可积核。恒等乘子的核为 \(\delta_0\)，Hilbert 乘子的核为主值分布。推导核估计时应先局部化频率，再在原点外重构普通函数；原点上的分布项要由已有的 \(L^2\) 定义处理。

\subsection{\texorpdfstring{\(L^2\)}{L2} 理论}
\label{b2:subsec:220102}

\begin{proposition}\label{b2:prop:22-l2}
可测 \(m\) 是 \(L^2\) 乘子当且仅当 \(m\in L^\infty\)，且
\(\|m\|_{M_2}=\|m\|_\infty\)。
\end{proposition}
\begin{proof}
Plancherel 给出充分性及范数上界。若算子范数为 \(C\)，在任意有限测度集合 \(E\) 上选 \(\widehat f=\mathbf1_E\)，得
\(\int_E|m|^2\le C^2|E|\)。若 \(|m|>C+\epsilon\) 的集合有正测度，取其中有限测度的正测度子集就矛盾。因此 \(|m|\le C\) 几乎处处，得到必要性和等号。
\end{proof}
对 \(|m|=1\) 的符号，算子在 \(L^2\) 上等距；这并不提供其他 \(p\) 上的范数控制。即使每个频率只发生相位旋转，不同频率叠加后的空间集中仍可能发生变化。

\ProseParagraphBreak
若 \(m_1,m_2\in M_p\)，则 \(m_1m_2\in M_p\)，且范数不超过两范数之积。先在 \(L^2\cap L^p\) 上验证 \(T_{m_1m_2}=T_{m_1}T_{m_2}\)，再用稠密性延拓即可。相似地，乘子空间对一致算子范数有界的点态极限封闭：在 \(L^2\cap L^p\) 上先取 \(L^2\) 极限及几乎处处子列，再用 Fatou，最后延拓。这是处理光滑截断极限的常用办法。

\subsection{\texorpdfstring{\(L^p\)}{Lp} 乘子问题}
\label{b2:subsec:220103}

若 \(m\in M_p\)、\(1<p<\infty\)，其 Banach 伴随在 \(L^{p'}\) 上的符号为 \(\overline m\)，范数相同。若 \(m\) 实值，就得到对偶指数上的相同范数。与 \(L^2\) 界插值，可控制 \(p\) 和 \(2\) 之间的全部指数，但不能越过尚未建立的端点。

\ProseParagraphBreak
缩放不改变乘子范数。若 \(D_rf(x)=f(rx)\)，则
\(T_{m(r\,\cdot)}=D_{1/r}T_mD_r\)，而两个缩放范数相消。
正交变换与频率平移也保持范数：前者对应空间坐标旋转，后者对应调制。符号乘上 \(e^{-2\pi\mathrm i y\cdot\xi}\) 则对应输出平移。这些等距对称使乘子条件应具有尺度不变性；要求所有导数在整个频率空间绝对有界，并不是适合齐次问题的条件。

\ProseParagraphBreak
一个不光滑却有界的例子是一维区间特征函数。半轴投影可用恒等算子与 Hilbert 变换表示，两个频率平移后的半轴投影之差就是区间投影，因此在 \(1<p<\infty\) 上一致有界。由此不能推断高维球的特征函数同样有界。符号边界的几何维数与曲率会参与问题，\chapref{b2:ch:24}将专门说明这一障碍。

% 扩写定位：b2:subsec:220103
Fourier 乘子并不是凭空选出的一类算子；在 \(L^2\) 上，它们正好是与所有平移交换的有界线性算子。
\begin{proposition}\label{b2:prop:22-translation-multiplier}
若 \(T:L^2(\mathbb R^d)\to L^2(\mathbb R^d)\) 有界，且 \(T\tau_y=\tau_yT\) 对全部 \(y\) 成立，其中 \(\tau_yf(x)=f(x-y)\)，则存在唯一的 \(m\in L^\infty\)，使 \(T=T_m\)，并且 \(\|T\|_{2\to2}=\|m\|_\infty\)。
\end{proposition}
\begin{proof}
令 \(A=\mathcal FT\mathcal F^{-1}\)。平移变成乘法 \(e^{-2\pi\mathrm iy\cdot\xi}\)，故 \(A\) 与所有这些乘法算子交换。对任意 \(\varphi\in C_c^\infty\)，Fourier 反演把 \(\varphi\) 写成这些指数乘法的 \(L^1\) 加权积分；其在 \(L^2\) 上的强积分与 \(A\) 交换。于是 \(A\) 与乘法 \(M_\varphi\) 交换。

由光滑函数的有界几乎处逼近，\(A\) 也与任意可测集合的投影 \(M_{\mathbf1_E}\) 交换。可以先对有界集合用 Lebesgue 正则性与卷积逼近得到 \(0\leq\varphi_n\leq1\)、\(\varphi_n\to\mathbf1_E\) 几乎处，支配收敛给出强算子收敛；再以球截断推广到一般集合。

取递增的有限测度立方体 \(Q_n\) 覆盖全空间。在 \(Q_n\) 上定义 \(m_n=A\mathbf1_{Q_n}\)。若 \(n<k\)，交换投影给出
\(\mathbf1_{Q_n}m_k=A(\mathbf1_{Q_n}\mathbf1_{Q_k})=m_n\)，所以拼成可测函数 \(m\)。对 \(E\subset Q_n\)，有
\[
 A\mathbf1_E=\mathbf1_E m,\qquad
 \int_E|m|^2\leq\|A\|^2|E|.
\]
若 \(\{|m|>\|A\|+\epsilon\}\cap Q_n\) 有正测度，将其代入即矛盾。因此 \(\|m\|_\infty\leq\|A\|\)。有限支集简单函数的稠密性说明 \(A=M_m\)，反向范数与唯一性由本节的 \(L^2\) 乘子结论得到。
\end{proof}
该结论只建立在 \(L^2\) 上。若算子只在 \(L^p\) 给定，还不能通过一次 Fourier 变换就把它视为 \(L^2\) 有界乘子；需先说明它在共同稠密类上的作用及其额外的平方界。另一方面，具有不连续符号的平移不变算子在 \(L^2\) 上完全没有困难；真正的乘子问题是同一个算子能否保持其他 \(L^p\)。

% !TeX root = ../../Book2.tex
\section{Mihlin–Hörmander 定理}
\label{b2:sec:2202}

导数条件的作用是使每个频率环带上的核有足够快的衰减，再让环带总和满足奇异积分的抵消估计。证明须同时使用符号的有界性来建立 \(L^2\) 起点；核的大小条件本身不能替代这个起点。

\subsection{符号导数}
\label{b2:subsec:220201}

设 \(N\) 为整数。常用的尺度不变符号界为
\[
 |\partial^\alpha m(\xi)|\le A|\xi|^{-|\alpha|},
 \qquad 0\le|\alpha|\le N,\quad \xi\ne0.
\]
零阶包含 \(|m|\le A\)。例如 \(\xi_j/|\xi|\) 的各阶导数都是相应次数的齐次函数，在单位球面连续有界，因此满足任意有限阶条件。相反，\(|\xi|^{-\alpha}\) 在零频无界，不能成为 \(L^2\) 上的有界零阶乘子。

\ProseParagraphBreak
若 \(m(\xi)=b(\log|\xi|)\)，而 \(b\) 的有限阶导数有界，则链式法则给出同样的尺度界。符号不需要在零点有极限，甚至可以沿尺度持续振荡。这里对零点的忽略不影响 \(L^p\) 算子，却提醒我们不能把定理误述成“符号在闭空间上光滑”。

\subsection{二进分解}
\label{b2:subsec:220202}

使用\chapref{b2:ch:19}的环带分割 \(\varphi(2^{-j}\xi)\)，令
\[
 m_j(\xi)=m(\xi)\varphi(2^{-j}\xi),\qquad
 k_j(x)=2^{jd}h_j(2^jx),\quad
 h_j=\mathcal F^{-1}\bigl[m(2^j\,\cdot)\varphi\bigr].
\]
在固定环带上，归一化符号的有限阶导数一致有界。对 Fourier 积分分部积分，并分别选择绝对值最大的坐标，可得
\[
 |h_j(x)|+|\nabla h_j(x)|\le C A(1+|x|)^{-N}.
\]
梯度估计只是把符号乘上有界环带中的 \(2\pi\mathrm i\xi_\ell\)，仍使用相同阶数的符号导数；不需要额外的一阶 \(m\) 导数。该估计对所有 \(j\) 一致。

\ProseParagraphBreak
当 \(N>d+1\)，在 \(x\ne0\) 上绝对求和：
\[
 \sum_j|k_j(x)|\le CA|x|^{-d},\qquad
 \sum_j|\nabla k_j(x)|\le CA|x|^{-d-1}.
\]
例如第一式把 \(2^j|x|\le1\) 与 \(>1\) 分开，分别求等比级数
\(\sum2^{jd}\) 和 \(|x|^{-N}\sum2^{j(d-N)}\)。
第二式同理需要 \(N>d+1\)。这些级数在离原点的紧集上一致收敛，从而其和 \(K\) 是那里的一阶光滑函数。

\subsection{Mihlin 乘子定理}
\label{b2:subsec:220203}

下面采用\term[mihlin-chengzi-dingli]{Mihlin 乘子定理}[Mihlin multiplier theorem]的下述版本。

\begin{theorem}[Mihlin 乘子定理的有限导数版本]\label{b2:thm:22-mihlin}
设 \(m\in C^{d+2}(\mathbb R^d\setminus\{0\})\)，且对
\(|\alpha|\le d+2\) 满足上述尺度导数界。则
\[
 \|T_mf\|_p\le C_{d,p}A\|f\|_p,\qquad1<p<\infty.
\]
该导数数目是本证明使用的方便充分条件，并不声称最优。
\end{theorem}
\begin{proof}
Plancherel 给出 \(L^2\) 范数至多 \(A\)。有限环带和的符号在几乎处处频率上趋于 \(m\)，且一致有界，因此对应算子在 \(L^2\) 中趋于 \(T_m\)。
取紧支撑 \(f\in L^2\)，其支撑之外的任意紧集到支撑的距离为正。前述核级数在那里一致绝对可积，可交换求和和积分，得到
\(T_mf(x)=\int K(x-y)f(y)\,\mathrm dy\)。
更准确地，先在两个不交支撑的测试函数之间配对，再取 \(L^2\) 极限，即确认这就是给定算子的支撑外核，而不是另行定义的算子。

核的大小与梯度界满足\chapref{b2:ch:18}的标准 Calderón--Zygmund 条件，且卷积核的两变量差分由同一个梯度界控制。\cref{b2:thm:18-cz-lp}因而给出所有 \(1<p<\infty\) 上的界。对 \(L^2\cap L^p\) 的界先建立，再以该空间在 \(L^p\) 中的稠密性定义唯一延拓。
\end{proof}
证明实际上还给出弱型 \((1,1)\)，但通常没有强型 \((1,1)\) 或 \(L^\infty\to L^\infty\) 界。Hilbert 变换的区间输入产生对数尾和边界对数奇性，已分别说明这两个端点为何不能由定理自动补齐。

\ProseParagraphBreak
有限频率环带上的导数条件也可使用：若 \(m\) 支撑在固定环带并有 \(d+1\) 阶导数，则它的核属于 \(L^1\)，从而包括端点在内都可有界。真正造成端点困难的是无限尺度叠加，不能把单个环带的 Young 估计在没有可求和常数时直接相加。

\subsection{Hörmander 条件}
\label{b2:subsec:220204}

导数的逐点控制可改为环带上的 Sobolev 控制。固定 \(\eta\in C_c^\infty\)，在 \(\varphi\) 支撑邻域等于 \(1\)，令
\[
 B_s(m)=\sup_{j\in\mathbb Z}\|m(2^j\,\cdot)\eta\|_{H^s(\mathbb R^d)}.
\]
这里 \(H^s\) 的变量是频率变量本身；此处的 Sobolev 范数不是输入 \(f\) 的范数。
下面采用\term[hormander-chengzi-dingli]{Hörmander 乘子定理}[Hörmander multiplier theorem]的下述版本。

\begin{theorem}[Hörmander 环带条件]\label{b2:thm:22-hormander}
若 \(s>d/2\)、\(m\in L^\infty\) 且 \(B_s(m)<\infty\)，则 \(m\in M_p\)、\(1<p<\infty\)，算子范数不超过
\(C_{d,p,s}(\|m\|_\infty+B_s(m))\)。
\end{theorem}
\begin{proof}
设 \(h_j\) 如上一小节。Sobolev 范数经 Plancherel 等价于
\(\|(1+|x|)^s h_j\|_2\)；固定光滑截断的乘法在 \(H^s\) 有界，故这些范数一致受 \(CB_s\) 控制。
证明这一乘法事实，只需把乘积的 Fourier 变换写为卷积，使用
\((1+|x|)^s\le C_s(1+|x-y|)^s(1+|y|)^s\) 及 Schwartz 核的带权 \(L^1\) 范数。
同样，乘上环带中的坐标 \(\xi_\ell\) 表明
\(\|(1+|x|)^s\nabla h_j\|_2\le CB_s\)。
由 Cauchy–Schwarz，
\[
 \|h_j\|_1+\|\nabla h_j\|_1\le CB_s,\qquad
 \int_{|x|>R}|h_j(x)|\,\mathrm dx\le CB_sR^{-(s-d/2)}
 \quad(R\ge1).
\]

记 \(\gamma=s-d/2>0\)。对位移 \(v\ne0\)，低频环带 \(2^j|v|\le1\) 满足
\(\|k_j(\,\cdot-v)-k_j\|_1\le CB_s2^j|v|\)。
高频环带 \(2^j|v|>1\) 在 \(|x|>2|v|\) 上用两个尾积分，得到
\[
 \int_{|x|>2|v|}|k_j(x-v)-k_j(x)|\,\mathrm dx
 \le CB_s(2^j|v|)^{-\gamma}.
\]
两段等比级数可求和，故每个有限环带和的核及其伴随都满足一致的积分 Hörmander 条件。

\chapref{b2:ch:18}弱型证明中，坏块的估计只用该积分条件：用零平均把核减去块中心的值，再把差分积分界乘以块的 \(L^1\) 范数。好部仍用一致的 \(L^2\) 范数。因此同一证明给出有限和的一致弱型界；Marcinkiewicz 插值和伴随对偶给出一致的 \(L^p\) 界。
有限和在 \(L^2\) 中趋于 \(T_m\)，取几乎处处子列并用 Fatou，便将 \(L^p\) 界传到极限，最后用稠密性延拓。
这条证明不要求极限核有逐点梯度界，因而确实比上一版本弱。
\end{proof}
严格条件 \(s>d/2\) 用于尾积分中的正指数 \(\gamma\)。在等号处，高频环带级数不再衰减；这说明本证明失效的位置，不能把失效本身当作等号处任何可能强化的反例。

% 扩写定位：b2:subsec:220204
有限阶条件中的尺度因子不可删除。若在频率 \(\xi\) 附近的变化尺度是 \(|\xi|\)，则一次求导自然带来 \(|\xi|^{-1}\)；Mihlin 假设以 \(|\xi|^{|\alpha|}\) 恢复这个尺度。普通的“全部导数有界”在无穷频率未必提供这种衰减，也就不能直接控制重缩放后的二进核。

\ProseParagraphBreak

例如 \(m(\xi)=e^{2\pi\mathrm ia\cdot\xi}\) 的乘子是平移，故在每个 \(L^p\) 上都是等距。但其导数在无穷处不衰减，除非 \(a=0\)，不满足 Mihlin 的尺度导数条件。这个例子表明定理提供的是充分条件，而不是所有有界乘子的分类。反过来，把一个任意有界但高度不规则的符号放入 \(L^\infty\)，只能保证 \(L^2\) 有界，不能得到其他指数。

\ProseParagraphBreak

还应区分单个乘子和一族乘子的界。若 \(m_\nu\) 的 Mihlin 常数一致，则标量算子范数一致；若要在位置上先对输出求平方和，需要更强的向量论证。不过在此处统一符号条件已经足以给出这一论证：对有限族取对角乘子 \(\operatorname{diag}(m_\nu)\)，归一化后的核的算子范数由各个标量核上界的上确界控制，\(L^2(\ell^2)\) 界由逐项 Plancherel 相加得到。\cref{b2:prop:18-hilbert-valued-cz}于是给出
\[
 \left\|\left(\sum_\nu|T_{m_\nu}f_\nu|^2\right)^{1/2}\right\|_p
 \leq C_p\left\|\left(\sum_\nu|f_\nu|^2\right)^{1/2}\right\|_p.
\]
该常数与族的大小无关。若只知道各个算子的 \(L^p\) 范数一致，而不知道统一核或其他向量结构，则不能凭空推出上述平方和界。这也是前面 Bessel 平方函数识别时应核对的额外条件。

% !TeX root = ../../Book2.tex
\section{应用}
\label{b2:sec:2203}

应用乘子定理时，先判断符号阶数，再选取匹配的函数空间。零阶符号可保持 \(L^p\)，负阶符号在不同频率区间可能改善正则性或改变可积指数，正阶符号则需输入具有相应导数。

\subsection{Riesz 变换}
\label{b2:subsec:220301}

Riesz 变换的符号 \(-\mathrm i\xi_j/|\xi|\) 满足任意阶的尺度导数条件，故\chapref{b2:ch:18}的 \(L^p\) 定理也可由 Mihlin 定理再次得到。两条路线揭示不同方面：奇异积分证明直接看到零平均和核尾部；乘子证明利用频率光滑性，以 \(L^2\) 自动建立算子起点。

\ProseParagraphBreak
还可处理任意光滑零次齐次符号 \(m\)，因为单位球面上的有限阶导数有界。零次齐次性并不要求球面平均为零；常数符号的核有一个原点质量。只有在提取原点外的奇异核时，才需要辨认其球面抵消，不能把抵消条件错误地加到所有零阶符号上。

\subsection{分数幂}
\label{b2:subsec:220302}

记
\[
 J^sf=\mathcal F^{-1}\bigl[(1+4\pi^2|\xi|^2)^{s/2}\widehat f\bigr],
 \qquad (-\Delta)^{s/2}f
 =\mathcal F^{-1}\bigl[(2\pi|\xi|)^s\widehat f\bigr].
\]
前者是不齐次分数阶算子，后者是齐次分数阶算子。对 \(s>0\)，\(J^{-s}\) 的符号满足 Mihlin 条件，故在 \(L^p\) 上有界；\((-\Delta)^{-s/2}\) 在零频发散，应在适当齐次空间或以\chapref{b2:ch:19}的分数积分解释，不能据此声称它在同一个 \(L^p\) 上有界。

\ProseParagraphBreak
对 \(1<p<\infty\)，定义 Bessel 位势空间
\[
 H^{s,p}=\{f\in\mathcal S':J^sf\in L^p\},\qquad
 \|f\|_{H^{s,p}}=\|J^sf\|_p.
\]
这是完整空间，因为 \(J^{-s}\) 把 \(L^p\) 连续送入温和分布，\(L^p\) 极限可经它反解。连续性的核对可对任意测试函数 \(\psi\) 使用
\(|\langle J^{-s}g,\psi\rangle|\le\|g\|_p\|J^{-s}\psi\|_{p'}\)；
任意实 \(s\) 时 \(J^{-s}\psi\) 仍为 Schwartz 函数。

\begin{proposition}\label{b2:prop:22-bessel-square}
对非齐次分割 \(P_0,\Delta_j\)（\(j\ge1\)），有
\[
 \|J^sf\|_p\asymp
 \left\|\left(|P_0f|^2+\sum_{j\ge1}2^{2js}|\Delta_jf|^2\right)^{1/2}\right\|_p.
\]
因而\chapref{b2:ch:19}定义的 \(F^s_{p,2}\) 与 \(H^{s,p}\) 相同，范数等价。
\end{proposition}
\begin{proof}
设 \(g=J^sf\)。第 \(j\) 个加权块的符号为
\(2^{js}\varphi(2^{-j}\xi)(1+4\pi^2|\xi|^2)^{-s/2}\)。
在归一化环带上，它及有限阶导数一致有界，核具有\chapref{b2:ch:19}向量值平方函数证明使用的一致大小、光滑与 \(L^2(\ell^2)\) 界。于是该证明给出
\(\|(\sum2^{2js}|\Delta_jf|^2)^{1/2}\|_p\le C\|g\|_p\)。
低频块同样由一个 Schwartz 卷积核控制。

反向利用环带上等于 \(1\) 的伴随截断 \(\widetilde\varphi_j\) 及分解恒等式。
与测试函数 \(h\in\mathcal S\) 配对，把 \(J^sf\) 写成块的重构，每项中的 \(J^s\) 除以 \(2^{js}\) 后又是一个一致零阶环带符号。逐点 Cauchy–Schwarz 与\chapref{b2:ch:19}相同的向量值界得到
\[
 |\langle J^sf,h\rangle|
 \le C\left\|\left(|P_0f|^2+\sum2^{2js}|\Delta_jf|^2\right)^{1/2}\right\|_p
          \|h\|_{p'}.
\]
先对有限和证明，再由平方函数尾部的支配收敛及分布重构取极限；对偶性给出反向界。这里没有用习题中的未证明等价，也没有把一般 \(F^s_{p,q}\) 的理论混进来。
\end{proof}
当 \(s\) 为非负整数，导数符号除以 \(J^s\) 满足 Mihlin 条件。反向则以有限个局部符号分解或椭圆符号恒等式控制 \(J^s\)，得到通常 Sobolev 导数范数的等价。仅在 \(p=2\) 时，这种等价可以完全不经过乘子理论而直接由 Plancherel 得到。

% 扩写定位：b2:subsec:220302
整数阶的识别可以把所需符号明确写出。设 \(k\geq1\) 为整数，\(1<p<\infty\)。正向符号
\[
 \frac{(2\pi\mathrm i\xi)^\alpha}{(1+4\pi^2|\xi|^2)^{k/2}},
 \qquad |\alpha|\leq k,
\]
满足 Mihlin 条件，所以 \(\sum_{|\alpha|\leq k}\|\partial^\alpha f\|_p\leq C\|J^kf\|_p\)。
反向先在 \(|\xi|\leq2\) 取一个低频截断，再将 \(|\xi|\geq1\) 的球面分成有限个坐标锥：每片有某个 \(|\xi_j|\geq c|\xi|\)。以光滑齐次分割 \(\chi_j\) 写成
\[
 (1-\chi_0(\xi))(1+4\pi^2|\xi|^2)^{k/2}
 =\sum_{j=1}^d a_j(\xi)(2\pi\mathrm i\xi_j)^k,
\]
其中 \(a_j=(1-\chi_0)\chi_j(1+4\pi^2|\xi|^2)^{k/2}/(2\pi\mathrm i\xi_j)^k\)，在其支集附近定义并平滑延拓为零。分母在该锥上有统一下界，每个 \(a_j\) 满足零阶导数条件。低频符号 \(\chi_0(1+4\pi^2|\xi|^2)^{k/2}\) 的核为 Schwartz 函数。Mihlin 与 Young 给出
\[
 \|J^kf\|_p\leq C\left(\|f\|_p+\sum_{j=1}^d\|\partial_j^kf\|_p\right).
\]
初始等式在 Schwartz 类上成立，再通过分布与光滑逼近得到 \(H^{k,p}=W^{k,p}\)。这里连纯坐标的最高阶导数就已足够控制全部混合导数；它是乘子分解的结论，而不是可以在任意区域或任意边界条件下直接套用的代数事实。

\ProseParagraphBreak

分数阶算子还具有一个反映非局部性的实空间公式。对 \(0<s<2\)，设
\[
 A_{d,s}=\int_{\mathbb R^d}\frac{1-\cos(2\pi z_1)}{|z|^{d+s}}\,\mathrm dz,
 \qquad c_{d,s}=\frac{(2\pi)^s}{A_{d,s}}.
\]
原点处 \(1-\cos(2\pi z_1)\leq C|z|^2\)，无穷处它有界，故 \(0<A_{d,s}<\infty\)；这同时解释范围 \(0<s<2\)。
\begin{proposition}\label{b2:prop:22-fractional-kernel}
对 \(f\in\mathcal S\)，
\[
 (-\Delta)^{s/2}f(x)=\frac{c_{d,s}}2
 \int_{\mathbb R^d}\frac{2f(x)-f(x+y)-f(x-y)}{|y|^{d+s}}\,\mathrm dy.
\]
并且
\[
 \langle(-\Delta)^{s/2}f,f\rangle
 =\frac{c_{d,s}}2
 \iint\frac{|f(x)-f(y)|^2}{|x-y|^{d+s}}\,\mathrm dx\,\mathrm dy.
\]
\end{proposition}
\begin{proof}
对称二阶差在原点是 \(O(|y|^2)\)，远场由函数有界和核尾部可积控制，故第一式绝对收敛。先在 \(\epsilon<|y|<R\) 上作 Fourier 变换，得到乘子
\[
 c_{d,s}\int_{\epsilon<|y|<R}
 (1-\cos(2\pi y\cdot\xi))|y|^{-d-s}\,\mathrm dy.
\]
旋转及换元 \(z=|\xi|y\) 把全积分变成
\(c_{d,s}|\xi|^s A_{d,s}=(2\pi|\xi|)^s\)。可用非负积分的单调收敛，或用 Schwartz 配对的支配收敛取极限，证明第一式。第二式由交换 \(x,y\)、展开平方得到，等价地也可由 Plancherel 与同一个非负乘子积分推导；非负性允许 Tonelli。
\end{proof}
若 \(f\geq0\) 且紧支集，在支集外有
\((-\Delta)^{s/2}f(x)=-c_{d,s}\int f(y)|x-y|^{-d-s}\,\mathrm dy<0\)，除非 \(f=0\)。
因此这个导数一般不保持支集，也不保持非负性；它的二次型却非负。这三个性质并不矛盾：非负的是总能量，而不是每一点的算子输出。与整数 Laplace 算子的局部性相比，这个尾部正是分数阶模型的本质区别。

频率定义中的 Bessel 势也有一个不振荡的实空间核。设 \(s>0\)，定义
\[
 G_s(x)=\frac1{\Gamma(s/2)}
 \int_0^\infty t^{s/2-1}e^{-t}
 (4\pi t)^{-d/2}e^{-|x|^2/(4t)}\,\mathrm dt.
\]
这是把热核作正权平均，权的总质量为一。它与齐次 Riesz 势的区别恰在额外因子 \(e^{-t}\)：大时间被截去，因而不会保留齐次核的长距离幂尾。
\begin{proposition}\label{b2:prop:22-bessel-kernel}
有 \(G_s\geq0\)、\(\int G_s=1\)、\(\widehat G_s(\xi)=(1+4\pi^2|\xi|^2)^{-s/2}\)。故 \(J^{-s}f=G_s*f\)，并在全部 \(1\leq p\leq\infty\) 上满足 \(\|J^{-s}f\|_p\leq\|f\|_p\)。
当 \(0<s<d\) 时，\(G_s(x)\asymp |x|^{s-d}\) 于 \(0<|x|\leq1\)；当 \(s=d\) 时量级为 \(1+\log(1/|x|)\)；当 \(s>d\) 时在原点附近有界。无穷远处三种情形均有指数衰减。
\end{proposition}
\begin{proof}
非负性与 Tonelli 给出总质量一。热核 Fourier 变换为 \(e^{-4\pi^2t|\xi|^2}\)，对 \(t\) 的 gamma 积分得到所述符号，Young 即给出压缩界。

记 \(r=|x|\)。在 \(0<r\leq1\) 时，积分的 \(t\geq1\) 部分有统一有限上界。对 \(0<t<1\)，用 \(e^{-1}\leq e^{-t}\leq1\)，再换元 \(u=r^2/(4t)\)，得到与
\[
 r^{s-d}\int_{r^2/4}^\infty
 u^{(d-s)/2-1}e^{-u}\,\mathrm du
\]
可比的表达。若 \(s<d\)，积分趋于一个正的 gamma 常数；若 \(s=d\)，零点附近是 \(\int_{r^2/4}^1\mathrm du/u\)；若 \(s>d\)，它的下限奇性恰好抵消外面的 \(r^{s-d}\)，结果有界。下界可以分别在固定 \(u\) 区间或 \(r^2/4<u<1\) 中取出，故这些不仅是上界。

当 \(r\geq1\)，在指数中使用 \(t+r^2/(4t)\geq r\)，抽出 \(e^{-r/2}\)。剩下的指数 \(e^{-t/2-r^2/(8t)}\) 把 \(t\to0\) 与 \(t\to\infty\) 的两端都控制；将 \(r\geq1\) 代入后，剩余积分有统一上界或至多一个固定多项式因子，这个因子可再吸收到较慢的指数 \(e^{-cr}\) 中。因此取得远场指数衰减。
\end{proof}
这个核解释了非齐次空间的低频项：常数频率附近的符号有界，所以 \(J^{-s}\) 保持同一 \(L^p\)；齐次符号 \(|\xi|^{-s}\) 在零频发散，改善可积性时必须改变目标指数或规定低频归一化。两者在高频上具有相同阶数，却不能在全空间上互相替代。

\ProseParagraphBreak

核估计还给出 Sobolev 嵌入的另一种表述。若 \(1<p<\infty\)、\(0<s<d/p\)，置 \(1/q=1/p-s/d\)。由 \(G_s(x)\leq C|x|^{s-d}\) 和 HLS，
\[
 \|f\|_q=\|G_s*(J^sf)\|_q\leq C\|J^sf\|_p.
\]
若 \(s>d/p\)，则 \(G_s\in L^{p'}\)：原点幂次条件正是 \((d-s)p'<d\)，临界对数在有限指数也可积，远场由指数衰减处理。因此
\(\|f\|_\infty\leq\|G_s\|_{p'}\|J^sf\|_p\)，并得到连续代表。连续性由 \(L^{p'}\) 中的核平移连续性给出，而且紧支集输入的逼近说明该代表在无穷远趋零。

\ProseParagraphBreak

在临界 \(s=d/p\) 时，核不属于 \(L^{p'}\)，不能原样保留无穷端点。取
\(h(x)=|x|^{-s}(\log(e/|x|))^{-\beta}\mathbf1_{\{|x|<1/2\}}\)，其中 \(1/p<\beta\leq1\)。径向积分给出 \(h\in L^p\)；正核下界则使 \(G_s*h(0)\) 的径向积分成为 \(\int_0^{1/2}r^{-1}(\log(e/r))^{-\beta}\,\mathrm dr=\infty\)。对 \(x\to0\)，Fatou 给出输出任意大的邻域值，故它没有本性有界的代表。这个反例不仅指出一个点的发散，也排除了通过修改零测集恢复 \(L^\infty\) 的可能。

\ProseParagraphBreak

最后，椭圆导数估计与这些嵌入处理不同问题。前者从方程右端恢复同指数的最高阶导数，后者把已经拥有的导数换成更高的可积指数或连续性。应用偏微分方程时，常见路线是先证明解处于 \(H^{s,p}\)，再按 \(s\) 与 \(d/p\) 的相对位置使用嵌入。若只看方程的阶数就声称解连续，会遗漏右端的可积指数及维数这两个决定性参数。

\ProseParagraphBreak

齐次导数与非齐次范数之间还有一个常用的精确比较。对 \(s>0\)、\(1<p<\infty\)，在 Schwartz 类上有
\[
 \|J^sf\|_p\asymp\|f\|_p+\|(-\Delta)^{s/2}f\|_p.
\]
正向用两个乘子：\(J^{-s}\) 以及
\((2\pi|\xi|)^s(1+4\pi^2|\xi|^2)^{-s/2}\)。
后者在原点不一定普通光滑，但定义域只要求 \(\xi\ne0\)，其尺度导数依然受控，故满足 Mihlin 条件。反向取低频截断 \(\chi\)，低频 \(J^s\chi\) 的核可积；高频符号
\[
 (1-\chi(\xi))\frac{(1+4\pi^2|\xi|^2)^{s/2}}{(2\pi|\xi|)^s}
\]
是零阶 Mihlin 符号，把它作用于齐次导数即可恢复高频 \(J^sf\)。两部分合并证明比较；连续延拓把 Schwartz 类的闭包识别为相同的图范数空间。

\ProseParagraphBreak

因此非齐次空间不是拒绝齐次导数，而是在其之外补上低阶 \(L^p\) 大小。高频上的两个定义等价，真正差别留在零频附近。使用能量方法时，可以在已说明的空间范围内自由改写这个图范数；使用全空间齐次逆、缩放不变估计或商掉多项式时，却必须保留齐次定义自己的归一化。一个范数比较不能自动解除所有齐次分布的低频歧义。

\subsection{椭圆型估计}
\label{b2:subsec:220303}

设 \(A=(a_{jk})\) 为实对称正定常矩阵，且
\(\nu|\xi|^2\le\xi^{\mathsf T}A\xi\le\Lambda|\xi|^2\)。
对于紧支撑光滑 \(u\)，令
\(Lu=-\sum a_{jk}\partial_j\partial_k u\)。则
\[
 \widehat{\partial_j\partial_k u}
 =-\frac{\xi_j\xi_k}{\xi^{\mathsf T}A\xi}\widehat{Lu},
 \qquad
 \|D^2u\|_p\le C_{d,p,\nu,\Lambda}\|Lu\|_p.
\]
商是零阶齐次符号，其球面导数由椭圆下界控制；Mihlin 定理给出结论。这是先验估计，既不要求 \(L^{-1}f\) 在 \(L^p\) 中存在，也不将低频问题忽略。

\ProseParagraphBreak
若改用 \(1+L\)，符号
\((1+4\pi^2\xi^{\mathsf T}A\xi)^{-1}\) 及其乘上 \(\xi_j\xi_k\) 的版本均满足零阶界，所以
\[
 \|u\|_p+\|D^2u\|_p\le C\|(1+L)u\|_p.
\]
这可以扩展为 \((1+L)^{-1}:L^p\to W^{2,p}\) 的有界解算子，分布方程成立，唯一性在这一空间中可通过温和分布 Fourier 变换验证。

\ProseParagraphBreak
对于变系数方程，不能把 \(a_{jk}(x)\) 当常数移到 Fourier 符号中，因为 Fourier 变换不与空间乘法交换。局部冻结系数可以形成一个近似常系数问题，但需要小振荡条件、交换子估计和覆盖拼接。\chapref{b2:ch:16}的椭圆正则理论与\appref{b2:app:H}的低正则理论分别处理这些额外困难，本章的常系数证明不替代它们。

常系数模型可以进一步解释定义域与低频唯一性的角色。对正定矩阵 \(A\)，符号 \(1+4\pi^2\xi^{\mathsf T}A\xi\) 没有零点，它的逆及各阶导数至多多项式增长或衰减，因而乘法保持 Schwartz 测试空间。若温和分布 \(u\) 满足 \((1+L)u=0\)，把 Fourier 方程乘以该逆立即得到 \(\widehat u=0\)，故唯一性甚至在整个温和分布类中成立。这里有一个真正可逆的低频项。

\ProseParagraphBreak

对没有常数项的 \(L\)，情况不同。任何常数乃至适当的调和多项式都是齐次方程的温和分布解；若以 \(D^2u\) 的先验界工作，这些低阶自由度不会被看到。在 \(L^p(\mathbb R^d)\)、有限 \(p\) 中，非零常数当然不属于空间，但是否存在具有所需衰减的逆解，仍要检查频率原点的奇性。二阶导数估计已经给出，并不意味着 \(L^{-1}\) 自动为同指数有界算子。

\ProseParagraphBreak

在对角例子 \(L=-\sum_j a_j\partial_j^2\)、\(a_j>0\) 中，换元 \(x_j=\sqrt{a_j}\,y_j\) 将它化为标准 \(-\Delta_y\)。各个导数却分别得到 \(a_j^{-1/2}\) 或 \(a_j^{-1}\) 因子，因此当某个 \(a_j\downarrow0\) 时，椭圆正则估计会退化。一般正定矩阵可先正交对角化，正交变换保持 Lebesgue 测度与欧氏导数长度；最小特征值控制的正是这种退化。记录 \(\nu,\Lambda\) 的依赖有实际意义，而不是可删去的形式装饰。

\ProseParagraphBreak

带参数的预解算子把同一估计组织成一个尺度一致的族。对 \(\lambda>0\)，令 \(u=(\lambda+L)^{-1}f\)。以 \(\xi=\sqrt\lambda\,\zeta\) 归一化后，三个符号
\[
 \frac{\lambda}{\lambda+4\pi^2\xi^{\mathsf T}A\xi},\quad
 \frac{\sqrt\lambda\,2\pi\mathrm i\xi_j}{\lambda+4\pi^2\xi^{\mathsf T}A\xi},\quad
 \frac{-4\pi^2\xi_j\xi_k}{\lambda+4\pi^2\xi^{\mathsf T}A\xi}
\]
都具有与 \(\lambda\) 无关的 Mihlin 常数。于是
\[
 \lambda\|u\|_p+\sqrt\lambda\|\nabla u\|_p+\|D^2u\|_p
 \leq C_{d,p,\nu,\Lambda}\|f\|_p.
\]
符号缩放保持同指数算子范数，这是统一性的关键；直接逐项对未归一化符号求导往往会出现不必要的 \(\lambda^{-1}\) 因子。这里导数和零阶项分别以自己的尺度加权，也正好满足微分阶数的平衡。

\ProseParagraphBreak

令 \(\lambda\downarrow0\) 时，最高阶导数界仍统一，而对 \(\|u\|_p\) 的估计允许增长为 \(\lambda^{-1}\)。因此它没有绕过齐次逆算子的低频障碍。若需要在零参数得到实际解，必须另外控制数据的低频或选择齐次空间；有界参数预解族并不自动给出一个同指数的无质量逆。

\subsection{球面乘子 预告}
\label{b2:subsec:220304}

径向符号只依赖 \(|\xi|\)，仍可能在球面上发生跳跃。
光滑环带截断 \(\chi(|\xi|)\) 的核属于 \(L^1\)；球特征函数 \(\mathbf1_{\{|\xi|\le1\}}\) 则不满足任何跨边界的 Mihlin 光滑条件。
把跳跃平滑成宽度 \(\epsilon\) 的过渡层时，第 \(N\) 阶导数通常增长如 \(\epsilon^{-N}\)。对每个固定 \(\epsilon\) 应用定理可得有界性，却得不到 \(\epsilon\downarrow0\) 的统一界。

\ProseParagraphBreak
Bochner–Riesz 平均以 \((1-|\xi|^2)_+^\delta\) 逐渐消减边界跳跃，其所需光滑阶数与 \(p\)、维数有关。Mihlin 定理给出方便的充分条件，但常不是锐利范围。\chapref{b2:ch:24}用径向核的振荡和曲率研究这一问题；\chapref{b2:ch:23}的限制估计控制频率薄壳上的集中，则说明为什么乘子问题会与曲面延拓相遇。

\ProseParagraphBreak
因此“径向”“有界”“紧支撑”三项都不能单独替代乘子条件。对核使用 Young 不等式、对符号使用尺度导数、对薄壳使用曲面几何，是三条不同证明路线。应用时应明确正在使用哪一条，以及它在端点或退化极限中会失去哪一个条件。

平滑乘子还有一个时间演化模型，可与频率硬截断比较。热算子的符号是 \(e^{-4\pi^2t|\xi|^2}\)。其空间核非负，积分为一，因而在全部 \(L^p\) 上压缩；对每个整数 \(k\geq0\)，对 Gaussian 核求导及换元给出
\[
 \|D^ke^{t\Delta}f\|_p\leq C_k t^{-k/2}\|f\|_p,
 \qquad 1\leq p\leq\infty,\quad t>0.
\]
导数核是一个多项式乘 Gaussian，其 \(L^1\) 范数有限；时间缩放产生 \(t^{-k/2}\)，故该证明连两个端点也包括。单纯的开区间 Mihlin 定理已经足以处理 \(1<p<\infty\)，但核方法在这个特定模型中提供了额外端点。

\ProseParagraphBreak

任意实 \(s\geq0\) 时同样有
\[
 \|e^{t\Delta}f\|_{H^{s,p}}\leq C_{s,p}(1+t^{-s/2})\|f\|_p,
 \qquad1<p<\infty.
\]
在 \(0<t\leq1\) 中，把 \(t^{s/2}(1+4\pi^2|\xi|^2)^{s/2}e^{-4\pi^2t|\xi|^2}\) 以 \(\zeta=\sqrt t\,\xi\) 重写，成为
\((t+4\pi^2|\zeta|^2)^{s/2}e^{-4\pi^2|\zeta|^2}\)。
各个尺度导数在原点附近由 \(|\zeta|^{|\alpha|}\) 抵消可能的负幂，在无穷处由 Gaussian 控制，所以 Mihlin 常数统一；\(t\geq1\) 则用半群性将一次固定时间的平滑与剩余压缩复合。这个推导说明分数阶平滑并不要求先把 \(s\) 取整。

\ProseParagraphBreak

当 \(t\downarrow0\)，正阶范数常数发散，与一般 \(L^p\) 数据可能没有任何正阶导数完全相容。对每个固定正时间得到光滑函数，不意味着这些函数在更强空间中一致趋向原数据。连续半群的初值收敛发生在原来给定的空间，提升正则性则付出时间尺度。

\ProseParagraphBreak

把时间倒过来，符号成为 \(e^{4\pi^2t|\xi|^2}\)，甚至失去平方乘子的本性有界性。取 \(\widehat g\) 支于单位球，并令 \(f_N(x)=e^{2\pi\mathrm iNx_1}g(x)\)。输入平方范数不变，而逆热像的平方范数至少为
\(e^{4\pi^2t(N-1)^2}\|g\|_2\)。适当缩小输入振幅，可以使初值在 \(L^2\) 中趋零而逆演化的像趋于无穷。这个明确的高频检验说明正向平滑不能通过形式上的指数取反得到稳定逆过程。

\ProseParagraphBreak

与球面乘子相比，热符号没有一条频率边界上的跳跃，所有尺度的导数也具有统一控制；与齐次分数积分相比，它在零频不发散。良好的端点压缩、正时间平滑和低频稳定是这些特殊性质共同产生的，不应当归给任意一个写成 \(m(\xi)\) 的符号。

\ProseParagraphBreak

频率过渡层还可以用一个明确的算子范数检验。取固定 \(\chi\in C_c^\infty\)，令 \(m_R(\xi)=\chi(\xi/R)\)。其核为 \(R^d\check\chi(Rx)\)，所以全部 \(R>0\) 的 \(L^1\) 核范数都等于 \(\|\check\chi\|_1\)，在每个 \(L^p\) 上有同一个界。若 \(\chi(0)=1\)，在 \(R\to\infty\) 时，有限指数下这族算子强趋向恒等：先对 Schwartz 输入用核的近似恒等性，再用一致范数和稠密性推广。将频率支撑推向无穷本身并不造成常数发散；关键是过渡轮廓经过尺度归一化后仍固定。

\ProseParagraphBreak

若改用 \(m_\epsilon(\xi)=\theta((|\xi|-1)/\epsilon)\)，其中 \(\theta\) 固定而过渡层紧贴单位球面，归一化整个频率空间不能同时固定球的半径和层的厚度。径向一阶导数在该层上达到 \(\epsilon^{-1}\)，而 \(|\xi|\) 仍约为一，尺度导数条件因此真实地增长。两个模型都叫作光滑截断，却有不同的无量纲参数：前者只有整体半径，后者还有厚度与半径之比。将其混同，会把 Mihlin 常数的一致性误判成一个无条件事实。

\ProseParagraphBreak

在平方空间中只需符号的本性上确界，故上述有界过渡层即便变得越来越薄，仍可具有统一平方界。对其他指数则必须利用核、导数或特殊几何。这个差别再次说明，平方乘子的全部刻画并没有自动延伸到任意 \(L^p\)。讨论逼近一个不连续乘子时，要分别检查固定参数有界、一致参数有界和极限算子的有界性；只有前两项共同存在，才可能通过稠密输入上的极限传递到后一项。


\begin{chaptersummary}
零阶符号的 \(L^2\) 起点由 Plancherel 提供，环带光滑性使核差分可控，Calderón–Zygmund 分解和插值再推广到 \(1<p<\infty\)。Hörmander 条件展示了证明真正需要的积分控制。低频奇性、高频阶数与符号边界几何各有独立作用，不能用一个笼统的“符号光滑”代替核验。
\end{chaptersummary}

\BookExercises

% 自拟：围绕本章已经建立的工具，改变条件或组合方法。
\begin{exercise}\label{b2:ex:22-new-1}
证明非零常数函数 \(m\) 的乘子范数为 \(|m|\)，并说明其核为何不属于普通函数意义下的 \(L^1\)。
\end{exercise}
\begin{solution}
算子为 \(mf\)，故每个 \(L^p\) 上范数均为 \(|m|\)。分布核为 \(m\delta_0\)。若它由 \(L^1\) 函数表示，其 Fourier 变换按 Riemann–Lebesgue 引理应在无穷远趋零，与非零常数矛盾。
\end{solution}

% 自拟：围绕本章已经建立的工具，改变条件或组合方法。
\begin{exercise}\label{b2:ex:22-new-2}
验证 \(m(\xi)=e^{\mathrm i\log|\xi|}\) 满足任意有限阶 Mihlin 条件。它能否连续延拓到原点？
\end{exercise}
\begin{solution}
该符号的模为 \(1\)。每次求导产生一个 \(|\xi|^{-1}\) 因子及球面上的光滑有界函数，归纳得 \(|\partial^\alpha m|\le C_\alpha|\xi|^{-|\alpha|}\)。沿径向趋零时相位不断旋转，无极限；这不影响 \(1<p<\infty\) 的乘子结论。
\end{solution}

% 自拟：围绕本章已经建立的工具，改变条件或组合方法。
\begin{exercise}\label{b2:ex:22-new-3}
证明符号缩放 \(m(r\xi)\) 保持 \(M_p\) 范数，并准确写出共轭公式。
\end{exercise}
\begin{solution}
对 \(D_rf(x)=f(rx)\)，有 \(\widehat{D_rf}(\xi)=r^{-d}\widehat f(\xi/r)\)。直接计算得 \(D_{1/r}T_mD_r=T_{m(r\,\cdot)}\)。两缩放算子的范数为 \(r^{-d/p}\) 与 \(r^{d/p}\)，给出一个方向；用 \(1/r\) 再做一次得到反向，故范数相等。
\end{solution}

% 自拟：围绕本章已经建立的工具，改变条件或组合方法。
\begin{exercise}\label{b2:ex:22-new-4}
设 \(m\) 支撑在固定环带，具有 \(N>d\) 阶连续导数。证明 \(T_m\) 在包括端点的所有 \(L^p\) 上有界。
\end{exercise}
\begin{solution}
分部积分给出 \(|\check m(x)|\le C(1+|x|)^{-N}\)，近原点直接由 \(\|m\|_1\) 控制。由于 \(N>d\)，核可积；Young 不等式给出 \(1\le p\le\infty\) 的界。这个证明只针对单个紧频率环带，并未对无限尺度求和。
\end{solution}

% 自拟：围绕本章已经建立的工具，改变条件或组合方法。
\begin{exercise}\label{b2:ex:22-new-5}
设 \(0<\alpha<d\)。说明为什么 \((2\pi|\xi|)^{-\alpha}\) 不是 \(L^2\) 乘子，却仍可定义分数积分。
\end{exercise}
\begin{solution}
符号在每个零点邻域本质无界，\cref{b2:prop:22-l2}排除 \(L^2\) 有界性。其逆变换是\chapref{b2:ch:19}的 Riesz 位势核，可以在 \(1<p<d/\alpha\) 时定义为 \(L^p\to L^q\)、\(1/q=1/p-\alpha/d\) 的有界算子。空间指数改变，所以两结论一致。
\end{solution}

% 自拟：围绕本章已经建立的工具，改变条件或组合方法。
\begin{exercise}\label{b2:ex:22-new-6}
对 \((1-\Delta)u=f\)，用乘子证明 \(\|u\|_p+\sum_{j,k}\|\partial_j\partial_ku\|_p\le C\|f\|_p\)。
\end{exercise}
\begin{solution}
解符号为 \((1+4\pi^2|\xi|^2)^{-1}\)，二阶导数符号为 \(-4\pi^2\xi_j\xi_k/(1+4\pi^2|\xi|^2)\)。其所有所需尺度导数由有理函数求导直接控制，原点附近有界，高频为零阶或负阶。分别应用 Mihlin 定理并对有限个导数相加即得。
\end{solution}

% 自拟：围绕本章已经建立的工具，改变条件或组合方法。
\begin{exercise}\label{b2:ex:22-new-7}
若 \(\eta\) 为固定光滑环带截断，解释为何 Hörmander 证明同时控制 \(\|h_j\|_1\) 与 \(\|\nabla h_j\|_1\)，而没有假定 \(m\) 可微。
\end{exercise}
\begin{solution}
\(H^s\) 控制转化为 \((1+|x|)^sh_j\) 的 \(L^2\) 控制。梯度对应归一化频率中的坐标乘法；该坐标先乘一个紧支撑光滑截断，在 \(H^s\) 上有界，所以梯度具有相同带权 \(L^2\) 界。因 \(s>d/2\)，逆权重属于 \(L^2\)，Cauchy–Schwarz 给出两项 \(L^1\) 界。它使用频率乘法，不使用 \(m\) 的逐点导数。
\end{solution}

% 自拟：围绕本章已经建立的工具，改变条件或组合方法。
\begin{exercise}\label{b2:ex:22-new-8}
证明一维 \(\mathbf1_{[a,b]}\)、\(a<b\)，是所有 \(1<p<\infty\) 上的乘子，且范数可独立于区间选择。
\end{exercise}
\begin{solution}
Hilbert 变换符号为 \(-\mathrm i\operatorname{sgn}\xi\)，故正半轴投影为 \((I+\mathrm iH)/2\)。调制将半轴边界移到 \(a\) 或 \(b\)，这些操作保持范数。两个相移半轴投影之差的符号为 \(\mathbf1_{[a,b]}\)（端点值无关），算子范数至多 \(1+\|H\|_{p\to p}\)，独立于区间。
\end{solution}

